\documentclass[10pt]{amsart}
\usepackage{amsmath,amssymb,amsthm,mathtools,mathrsfs}
\usepackage[a4paper,hmargin=0.76in,vmargin=0.60in]{geometry}
\usepackage{booktabs}
\usepackage{array}
\usepackage{longtable}
\usepackage{tabularx}
\usepackage[hypertexnames=false]{hyperref}
\hypersetup{hidelinks}
\usepackage[final,expansion=false]{microtype}
\usepackage{url}
\emergencystretch=2.5em

\newtheorem{theorem}{Theorem}[section]
\newtheorem{proposition}[theorem]{Proposition}
\newtheorem{lemma}[theorem]{Lemma}
\newtheorem{corollary}[theorem]{Corollary}
\theoremstyle{definition}
\newtheorem{definition}[theorem]{Definition}
\newtheorem{problem}[theorem]{Problem}
\theoremstyle{remark}
\newtheorem{remark}[theorem]{Remark}

\newcommand{\C}{\mathbb C}
\newcommand{\R}{\mathbb R}
\newcommand{\N}{\mathbb N}
\newcommand{\Z}{\mathbb Z}
\newcommand{\Q}{\mathbb Q}
\newcommand{\Rel}{\operatorname{Re}}
\newcommand{\Iml}{\operatorname{Im}}
\newcommand{\supp}{\operatorname{supp}}
\newcommand{\TV}{\operatorname{TV}}
\newcommand{\Kp}{\mathcal K^{\mathrm{prim}}_{2T}}
\newcommand{\Kw}{\mathcal K_{2T}}
\newcommand{\PT}{\mathcal P_T}
\newcommand{\IT}{I_T}
\newcommand{\lamP}{\lambda^{\mathrm{prim}}_T}
\newcommand{\lamW}{\lambda^{\mathrm W}_T}
\newcommand{\lamA}{\lambda_{\mathrm A}(T)}
\newcommand{\sT}{\sigma_T}
\newcommand{\gh}[1]{\widehat g_{#1}}
\newcommand{\nuh}{\widehat\nu}
\newcommand{\wh}{\widehat w}
\newcommand{\Fh}{\widehat F}
\newcommand{\cF}{c_F}
\newcommand{\Lev}{\mathfrak E}
\newcommand{\Wcl}{\mathcal W_T}
\newcommand{\NT}{\mathcal N_T}
\newcommand{\Aconst}{\mathsf A_T}
\newcommand{\Ainf}{\mathcal A_\infty}
\newcommand{\PiT}{\Pi_T}
\newcommand{\mT}{\mathfrak m_T}
\newcommand{\laminf}{\lambda^{\infty}_{\min}}
\newcommand{\laminfP}{\lambda^{\infty,\mathrm P}_{\min}}
\newcolumntype{Y}{>{\raggedright\arraybackslash}X}

\providecommand{\anPa}{\Phi_a}
\providecommand{\anPb}{\Phi_b}
\providecommand{\aneu}{\mathrm{eu}_d}
\providecommand{\auAut}{\operatorname{Aut}}
\providecommand{\auAz}{\mathcal A_0}
\providecommand{\auHom}{\operatorname{Hom}}
\providecommand{\auhdir}{h_{\mathrm{dir}}}
\providecommand{\pao}{\mathfrak o}
\providecommand{\pcMat}{\mathrm M}
\providecommand{\pcR}{R}
\providecommand{\pcm}{\mathfrak m}
\providecommand{\pgsdu}{\mathrm{du}}
\providecommand{\pkiC}{\mathcal C}
\providecommand{\pkiWin}{\operatorname{Win}}
\providecommand{\pkilog}{\operatorname{lg}}
\providecommand{\pkjG}{G_{\Q_\ell}}
\providecommand{\pkjPb}{\mathsf P_b}
\providecommand{\pkjPl}{\mathsf P}
\providecommand{\pkwB}{B_{\mathrm{cris}}}
\providecommand{\pkwD}{D_{\mathrm{pcris}}}
\providecommand{\plone}{\mathbf 1}
\providecommand{\poA}{\mathcal A}
\providecommand{\psHS}{\mathcal H}
\providecommand{\psM}{\mathsf M}
\providecommand{\psZ}{Z}
\providecommand{\psrho}{\varrho}
\providecommand{\psvac}{\mathbf v}
\newenvironment{quotedthm}[2]{\par\medskip\noindent\begingroup
  \textbf{#1}#2.\ \itshape\ignorespaces}{\endgroup\par\medskip}
\newenvironment{quoteddef}[2]{\par\medskip\noindent\begingroup
  \textbf{#1}#2.\ \ignorespaces}{\endgroup\par\medskip}
\usepackage{xurl}
\begin{document}
\addtolength{\abovedisplayskip}{-3.4pt}
\addtolength{\belowdisplayskip}{-3.4pt}
\addtolength{\belowdisplayshortskip}{-2pt}
\title[The Covolume Theorem: regular families of twists]{The Covolume Theorem: regular families of twists, Selmer groups and Ribet's method}
\author{David Alejandro Trejo Pizzo}
\date{}
\begin{abstract}
Let $E$ be a CM field and consider the theta lifts, between unitary groups over $E$, that occur
in the Covolume Theorem of the main paper of this series. We prove that for every prime $\ell$
there is an anticyclotomic character of $E$ of $\ell$-power order, tamely ramified at auxiliary
primes, such that after twisting by its first $m$ powers the ratios of twisted to untwisted
Hecke $L$-values have bounded $\ell$-adic valuation and the congruences modulo powers of $\ell$
between theta lifts and other forms are bounded, with bounds whose sum over all primes is
controlled. The twist is chosen so that Bloch--Kato Selmer groups of the Galois representations
attached to the lifts are small; the passage from Selmer groups to congruences is Ribet's method,
carried out modulo powers of every prime $\ell$ with the help of auxiliary places.
\end{abstract}
\maketitle
\setcounter{tocdepth}{1}\tableofcontents

\section{Introduction}\label{C3:intro}

The main paper~\cite{P46} of this series deduces the Riemann hypothesis for Dirichlet
$L$-functions from three theorems on integral algebraic modular forms on a chain of $m$ theta
lifts between unitary groups over a CM field $E$: the Short Forms Theorem, the Spanning Theorem
and the Covolume Theorem. The Covolume Theorem is a lower bound obtained from the product
formula, in which the integral structure at a prime $\ell$ is tested after twisting the lifts by
a character that may depend on~$\ell$. Its proof requires, for each $\ell$, a twist for which the
quantities that enter have the least possible divisibility by~$\ell$. The precise requirement is
that of a \emph{regular family} of twists ({\cite[Definition~14.39]{P46}}), and the purpose of
this paper is to prove that regular families exist ({\cite[Proposition~14.40]{P46}}).

We state the result for one prime. The notation is that of
{\cite[\S\S14.1--14.10]{P46}}. A \emph{twist} is a unitary automorphic character
$\vartheta$ of $\mathrm U(1)$ over the maximal totally real subfield $E^+$, trivial at the
infinite places and tamely ramified at a finite set $Q_\vartheta$ of primes of $E$, disjoint
from its conjugate and prime to the level; it is \emph{$m$-fold} if the first $m$ powers of its
tame parts are non-trivial. For $1\le j\le m$ the $j$-th lift is the theta lift, formed with the
character $\mu_j^0\vartheta$, from good forms on the unitary group of a hermitian space $W_j$ of
rank $r$ to algebraic modular forms on the definite unitary group of rank $r+t$; we also
consider, for certain pairs $(\nu,z)$ of a character and an integer $m<z\le r$ (the \emph{main
pairs}), the lift of $\nu\vartheta$ from $\mathrm U(1)$ to the definite unitary group of rank
$2z$. The denominators of the projection to a space of lifts are measured on the integral forms
orthogonal to the false lifts (Definition~\ref{pg:def-den}). Fix a prime $\ell$ and an isomorphism
$\iota\colon\mathbb C\to\overline{\mathbb Q}_\ell$, and write $v$ for the valuation of
$\overline{\mathbb Q}_\ell$ with $v(\ell)=1$, transported to $\mathbb C$ by $\iota$.
Theorem~\ref{pg:main} asserts that there is an $m$-fold twist $\tau$ with
$\ell\nmid\operatorname{Nm}Q_\tau$ such that, for every ramified critical pair $(\nu,z)$, every
$1\le n\le m$ and both signs,
\[
v\Bigl(\frac{|L(z,(\nu\tau^{\pm n})^{-1})|}{|L(z,\nu^{-1})|}\Bigr)\le e'_\ell ,
\]
such that for $1\le j\le m$ the projection to the space of $j$-th lifts formed with $\tau^{\,j}$
has denominator at most $\beta_\ell$ at $\iota$, and such that the corresponding bound holds for
the lines of the lifts of the main pairs. The constants are explicit:
\[
e'_\ell=(4f+3)\,r\,v(N_{\mathrm{ch}})+2^{-\ell},\qquad
\beta_\ell=\mathrm{loc}_\ell+(1+2t)\,e'_\ell+e_\ell+v(N_{\mathrm{ch}}),\qquad
\mathrm{loc}_\ell=4f\cdot10^6\,\frac mk\bigl(b\,v(M)+\lambda_\ell\bigr),
\]
where $2f=[E:\mathbb Q]$, $N_{\mathrm{ch}}$ is the product of the orders of the groups
$\mathrm{GL}_{rt}$ over the residue fields of the chosen places, $e_\ell$ vanishes unless
$\ell$ divides one of two fixed integers, and
$\lambda_\ell$ is the valuation of an explicit product of integers not exceeding $2b$, where $b$
is the largest weight (see \eqref{pg:eq-beta}). What matters is their sum: by Lemma~\ref{pg:sums},
$\sum_\ell e'_\ell\log\ell\le x_1$ and $\sum_\ell(\beta_\ell+v(h_1))\log\ell\le
f\,a\,(\log b)^3/(mk)$, which is the budget allowed in~\cite{P46}. A family of twists
$\theta_\iota=\tau$ chosen in this way for each $\iota$ is regular.

The statement and its proof are the same for every prime: for $\ell=2$, for the primes that
ramify in $E$, for those inert or split in the imaginary quadratic subfield $K$, and for the
primes of the level. This uniformity is necessary, since the primes below $x_1$ that are inert
in $K$ carry the prime sum that~\cite{P46} estimates, and it excludes an appeal to the published theorems on
the non-vanishing modulo $\ell$ of Hecke $L$-values of anticyclotomic twists, for instance
\cite{Finis,HidaNonvanishing,Prasanna}, which assume $\ell$ odd, unramified and split in the CM
extension. Three devices take their place. The twist has $\ell$-power order and conductor prime to~$\ell$, so that it
is a deformation of the trivial twist, and only upper bounds for valuations at one point near
the trivial character are required. The types of the forms at the two primes $p_0$, $p_1$ of the
level are chosen so that the Galois representations in question are
residually absolutely irreducible for every~$\ell$. And above $\ell$ no ordinarity or
Fontaine--Laffaille condition is imposed: the failure of torsion classes to be crystalline is
bounded by the term $\mathrm{loc}_\ell$.

The proof has three parts (\S\ref{pg:ss-reduction}). The first attaches Galois representations.
By the theorems on Galois representations of automorphic representations of unitary groups
\cite{Labesse,BLGGT,HarrisTaylor,Shin,ChenevierHarris,CaraianiDuke,CaraianiANT,BLGGTlp1,BLGGTlp2},
an eigenform $y$ on the definite group that has the Hecke polynomials of a $j$-th lift has a
Galois representation whose trace, after a normalising twist, is
\[
\operatorname{tr}\bigl(S_y\otimes\eta\bigr)(g)=\psi_\vartheta(g)\,\operatorname{tr}R(g)
+\sum_{e=0}^{t-1}\chi(g)^e ,
\]
where $\psi_\vartheta$ is the anticyclotomic character of $G_E$ attached to $\vartheta$, $\chi$
is the inverse of the cyclotomic character, and $R$ is a \emph{main summand}
(Definition~\ref{pg:def-summand}): a semisimple representation of rank $r$, conjugate self-dual up to a power of $\chi$,
crystalline or potentially crystalline above $\ell$ with Hodge--Tate weights in prescribed
windows, of which at most the proportion $m/k$ are small, and all of whose lattices are
residually absolutely irreducible. This is Theorem~\ref{pg:I}; the properties of the Galois
representations that it uses are derived from the cited theorems in Section~\ref{C3:g5}.

The second part is the twisting theorem (Theorem~\ref{pg:II}): given finitely many main summands
$R_i$, there is an $m$-fold twist $\tau$, ramified at primes of prime norm away from~$\ell$, such
that for $\vartheta=\tau^{\,j}$, $1\le j\le m$, each $\psi_\vartheta R_i$ has a lattice whose
Selmer group with coefficients in $K/\mathcal O$, with the Bloch--Kato condition above $\ell$
and ramification allowed at $Q_\tau$, has length at most $e'_\ell$, and such that the ratios of
$L$-values satisfy the bound displayed above. The Selmer groups of all the twists ramified at a
given set of auxiliary primes are described by one system of linear equations
(Lemma~\ref{pg:minor}); the twist is then produced by a count of Greenberg--Wiles
type~\cite{WilesFLT} (Lemma~\ref{pg:count}), in which the local conditions above $\ell$ enter
only through the dimension formula of Bloch and Kato~\cite{BlochKato}, and by a descent in which
auxiliary primes are added one at a time (Proposition~\ref{pg:step}). The bound for the $L$-values is obtained in the same framework from one fact on
theta lifts, taken from~\cite{P54}: the normalised norms of the lifts of characters are integral
linear combinations of the values of the twist. No integral $\ell$-adic Hodge theory enters
here.

The third part passes from Selmer groups to denominators (Theorem~\ref{pg:III}). A denominator
$\ell^\nu$ in the projection of an integral form to the space of lifts is a congruence modulo
$\ell^\nu$ between a lift and eigenforms whose normalised Galois representations do not have the
eigenvalue $1$ identically, and the method of Ribet~\cite{Ribet1976} extracts from it a non-split
extension, modulo a power of $\ell$, between two of the $t+1$ summands on the right of the trace
formula above. Extensions between $\psi_\vartheta R$ and a character are bounded by the Selmer
groups of the second part. Extensions between two powers of $\chi$ are classes in Tate twists
of the trivial module over $E$. The classes that change sign under complex conjugation are bounded,
for every $\ell$, by an integer $c_{\mathrm{odd}}$ with $\log c_{\mathrm{odd}}\le x^{800}$
(Proposition~\ref{pri:odd}); we prove this in \S\ref{pod:sec} by Kolyvagin's Euler system
of Gauss sums in the form given by Rubin~\cite{RubinTexel}, made effective at finite level and
combined with a count from class field theory for the small primes. The
invariant classes are governed by class groups of real cyclotomic fields and cannot be bounded.
We remove them by passing to forms of auxiliary level at places chosen in a triangular order, at
each of which one such class is killed (Proposition~\ref{pri:triangle}, which rests on a count
of Greenberg--Wiles type for Kummer maps, \S\ref{pku:sec}), and by carrying out Ribet's
method to first order in parameters that tie tame inertia at these places to Frobenius
(Proposition~\ref{pri:ribet}). The operators that lead back from auxiliary level to the original forms
(Proposition~\ref{pri:norm}) depend on a local statement on unitary representations of $\mathrm{GL}_N$ of a $p$-adic field, the spherical
chain lemma (Theorem~\ref{pay:chain}), which we prove directly.

Ribet's method produces cocycles modulo $\ell^\nu$ that satisfy all linear relations satisfied by
the Galois representations of the forms, and these need not satisfy the Bloch--Kato condition
above~$\ell$. The length of the module of such cocycles modulo those that do is the
\emph{local loss} (Proposition~\ref{pg:loss}). In Section~\ref{pll:sec} we reduce it, by local
Tate duality and the results of Bloch and Kato on $H^1_f$~\cite{BlochKato,NSW}, to a dual
quantity defined by cup products, and then to the existence of an exponent of $\ell$ that
carries the crystalline cocycles of the Tate dual into the classes orthogonal to a given class
of crystalline lattices (Proposition~\ref{pll:dual}, Lemma~\ref{pll:exp}). The exponents needed
are stated in Theorem~\ref{pll:exponent} and proved in~\cite{P56}; they give the bound
$\mathrm{loc}_\ell$ for the loss per unit of rank of the main summand, in which the factor $m/k$
reflects the small number of low weights, and there is no loss when $\ell>2b$ and $\ell$ does
not divide $M=p_0p_1$. This is the only use of integral $\ell$-adic Hodge theory in the proof.

The notation and the data are those of~\cite{P46}. Theorem~\ref{pll:exponent} is stated in
Section~\ref{pll:sec} and proved in~\cite{P56}; the deduction of the local loss from it (Remark~\ref{pll:status}) also quotes from~\cite{P56} the chart of the data ({\cite[Theorem~7.9]{P56}}, {\cite[Definition~2.4]{P56}}). That paper takes from this one the definitions of
Section~\ref{pll:sec}, the elementary Lemma~\ref{pll:exp}, and the properties of the Galois
representations of the lifts above $\ell$ and at the primes of the level
(Definition~\ref{pg:def-summand}, Proposition~\ref{pga:G45}, Lemma~\ref{pga:depth0},
Proposition~\ref{pri:BL}); none of these depends on Theorem~\ref{pll:exponent} or on the local loss, so the two papers do not form
a circle. From~\cite{P54} we take the statement (IM) on the norms of the lifts of characters
({\cite[Propositions~3.8 and~3.12]{P54}}) and the equivariance and the Hecke equation of a
lift ({\cite[Lemma~2.2]{P54}}). They are proved there without regular families, which~\cite{P54}
uses only as the hypothesis of its bound for the heights of Hecke values. We also use the
equality of multiplicities of~\cite{P51} and some properties of the lifts from~\cite{P49}
and~\cite{P53}. The statements quoted are reproduced in Appendix~\ref{C3:quoted}.

The paper is organised as follows. Section~\ref{pg:sec} states Theorem~\ref{pg:main}, sums the
constants over the primes, reduces the theorem to Theorems~\ref{pg:I}, \ref{pg:II}
and~\ref{pg:III}, and proves the parts of Theorems~\ref{pg:I} and~\ref{pg:II} that do not require
further input. Section~\ref{pse:sec} completes the proof of the twisting theorem.
Section~\ref{pri:sec} proves Theorem~\ref{pg:III}, and Section~\ref{pll:sec} treats the local
loss. Section~\ref{C3:g5} supplies the properties of Galois representations used in
Theorem~\ref{pg:I}. The last two sections prove the inputs of Ribet's method: the odd classes and
Kummer maps in Section~\ref{C3:g6}, and in Section~\ref{C3:g7} auxiliary statements on forms and
on lattices (\S\ref{pax:sec}), the spherical chain lemma (\S\ref{pay:sec}), and the characters,
the norm induction and the base level (\S\ref{prm:sec}).


\providecommand{\auO}{\mathcal O}
\providecommand{\auI}{\mathcal I}
\providecommand{\auQlb}{\overline{\Q}_\ell}
\providecommand{\auPl}{\Phi}
\providecommand{\auPb}{\Phi_b}
\providecommand{\auT}{\mathcal T}
\providecommand{\auG}{\overline G}
\providecommand{\audT}{d_{\auT}}
\providecommand{\auNm}{\operatorname{Nm}}
\providecommand{\auU}{\mathrm U}
\providecommand{\auGL}{\mathrm{GL}}
\providecommand{\auA}{\mathbb A}
\providecommand{\auF}{\mathbb F}
\providecommand{\aurd}{\mathrm{rd}_E}
\providecommand{\aueu}{\varepsilon_{\mathrm u}}
\providecommand{\aubud}{\mathrm{bud}}
\providecommand{\aubar}[1]{\overline{#1}}
\providecommand{\auf}[1]{\texttt{#1}}
\providecommand{\pgZb}{\overline{\Z}_\ell}
\providecommand{\pgv}{v}
\providecommand{\pgNch}{N_{\mathrm{ch}}}
\providecommand{\pgcodd}{c_{\mathrm{odd}}}
\providecommand{\pgloc}{\mathrm{loc}}
\providecommand{\pgG}{G_E}
\providecommand{\pgSel}{\operatorname{Sel}}
\providecommand{\pglen}{\operatorname{len}}
\providecommand{\pgtr}{\operatorname{tr}}
\providecommand{\pgFrob}{\operatorname{Frob}}
\providecommand{\pgHom}{\operatorname{Hom}}
\providecommand{\pgS}{S_{\mathrm{bad}}}
\providecommand{\pgcyc}{\varepsilon_\ell}
\providecommand{\pgpr}{\operatorname{pr}}

\section{Regular families of twists}\label{pg:sec}

This section concerns {\cite[Proposition~14.40]{P46}}: for all admissible data there is a regular
family of $m$-fold twists. We state exactly what is asserted for one prime $\ell$ and one
isomorphism $\iota\colon\C\to\auQlb$ (Theorem~\ref{pg:main}) and reduce it to three theorems
(Theorems~\ref{pg:I}, \ref{pg:II}, \ref{pg:III}).

\emph{What is proved and what is not.} Proved in this section are: the deduction of
{\cite[Proposition~14.40]{P46}} from Theorem~\ref{pg:main}, including all sums over $\ell$
(Lemma~\ref{pg:sums}); the deduction of Theorem~\ref{pg:main} from Theorems~\ref{pg:I}--\ref{pg:III}
(\S\ref{pg:ss-reduction}); the part of Theorem~\ref{pg:I} that makes the argument uniform in $\ell$
(Lemmas~\ref{pg:cyclic} and~\ref{pg:burnside}: the labels at $p_0$ and $p_1$ force every lattice in
a main summand to be residually absolutely irreducible, for every $\ell$); and, inside
Theorem~\ref{pg:II}, the choice of the roots of unity (Proposition~\ref{pg:roots}), the reduction of
all twisted Selmer groups to one system of linear equations (Lemma~\ref{pg:minor}), the relation
between tame inertia and Frobenius (Lemma~\ref{pg:tame}), the count of Greenberg--Wiles type
(Lemma~\ref{pg:count}) and the bound for the constant of a Selmer datum (Lemma~\ref{pg:beta}). The
remaining statements are proved in later sections: Theorem~\ref{pg:I} in \S\ref{pga:sec} from
published theorems on Galois representations; Propositions~\ref{pg:mirror}, \ref{pg:step},
\ref{pg:stick} and statement (A) in \S\ref{pse:sec}, with the input (IM) proved in {\cite[\S3.1]{P54}};
Theorem~\ref{pg:III} (Ribet's method with auxiliary primes, Propositions~\ref{pg:ribet},
\ref{pg:triangle}, \ref{pg:odd}, \ref{pg:norm}) in \S\S\ref{pri:sec}, \ref{pod:sec}, \ref{pku:sec},
\ref{pax:sec}, \ref{pay:sec}, \ref{prm:sec}. The only use of integral $p$-adic Hodge theory,
Proposition~\ref{pg:loss}, is proved in \S\ref{pll:sec} from Theorem~\ref{pll:exponent}, which is
proved in {\cite[\S2.3]{P56}} (for $\ell\ge4b+2$, $\ell\nmid q_1dM$ see also
Remark~\ref{pll:FL} for another route). \S\ref{pg:ss-summary}
lists the inputs.

\emph{Conventions.} $\ell$ is a prime and $\iota\colon\C\to\auQlb$ an isomorphism; $\pgZb$ is the
ring of integers of $\auQlb$, and $\pgv$ the valuation of $\auQlb$ with $\pgv(\ell)=1$; for
$z\in\C$ we write $\pgv(z)$ for $\pgv(\iota z)$. A \emph{coefficient ring} is the ring of integers
$\auO$ of a finite Galois extension of $\Q_\ell$ inside $\auQlb$; ``for all large $\auO$'' means
for all $\auO$ containing a given one. Lengths of $\auO$-modules are \emph{normalised} so that
$\auO/\ell$ has length $1$; thus $\pglen\auO/a=\pgv(a)$. $\pgG=\operatorname{Gal}(\overline\Q/E)$,
$c$ is a fixed complex conjugation, $g^c=cgc$ and $g^{-c}=(g^c)^{-1}$ for $g\in\pgG$;
$\pgcyc\colon\pgG\to\Z_\ell^\times$ is the cyclotomic character and $\chi=\pgcyc^{-1}$. For a
finite place $w$ of $E$, $D_w\supseteq I_w$ are a decomposition and an inertia group at a fixed
place of $\overline\Q$ above $w$, and, if $w\nmid\ell$, $\gamma_w\in I_w$ is a \emph{tame
generator}: an element whose image generates the pro-$\ell$ quotient of tame inertia; at the
places $w$ above $p_0p_1$ (the chosen places and their conjugates) we require more, namely that the
image of $\gamma_w$ is a topological generator of the whole tame inertia group $I_w/P_w$ ($P_w$ the
wild inertia group), so that its image under the tame character of every level $d$ generates
$\auF_{q^d}^\times$, $q=\auNm w$ (this is what {\rm(G4)} and {\rm(M3)--(M5)} below use). These
elements are fixed once and for all; $\pgFrob_w$ denotes a geometric Frobenius element in $D_w$. The \emph{bad places} are
the places of $E$ dividing $\ell p_0p_1N$; their set is $\pgS=\pgS(\ell)$. The \emph{chosen places}
are the $2f$ places dividing $\pi_M$.

\subsection{The statement for one prime}\label{pg:ss-statement}

\subsubsection*{Pairs, forms, lifts}
A \emph{pair} $O=(W,V,\mu,Q)$ consists of hermitian spaces $W$, $V$ of ranks $n_W<n_V$ with $V$
positive definite and standard, an automorphic character $\mu$ of $\auU(1)$ tamely ramified at the
chosen places and at a finite set $Q$ of places as in the definition of a twist, and the weight
and level of the forms on $\auU(V)$ that the theta lift for $(\auU(W),\auU(V))$ with the character
$\mu$ produces. We write $\mathcal F(O)$ for this space of algebraic modular forms on $\auU(V)$,
with its inner product, and $t(O)=n_V-n_W$. Two families of pairs occur:
\begin{itemize}
\item for $1\le j\le m$ and a twist $\vartheta$, the pair $O_j(\vartheta)=(W_j,V_j,\mu_j^0\vartheta,
Q_\vartheta)$ of the $j$-th lift, with $n_W=r$, $t(O)=t$; the lift is the map $u$ of
{\cite[Proposition~14.37]{P46}}, restricted to $e_0\mathcal A(W_j)$, and $\mathcal Z_j(\vartheta)\subseteq
\mathcal F(O_j(\vartheta))$ is its image, the \emph{space of lifts};
\item for a main pair $(\nu,z)$ for $j$ (a critical pair with $m<z\le r$, tame part $\omega_w^{\,j}$
at each chosen place $w$, where $\omega_w$ is the character of $\auF_w^\times$ of order $q_2$ fixed in
{\cite[\S14.3]{P46}} (so that $\mu_j^0$ has tame part $\omega_w^{\,j}$), and weight $-c$
with $\kappa_j(\tau)+r-t\le c(\tau)\le\kappa_j(\tau)$) and a twist $\vartheta$, the pair
$O_{\nu,z}(\vartheta)=(W,V,\nu\vartheta,Q_\vartheta)$ with $n_W=1$, $n_V=2z$; the space of lifts
is the line $\C\,L_{\nu,z}(\vartheta)$ spanned by the lift of the constant function.
\end{itemize}
The product $\auG_O=\prod_w\auGL_{n_W}(\auF_w)$ over the chosen places acts on $\mathcal F(O)$
(twisted by the tame part of $\mu$, as in {\cite[Proposition~14.27]{P46}}); for a family $c=(c_w)$ of
labels of degree $n_W$, $e_c$ is the corresponding central idempotent. For the pairs $O_j$ we use
the idempotent $e_0$; for the pairs $O_{\nu,z}$ the idempotent $e_1$ of the labels $X-1$ (the mean
over $\prod_w\auF_w^\times$).

Let $w$ be a place of $E$ of degree one with $w\ne\aubar w$ such that neither $w$ nor $\aubar w$
divides $\ell\,M\,\auNm Q$ (a \emph{split place off the level}). The spherical Hecke algebra of
$\auGL_{n_V}(E_w)$ acts on $\mathcal F(O)$; let $T_{w,i}$, $0\le i\le n_V$, be the operator of the
double coset of $\operatorname{diag}(\varpi_w,\dots,\varpi_w,1,\dots,1)$ with $i$ entries
$\varpi_w$, and
\[
H_w(X)=\sum_{i=0}^{n_V}(-1)^i\,(\auNm w)^{i(i-1)/2}\,T_{w,i}\,X^{n_V-i}
\]
the \emph{Hecke polynomial}, a polynomial with coefficients in commuting self-adjoint-up-to-scalar
operators. For an eigenform $y$, $H_w(y;X)\in\C[X]$ is the polynomial of its eigenvalues. An
eigenform $y\in\mathcal F(O)$ \emph{has the Hecke polynomials of a lift of the system}
$q=(q_{w,i})$ (a system of eigenvalues of the operators $T_{w,i}$, $i\le n_W$, on forms on
$\auU(W)$) if for every split place $w$ off the level
\begin{equation}\label{pg:eq-liftpoly}
H_w(y;X)=\Bigl(\text{the polynomial of }q\text{ at }w,\text{ rescaled by }\mu\Bigr)\cdot
\prod_{i=0}^{t(O)-1}\bigl(X-\lambda_w(\auNm w)^{i}\bigr),
\end{equation}
where $\lambda_w=\lambda_w(O)$, the \emph{lift root}, is an explicit number depending on $\mu$ and
$w$ only, and the first factor is the Hecke polynomial of $q$ with its roots multiplied by a
number depending on $\mu$ and $w$ only. (The exact normalisations are those of the unramified
theta correspondence; we do not need them. That theta lifts of eigenforms have these Hecke
polynomials is statement (L) of \S\ref{pg:ss-galois}.)

\begin{definition}[False lifts; denominators]\label{pg:def-den}
Let $O$ be a pair, $e$ a central idempotent of $\C[\auG_O]$, and $\mathcal Z\subseteq\mathcal F(O)$
a subspace.
\begin{enumerate}
\item[(a)] For a split place $w$ off the level, the \emph{false lifts at $w$} are the elements of
\[
\mathcal F^{\mathrm{false}}_w=e\,\mathcal F(O)\ \cap\ \ker H_w(\lambda_w)\ \cap\ \mathcal Z^{\perp},
\]
and $y\in\mathcal F(O)$ is \emph{orthogonal to the false lifts} if $y\perp\mathcal F^{\mathrm{
false}}_w$ for every such $w$.
\item[(b)] Let $\beta\ge0$. The projection to $\mathcal Z$ \emph{has denominator at most $\beta$ at
$\iota$} if there is $\Lambda\in\operatorname{End}(\mathcal Z)$ with
\[
\pgv(\det\Lambda)\le\dim\mathcal Z\cdot n_W\cdot\beta
\]
such that for every $y\in\mathcal F(O)$ that is integral at $\iota$ and orthogonal to the false
lifts, the form $\Lambda\,\pgpr_{\mathcal Z}(y)$ is integral at $\iota$, and
$\langle y',\Lambda\,\pgpr_{\mathcal Z}(y)\rangle\in\pgZb$ for every $y'$ that is integral at every
isomorphism $\C\to\auQlb$.
\end{enumerate}
\end{definition}

\subsubsection*{The constants}
Put
\[
\pgNch=\prod_{w\text{ chosen}}\#\auGL_{rt}(\auF_w),\qquad h_1=\#\bigl(\auU(1)(E^+)\backslash
\auU(1)(\auA_{E^+,f})/J_1\bigr),
\]
where $J_1$ is the group of integral elements that are $\equiv1$ at the chosen places. Let
$\pgcodd$ be a positive integer such that, for every $\ell$, all large $\auO$, all $0\le e<t$ and
all $\nu\ge1$, the group of \emph{odd classes} with coefficients in $(\auO/\ell^\nu)(\chi^e)$
(Proposition~\ref{pg:odd}) has length at most $\pgv(\pgcodd)$. Let $h=r^8=\Delta/2$, and for a prime
$\ell$
\begin{align*}
\lambda_\ell&=\pgv\Bigl((2b)!\cdot\Bigl(\prod_{b-30h^3<n\le b+30h^3}n\Bigr)^{h^3}\Bigr),&
\pgloc_\ell&=4f\cdot10^6\,\frac mk\bigl(b\,\pgv(M)+\lambda_\ell\bigr),\\
e'_\ell&=(4f+3)\,r\,\pgv(\pgNch)+2^{-\ell},&
e_\ell&=t\,\pgv(\pgcodd)+\bigl(8ft+2+(4t+1)h^{22}+h\bigr)\pgv(\pgNch),
\end{align*}
\begin{equation}\label{pg:eq-beta}
\beta_\ell=\pgloc_\ell+(1+2t)\,e'_\ell+e_\ell+\pgv(\pgNch).
\end{equation}
The four terms are: the \emph{local loss} above $\ell$; the allowance for the Selmer groups and
for the ramification of the twist; an error term for the class group, the level and the
congruences between powers of the cyclotomic character; and the order of the group $\auG_O$.

\begin{theorem}[A twist with few congruences at one prime]\label{pg:main}
Let $D$ be admissible data, $\ell$ a prime and $\iota\colon\C\to\auQlb$ an isomorphism. There is an
$m$-fold twist $\tau$ such that:
\begin{enumerate}
\item[(a)] $\ell\nmid\auNm Q_\tau$;
\item[(b)] for every ramified critical pair $(\nu,z)$, every $1\le n\le m$ and both signs,
\[
\pgv\Bigl(\frac{|L(z,(\nu\tau^{\pm n})^{-1})|}{|L(z,\nu^{-1})|}\Bigr)\le e'_\ell;
\]
\item[(c)] for $1\le j\le m$, the projection to $\mathcal Z_j(\tau^j)$ has denominator at most
$\beta_\ell$ at $\iota$, for the idempotent $e_0$;
\item[(d)] for $1\le j\le m$, every main pair $(\nu,z)$ for $j$ and all $\Xi,\Xi'\in\mathcal
F(O_{\nu,z}(\tau^j))$ that are integral at every isomorphism $\C\to\auQlb$ and orthogonal to those
false lifts (for $e_1$) that lie in the kernel of the adjoint of the lift from $\auU(1)$
(\emph{test forms}),
\[
\pgv\bigl(\langle\pgpr\,\Xi,\pgpr\,\Xi'\rangle\bigr)\ge-\beta_\ell-\pgv(h_1),
\]
where $\pgpr$ is the orthogonal projection to the line $\C L_{\nu,z}(\tau^j)$.
\end{enumerate}
\end{theorem}

\begin{remark}[The exact form of (R3)]\label{pg:rem-R3}
Clauses (c), (d) are the exact content of ``$\theta_\iota^{\,j}$ has denominators at most
$\beta$ at $\iota$ for the $j$-th lift'' in {\cite[Definition~14.39(R3)]{P46}}, with
$\beta=(\beta_\ell+\pgv(h_1))\log\ell$. The \emph{false lifts} of (R3)(i) are those of
Definition~\ref{pg:def-den}(a) (forms in $e_0\mathcal F$, orthogonal to $\mathcal Z$, killed by
$H_w(\lambda_w)$ at \emph{one} split place $w$; one is orthogonal to them for every $w$); and
(R3)(ii) is (d): a bound $e^{\beta}$ for the absolute value at $\iota$ of
$\langle\pgpr\Xi,\pgpr\Xi'\rangle$ for test forms $\Xi,\Xi'$. The condition in (R3) on the sum is
$\sum_{\ell\in T}\beta(\ell)\le\aubud$ with $\beta(\ell)$ already containing the factor
$\log\ell$. This is the form in which {\cite[Propositions~14.41(b), 14.42 and~14.43]{P46}} quote (R3).
\end{remark}

\begin{remark}[How the primes are treated]\label{pg:rem-primes}
The statement and the argument are the same for every $\ell$: for $\ell=2$, for $\ell\mid q_1d$
(ramified in $E$), for $\ell$ inert or split in $K$, and for $\ell\in\{p_0,p_1,N\}$. No hypothesis
``$\ell$ odd, unramified, ordinary'' is made, and none of the published non-vanishing theorems
modulo $\ell$ for anticyclotomic twists is used. The dependence on $\ell$ is entirely in the
constants: $\pgv(\pgNch)$, $\pgv(\pgcodd)$ and $\pgv(h_1)$ vanish unless $\ell$ divides one of three
fixed integers; $2^{-\ell}$ is summable; and $\pgloc_\ell$ vanishes for $\ell>2b$,
$\ell\nmid M$, while for small $\ell$ it is of size $f(m/k)\,b/(\ell-1)$. The three devices that
make this possible are: (i) the twist has $\ell$-power order and conductor prime to $\ell$, so that
twisting is a deformation of the trivial twist and only upper bounds for valuations at one nearby
point are needed (\S\ref{pg:ss-twisting}); (ii) the labels at the two level primes make the main
summands residually irreducible for all $\ell\ne p_i$, and for $\ell=p_0$ one uses $p_1$
(Lemmas~\ref{pg:cyclic}, \ref{pg:burnside}); (iii) above $\ell$ no ordinarity is used; the failure
of torsion crystalline classes to lift is bounded by $\pgloc_\ell$ (Proposition~\ref{pg:loss}), and
this is what the budget pays for.
\end{remark}

\subsection{The sums over the primes}\label{pg:ss-sums}

\begin{lemma}\label{pg:sums}
Assume that $\pgcodd$ can be taken with $\log\pgcodd\le x^{800}$ (Proposition~\ref{pg:odd}). Then
for every finite set $T$ of primes
\[
\sum_{\ell\in T}e'_\ell\log\ell\le x_1,\qquad
\sum_{\ell\in T}\bigl(\beta_\ell+\pgv(h_1)\bigr)\log\ell\le\aubud=\frac{f\,a\,(\log b)^3}{m\,k}.
\]
\end{lemma}
\begin{proof}
For a positive integer $n$, $\sum_{\ell\in T}\pgv(n)\log\ell\le\log n$; and
$\sum_\ell2^{-\ell}\log\ell\le\sum_{n\ge1}n2^{-n}=2$. We bound the logarithms of the integers that
occur, using {\cite[(36) and~(38)]{P46}} and $\log b\le x^{1/100}$.

\emph{$\pgNch$.} $\#\auGL_n(\auF_p)<p^{n^2}$, there are $2f$ chosen places, $p_i\le b^{11}$ and
$rt=5r^2\le45x^{10}$; so $\log\pgNch\le2f(rt)^2\cdot11\log b\le x^{34}$.

\emph{$h_1$.} Since $E/E^+$ is unramified at all finite places, every element of
$\auU(1)(\auA_{E^+,f})$ is of the form $y/\aubar y$ with $y\in\auA_{E,f}^\times$, and the integral
elements are those with $y$ a unit idele (Hilbert's Theorem~90 at each place, and the vanishing of
$H^1$ of the units in an unramified extension). Hence the class group of $\auU(1)$ at the maximal
level is a quotient of the class group of $E$, and $h_1\le h_E\cdot\prod_{w\text{ chosen}}
(\auNm w-1)$. By Minkowski's bound every ideal class contains an integral ideal of norm at most
$|d_E|^{1/2}$, and the number of integral ideals of norm at most $X$ is at most
$X(1+\log X)^{2f-1}$; with $\log|d_E|=2f\log\aurd\le24x^{12}\log x$ this gives $\log h_E\le
x^{14}$, and $\log h_1\le x^{14}+22f\log b\le x^{15}$.

\emph{The local loss.} $\sum_{\ell\in T}\lambda_\ell\log\ell\le\log(2b)!+h^3\cdot60h^3\log(b+30h^3)
\le2b\log(2b)+60h^6\log(2b)\le3b\log b$, and $\sum_\ell b\,\pgv(M)\log\ell\le b\log M\le22\,b\log
b$. So $\sum_{\ell\in T}\pgloc_\ell\log\ell\le10^8f(m/k)\,b\log b$.

\emph{Conclusion.} $\sum_T e'_\ell\log\ell\le(4f+3)r\,x^{34}+2\le x^{52}\le x_1$. Next,
$\sum_Te_\ell\log\ell\le t\,x^{800}+(8ft+2+(4t+1)h^{22}+h)\,x^{34}\le x^{1000}$, since
$h^{22}=r^{176}\le(3x^5)^{176}$. Therefore
\[
\sum_{\ell\in T}\bigl(\beta_\ell+\pgv(h_1)\bigr)\log\ell\le10^8f\frac mk\,b\log b+(1+2t)x^{52}
+x^{1000}+x^{34}+x^{15}\le2\cdot10^8f\frac mk\,b\log b,
\]
because $b\ge x_1^{10^6}\ge x^{10^{12}}$. Finally $2\cdot10^8\,m^2\log b\le8\cdot10^8x\log b\le
x(\log b)^3$, as $m^2\le4x$ and $(\log b)^2\ge8\cdot10^8$; that is,
$2\cdot10^8f(m/k)b\log b\le f\,xb(\log b)^3/(mk)=\aubud$.
\end{proof}

\begin{proof}[Proof of {\cite[Proposition~14.40]{P46}} from Theorem~\ref{pg:main} and
Proposition~\ref{pg:odd}]
For each $\iota\in\auI$, with prime $\ell$, let $\theta_\iota$ be the twist $\tau$ of
Theorem~\ref{pg:main}. (R1) is (a). For (R2), let $(\nu,z)$ be a ramified critical pair,
$1\le n\le m$ and a sign be given, and let $y_\iota=|L(z,\nu^{-1})|/|L(z,(\nu\theta_\iota^{\pm
n})^{-1})|$. By (b), $|y_\iota|_\iota\le\ell^{e'_\ell}$; so $y$ has height at most $x_1$ with the
multiplier $H=1$, $\alpha=0$ and the exponents $h_\ell=e'_\ell$, by the first inequality of
Lemma~\ref{pg:sums}. For (R3), put $\beta(\ell)=(\beta_\ell+\pgv(h_1))\log\ell$; the second
inequality of Lemma~\ref{pg:sums} is the condition on the sum, (c) gives the operator $\Lambda$
with $|\det\Lambda|_\iota\ge\exp(-\dim\mathcal Z\cdot r\,\beta_\ell\log\ell)\ge\exp(-\dim\mathcal
Z\cdot r\,\beta(\ell))$, and (d) is the bound for the main pairs (Remark~\ref{pg:rem-R3}).
\end{proof}

\subsection{Reduction to three theorems}\label{pg:ss-reduction}

\subsubsection*{Characters}
For a twist $\vartheta$ let $\psi_\vartheta\colon\pgG\to\pgZb^\times$ be the character of finite
order that corresponds by class field theory, through $\iota$, to the Hecke character
$y\mapsto\vartheta^{-1}(y/\aubar y)$ of $E$. It is \emph{anticyclotomic}:
$\psi_\vartheta(g^c)=\psi_\vartheta(g)^{-1}$; it is unramified outside $Q_\vartheta\cup\aubar{
Q_\vartheta}$, and its restriction to $I_w$, $w\in Q_\vartheta$, is given by the tame part of
$\vartheta$. Conversely, every continuous character $\psi\colon\pgG\to\pgZb^\times$ of finite order
with $\psi(\mathrm{Art}(y))=\eta(y/\aubar y)$ for a character $\eta$ of $\auU(1)(E^+)\backslash
\auU(1)(\auA_{E^+,f})$ that is tamely ramified on a set $Q$ with $Q\cap\aubar Q=\emptyset$ of
places prime to $MN$, and unramified elsewhere, is $\psi_\vartheta$ for the twist
$\vartheta=\iota^{-1}\circ\eta^{-1}$ (trivial at the infinite places); if moreover $\psi(\gamma_w)$
has order greater than $m$ for all $w\in Q$, then $\vartheta$ is $m$-fold.

\begin{definition}[Main summands]\label{pg:def-summand}
Let $n\ge1$ and $t'\ge1$. A continuous representation $R\colon\pgG\to\auGL_n(\auQlb)$ is a
\emph{main summand of rank $n$ and string length $t'$} if:
\begin{enumerate}
\item[(M1)] $R$ is semisimple and $\pgtr R(g^c)=\chi(g)^{t'-1}\,\pgtr R(g^{-1})$ for all $g$;
\item[(M2)] $R$ is unramified outside $\ell p_0p_1$; at $v\mid\ell$ it is crystalline if
$v\nmid p_0p_1$ and potentially crystalline otherwise; its Hodge--Tate weights $J_\sigma$, for
the embeddings $\sigma\colon E\to\C$ (transported by $\iota$), satisfy:
$1\le n\le r$, $t'\le t$; for one $j\le m$ every $u\in J_\sigma$ is \emph{low},
$\Delta/4\le|u|\le2m\Delta$, or \emph{high}, $|u\pm\kappa_j(\sigma)|\le t$ with the sign $+$ for
$\sigma$ in the CM type of the data and $-$ otherwise; at most $nm/k$ elements of $J_\sigma$ have
$|u|\le2h^2$; and $J_\sigma$ and $J_{\aubar\sigma}$ together have at most $n$ negative elements;
\item[(M3)] at every $v\mid p_0p_1$, $v\nmid\ell$: there is $N_v$ prime to $v$ with $R(s^{N_v})=1$
for $s\in I_v$, and an element $\gamma\in I_v$, equal to $\gamma_v$ if $v$ is chosen, with
$\pgv(\det(R(\gamma)-1))\le n\,\pgv(\pgNch)$ and $\pgv(\det(R(\gamma^{-c})-1))\le n\,\pgv(\pgNch)$;
\item[(M4)] for all large $\auO$ and every $\pgG$-stable $\auO$-lattice, with matrices $u(g)$: the
$u(g)$ span $M_n(\auO)$ over $\auO$, and there are $g\in\pgG$ and $i\in\{0,1\}$ with
$\pgcyc(g)=p_i^{\,n}$ and $u(g)$ a scalar;
\item[(M5)] there are a place $v\mid p_0p_1$, $v\nmid\ell$, and $s\in I_v$ such that $R(s)$ is a
scalar $y$ with $y^{\ell^aj}-1$ a unit for all $a\ge0$ and $1\le j\le4n+2$.
\end{enumerate}
\end{definition}

The inequality on negative weights in (M2) says that the Bloch--Kato local conditions are
\emph{balanced}: $\sum_{v\mid\ell}\dim H^1_f(E_v,R)-\dim H^0(E_v,R)\le f\,n$, half of
$[E:\Q]\,n$. Condition (M4) says that every lattice in $R$ is residually absolutely irreducible
(Burnside), so that all lattices are homothetic.

\begin{definition}[Selmer groups]\label{pg:def-sel}
Let $u\colon\pgG\to\auGL_n(\auO)$ be continuous, $Q$ a finite set of places and $\gamma\in\pgG$.
For $v\mid\ell$ let $Z_v(u)$ be the $\auO$-module of \emph{crystalline cocycles}: the continuous
cocycles $D_v\to\auO^n$ whose class in $H^1(E_v,\auO^n\otimes\Q)$ lies in the subspace $H^1_f$ of
Bloch and Kato~\cite[(3.7.2)]{BlochKato}. A function $z\colon\pgG\to B^n$, $B$ an $\auO$-module,
\emph{obeys the crystalline relations} if for every $v\mid\ell$ its restriction to $D_v$ satisfies
every $\auO$-linear relation $\sum_ia_i\,z(g_i)=0$ ($a_i\in M_n(\auO)$, $g_i\in D_v$) that holds
for all $z\in Z_v(u)$. The \emph{Selmer group} $\pgSel_Q(u,\gamma;B)$ is the module of locally
constant cocycles $z\colon\pgG\to B^n$ for $u$ that are unramified outside $\ell p_0p_1$,
$Q\cup\aubar Q$ (that is, $z(I_w)=0$ elsewhere), obey the crystalline relations and vanish at
$\gamma$. We use $B=\auO/a$ and $B=K/\auO$, $K=\operatorname{Frac}\auO$.
\end{definition}

If $u(\gamma)$ is a scalar $\zeta$ with $\zeta-1$ a unit, every class has exactly one cocycle
vanishing at $\gamma$ and $H^0=0$, so $\pgSel_Q(u,\gamma;B)$ is a group of cohomology classes. The
relations define the local condition for every coefficient module at once; for $B=\auO/\ell^M$
the cocycles obeying them contain the reductions of the elements of $Z_v(u)$, and the two coincide
when $Z_v(u)\otimes\auO/\ell^M$ is all of the torsion condition, for instance in the range of
Fontaine and Laffaille.

\begin{theorem}[Main summands exist]\label{pg:I}
Let $\ell$, $\iota$ be given.
\begin{enumerate}
\item[(a)] For $1\le j\le m$ and every system $q$ of Hecke eigenvalues that occurs in
$e_0\mathcal A(W_m)$ there is a main summand $R_{j,q}$ of rank $r$ and string length $t$ such
that, for every twist $\vartheta$, every eigenform $y\in\mathcal F(O_j(\vartheta))$ with the
Hecke polynomials of a lift of $q$ and every Galois representation $S_y$ attached to $y$
(\S\ref{pg:ss-galois}, (G1)),
\begin{equation}\label{pg:eq-trace}
\pgtr\bigl(S_y\otimes\eta_O\bigr)(g)=\psi_\vartheta(g)\,\pgtr R_{j,q}(g)+\sum_{e=0}^{t-1}\chi(g)^e
\qquad(g\in\pgG),
\end{equation}
where $\eta_O$ is the $\ell$-adic character of an explicit algebraic Hecke character depending on
the pair only (the \emph{normalising character}).
\item[(b)] For every $j$ and every main pair $(\nu,z)$ for $j$ there is a main summand
$R_{\nu,z}$ of rank $1$ and string length $2z-1$ with the same property for the pairs
$O_{\nu,z}(\vartheta)$ and the eigenforms with the Hecke polynomials of a lift of the trivial
system.
\end{enumerate}
\end{theorem}

\begin{theorem}[The twisting theorem]\label{pg:II}
Let $\ell$, $\iota$ be given and let $R_1,\dots,R_s$ be main summands, of ranks $n_i\le r$. There
is an $m$-fold twist $\tau$ such that for $1\le j\le m$, with $\vartheta=\tau^j$,
$\psi=\psi_\vartheta$ and $Q=Q_\tau$:
\begin{enumerate}
\item[(a)] the places of $Q$ have prime norm, and $Q$, $\aubar Q$ and $\pgS(\ell)$ are pairwise
disjoint; $\psi$ is trivial on $D_v$ for $v\mid\ell$; and $\sum_{w\in Q}\pgv(\psi(\gamma_w)-1)\le
e'_\ell$;
\item[(b)] for each $i$ there are $\gamma\in I_v$ for some $v\nmid\ell$ and $\zeta\in\pgZb$ with
$\zeta-1$ a unit such that, for all large $\auO$, there is a lattice $u$ in $R_i$ with
$\psi(\gamma)u(\gamma)=\zeta$ and $\pglen\pgSel_Q(\psi u,\gamma;K/\auO)\le e'_\ell$;
\item[(c)] $\pgv\bigl(|L(z,\nu^{-1})|/|L(z,(\nu\vartheta^{\pm1})^{-1})|\bigr)\ge-e'_\ell$ for every
ramified critical pair $(\nu,z)$ and both signs.
\end{enumerate}
\end{theorem}

\begin{theorem}[From Selmer groups to denominators]\label{pg:III}
Let $\ell$, $\iota$ be given, let $\vartheta$ be a twist satisfying
Theorem~\ref{pg:II}(a), $\psi=\psi_\vartheta$, and let $O$ be one of the pairs
$O_j(\vartheta)$, $O_{\nu,z}(\vartheta)$, with its idempotent $e$. Let $\mathcal Z\subseteq
\mathcal F(O)$ be spanned by eigenforms $y$ with the following properties: $y$ lies in
$e_c\mathcal F(O)\subseteq e\,\mathcal F(O)$ for a family $c$ of labels in which, above each of
$p_0$, $p_1$, one label is primitive (if $n_W=r$) or all are $X-1$ (if $n_W=1$), whose rescaled labels do not have the root $1$, and with $n_W$ and $t(O)$ odd (for the data, $r$ is a power of $3$, {\cite[Lemma~13.4(S7)]{P46}}); $y$ has the Hecke polynomials of a lift of a system
$q$; and there is a main summand $R$ for which \eqref{pg:eq-trace} holds for all twists, and for
which the conclusion of Theorem~\ref{pg:II}(b) holds for $\psi$. Then the projection to
$\mathcal Z$ has denominator at most $\beta_\ell$ at $\iota$ (Definition~\ref{pg:def-den}).
\end{theorem}

\begin{lemma}[The line of a main pair]\label{pg:line}
Let $O=O_{\nu,z}(\vartheta)$, $L=L_{\nu,z}(\vartheta)$, and suppose that the projection to $\C L$
has denominator at most $\beta$ at $\iota$, for $e_1$. Assume statement \emph{(C)} below. Then
$\pgv(\langle\pgpr\Xi,\pgpr\Xi'\rangle)\ge-\beta-\pgv(h_1)$ for all $\Xi,\Xi'$ as in
Theorem~\ref{pg:main}(d).
\end{lemma}
\smallskip\noindent\textbf{(C)} For $s\in\auU(1)(\auA_{E^+,f})$ let $\varrho(s)$ be the translation of forms on $\auU(V)$ by the
central element $s$, multiplied by the value at $s$ of the character that makes $L$ invariant.
Then $\varrho(s)L=L$, $\langle\varrho(s)y,y'\rangle=\langle y,\varrho(s)^{-1}y'\rangle$; if $s$ is
integral above $\ell$ then $\varrho(s)^{\pm1}$ preserves the forms integral at each isomorphism
$\C\to\auQlb$; every class of $\auU(1)(E^+)\backslash\auU(1)(\auA_{E^+,f})/J_1$ has a
representative integral above $\ell$; and if $\Xi'$ is orthogonal to the false lifts in the kernel
of the adjoint of the lift, and $s_1,\dots,s_{h_1}$ represent the classes, then
$\sum_c\varrho(s_c)^{-1}\Xi'$ is orthogonal to the false lifts in $(\C L)^\perp$. The level $J_1$ is
free by {\cite[Lemma~14.5]{P46}}; this is used in the last assertion. This is proved in Proposition~\ref{pga:C}.\par\smallskip
\begin{proof}
We may assume $L\ne0$. Let $s_1,\dots,s_{h_1}$ be representatives integral above $\ell$ and
$\Xi''=\sum_c\varrho(s_c)^{-1}\Xi'$. By (C), $\Xi''$ is integral at every isomorphism and
orthogonal to the false lifts, and $\langle L,\Xi''\rangle=\sum_c\langle\varrho(s_c)L,\Xi'\rangle
=h_1\langle L,\Xi'\rangle$, so $\pgpr\Xi''=h_1\pgpr\Xi'$. On the line, $\Lambda$ is a scalar
$\mu$ with $\pgv(\mu)\le\beta$ (here $\dim\mathcal Z=n_W=1$). By Definition~\ref{pg:def-den}(b)
applied to $y=\Xi''$ and $y'=\Xi$, $\mu\,h_1\langle\Xi,\pgpr\Xi'\rangle\in\pgZb$. As $\pgpr$ is
self-adjoint and idempotent, $\langle\pgpr\Xi,\pgpr\Xi'\rangle=\langle\Xi,\pgpr\Xi'\rangle$, and
the claim follows.
\end{proof}

\begin{proof}[Proof of Theorem~\ref{pg:main} from Theorems~\ref{pg:I}, \ref{pg:II}, \ref{pg:III},
{\cite[Proposition~14.36]{P46}} and statements \emph{(L)}, \emph{(C)}]
The set $\mathcal Q$ of systems of Hecke eigenvalues occurring in $e_0\mathcal A(W_m)$ is finite,
and so is the set of main pairs (the weights and the tame parts are bounded and the class group
of $\auU(1)$ at a fixed level is finite). Apply Theorem~\ref{pg:II} to the finitely many main
summands $R_{j,q}$ ($1\le j\le m$, $q\in\mathcal Q$) and $R_{\nu,z}$ of Theorem~\ref{pg:I}. Let
$\tau$ be the resulting twist. (a) holds because $Q_\tau$ is disjoint from the places above
$\ell$. (b) is Theorem~\ref{pg:II}(c) for the powers $\tau^n$.

(c) Fix $j$ and put $\vartheta=\tau^j$. We check that $\mathcal Z=\mathcal Z_j(\vartheta)$
satisfies the hypothesis of Theorem~\ref{pg:III}. The space $e_0\mathcal A(W_j)$ is spanned by
Hecke eigenforms $y_0$ lying in $e_c\mathcal A(W_j)$ for good families $c$ of labels
(definition of $e_0$ in {\cite[\S14.3]{P46}}; the Hecke operators at the split places off the
level are normal and commute with $\auG$); put $y=u(y_0)$. A good family has a primitive label above each of $p_0$ and $p_1$ (its root generates
$\auF_{p_i^r}^\times$), and its rescaled labels do not have the root $1$: the rescaling multiplies
the roots of $c_w$ by the element $z_w\in\auF_w^\times$ that names the inverse of the tame part of
$\mu_j^0\vartheta$ at $w$, which is $\omega_w^{-j}$ and has order dividing $q_2$, the order of
$\omega_w$; a rescaled root equal to $1$ would be a root $z_w^{-1}\in\auF_w$ of $c_w$ that is a
$q_2$-th root of unity, which a good family excludes. By (L), $y$ is in $e_c\mathcal F(O)$ and has
the Hecke polynomials of a lift of the system $q$ of $y_0$; by {\cite[Proposition~14.36]{P46}} (equality
of multiplicities) the system $q$ occurs in $e_0\mathcal A(W_m)$, so
$q\in\mathcal Q$. The main summand $R_{j,q}$ satisfies \eqref{pg:eq-trace} by
Theorem~\ref{pg:I}(a) and the conclusion of Theorem~\ref{pg:II}(b) by the choice of $\tau$. Hence
Theorem~\ref{pg:III} applies and gives (c).

(d) In the same way Theorem~\ref{pg:III} applies to the line $\C L_{\nu,z}(\vartheta)$ of a main
pair for $j$, with the labels $X-1$ and the summand $R_{\nu,z}$: the lift of the constant
function is fixed by $\prod_w\auF_w^\times$ up to the tame part of the character, has the Hecke
polynomials of a lift of the trivial system by (L), and the rescaled label is $X-\xi$ with $\xi$
the element naming $\omega_w^{-j}$ at $w$, which is not $1$ because $\omega_w^{\,j}\ne1$ for $1\le
j\le m$ ($\omega_w$ has order $q_2>m$). Lemma~\ref{pg:line} gives (d).
\end{proof}

\subsection{Galois representations: what is used, and Theorem~\ref{pg:I}}\label{pg:ss-galois}

Let $V$ be a standard positive definite hermitian space of rank $n$ over $E$ and $y$ an algebraic
modular form on $\auU(V)$ of weight $\lambda$ that is an eigenform of the operators $T_{w,i}$ at
almost all split places. The following properties of Galois representations are used in
Theorems~\ref{pg:I} and~\ref{pg:III}, and no others.
\begin{enumerate}
\item[(G1)] \emph{Existence.} There is a continuous semisimple $S_y\colon\pgG\to\auGL_n(\auQlb)$,
unramified at almost all places, such that for almost all split places $w$ at which $y$ is
spherical, $\det(X-S_y(\pgFrob_w))=\iota H_w(y;X)$.
\item[(G2)] \emph{Good places.} If $y$ is fixed by a hyperspecial subgroup at $v$ (for instance by
$\auGL_n(\auO_{E,w})$ at a split place $w$) then $S_y$ is unramified at the places of $E$ above $v$
not dividing $\ell$, with the Frobenius relation of (G1) at split $w$, and crystalline at those
dividing $\ell$.
\item[(G3)] \emph{Weights.} At every place above $\ell$, $S_y$ is de Rham with the Hodge--Tate
weights determined by $\lambda$: at $\sigma\colon E\to\C$ in the CM type the numbers
$i-1-\lambda_{\sigma,i}$ and at $\aubar\sigma$ the numbers $\lambda_{\sigma,i}+n-i$, $1\le i\le n$.
\item[(G4)] \emph{Tame level.} Let $w$ be a split place of residue characteristic $p$ at which $y$
is fixed by the principal congruence subgroup and lies in the isotypic part, for
$\auGL_n(\auF_w)$, of the character named by a label $c$. Then $S_y|_{I_w}$ is tame, of finite
order if $w\nmid\ell$ and potentially crystalline with the corresponding inertial type if
$w\mid\ell$, and its eigenvalues on a generator of tame inertia are the images of the roots of
$c$ under a fixed injective character $\overline{\auF}_p^\times\to\auQlb^\times$ (with the sign
convention $\epsilon=1$ of Proposition~\ref{pga:G45}).
\item[(G5)] \emph{Parahoric level with a character.} If at a split place $w\in Q$, $w\nmid\ell$,
the form $y$ is an eigenvector of a parahoric subgroup of type $(n_W,t)$ with a tame character
$\psi_w$ on the first block and the trivial one on the second, then $(S_y(s)-\psi_w(s))(S_y(s)-1)
=0$ for $s\in I_w$ (after the normalisation; in particular there is no monodromy; with the sign
convention $\epsilon=1$ of Proposition~\ref{pga:G45}).
\end{enumerate}

\emph{Published sources.} By Labesse's base change~\cite{Labesse} (which applies because $E/E^+$
is unramified at all finite places, $[E^+:\Q]\ge2$ and the group is definite), $y$ generates
automorphic representations $\pi$ of $\auU(V)$ with a base change $\Pi$ to $\auGL_n(\auA_E)$, an
isobaric sum of conjugate self-dual discrete representations, compatible with the local base
change at all split places and at all places where $\pi$ is unramified. The Galois
representations attached to the cuspidal constituents of $\Pi$ are those of Harris--Taylor, Shin,
Chenevier--Harris, with local--global compatibility at all finite places, including those above
$\ell$, by Caraiani; all this is collected, for every prime $\ell$ and every regular algebraic
conjugate self-dual cuspidal representation, in~\cite[Theorem~2.1.1]{BLGGT} (above $\ell$ that
theorem covers only components with an Iwahori-fixed vector; above $\ell$ in general we use
\cite[Theorem~A]{BLGGTlp2}, which gives the semisimplified identity and a bound on the monodromy,
and which already yields potential crystallinity with the expected inertial type when the
local parameter has no monodromy, see Proposition~\ref{pga:G45}; see also
\cite{BLGGTlp1,CaraianiANT}). Thus (G1)--(G3)
follow for every $\ell$, including $\ell=2$ and $\ell$ ramified in $E$; (G4) and (G5) follow from
the compatibility at $w$ together with the description of the local Langlands correspondence for
representations of depth zero (a representation containing the inflation of the cuspidal
representation of $\auGL_n(\auF_w)$ of a primitive label corresponds to the representation induced
from the character of the inertia group of the unramified extension of degree $n$ given by a root
of the label) and for representations with Iwahori or parahoric fixed vectors with a tame
character (principal series and their subquotients). \emph{We have not found one published
statement that contains {\rm(G1)--(G5)} for algebraic modular forms on definite unitary groups in
this generality}; each item is a consequence of the cited theorems, but the deduction (passage
from $\Pi$ to its cuspidal constituents, including possible Speh constituents; places of $E^+$
inert in $E$; the normalisations in (G3)) is written in \S\ref{pga:sec}, Propositions~\ref{pga:G123} and~\ref{pga:G45}, with
the twist of Proposition~\ref{pax:sign} for Speh constituents.



\smallskip
The following statement about theta lifts is used; it is proved in Proposition~\ref{pga:L}.
\smallskip\noindent\textbf{(L)} If $y_0$ is a Hecke
eigenform on $\auU(W)$ with system $q$ and its theta lift $y$ to $\auU(V)$ for the character $\mu$
is not zero, then $y$ is an eigenform and has the Hecke polynomials \eqref{pg:eq-liftpoly} of a
lift of $q$ at every split place off the level; and the lift is equivariant for the groups
$\auG_O$ as in {\cite[Proposition~14.27]{P46}}, so that $y\in e_c\mathcal F(O)$ if $y_0$ is of type
$e_c$. This is proved in Proposition~\ref{pga:L}, for the Weil representation of {\cite[\S2.1]{P54}}.\par\smallskip

\subsubsection*{The image of inertia at the level primes}

\begin{lemma}\label{pg:cyclic}
Let $p\equiv1\pmod3$ be a prime, $r$ a power of $3$, $\xi\in\overline{\auF}_p^\times$ of order
$p^r-1$, $z\in\auF_p^\times$ and $\xi'=\xi z$. For every prime $\ell\ne p$ and every $u\ge0$ the
elements $(\xi'^{\,p^a})^{\ell^u}$, $0\le a<r$, are pairwise distinct.
\end{lemma}
\begin{proof}
If two of them coincide, for $a<b$, then $\xi^{\ell^up^a(p^{b-a}-1)}=1$, because
$z^{p^a}=z^{p^b}$. So the order $o$ of $\xi^{\ell^u}$, which is $(p^r-1)/\gcd(\ell^u,p^r-1)$,
divides $p^d-1$ with $0<d=b-a<r$, hence divides $p^{\gcd(d,r)}-1$, hence $x-1$ with $x=p^{r/3}$.
(Here we used that $o$ is prime to $p$ and divides $p^r-1$, and that a proper divisor of $r$ divides
$r/3$.) From $o\mid x-1$ and $p^r-1=x^3-1$ we get $x^3-1\mid(x-1)\ell^u$, that is,
$x^2+x+1\mid\ell^u$. But $x\equiv1\pmod3$, so $3\mid x^2+x+1$ and $x^2+x+1\equiv3\pmod9$; a divisor
of $\ell^u$ that is divisible by $3$ and not by $9$ is $3$, and $x^2+x+1>3$. Contradiction.
\end{proof}

\begin{lemma}\label{pg:burnside}
Let $R\colon\pgG\to\auGL_n(\auQlb)$ be continuous, $w\nmid\ell$ a place of $E$ of degree one and
norm $p$, and $s\in I_w$ such that $R|_{I_w}$ is tame of finite order and the eigenvalues of $R(s)$
are $\zeta^{p^a}$, $0\le a<n$, for a root of unity $\zeta$ with the $\zeta^{p^a}$ pairwise
incongruent modulo the maximal ideal of $\pgZb$. Then for every coefficient ring $\auO$ over which
$R$ and $\zeta$ are defined and every stable lattice, with matrices $u(g)$: the $u(g)$ span
$M_n(\auO)$, and $u(\phi^n)$ is a scalar for every arithmetic Frobenius element $\phi\in D_w$;
moreover $\pgcyc(\phi^n)=p^n$.
\end{lemma}
\begin{proof}
Let $k$ be the residue field of $\auO$ and $\bar u$ the reduction. Since the eigenvalues of
$u(s)$ are distinct modulo the maximal ideal, the lattice is the direct sum of the eigenlines
$\Lambda_a$ of $u(s)$, and $k^n$ is the direct sum of the eigenlines $\bar\Lambda_a$ of
$\bar u(s)$ for the distinct eigenvalues $\bar\zeta^{p^a}$. In the tame quotient $\phi
s\phi^{-1}=s^{p}$, and $R$ is trivial on wild inertia; so $u(\phi)$ maps the eigenline of $u(s)$
for an eigenvalue $\alpha$ to the eigenline of $u(s^{p})$ for $\alpha$, which is the eigenline of
$u(s)$ for the eigenvalue $\alpha'$ with $\alpha'^{\,p}=\alpha$. (the $p$-th power map is injective on the eigenvalues, which have order prime to $p$). So
$u(\phi)$ permutes the eigenlines, and the set of eigenvalues is stable under $\alpha\mapsto
\alpha'$; for $\alpha=\zeta$ this gives $\zeta=\zeta^{p^{b+1}}$ with $b<n$, hence $b=n-1$ since
the $\zeta^{p^a}$ are distinct. Thus $u(\phi)$ permutes the $\Lambda_a$ cyclically. A $\bar u$-stable
subspace of $k'^n$, $k'\supseteq k$, is a sum of lines $\bar\Lambda_a\otimes k'$ and is stable
under $\bar u(\phi)$, so it is $0$ or everything: $\bar u$ is absolutely irreducible. By Burnside's
theorem the $\bar u(g)$ span $M_n(k)$, and by Nakayama's lemma the $u(g)$ span $M_n(\auO)$.
Finally $u(\phi^n)$ preserves each $\Lambda_a$ and commutes with $u(\phi)$, which permutes them
transitively; so its eigenvalues on the $\Lambda_a$ are equal, and it is a scalar.
\end{proof}

\begin{proof}[Proof of Theorem~\ref{pg:I}, given {\rm(G1)--(G4)} (Propositions~\ref{pga:G123}, \ref{pga:G45}) and {\rm(L)} (Proposition~\ref{pga:L})]
(a) Let $q\in\mathcal Q$, let $y_0\in e_c\mathcal A(W_m)$ be an eigenform with system $q$, for a
good family $c$ (recall that $\auU(W_m)$ is definite), and let $S_q$ be attached to $y_0$ by (G1).
Put $R_{j,q}=S_q\otimes\eta_j$, where $\eta_j$ is the $\ell$-adic character of the algebraic Hecke
character $y\mapsto(\mu_j^0)^{-1}(y/\aubar y)\,|y|^{(r-t)/2}$ (note that $r\equiv t\pmod2$).

\emph{The trace identity.} Let $\vartheta$, $y$, $S_y$ be as in the statement. By (G1) and
\eqref{pg:eq-liftpoly}, for almost all split places $w$ the characteristic polynomial of
$S_y(\pgFrob_w)$ is the product of that of $S_q(\pgFrob_w)$, with roots rescaled by a number
depending on $\mu_j^0\vartheta$ and $w$, and of $\prod_{i<t}(X-\lambda_w(\auNm w)^i)$. The
rescaling numbers and $\lambda_w$ are values at $\pgFrob_w$ of $\ell$-adic characters of algebraic
Hecke characters built from $\mu_j^0\vartheta$; after multiplication by $\eta_O(\pgFrob_w)$ the
second factor becomes $\prod_{i<t}(X-(\auNm w)^i)=\prod_{i<t}(X-\chi(\pgFrob_w)^i)$ and the first
the characteristic polynomial of $(\psi_\vartheta\otimes R_{j,q})(\pgFrob_w)$; this computation is
carried out in Proposition~\ref{pga:trace} (with $\eta_O$ the character of $(\chi_W(y/\aubar y)\upsilon(y)^{-r}|y|^{-r/2})^{-1}$). The split
places of degree one have density one, so by the Chebotarev density theorem and continuity
\eqref{pg:eq-trace} holds for all $g$.

\emph{{\rm(M1)}.} $S_q$ is semisimple, and $S_q^c\cong S_q^\vee\otimes\chi^{r-1}$
by~\cite[Theorem~2.1.1]{BLGGT} applied to the constituents of the base change (or directly: the
Hecke operators at $w$ and $\aubar w$ are dual to each other, which gives the identity of traces
at Frobenius elements, and then Chebotarev). Twisting by $\eta_j$, which satisfies
$\eta_j^c=\eta_j^{-1}\chi^{t-r}$ (its $\mu$-part is anticyclotomic and $|y|^{(r-t)/2}$
corresponds to $\chi^{(t-r)/2}$), gives $\pgtr R(g^c)=\chi(g)^{t-1}\pgtr R(g^{-1})$.

\emph{{\rm(M2)}.} The forms in $\mathcal A(W_m)$ have level $J_M$: hyperspecial outside $M$ and
principal at the chosen places. So $R_{j,q}$ is unramified outside $\ell p_0p_1$ and crystalline
above $\ell$ if $\ell\nmid M$ by (G2), and potentially crystalline above $\ell=p_i$ by (G4). The
weights: by (G3) and the definition of the parameter $\varsigma_\tau$ of the target forms, the
Hodge--Tate weights of $R_{j,q}$ at $\sigma$ are, up to a shift bounded by $t$ in absolute value
and with $\epsilon=\pm1$ according as $\sigma$ is in the CM type or not,
\[
-\epsilon\bigl(\kappa_j(\sigma)+O(k)\bigr)\quad(k-1\text{ weights}),\qquad
-\epsilon\,(j-i+\tfrac12)\Delta\quad(1\le i\le m).
\]
The first are high. The second are low: $\Delta/4\le|(j-i+\frac12)\Delta|-t$ and $|(j-i+\frac12)\Delta|+t\le2m\Delta$ (as $t=5r\le\Delta/4$). Exactly
$m$ weights have absolute value at most $2h^2$, and $m\le rm/k$. For $\sigma$ in the CM type the
negative weights are the $k-1$ high ones and the $j$ low ones with $i\le j$; for $\aubar\sigma$
they are the $m-j$ low ones with $i>j$; together $k-1+m=r$. The exact weights are computed in Proposition~\ref{pga:weights}: they are
$\frac{t-1}2-\epsilon_\sigma(\varsigma_{\sigma,i}+\kappa_j(\sigma))$.

\emph{{\rm(M3)}, {\rm(M4)}.} Let $i\in\{0,1\}$ with $p_i\ne\ell$ and let $w\mid p_i$ be a chosen
place at which the label $c_w$ is primitive. By (G4), $R_{j,q}|_{I_w}$ is tame of finite order
dividing $p_i^r-1$, and the eigenvalues of $R_{j,q}(\gamma_w)$ are the images of the roots
$\xi'^{\,p^a}$ of the label rescaled by the element $z\in\auF_{p_i}^\times$ naming the restriction
of $\eta_j$ to $I_w$. As $p_i\equiv1\pmod{q_2^2}$ and $3\mid q_0\mid q_2$, Lemma~\ref{pg:cyclic}
applies: for every $u$ the $\ell^u$-th powers of these $r$ roots of unity are distinct, that is, no
quotient of two of them has $\ell$-power order, so they are pairwise incongruent modulo the
maximal ideal. Lemma~\ref{pg:burnside} gives (M4) with $g=\phi^r$,
$\pgcyc(g)=p_i^{\,r}$. For (M3): the places above $p_0p_1$ are the chosen places and their conjugates. At a chosen
$v\nmid\ell$ the eigenvalues of $R(\gamma_v)$ are, by (G4), the images of the roots of the
rescaled label: roots of unity $\omega\ne1$ of order dividing $p_i^r-1$ (not $1$ because the
rescaled labels do not have the root $1$: this follows from the definition of a good family alone, by the computation in part (c) of the proof of Theorem~\ref{pg:main}, which uses neither that theorem nor the present one). By (M1),
$R(g^c)$ is conjugate to $\chi(g)^{t-1}\,{}^{\mathrm t}R(g)^{-1}$, so $R(\gamma_v^{-c})$ has the
same eigenvalues as $R(\gamma_v)$; this treats the conjugate places. If
$\pgv(\omega-1)>0$ then $\omega$ has order a power $\ell^a>1$ of $\ell$, so
$\pgv(\omega-1)\le1/(\ell-1)\le1$ and $\ell\mid p_i^r-1\mid\pgNch$. Hence
$\pgv(\det(R(\gamma)-1))\le r\le r\,\pgv(\pgNch)$ if some eigenvalue is not a unit, and it is $0$
otherwise.

\emph{{\rm(M5)}.} This is Proposition~\ref{pga:M5}: with $e=(p_i^r-1)/(p_i-1)$, $s=\gamma_w^{\,e}$
acts by a scalar $y$ (the image of the norm $y_0$ of $\xi'$ to $\auF_{p_i}$) whose order is
divisible by $q_2^2$, because $p_i\equiv1\pmod{q_2^2}$; at a chosen place with label $X-1$ (the
case of the twist $\eta_j$, $1\le j\le m$) the inertia generator acts by a scalar of order
$q_2/\gcd(q_2,j)$. In both cases $y_0^x=1$ forces $q_2\mid jx$, and Lemma~\ref{pga:q2} (which
uses $q_2=(q_0d)^{s+6}$ and $2^s=m^{22}$) excludes $q_2\mid j\,u\,\ell^{a'}$ for $1\le u\le4r+2$,
$1\le j\le m$, $a'\ge0$. Hence $y^{\ell^a u}-1$ is a unit for all $a\ge0$ and $1\le u\le4r+2$,
which is (M5); the details are in Proposition~\ref{pga:M5}.

(b) For a main pair, $W$ has rank one, the eigensystem is trivial and $R_{\nu,z}$ is the
$\ell$-adic character of the algebraic Hecke character $y\mapsto\nu^{-1}(y/\aubar y)|y|^{1-z}$. It
is crystalline above $\ell\nmid M$ and potentially crystalline above $\ell\mid M$, with the single
weight $z-1-\epsilon\,c(\sigma)$ at $\sigma$, which is high because $|c-\kappa_j|\le t$; exactly
one of the weights at $\sigma$, $\aubar\sigma$ is negative. (M4) is empty in rank one except for
the scalar, and (M3), (M5) hold as above, the eigenvalue of inertia at a chosen place being the
image of the name of $\omega_w^{-j}$. The trace identity is proved as in (a).
\end{proof}

\noindent Theorem~\ref{pg:I} is proved, modulo the published theorems quoted with their hypotheses in
Theorem~\ref{pga:lit} and Lemmas~\ref{pga:depth0}, \ref{pga:parahoric}; the twist for Speh
constituents is Proposition~\ref{pax:sign} and the convention check is Lemma~\ref{pax:conv}.\par\smallskip

\subsection{The twisting theorem}\label{pg:ss-twisting}

\subsubsection*{Roots of unity}
For a polynomial $g=\sum_iy_iY^i\in\pgZb[Y]$ and $\rho\ge0$ let $\pgv_\rho(g)=\min_i(\pgv(y_i)+
\rho i)$ be the Gauss valuation. It is multiplicative, nondecreasing in $\rho$, and
$\pgv(g(y))\ge\pgv_{\rho}(g)$ if $\pgv(y)\ge\rho$. Put $\rho_h=1/\varphi(\ell^h)$, the valuation
of $\zeta-1$ for $\zeta$ of order $\ell^h$.

\begin{lemma}\label{pg:roots1}
Let $0<\delta<1$, $\rho>0$, $B\ge0$, and let $h_0$ be such that $\delta\sum_h(h+1)^{-1}>B$, the
sum over the $h$ with $1\le h\le h_0$ and $\rho_h\le\rho$. Let $g\in\pgZb[Y]$ with
$\pgv_\rho(g)\le B$. Then there are such an $h$ and a root of unity $\zeta$ of order $\ell^h$ with
$\pgv(g(\zeta-1))\le\pgv_{\rho_h}(g)+\delta$.
\end{lemma}
\begin{proof}
We may assume $g\ne0$; write $g=c\prod_i(Y-\alpha_i)$ over $\auQlb$. By multiplicativity,
$\pgv_{\rho'}(g)=\pgv(c)+\sum_i\min(\rho',\pgv(\alpha_i))$ for every real $\rho'$. For $\pgv(y)=\rho'$ one has $\pgv(y-\alpha_i)=\min(\rho',\pgv(
\alpha_i))$ unless $\pgv(\alpha_i)=\rho'$; hence
\begin{equation}\label{pg:eq-newton}
\pgv(g(y))-\pgv_{\rho'}(g)=\sum_{\alpha\in R(\rho')}\bigl(\pgv(y-\alpha)-\rho'\bigr),
\end{equation}
where $R(\rho')$ is the multiset of roots of valuation exactly $\rho'$. Since $g$ is integral, $\pgv_0(g)\ge0$, so $\pgv_\rho(g)\ge\pgv_\rho(g)-\pgv_0(g)\ge
\sum_{i:\pgv(\alpha_i)>0}\min(\rho,\pgv(\alpha_i))\ge\sum_h\#R(\rho_h)\rho_h$, the last sum over
the $h$ of the statement (the $\rho_h$ are distinct). Hence this sum is at most $B$; so for one
of them $\#R(\rho_h)\,\rho_h\le\delta/(h+1)$. Fix it. For $\alpha\in R(\rho_h)$ let $\zeta_0$ be a
root of unity of order $\ell^h$ nearest to $1+\alpha$. For the other $\zeta$ of order $\ell^h$,
$\pgv(\zeta-\zeta_0)\ge\min(\pgv(\zeta-1-\alpha),\pgv(\zeta_0-1-\alpha))=\pgv(\zeta-1-\alpha)$, and $\pgv(\zeta-\zeta_0)=\rho_{h'}$ if
$\zeta/\zeta_0$ has order $\ell^{h'}$; there are at most $\varphi(\ell^{h'})$ such $\zeta$. So
\[
\sum_{\zeta}\min\bigl(\pgv(\zeta-1-\alpha),1\bigr)\le1+\sum_{h'=1}^{h}\varphi(\ell^{h'})\rho_{h'}
=1+h .
\]
Summing over $\alpha$ and dividing by the number $\varphi(\ell^h)=1/\rho_h$ of roots of unity of
order $\ell^h$, some $\zeta$ has $\sum_{\alpha\in R(\rho_h)}\min(\pgv(\zeta-1-\alpha),1)\le
\#R(\rho_h)\rho_h(h+1)\le\delta<1$. Then no term is $1$, and by \eqref{pg:eq-newton} with
$y=\zeta-1$, $\pgv(g(\zeta-1))-\pgv_{\rho_h}(g)\le\delta$.
\end{proof}

\begin{proposition}[Roots of unity]\label{pg:roots}
Let $s_0,m\ge1$ be integers, $0<\rho\le1$, $\beta\ge0$ and $\varepsilon>0$. There is $h_0$ with the
following property. Let $P_1,\dots,P_K\in\pgZb[X_1,\dots,X_s]$, $s\le s_0$, and
$x_1,\dots,x_K\in\pgZb^{\,s}$ with coordinates of valuation at least $\rho$ and
$\sum_i\pgv(P_i(x_i))\le\beta$. Then there are roots of unity $\zeta_1,\dots,\zeta_s$ of orders
dividing $\ell^{h_0}$ and greater than $m$ such that $\pgv(\zeta_u-1)\le\rho/m$ and, for
$1\le j\le m$ and all $i$,
\[
\pgv\bigl(P_i(\zeta_1^{\,j}-1,\dots,\zeta_s^{\,j}-1)\bigr)\le\pgv(P_i(x_i))+\varepsilon .
\]
\end{proposition}
\begin{proof}
\emph{Reduction to $m=1$.} For $1\le j\le m$ let $P_{i,j}(X)=P_i((1+X_1)^j-1,\dots,(1+X_s)^j-1)$,
and let $x_{i,j}$ have coordinates $y$ with $(1+y)^j=1+x$, $x$ the coordinate of $x_i$, where
$y$ is chosen of maximal valuation among the $j$ roots of $(1+Y)^j-(1+x)$. This polynomial is
monic of degree $j$ with constant term $-x$, so the product of its roots is $\pm x$ and
$\pgv(y)\ge\pgv(x)/j\ge\rho/m$. Then $P_{i,j}(x_{i,j})=P_i(x_i)$. Apply the case $m=1$ to the $Km$
polynomials $P_{i,j}$, the radius $\rho/m$ and the bound $m\beta$. Replacing $\rho$ by $\min(\rho,1/m)$ at the start, we get $1/\varphi(\ell^h)=\pgv(\zeta_u-1)\le
1/m^2$ for the order $\ell^h$ of $\zeta_u$, so this order exceeds $m$.

\emph{The case $m=1$.} Let $\delta=\min(\varepsilon,1)/(2s_0)$ and let $h_0$ be as in
Lemma~\ref{pg:roots1} for $\delta$, $\rho$ and $B=\beta+1$. We replace the coordinates one at a
time. Suppose the first $u-1$ coordinates of all $x_i$ have been replaced by $\zeta_1-1,\dots,
\zeta_{u-1}-1$, giving points $x_i'$ with $\sum_i\pgv(P_i(x'_i))\le\beta+(u-1)\delta$ and
$\pgv(P_i(x_i'))\le\pgv(P_i(x_i))+(u-1)\delta$. Let $g_i(Y)$ be $P_i$ with the $u$-th variable
equal to $Y$ and the others equal to the coordinates of $x_i'$, and $g=\prod_ig_i$. Then
$\pgv_\rho(g)\le\sum_i\pgv(g_i(x'_{i,u}))\le\beta+1$, so Lemma~\ref{pg:roots1} gives $h$ and
$\zeta_u$ of order $\ell^h$, $\rho_h\le\rho$, with $\sum_i\pgv(g_i(\zeta_u-1))\le\sum_i\pgv_{
\rho_h}(g_i)+\delta$. Since $\pgv(g_i(\zeta_u-1))\ge\pgv_{\rho_h}(g_i)$ for each $i$, each
difference is at most $\delta$; and $\pgv_{\rho_h}(g_i)\le\pgv_{\rho}(g_i)\le\pgv(g_i(x'_{i,u}))$.
So the replacement of the $u$-th coordinate raises each valuation by at most $\delta$. After $s$
steps the total increase is at most $s\delta\le\varepsilon$.
\end{proof}

\subsubsection*{Selmer data}
Fix $S_0=\pgS(\ell)$, a chosen place $w_0\nmid\ell$ and $g_0=\gamma_{w_0}$.

\begin{definition}\label{pg:def-datum}
A \emph{Selmer datum} over $\auO$ is a continuous $u\colon\pgG\to\auGL_n(\auO)$, unramified outside
a set $S\subseteq S_0$ containing the places above $\ell$, with: for $v\mid\ell$ a finitely
generated module $Z_v$ of cocycles $D_v\to\auO^n$ containing the coboundaries, with
$\sum_{v\mid\ell}(\operatorname{rk}Z_v-n)\le f\,n$; the $u(g)$ span $M_n(\auO)$; and an element
$\gamma_*\in I_v$, $v\in S$, $v\nmid\ell$, such that $u(\gamma_*)=u(\gamma_*^{-c})=\zeta_*$ is a
scalar with $\zeta_*^2-1$ a unit. Its \emph{constant} is $\beta_T=\mathrm{gl}+\mathrm{pl}+
\mathrm{fp}$, where
\begin{align*}
\mathrm{gl}&=\sup_{a\ne0}\pglen\pgHom_{\pgG}\bigl(T/a,(T/a)^{*c}\bigr),\qquad
\mathrm{pl}=\sum_{v\in S,\,v\nmid\ell}\ \sup_{a\ne0}\pglen H^0\bigl(E_v,(T/a)^*\bigr),\\
\mathrm{fp}&=\pgv\bigl(\det(u(g_0)-1)\bigr)+\pgv\bigl(\det(u(g_0^{-c})-1)\bigr);
\end{align*}
here $T=\auO^n$, $(T/a)^*=\pgHom(T/a,\Q/\Z)(1)$, and $^{c}$ means that $g$ acts as $g^c$. The
local conditions are the relations that hold on the $Z_v$, as in Definition~\ref{pg:def-sel}.
\end{definition}

\begin{lemma}\label{pg:beta}
Let $R$ be a main summand of rank $n$. For all large $\auO$ there is a lattice $u$ in $R$ that,
with $S$ the set of places above $\ell p_0p_1$ and $Z_v=Z_v(u)$ the crystalline cocycles, is a
Selmer datum with $\beta_T\le(4f+3)\,n\,\pgv(\pgNch)$.
\end{lemma}
\begin{proof}
\emph{The rank condition.} For a de Rham representation $V$ of $\operatorname{Gal}(\overline{E_v}
/E_v)$, $\dim H^1_f=\dim D_{\mathrm{dR}}(V)/\mathrm{Fil}^0+\dim H^0$ (Bloch--Kato~\cite[Corollary~3.8.4]{BlochKato}), so $\operatorname{rk}Z_v-n=\dim H^1_f-\dim H^0=\dim_{\Q_\ell}D_{\mathrm{dR}}/
\mathrm{Fil}^0$ counted over $K$, which is the number of negative Hodge--Tate weights at the
embeddings of $E_v$ (with the convention of (G3)). Summing over $v\mid\ell$ gives the number of
negative weights at all $2f$ embeddings of $E$, at most $f\,n$ by (M2).

\emph{The scalar element.} Take $\gamma_*=s$ from (M5), $\zeta_*=y$; $y^2-1$ is a unit. By (M1)
and semisimplicity, $g\mapsto R(g^c)$ is isomorphic to $g\mapsto\chi(g)^{t'-1}\,{}^{\mathrm
t}R(g)^{-1}$; so $R(g^{-c})$ is conjugate, by a matrix independent of $g$, to
$\chi(g)^{1-t'}\,{}^{\mathrm t}R(g)$. As $\chi(\gamma_*)=1$, $R(\gamma_*^{-c})$ is the scalar
$y$.

\emph{The global term.} By (M4) there is $g$ with $\pgcyc(g)=p_i^{\,n}$ and $u(g)=\alpha$ a
scalar; by the same remark, $u(g^{-c})$ is the scalar $\alpha'=\alpha\,\chi(g)^{1-t'}$. On $(T/a)^{*c}$ the
element $g$ acts by $\pgcyc(g^c)\alpha'=\pgcyc(g)\alpha'$. So $\alpha-\pgcyc(g)\alpha'=\alpha(1-
\pgcyc(g)^{t'})=\alpha(1-p_i^{\,nt'})$ kills $\pgHom_{\pgG}(T/a,(T/a)^{*c})$. This module is
cyclic: both $T/a$ and $(T/a)^{*c}$ are free of rank $n$ over $\auO/a$ with residually absolutely
irreducible actions; if $\eta=\varpi^k\eta_0$ is an equivariant map with $k$ least, then $\eta_0$,
an equivariant map modulo $a/\varpi^k$, is not zero modulo $\varpi$, hence an isomorphism, and
every equivariant $\eta'$, being divisible by $\varpi^k$, differs from $\eta$ by an endomorphism
of $T/(a/\varpi^k)$ commuting with $M_n(\auO)$, a scalar. Hence $\mathrm{gl}\le\pgv(p_i^{\,nt'}-1)\le\pgv(\pgNch)$, because
$p^{nt'}-1$ divides $\#\auGL_{rt}(\auF_p)$ for $nt'\le rt$.

\emph{The local terms.} For $v\mid p_0p_1$, $v\nmid\ell$, let $\gamma\in I_v$ be as in (M3). Then
$H^0(E_v,(T/a)^*)$ is contained in the kernel of $u(\gamma)^{-1,\mathrm t}-1$ on
$\pgHom(T/a,\Q/\Z)$ (inertia acts trivially on the Tate twist), whose length is at most
$\pgv(\det(u(\gamma)-1))\le n\,\pgv(\pgNch)$. There are $4f$ such places at most, so
$\mathrm{pl}\le4f\,n\,\pgv(\pgNch)$; and $\mathrm{fp}\le2n\,\pgv(\pgNch)$ by (M3) for the chosen
place $w_0$. If $\pgv(\pgNch)=0$ all terms vanish; otherwise the sum is at most $(1+4f+2)\,n\,
\pgv(\pgNch)$.
\end{proof}

\subsubsection*{One system of equations for all twists}
Let $N\ge M\ge1$. Let $J$ be a finite set of places of $E$ of prime norm $\equiv1\pmod{\ell^N}$
such that $J$, $\aubar J$, $S_0$ are pairwise disjoint, and for $v\in J$ let $\lambda_v\colon\pgG
\to\Z/\ell^N$ be a homomorphism that is anticyclotomic, unramified outside $v,\aubar v$, trivial
on $D_w$ for $w\mid\ell$, with $\lambda_v(\gamma_v)=1$. For a point $x=(x_v)_{v\in J}$ with
coordinates in the ring of integers $\auO'$ of a finite extension of $K$ in $\auQlb$ put
$\psi_x(g)=\prod_v(1+x_v)^{\lambda_v(g)}$ (exponents in $[0,\ell^N)$). If $\pgv(x_v)\ge\rho$ with
$M\le N\rho$, $\rho\le1$, then $(1+x_v)^{\ell^N}\equiv1\pmod{\ell^M}$ (each term
$\binom{\ell^N}{i}x^i$, $i=\ell^ai'$, has valuation at least $N-a+\ell^a\rho\ge N\rho$), so
$\psi_x$ is a character modulo $\ell^M$. For a Selmer datum $T$ let $\pgSel(x)$ be the module of
locally constant cocycles $\pgG\to(\auO'/\ell^M)^n$ for $\psi_xu$ that are unramified outside
$S\cup J\cup\aubar J$, obey the local relations of $T$ and vanish at $\gamma_*$. If $x_v=\zeta_v-1$
for $\ell^N$-th roots of unity $\zeta_v$, then $\psi_x=\prod\zeta_v^{\lambda_v}$ is a character
$\psi$ of $\pgG$ and $\pgSel(x)$ is the Selmer group of the twist $\psi u$ modulo $\ell^M$ (the
local relations are unchanged because $\psi$ is trivial on $D_w$, $w\mid\ell$).

\begin{lemma}\label{pg:minor}
There are a finite set $\mathcal P$ of polynomials in $\auO[X_v:v\in J]$ with the following
properties, for every $\auO'$ and every point $x$ as above. {\rm(i)} For every $P\in\mathcal P$:
$\pglen\pgSel(x)\le\pgv(P(x))$. {\rm(ii)} If $\pglen\pgSel(x)<M$ then some $P\in\mathcal P$ has
$\pgv(P(x))=\pglen\pgSel(x)$. {\rm(iii)} If $x=(\zeta_v-1)$ and $\pgv(P(x))<M$ then, for all large
coefficient rings, $\pglen\pgSel_J(\psi u,\gamma_*;K/\auO)\le\pgv(P(x))$.
\end{lemma}
\begin{proof}
Let $U\subseteq\pgG$ be the kernel of the reduction of $u$ modulo $\ell^M$ and of all $\lambda_v$.
The restriction to $U$ of a cocycle of $\pgSel(x)$ is a homomorphism $U\to(\auO'/\ell^M)^n$
unramified outside $S\cup J\cup\aubar J$, so it factors through the Galois group of the maximal
abelian extension of exponent $\ell^M$ of the fixed field of $U$ unramified outside these places,
which is finite. So there is a finite quotient $\Gamma$ of $\pgG$, independent of $x$ and
$\auO'$, such that $\pgSel(x)$ is the module of functions $z\colon\Gamma\to(\auO'/\ell^M)^n$
satisfying: $z(gh)=z(g)+\psi_x(g)u(g)z(h)$; $z(g)=0$ for $g$ in the images of the inertia groups
outside $S\cup J\cup\aubar J$ and for $g=\gamma_*$; and the local relations, of which finitely
many suffice because the $Z_v$ are finitely generated and only the values on $\Gamma$ matter. These
are linear equations whose coefficients are values at $x$ of polynomials over $\auO$: $\pgSel(x)=
\ker A(x)$ on $(\auO'/\ell^M)^{d}$, $d=n\#\Gamma$, for a matrix $A$ over $\auO[X]$ with $d$
columns. Let $\mathcal P$ be the set of $d\times d$ minors of $A$. (i) The kernel of $A(x)$ is
contained in the kernel of a square submatrix $B$, and for a square matrix over a discrete
valuation ring the kernel of $B$ on $(\auO'/a)^d$ has length $\sum_i\min(\pgv(a),\pgv(d_i))\le
\pgv(\det B)$, the $d_i$ being the elementary divisors. (ii) Let $d_1\mid d_2\mid\dots$ be the
elementary divisors of $A(x)$ over $\auO'$; the length of the kernel is
$\sum_{i\le d}\min(M,\pgv(d_i))$ (with $d_i=0$ beyond the rank). If it is less than $M$ then all
$\pgv(d_i)<M$, the rank is $d$ and the length is $\pgv(d_1\cdots d_d)$, the least valuation of a
$d\times d$ minor. (iii) Multiplication by $\ell^M$ on $\pgSel_J(\psi u,\gamma_*;K/\auO)$ has
kernel $\pgSel(x)$, of length less than $M=\pglen\auO/\ell^M$; if some element of the group were not killed by $\ell^M$, the cyclic submodule it
generates would have $\ell^M$-torsion of length $M$. So the group equals its $\ell^M$-torsion.
\end{proof}

\begin{lemma}[Tame inertia and Frobenius]\label{pg:tame}
Let $v\in J\cup\aubar J$, $\gamma$ a tame generator and $\phi$ a Frobenius element at $v$, and let
$z$ be a cocycle of $\pgSel(x)$, where $0<\rho\le1$, $\pgv(x_w)\ge\rho$ for all $w$ and $M\le N\rho$. Then
\[
\bigl(\psi_x(\phi)u(\phi)-1\bigr)z(\gamma)=\bigl(\psi_x(\gamma)-1\bigr)z(\phi).
\]
\end{lemma}
\begin{proof}
$u$ is unramified at $v$, so $\gamma$ acts by the scalar $\alpha=\psi_x(\gamma)$, an $\ell^N$-th
root of unity modulo $\ell^M$, with $\sum_{j<\ell^N}\alpha^j\equiv0\pmod{\ell^M}$. Indeed, write
$\alpha=1+y$; then $\pgv(y)\ge\rho$, because $\alpha$ is a product of integer powers of the
$1+x_w$, each congruent to $1$ modulo elements of valuation at least $\rho$. Then
$\sum_{j<\ell^N}\alpha^j=\sum_{i=1}^{\ell^N}\binom{\ell^N}{i}y^{i-1}$, and for $i=\ell^ai'$ with
$\ell\nmid i'$ (so $a\le N$) the $i$-th term has valuation at least
$N-a+(\ell^a-1)\rho\ge N\rho\ge M$, since $N-a+(\ell^a-1)\rho-N\rho=(N-a)(1-\rho)+(\ell^a-1-a)\rho\ge0$
for $0<\rho\le1$. Here $\phi$ may be an arithmetic or a geometric Frobenius element (both
occur below, e.g.\ $\phi_v^{-c}$ at $\aubar v$). The cocycle $z$ vanishes on wild inertia (a homomorphism from a pro-$p$ group to an
$\ell$-group), and $\phi\gamma\phi^{-1}=\gamma^{q}$ modulo wild inertia, where $q$ is a positive integer
congruent to $(\auNm v)^{\pm1}$ modulo $\ell^N$ (the sign depending on whether $\phi$ is
arithmetic or geometric); as $\auNm v\equiv1\pmod{\ell^N}$, $q\equiv1\pmod{\ell^N}$. Since
$z(\gamma^{k+\ell^N})=z(\gamma^k)+\alpha^k\bigl(\sum_{j<\ell^N}\alpha^j\bigr)z(\gamma)=z(\gamma^k)$
and $\alpha^{\ell^N}=1$, the value $z(\gamma^k)$ and $\alpha^k$ depend only on $k$ modulo
$\ell^N$. Hence $z(\gamma^q)=z(\gamma)$ and $\alpha^q=\alpha$. Comparing $z(\phi\gamma)=z(\phi)+\psi_x(\phi)u(\phi)z(\gamma)$ with
$z(\gamma^q\phi)=z(\gamma)+\alpha z(\phi)$ gives the relation.
\end{proof}

\subsubsection*{Tied Selmer groups and the descent}
A \emph{set of tied primes} is a set $J$ as above with a \emph{first prime} $P_0\in J$, Frobenius
elements $\phi_v$ and, for each datum $T_i$ of a finite family, \emph{slopes} $s_{v,i}\in\auO$
($v\in J$), such that: $s_{P_0,i}=0$; $u_i(\phi_{P_0})\equiv u_i(g_0)$ and
$u_i(\phi_{P_0}^{-c})\equiv u_i(g_0^{-c})$ modulo $\ell^M$; for $v\ne P_0$,
$\lambda_{P_0}(\phi_v)=1$ and $\lambda_w(\phi_v)=0$ for $w\in J\setminus\{P_0\}$; and
$u_i(\phi_v)\equiv u_i(\phi_v^{-c})\equiv1\pmod{\ell^M}$ whenever $s_{v,i}\ne0$. The \emph{tied
Selmer group} $\pgSel^{\mathrm{tie}}_i(J;a)$, $a\mid\ell^M$, is the module of cocycles
$z\colon\pgG\to(T_i/a)$ for $u_i$ (no twist), unramified outside $S\cup J\cup\aubar J$, obeying
the local relations, vanishing at $\gamma_*$, with the \emph{ties}
\[
z(\gamma_v)=s_{v,i}\,z(\phi_v),\qquad z(\gamma_v^{-c})=s_{v,i}\,z(\phi_v^{-c})\qquad(v\in J).
\]
A tie with slope $0$ says that $z$ is unramified at $v$ and $\aubar v$; a tie with any slope is a
local condition at $v$ of the same size as the unramified one, since a cocycle of the tame
quotient is determined by its values at $\gamma_v$ and $\phi_v$.

\begin{lemma}[The count]\label{pg:count}
Let $J$ be a set of tied primes and $a\mid\ell^M$. Then
\[
\pglen\pgSel^{\mathrm{tie}}_i(J;a)\le\pglen\pgSel^{\mathrm{tie},*}_i(J;a)+\mathrm{pl}(T_i),
\]
where $\pgSel^{\mathrm{tie},*}$ is the Selmer group of $(T_i/a)^*$ with the local conditions
orthogonal, under local Tate duality, to those of $\pgSel^{\mathrm{tie}}$ (no condition at the
places of $S$ prime to $\ell$ becomes the zero condition there).
\end{lemma}
\begin{proof}
By the formula of Greenberg and Wiles for Selmer groups of a finite module $X=T_i/a$ with
arbitrary local conditions $L_v\subseteq H^1(E_v,X)$ (\cite[Proposition~1.6]{WilesFLT} over
$\Q$ for $\ell$ odd; the proof, by the Poitou--Tate sequence and the global Euler characteristic formula, is the
same over a number field, see~\cite[Theorem~8.7.9]{NSW}),
\[
\pglen H^1_L-\pglen H^1_{L^\perp}=\pglen H^0(E,X)-\pglen H^0(E,X^*)+\sum_v\bigl(\pglen L_v-\pglen
H^0(E_v,X)\bigr),
\]
the sum over the places of $S\cup J\cup\aubar J$ and the $f$ complex places, where $L_v=0$ and
the term is $-\pglen X$. Here $H^0(E,X)=0$ because $\zeta_*-1$ is a unit, and the classes are
the cocycles vanishing at $\gamma_*$. We bound the local terms. For $v\mid\ell$: let $C_v$ be the
module of functions $D_v\to X$ obeying the relations of $Z_v$; they are cocycles, because the
cocycle rule is among the relations, and $C_v$ contains the coboundaries, because $Z_v$ does; $L_v$
is the image of $C_v$. Let $\Phi$ be the free $\auO$-module on the symbols $(g,i)$, $g\in D_v$,
$1\le i\le n$, so that functions $D_v\to B^n$ are the linear maps $\Phi\to B$, and let
$N\subseteq\Phi$ be the common kernel of the elements of $Z_v$. Then $\Phi/N$ embeds in
$\pgHom_{\auO}(Z_v,\auO)$, so it is free of rank at most $\operatorname{rk}Z_v$, and
$C_v=\pgHom_{\auO}(\Phi/N,\auO/a)$ has length at most $\operatorname{rk}Z_v\cdot\pgv(a)$. The
coboundaries have length $\pglen X-\pglen H^0(E_v,X)$; so the term is at most
$(\operatorname{rk}Z_v-n)\pgv(a)$, and the sum over $v\mid\ell$ is at most $f\,n\,\pgv(a)=f\pglen
X$, which cancels the complex places. For $v\in S$, $v\nmid\ell$, with no condition:
$\pglen H^1(E_v,X)-\pglen H^0(E_v,X)=\pglen H^2(E_v,X)=\pglen H^0(E_v,X^*)$ by the local Euler
characteristic formula and local duality; the sum is at most $\mathrm{pl}(T_i)$. For $v\in
J\cup\aubar J$: a cocycle of $D_v$ with values in $X$ is a homomorphism on $I_v$ and factors
through the quotient topologically generated by $\gamma_v$ and $\phi_v$; a tied one is determined
by $z(\phi_v)$, so the tied cocycles have length at most $\pglen X$. They contain the
coboundaries (which vanish at $\gamma_v$, and at $\phi_v$ when the slope is not $0$). So
$\pglen L_v\le\pglen X-(\pglen X-\pglen H^0(E_v,X))$ and the term is at most $0$. Finally
$\pglen H^0(E,X^*)\ge0$.
\end{proof}

\begin{proposition}[The dual of a tied Selmer group]\label{pg:mirror}
Let $J$ be a set of tied primes, $a=\varpi^k\mid\ell^M$, $M<N$, and suppose that
$\pgSel_i^{\mathrm{tie}}(J;a)$ contains a cocycle $z$ one of whose values has a unit coordinate.
Suppose that $z$ is balanced at every place $P$ admissible for $T_i$ (\S\ref{pse:ss-aux}): the pair
$(z(\phi_P),z(\phi_P^{-c}))$ is the pair of values at $(\gamma_P,\gamma_P^{-c})$ of a cocycle $z'$
with the local conditions and ties of $J$ that may ramify at $P$ and $\aubar P$. Then
$\pglen\pgSel_i^{\mathrm{tie},*}(J;a)\le\mathrm{gl}(T_i)$.
\end{proposition}
\smallskip\noindent Proved in Proposition~\ref{pse:mirror}.\par\smallskip

\begin{proposition}[One step of the descent]\label{pg:step}
Let $J$ be a set of tied primes for data $T_1,\dots,T_s$ over $\auO$ with uniformiser $\varpi$,
$\varpi^e\mid\ell^M$, and suppose that the residue field of $\auO$ has more than $2n_ie+1$ elements
for all $i$. If for some $i$, some $1\le k\le e$ and some place $P$ that can be added to $J$, an
element of $\pgSel^{\mathrm{tie}}_i(J;\varpi^k)$ is not balanced at $P$, then $P$ can be added to
$J$ with a unit slope for $T_i$ and slope $0$ for the other data so that
$\sum_{i'}\sum_{k'=1}^e\pglen\pgSel_{i'}^{\mathrm{tie}}(J;\varpi^{k'})$ decreases.
\end{proposition}
\smallskip\noindent Proved in Proposition~\ref{pse:step}.\par\smallskip

\begin{proposition}[Polynomials and points]\label{pg:poly}
Assume Propositions~\ref{pg:mirror}, \ref{pg:step} and statement \emph{(A)} below. Let $T_1,\dots,T_s$ be Selmer data over
$\auO$, let $M>\max_i\beta_{T_i}+1$, and suppose the residue field of $\auO$ has more than
$2n_i\pgv(\ell^M)/\pgv(\varpi)+1$ elements. There is an integer $s_0\ge1$ such that for every
rational $0<\rho<1/s_0$ and every large $N$ there are a set $J$ with characters $\lambda_v$ as
above, $2\#J\le s_0$, and for each $i$ a polynomial $P_i\in\pgZb[X_v:v\in J]$ and a point $x_i$
with coordinates of valuation at least $\rho$, such that $P_i$ has the properties {\rm(i)},
{\rm(iii)} of Lemma~\ref{pg:minor} for $T_i$ and $\pgv(P_i(x_i))\le\beta_{T_i}+s_0\rho$.
\end{proposition}
\smallskip\noindent\textbf{(A)} is Proposition~\ref{pse:A}; ``a place that can be added to $J$'' means a place admissible in the sense of \S\ref{pse:ss-aux}.\par\smallskip
\begin{proof}
\emph{The descent.} Start with $J=\{P_0\}$, where $P_0$ is a place of prime norm $\equiv1
\pmod{\ell^N}$ outside $S_0\cup\aubar{S_0}$ whose Frobenius element is so near $g_0$ that it acts
as $g_0$, and $\phi_{P_0}^{-c}$ as $g_0^{-c}$, on all $T_i/\ell^M$, with an anticyclotomic
character $\lambda_{P_0}$ of level $P_0$, $\lambda_{P_0}(\gamma_{P_0})=1$; such places and characters exist by Proposition~\ref{pse:A} (see Remark~\ref{pse:rem-poly}). All slopes are $0$, and the tied groups are the Selmer
groups without auxiliary primes, of total length $b_0$ over the levels $\varpi^k$, $k\le e$,
independent of $N$; here, and only here, $b_0$ is the \emph{unnormalised} length (the number of
composition factors $\auO/\varpi$), an integer, so that a strict decrease is a decrease by at least
$1$. Apply Proposition~\ref{pg:step} as long as it applies; each step lowers the
total length, so after at most $b_0$ steps every element of every tied group modulo every
$\varpi^k$ is balanced at every admissible $P$. Then for each $i$: by passage to the exponent,
there is $k'\le e$ with $\pglen\pgSel^{\mathrm{tie}}_i(J;\varpi^e)\le\pglen\pgSel^{\mathrm{tie}}_i
(J;\varpi^{k'})$ and, if $k'\ne0$, an element of the latter with a unit coordinate (take $k'$
least such that $\varpi^{k'}$ kills the group modulo $\varpi^e$; its elements have values in
$\varpi^{e-k'}T/\varpi^e\cong T/\varpi^{k'}$, and one not killed by $\varpi^{k'-1}$ has, there, a
value outside $\varpi T/\varpi^{k'}$). By Lemma~\ref{pg:count} and Proposition~\ref{pse:mirror},
$\pglen\pgSel^{\mathrm{tie}}_i(J;\ell^M)\le\mathrm{gl}(T_i)+\mathrm{pl}(T_i)$. Put $s_0=2\max(n,2)(1+b_0)$,
$n=\max n_i$.

\emph{Release of the first prime.} Let $\pgSel'_i$ be the group with the ties at $J\setminus
\{P_0\}$ only. By Lemma~\ref{pg:tame} at the point $x=0$ and $v=P_0$, every $z\in\pgSel_i'$ has
$(u_i(g_0)-1)z(\gamma_{P_0})=0$, and likewise at $\aubar{P_0}$; so the quotient
$\pgSel'_i/\pgSel_i^{\mathrm{tie}}$ embeds in the product of the kernels of $u_i(g_0)-1$ and
$u_i(g_0^{-c})-1$ on $T_i/\ell^M$ and has length at most $\mathrm{fp}(T_i)$. Thus
$\pglen\pgSel_i'\le\beta_{T_i}<M$.

\emph{The line.} Fix $i$ and consider the points $x(y)$ with $x_{P_0}=y$, $x_v=s_{v,i}y$
($v\ne P_0$), $\pgv(y)\ge\rho$. For $v\ne P_0$ in $J$ we have $\psi_{x(y)}(\gamma_v)=1+s_{v,i}y$
and $\psi_{x(y)}(\phi_v)=1+y$. Let $z\in\pgSel(x(y))$. By Lemma~\ref{pg:tame},
$((1+y)u_i(\phi_v)-1)z(\gamma_v)=s_{v,i}y\,z(\phi_v)$, and therefore
\[
\bigl((1+y)u_i(\phi_v)-1\bigr)\bigl(z(\gamma_v)-s_{v,i}z(\phi_v)\bigr)=-s_{v,i}(1+y)\bigl(u_i(
\phi_v)-1\bigr)z(\phi_v)=0,
\]
because $u_i(\phi_v)\equiv1$ if $s_{v,i}\ne0$; the same holds at $\aubar v$. Let
$\pgSel^{\mathrm{tie}}(y)\subseteq\pgSel(x(y))$ be the submodule on which the ties at
$J\setminus\{P_0\}$ hold. The quotient embeds in the product over $v\in J\setminus\{P_0\}$ and
$g\in\{\phi_v,\phi_v^{-c}\}$ of the kernels of $(1+y)u_i(g)-1$, so its length is at most
$\pgv(P^1(y))$, $P^1(Y)=\prod_{v,g}\det((1+Y)u_i(g)-1)$. Each factor has the unit coefficient
$\det u_i(g)$ in degree $n_i$; hence $\pgv_\rho(P^1)\le2n_i\#J\rho\le s_0\rho$. The modules
$\pgSel^{\mathrm{tie}}(y)$ are the kernels of one matrix over $\auO[Y]$, as in
Lemma~\ref{pg:minor}; at $y=0$ the module is $\pgSel'_i$, of length less than $M$, so by
Lemma~\ref{pg:minor}(ii) there is a minor $P^0\in\auO[Y]$ with $\pgv(P^0(0))=\pglen\pgSel'_i\le
\beta_{T_i}$ and $\pglen\pgSel^{\mathrm{tie}}(y)\le\pgv(P^0(y))$ for all $y$. In particular
$\pgv_\rho(P^0)\le\pgv(P^0(0))\le\beta_{T_i}$. Choose $y$ of valuation exactly $\rho$ (in a
ramified extension) with $\pgv(P^0P^1(y))=\pgv_\rho(P^0P^1)$; this excludes finitely many residue
classes of $y/\varpi'^{\,\rho}$. Then
\[
\pglen\pgSel(x(y))\le\pgv(P^0(y))+\pgv(P^1(y))\le\beta_{T_i}+s_0\rho<M .
\]
By Lemma~\ref{pg:minor}(ii) there is $P_i\in\mathcal P$ with $\pgv(P_i(x(y)))=\pglen\pgSel(x(y))$;
take $x_i=x(y)$. All its coordinates have valuation at least $\rho$.
\end{proof}

\subsubsection*{The $L$-values}

\begin{proposition}[Polynomials for the ratios of $L$-values]\label{pg:stick}
In the situation of Proposition~\ref{pg:poly}, let finitely many pairs of a ramified critical
pair $(\nu,z)$ and a sign $\delta=\pm1$ be given, with $M>\beta_Y+1$ for each ($\beta_Y$ as in
\S\ref{pse:ss-stick}). Then, after enlarging $s_0$, for every rational $0<\rho<1/s_0$ and every
large $N$ the set $J$ with its characters can be chosen so that moreover, for each of the given
pairs, there are a polynomial $P\in\pgZb[X_v:v\in J]$ and a point $x$ with coordinates of valuation
at least $\rho$ with
\[
\pgv(P(x))\le v_0+4\pgv(\auNm w_0-1)+s_0\rho ,
\]
where $v_0$ is the valuation at $\iota$ of the value at the trivial twist of the measure of
statement (IM) of \S\ref{pse:ss-stick} and $w_0$ is a chosen place, and such that, for all
$\ell^N$-th roots of unity $\zeta_v\ne1$, with $\psi=\prod_v\zeta_v^{\lambda_v}=\psi_\vartheta$,
\[
\pgv\Bigl(\frac{|L(z,(\nu\vartheta^{\delta})^{-1})|}{|L(z,\nu^{-1})|}\Bigr)\le
\max\bigl(0,\ \pgv\bigl(P((\zeta_v-1)_v)\bigr)-v_0\bigr).
\]
\end{proposition}
\smallskip\noindent This is Proposition~\ref{pse:stick}, proved there from statement (IM), which is
proved in {\cite[Propositions~3.8 and~3.12]{P54}}.\par\smallskip

\begin{proof}[Proof of Theorem~\ref{pg:II} from Lemma~\ref{pg:beta} and
Propositions~\ref{pg:roots}, \ref{pg:poly}, \ref{pg:stick}]
Let $c=(4f+3)\,r\,\pgv(\pgNch)$ and $\varepsilon=2^{-\ell}$. By Lemma~\ref{pg:beta} there is
$\auO$ over which each $R_i$ has a lattice $u_i$ that is a Selmer datum $T_i$ with $\beta_{T_i}\le
c$; enlarging $\auO$ (the constants do not grow, lengths being normalised) we may assume its
residue field is as large as Proposition~\ref{pg:poly} requires for an integer $M>\max(c,\max_Y\beta_Y)+1$ (\S\ref{pse:ss-end}). Let $s_0$
be given by Propositions~\ref{pg:poly}, \ref{pg:stick}, choose a rational $\rho$ with
$s_0\rho\le\varepsilon/2$, let $h_0$ be given by Proposition~\ref{pg:roots} for $s_0$, $m$,
$\rho$, the bound $\beta=s(c+1)+\sum_Y(v_0(Y)+4\pgv(\auNm w_0-1)+1)$ and
$\varepsilon/2$, and let $N\ge\max(h_0,M/\rho)$ be large. Take $J$, the $\lambda_v$, the
polynomials and the points of Propositions~\ref{pg:poly} and~\ref{pg:stick}, and let
$\zeta_v$ ($v\in J$) be the roots of unity of Proposition~\ref{pg:roots}; their orders divide
$\ell^{h_0}\mid\ell^N$ and exceed $m$. Let $\psi_1=\prod_v\zeta_v^{\lambda_v}$. It is a character of
$\pgG$ of $\ell$-power order, anticyclotomic, unramified outside $J\cup\aubar J$, trivial on the
decomposition groups above $\ell$, tamely ramified at $v\in J$ with $\psi_1(\gamma_v)=\zeta_v$ of
order greater than $m$. By the remarks at the beginning of \S\ref{pg:ss-reduction},
$\psi_1=\psi_\tau$ for an $m$-fold twist $\tau$ with $Q_\tau=J$, and $\psi_{\tau^j}=\psi_1^{\,j}$
corresponds to the point $(\zeta_v^{\,j}-1)_v$.

(a) $J$, $\aubar J$, $S_0$ are pairwise disjoint and the norms are prime; and
$\sum_{v\in J}\pgv(\zeta_v^{\,j}-1)\le\#J\cdot\rho\le s_0\rho\le\varepsilon$ (if $\zeta$ has order $\ell^h>m\ge j$ then
$\zeta^j$ has order $\ell^{h-a}$, $\ell^a\mid j$, and $\pgv(\zeta^j-1)=\ell^a\pgv(\zeta-1)\le
j\,\rho/m\le\rho$). The same bound holds at the conjugate places.

(b) For each $i$ and $j$: $\pgv(P_i((\zeta_v^{\,j}-1)_v))\le\pgv(P_i(x_i))+\varepsilon/2\le
\beta_{T_i}+\varepsilon<M$; by Lemma~\ref{pg:minor}(iii) the Selmer group of $\psi_1^{\,j}u_i$ with
coefficients in $K/\auO$ has length at most $c+\varepsilon=e'_\ell$ for all large coefficient
rings. Its local conditions are those of Definition~\ref{pg:def-sel}, and $\gamma=\gamma_*$,
$\zeta=\zeta_*$ serve, since $\psi_1(\gamma_*)=1$.

(c) By Proposition~\ref{pse:stick} and \S\ref{pse:ss-end}(4), the valuation of each ratio at $\vartheta=\tau^j$ is at most $4\pgv(\pgNch)+\varepsilon\le c+\varepsilon$; this is the lower bound $-e'_\ell$ for the
valuation of the inverse ratio.
\end{proof}

\subsection{From Selmer groups to denominators}\label{pg:ss-ribet}

Theorem~\ref{pg:III} is proved in \S\ref{pri:sec} (Theorem~\ref{pri:III}) from Propositions~\ref{pri:loss}, \ref{pri:odd}, \ref{pri:kummer}, \ref{pri:dict}, \ref{pri:BL} and Lemmas~\ref{pri:chain}, \ref{pri:block}, \ref{pri:forms}; the propositions below are restated there in exact form. Fix $\ell$, $\iota$, $\vartheta$, $O$, $\mathcal Z$ as in the theorem,
$n=n_V$, $r'=n_W$, $t'=t(O)$, and one generator $y_0$ of $\mathcal Z$, with system $x_0$ and main
summand $R$.

\emph{The mechanism.} Let $y$ be integral at $\iota$ and orthogonal to the false lifts. Its
component orthogonal to $\mathcal Z$ lies, in each type $e_c$, in the span of the eigenforms that
are not \emph{weak lifts} (an eigenform is a weak lift if $\det(1-(S\otimes\eta_O)(g))=0$ for all
$g$); this is where orthogonality to the false lifts is used (Lemma~\ref{pri:weak}). A denominator $\ell^\nu$ in
$\pgpr_{\mathcal Z}(y)$ at the component of $y_0$ is therefore a congruence modulo $\ell^\nu$
between $y_0$ and eigenforms $y_x$ whose normalised Galois representations $S_x$ do not have the
eigenvalue $1$ identically, while $S_{x_0}=\psi R\oplus\bigoplus_{e<t'}\chi^e$ on traces. Ribet's
method produces from such a congruence a non-split extension between two of the $t'+1$ summands
modulo a power of $\ell$. Extensions between $\psi R$ and a character are bounded by the Selmer
group of $\psi R$ (Theorem~\ref{pg:II}(b)), up to the local loss above $\ell$. Extensions between
two characters $\chi^e$, $\chi^{e'}$ are classes in Tate twists of the trivial module: the
\emph{odd} ones are bounded by $\pgv(\pgcodd)$, and for the \emph{even} ones, whose order is
governed by class groups of real cyclotomic fields and cannot be bounded, one passes to
forms of auxiliary level at Taylor--Wiles places $w_j$ chosen in a triangular way, where the
even classes are killed one after another, and works to first order in parameters $\alpha_j$ that
tie tame inertia at $w_j$ to Frobenius.

\begin{proposition}[Ribet's method to first order]\label{pg:ribet}
Let $G$ be a group, $\auO$ a discrete valuation ring with uniformiser $\varpi$, $p=\ell^\nu$, and
$u_x\colon G\to\auGL_n(\auO)$, $x\in Y$ finite, representations such that for $x$ in a subset
$\mathcal X$ the representation $u_x$ is block diagonal with a block $\rho_x$ of rank $r'$ and
the characters $\xi_e$, $e<t'$, $\xi_0=1$. Let $\tau_j,\phi_j\in G$ ($j\in I$), polynomials
$H_j$ with $H_j(1)=1$, and let $\mathfrak a$ be the left ideal of $\auO[X_j:j\in I][G]$ generated
by the elements killed by all $u_x$, by $p$ times the augmentation ideal, and by the \emph{tied
relations} $\tau_j-1-X_j(\phi_j-1)H_j(\phi_j)$. Assume (in the exact form of
Proposition~\ref{pri:ribet}): a triangle
cost $b_e$ for the cocycles of each pair of characters, a bound $b_0$ for the length of the
cocycles of the block with coefficients in $K/\auO$ obeying the relations among the $u_x(g)-1$ and
vanishing at the $\tau_j$, and an element separating the characters up to a cost $C$. Then there
is an element $\mathcal D=\sum_{S\subseteq I}\mathcal D_SX_S$ of the algebra generated over
$\auO[X]$ by the coefficients of the characteristic polynomials of the $u_x(g)$ (as functions of
$x$) such that $\mathcal D\cdot\det(1-u_x(g))$ lies, for every $g$, in the ideal of traces of
$\mathfrak a$; $\mathcal D_S$ vanishes on $\mathcal X$ for $S\ne I$; and $\mathcal D_I$ takes on
$\mathcal X$ a value in $\auO$ of valuation at most $b_0+\sum_eb_e+C$.
\end{proposition}
\smallskip\noindent Exact statement and proof: Proposition~\ref{pri:ribet}.\par\smallskip

\begin{proposition}[Taylor--Wiles places in a triangle]\label{pg:triangle}
For every $\nu\ge1$ and every prescribed open condition on Frobenius elements that is satisfied by
the elements of a non-empty open subset of the kernel of the cyclotomic character modulo
$\ell^\nu$, there are places $w_j\nmid\ell$ of degree one, $j$ in a finite ordered set $I$ whose
size does not depend on $\nu$, with $\auNm w_j\equiv1\pmod{\ell^\nu}$, Frobenius elements $\phi_j$
satisfying the condition, and exponents $0<\mathrm{ty}_j<t'$, such that every even class in
$(\auO/\ell^\nu)(\chi^{\mathrm{ty}_j})$ that vanishes at the earlier $\tau_i$ and at the earlier
$\phi_i$ of the same exponent has its value at $\tau_j$ killed by $\ell^{h'}$, where $h'\le
h\,\pgv(\pgNch)$ and $h'=0$ if $\ell>h^2$. Here a class is \emph{even} if it is invariant under
$g\mapsto g^c$, unramified outside $\ell$ and the $w_j$, and vanishes on the elements of trivial cyclotomic character in the decomposition groups above $\ell$.
\end{proposition}
\smallskip\noindent Proved in Proposition~\ref{pri:triangle}, with Proposition~\ref{pri:kummer} proved in
\S\ref{pku:sec}.\par\smallskip

\begin{proposition}[Odd classes]\label{pg:odd}
There is a positive integer $\pgcodd\le2^{4(q_1d+t+3)^{27}}h_E$ such that for every prime $\ell$,
all large $\auO$, all $0\le e<t$ and $\nu\ge1$ the module of cocycles $\pgG\to(\auO/\ell^\nu)(
\chi^e)$ that, for $e\ne0$, are unramified outside $\ell$, satisfy the local condition above $\ell$ of
Proposition~\ref{pg:triangle} and change sign under $g\mapsto g^c$, and, for $e=0$, the module of
everywhere unramified homomorphisms $\pgG\to\auO/\ell^\nu$, has length at
most $\pgv(\pgcodd)$. In particular $\log\pgcodd\le x^{800}$ (Remark~\ref{pod:rem-size}).
\end{proposition}
\smallskip\noindent Exact statement: Proposition~\ref{pri:odd}; proved in Proposition~\ref{pod:odd}.\par\smallskip

\begin{proposition}[The local loss]\label{pg:loss}
Let $v\mid\ell$, let $\auO$ be large, and let $\mathcal C_v$ be a class of lattices in potentially
crystalline representations of $D_v$ of the following kind: crystalline with Hodge--Tate weights in
prescribed sets, or, if $\ell=p_i$, potentially crystalline with tame inertial type of the shape
allowed by (G4). Let $L$ be a lattice of $\mathcal C_v$ that is the restriction of a lattice in a
main summand twisted as in Theorem~\ref{pg:III}, or of a character $\chi^e$, $e<t'$, where all
weights of all members of $\mathcal C_v$ are weights of these summands. Then for every finite family $u_x$ of global lattices whose restrictions to $D_v$ lie in $\mathcal C_v$, the module of cocycles $\pgG\to L/\ell^M$ obeying all $\auO$-linear relations among the families $(u_x(g)-1)_x$ has, modulo those obeying the crystalline relations of $L$ at $v$, length at most the \emph{defect} $\delta_v(L)$ (the length of the crystalline cocycles of the Tate dual of $L$ modulo those cup-orthogonal to $\mathcal C_v$; see Proposition~\ref{pri:loss}),
and
\[
\sum_{v\mid\ell}\ \sum_{L}\delta_v(L)\le r'\cdot\pgloc_\ell ,
\]
the inner sum over the $t'+1$ summands. In particular there is no loss if $\ell>2b$ and $\ell\nmid M$ (Theorem~\ref{pll:exponent}(d); for $\ell\ge4b+2$, $\ell\nmid q_1dM$ see also Remark~\ref{pll:FL} for another route).
\end{proposition}
\smallskip\noindent\emph{Proof.} The exact form of the statement is Proposition~\ref{pll:main}. The
restrictions to $D_v$ of the cocycles in question come from $\mathcal C_v$ by Lemma~\ref{pll:ribet}
(with $\Gamma=D_v$ inside $\pgG$; its hypothesis is exactly that they obey the $\auO$-linear relations
among the $(u_x(g)-1)_x$), so the first assertion is Proposition~\ref{pll:dual} with $T=L\in\mathcal C_v$. Its
first inequality is Proposition~\ref{pll:dual}; the second is Theorem~\ref{pll:exponent} with
Lemma~\ref{pll:exp} (Remark~\ref{pll:status}), Theorem~\ref{pll:exponent} being proved in
{\cite[\S2.3]{P56}} from {\cite[Propositions~2.5--2.10]{P56}}, which are proved in
{\cite[\S\S7, 8.1 and~8.2]{P56}}; the sum over places and summands is
Lemma~\ref{pll:sum}. The clause ``no loss if $\ell>2b$ and $\ell\nmid M$'' is
Theorem~\ref{pll:exponent}(d) (Remark~\ref{pll:status}); for $\ell\ge4b+2$, $\ell\nmid q_1dM$ see
also Remark~\ref{pll:FL} for another route.\par\smallskip

\begin{proposition}[The operators at the auxiliary places]\label{pg:norm}
Let $w_j$, $j\in I$, be places as in Proposition~\ref{pg:triangle} off the level, and consider
the forms of principal level at the $w_j$. Let $\mathbb T$ be the algebra of functions on Hecke
eigensystems that preserve the forms integral at $\iota$. Suppose that at each $w_j$ the Hecke
polynomial of every eigensystem of base level congruent to $x_0$ on which a given $\delta\in
\mathbb T$ does not vanish is, after normalisation, $\prod_{i<t'}(X-(\auNm w_j)^i)$ times a divisor
of $H_j$. Then there are integral operators $\mathcal N_S$, $S\subseteq I$, commuting with
$\mathbb T$, with $\mathcal N_\emptyset$ the sum over $\prod_j\auGL_n(\auF_{w_j})$, such that
$\mathcal N=\sum_S\mathcal N_SX_S$ kills the traces of the ideal $\mathfrak a$ of
Proposition~\ref{pg:ribet} modulo $\ell^\nu/(2\,n!)$ and $\delta^{\#I}\mathcal N$ factors through
$\mathcal N_\emptyset$; consequently, if $\mathcal D$ is as in Proposition~\ref{pg:ribet}, the top
coefficient $\mathcal D_I$, times the least common multiple of the orders of the stabilisers and
times the idempotent of the congruence class of $x_0$, maps $\pgpr_{\mathcal Z}(y)$ to a form
integral at $\iota$, for every $y$ as in Definition~\ref{pg:def-den}.
\end{proposition}
\smallskip\noindent Proved in Lemma~\ref{pri:oneplace} and Proposition~\ref{pri:norm} from
Proposition~\ref{pri:dict} and Lemma~\ref{pri:forms}, which are proved in \S\S\ref{pax:ss-dict},
\ref{pay:ss-cons}, \ref{pax:ss-forms}, \ref{pax:ss-pairing}.\par\smallskip

\emph{How the bound $\beta_\ell$ arises.} In the proof of Theorem~\ref{pri:III}, the eigenvalue at
$y_0$ of the operator
that clears the denominator has valuation at most
\[
\underbrace{e'_\ell}_{\text{Selmer}}+\underbrace{2t'\,e'_\ell}_{\text{ramification of }\psi
\text{ at }Q}+\underbrace{r'\,\pgloc_\ell}_{\text{Proposition \ref{pg:loss}}}+\underbrace{e_\ell}_{
\text{odd classes, level, triangle, characters}},
\]
and one more factor $\pgNch$ accounts for the idempotents; with $r'\ge1$ this is at most
$r'\beta_\ell$, and the determinant of the operator on $\mathcal Z$ has valuation at most
$\dim\mathcal Z\cdot r'\beta_\ell$ (the operator is built generator by generator and is diagonal
in a basis of eigenforms; Steps~1--3 of the proof of Theorem~\ref{pri:III}).

\subsubsection*{The $p$-adic Hodge theoretic input}
The local conditions above $\ell$ enter in exactly two places. (1) In the twisting theorem only
through the \emph{rank} of the module of crystalline cocycles (Lemma~\ref{pg:beta}); this is the
formula of Bloch and Kato and needs nothing integral. (2) In Theorem~\ref{pg:III}: the cocycles
modulo $\ell^M$ that Ribet's method extracts from a congruence are the Ribet cocycles of
Proposition~\ref{pg:loss}, and one needs that they obey the crystalline relations of the
summand, up to a bounded defect. If $\ell$ is unramified in $E$, $\ell\nmid M$ and all Hodge--Tate
weights of all representations involved lie in an interval of length at most $\ell-2$, one
expects defect $0$ from the theory of Fontaine and Laffaille~\cite{FontaineLaffaille} together with
\cite[Lemmas~4.4--4.5]{BlochKato}; this route needs an integral statement on lattices that we do
not claim (Remark~\ref{pll:FL}), and the paper does not use it. Since the weights here are bounded by $2b$ in absolute value, this would concern
the primes $\ell\ge4b+2$ that are prime to $q_1dM$. For the remaining primes ($\ell<4b+2$ or $\ell\mid q_1dM$), in particular all $\ell\le2b$, which include $\ell=2$ and the primes inert in $K$ below $x_1$ that carry the prime
sum, no published result gives the bound: the weights (of size $b$ and $a/x_1$, with gaps
$\Delta$) are far outside the Fontaine--Laffaille range, and results on torsion crystalline
representations in large weight (Liu) give a defect bounded in terms of the weights and the
ramification but not the explicit bound $\pgloc_\ell$, whose shape $b/(\ell-1)$ per unit of rank,
times $m/k$, is what the budget $\aubud$ can afford (Lemma~\ref{pg:sums}). The argument is an
analysis of Kisin modules of the lattices involved (duality: the defect is the length of the
crystalline classes of the Tate dual modulo those orthogonal under cup product to the class
$\mathcal C_v$, Proposition~\ref{pll:dual}; then an exponent of $\ell$ that carries crystalline
cocycles into the orthogonal, Lemma~\ref{pll:exp}, computed on Breuil--Kisin modules, using that the
weights are either small or within a window near $\pm\kappa_j$ and that few are small; the factor
$m/k$ comes from the sparseness of $\auPb$). This is Theorem~\ref{pll:exponent}, proved in
{\cite[\S2.3]{P56}}. \emph{The $p$-adic Hodge theoretic input is not replaced by a citation for any $\ell$}; for
$\ell\ge4b+2$, $\ell\nmid q_1dM$, Remark~\ref{pll:FL} describes a possible alternative route.

\subsection{The inputs of \texorpdfstring{{\cite[Proposition~14.40]{P46}}}{[Row D, Proposition~14.40]}}\label{pg:ss-summary}

{\cite[Proposition~14.40]{P46}} follows from Theorem~\ref{pg:main} by the sums over $\ell$ (Lemma~\ref{pg:sums}),
given $\log\pgcodd\le x^{800}$; Theorem~\ref{pg:main} follows from Theorems~\ref{pg:I}--\ref{pg:III} and
{\cite[Proposition~14.36]{P46}}. Theorem~\ref{pg:I} is proved in \S\ref{pga:sec} (with Proposition~\ref{pax:sign} and
Lemma~\ref{pax:conv}) from published results; statements (L), (C) are Propositions~\ref{pga:L}, \ref{pga:C}.
Theorem~\ref{pg:II} uses Propositions~\ref{pse:mirror}, \ref{pse:step}, \ref{pse:A}, \ref{pse:stick} of
\S\ref{pse:sec} and (IM), which is {\cite[Propositions~3.8 and~3.12]{P54}}, besides published results
(\cite{BlochKato}, \cite{NSW}). Theorem~\ref{pg:III} is proved in \S\S\ref{pri:sec}, \ref{prm:sec} from
Proposition~\ref{pg:loss} (\S\ref{pll:sec}, with Theorem~\ref{pll:exponent} proved in {\cite[\S2.3]{P56}}) and
Proposition~\ref{pg:odd} (Proposition~\ref{pod:odd}); see \S\ref{pri:ss-summary}.




\providecommand{\auO}{\mathcal O}
\providecommand{\auNm}{\operatorname{Nm}}
\providecommand{\auU}{\mathrm U}
\providecommand{\auA}{\mathbb A}
\providecommand{\auF}{\mathbb F}
\providecommand{\aubar}[1]{\overline{#1}}
\providecommand{\auf}[1]{\texttt{#1}}
\providecommand{\panot}[1]{\par\smallskip{\raggedright\noindent\textbf{[Not proved here:} #1\textbf{]}\par}\smallskip}
\providecommand{\pgZb}{\overline{\Z}_\ell}
\providecommand{\pgv}{v}
\providecommand{\pgNch}{N_{\mathrm{ch}}}
\providecommand{\pgG}{G_E}
\providecommand{\pgSel}{\operatorname{Sel}}
\providecommand{\pglen}{\operatorname{len}}
\providecommand{\pgHom}{\operatorname{Hom}}
\providecommand{\pgS}{S_{\mathrm{bad}}}
\providecommand{\pgcyc}{\varepsilon_\ell}
\providecommand{\pseinv}{\operatorname{inv}}
\providecommand{\pseev}{\operatorname{ev}}
\providecommand{\pseCl}{\mathrm{Cl}}
\providecommand{\pseArt}{\operatorname{Art}}

\section{The twisting theorem: Selmer groups and \texorpdfstring{$L$}{L}-values}\label{pse:sec}

This section completes the proof of Theorem~\ref{pg:II}. In \S\ref{pg:ss-twisting} that theorem
was deduced from Propositions~\ref{pg:mirror}, \ref{pg:step}, \ref{pg:stick} and statement~(A).
We prove (A) (\S\ref{pse:ss-aux}), Proposition~\ref{pg:mirror} in the form of Proposition~\ref{pse:mirror}
(\S\ref{pse:ss-mirror}), Proposition~\ref{pg:step} (\S\ref{pse:ss-step}), and
Proposition~\ref{pg:stick} in the form of Proposition~\ref{pse:stick} (\S\ref{pse:ss-stick}) \emph{from one input on
theta lifts}, statement~(IM) of \S\ref{pse:ss-stick}, which is proved in {\cite[\S3.1]{P54}}
({\cite[Propositions~3.8 and~3.12]{P54}}). \S\ref{pse:ss-end} collects the choices of constants used in the proof of
Theorem~\ref{pg:II} in \S\ref{pg:ss-twisting}; Theorem~\ref{pg:II} depends on published results and
on~(IM), and on nothing else.

Notation is that of \S\ref{pg:ss-twisting}: $E$ is the CM field, $c$ the fixed complex conjugation
of $\overline{\Q}$, $g^c=cgc$ and $g^{-c}=(g^c)^{-1}$ for $g\in\pgG$; $S_0=\pgS(\ell)$; $w_0\nmid\ell$
is a chosen place and $g_0=\gamma_{w_0}$ its tame generator. Lengths of $\auO$-modules are
normalised as in \S\ref{pg:ss-twisting}. For a finite $\auO$-module $X$ with $\pgG$-action,
$X^*=\pgHom(X,\Q/\Z)(1)$; $g$ acts on $f\in X^*$ by $(gf)(x)=\pgcyc(g)f(g^{-1}x)$.
For a place $B$ of $\overline{\Q}$ over a finite place $P\nmid\ell$ of $E$ we write $I_B\subseteq
D_B$ for the inertia and decomposition groups.

\subsection{Tame generators, and the local pairing at a split prime}\label{pse:ss-tame}

Fix a compatible system $(\xi_m)_{m\ge1}$ of primitive $m$-th roots of unity in $\overline{\Q}$.
Let $P\nmid\ell$ and $B\mid P$. An element $\delta\in I_B$ is a \emph{normalised tame generator for
$m$} (prime to $P$) if for some uniformiser $\varpi$ of $E_P$ and some $m$-th root $\varpi^{1/m}$ in
$\overline{\Q}_B$ one has $\delta(\varpi^{1/m})=\xi_m\varpi^{1/m}$. The condition does not depend on
the uniformiser: for $u\in\mathcal O_{E_P}^\times$ the roots $u^{1/m}$ generate an unramified
extension of $E_P$ ($m$ is prime to $P$), so $\delta$ fixes them and $\delta((u\varpi)^{1/m})/
(u\varpi)^{1/m}=\delta(\varpi^{1/m})/\varpi^{1/m}$. In particular we may take $\varpi\in E$. Normalised tame generators
exist (the tame inertia group $I_B/I_B^{\mathrm{wild}}$ is $\prod_{p\ne q}\Z_p(1)$, $q$ the residue
characteristic, and maps onto $\mu_m$ by $\delta\mapsto\delta(\varpi^{1/m})/\varpi^{1/m}$, independently
of the choice of $\varpi^{1/m}$).

\begin{lemma}\label{pse:conjgen}
If $\delta$ is a normalised tame generator for $m$ at $B$, then $\delta^{-c}$ is one at $cB$, a
place over $\aubar P$.
\end{lemma}
\begin{proof}
$\delta^{-c}=c\delta^{-1}c$ lies in $I_{cB}$. Take the uniformiser $\varpi\in E$ (allowed by the remark after the
definition) and put $\alpha=c(\varpi^{1/m})$, an $m$-th root of the
uniformiser $\aubar\varpi$ of $E_{\aubar P}$. Then $\delta^{-c}(\alpha)=c\,\delta^{-1}(\varpi^{1/m})
=c(\xi_m^{-1}\varpi^{1/m})=\xi_m\alpha$, because $c$ inverts roots of unity.
\end{proof}

\begin{lemma}\label{pse:kpower}
Let $P\nmid\ell$, $B\mid P$, $\gamma\in I_B$ any element, $\delta$ a normalised tame generator for
$\ell^N$ at $B$, and $X$ a finite $\auO$-module killed by $\ell^N$. There is an integer $k$ such that
for every continuous map $z\colon\pgG\to X$ whose restrictions to $I_B$ and to $I_{cB}$ are
homomorphisms one has $z(\gamma)=kz(\delta)$ and $z(\gamma^{-c})=kz(\delta^{-c})$.
\end{lemma}
\begin{proof}
A continuous homomorphism $I_B\to X$ kills wild inertia (a pro-$q$ group, $q\ne\ell$) and factors
through the maximal quotient of exponent $\ell^N$ of tame inertia, which is cyclic of order
$\ell^N$ and generated by the image of $\delta$. Hence $\gamma\equiv\delta^k$ modulo the common
kernel $K_B$ of these homomorphisms, for some $k$. Conjugation by $c$ carries $I_B$, $K_B$ onto
$I_{cB}$, $K_{cB}$, so $\gamma^{c}\equiv(\delta^{c})^k$ modulo $K_{cB}$ and, inverting in the abelian
quotient, $\gamma^{-c}\equiv(\delta^{-c})^k$.
\end{proof}

\begin{lemma}[The pairing at a split prime]\label{pse:pairing}
Let $P\nmid\ell$, $B\mid P$, and let $X$ be a finite $\auO$-module killed by $a\mid\ell^M$ on which
$D_B$ acts trivially, with $\auNm P\equiv1\pmod{\ell^M}$, so that $D_B$ acts trivially on $X^*$ as
well. Let $\phi\in D_B$ be a Frobenius element and $\delta$ a normalised tame generator for
$\ell^M$ at $B$. For a cocycle (here: a homomorphism) $z\colon D_B\to X$ and an unramified
homomorphism $y\colon D_B\to X^*$,
\[
\pseinv_B(z\cup y)=\epsilon\,\bigl\langle z(\delta),y(\phi)\bigr\rangle ,
\]
where $\langle\ ,\ \rangle\colon X\times X^*\to\mu_{\ell^M}\cong\frac1{\ell^M}\Z/\Z$ is the
evaluation (the isomorphism given by $\xi_{\ell^M}\mapsto1/\ell^M$), and $\epsilon\in\{\pm1\}$
depends only on the conventions for cup products and for the local invariant, not on $P$, $B$,
$X$, $z$, $y$.
\end{lemma}
\begin{proof}
Both sides are biadditive in $(z,y)$ and compatible with decompositions of $X$ into cyclic groups;
so we may assume $X=\Z/\ell^s$, $s\le M$, $X^*=\mu_{\ell^s}$. Then $z$ is a character
$\chi\colon D_B\to\Z/\ell^s$, and by Kummer theory $y$ is the class of some $b\in E_P^\times$,
a unit up to $\ell^s$-th powers because $y$ is unramified, with $y(\phi)=\phi(b^{1/\ell^s})/
b^{1/\ell^s}=\bar b^{(\auNm P-1)/\ell^s}$ (residue in $\mu_{\ell^s}$). By local class field theory
$\pseinv(\chi\cup\partial b)=\pm\chi(\pseArt(b))$, and for a unit $b$ the element $\pseArt(b)$ is the
element of tame inertia sending $\varpi^{1/\ell^s}$ to $\bar b^{\mp(\auNm P-1)/\ell^s}
\varpi^{1/\ell^s}$ (the tame symbol); that is, $\pseArt(b)\equiv\delta^{k}$ with $\xi_{\ell^s}^k=
y(\phi)^{\mp1}$. So $\pseinv(z\cup y)=\pm k\chi(\delta)=\pm\langle z(\delta),y(\phi)\rangle$, with a
sign fixed by the conventions.
\end{proof}
\noindent References: local reciprocity and the cup product with a Kummer class, and the tame
symbol: \cite[XIV, \S1, Proposition~3; \S2, Proposition~6; \S3, Proposition~8 and its Corollary]{SerreLF}; \cite[Ch.~I, Remark~2.5(b)]{MilneADT} (the tame symbol).

\subsection{Auxiliary primes: statement (A)}\label{pse:ss-aux}

\begin{proposition}[Statement (A)]\label{pse:A}
Let $N\ge1$, let $S$ be a finite set of finite places of $E$ containing those above $\ell$ and
stable under $c$, let $R$ be a finite set of finite places, let $Z$ be a finite set of finite places
outside $S$ and prime to $\ell$, and let $z$ lie in the normal subgroup of $\pgG$ generated by the
inertia groups $I_B$, $B$ over a place of $Z$. Let $X\subseteq\pgG$ be open with $z\in X$. Then
there are a place $P\notin R\cup\aubar R$ of $E$ of prime norm $q\equiv1\pmod{\ell^N}$ with
$P\ne\aubar P$, a place $B\mid P$, a normalised tame generator $\delta$ for $\ell^N$ at $B$, a
Frobenius element $\phi\in X$ at $B$, and a continuous homomorphism $\lambda\colon\pgG\to\Z/\ell^N$
such that:
\begin{enumerate}
\item[(i)] $\lambda(\pseArt(y))=\theta(y/\aubar y)$ for a character $\theta$ of the finite ideles
trivial on $E^\times$ and on $E_w^\times$ for $w\in S$; hence $\lambda$ is anticyclotomic
($\lambda(g^c)=-\lambda(g)$), unramified outside $P,\aubar P$, and zero on $D_{B'}$ for every
$B'$ over a place of $S$;
\item[(ii)] $\lambda(\delta)=1$ and $\lambda(\phi)=0$;
\item[(iii)] every Frobenius element at $B$ lies in $X$.
\end{enumerate}
\end{proposition}

\begin{proof}
\emph{Step 1: a set $T$.} Choose a finite $c$-stable set $T\supseteq S$ of finite places, disjoint
from $Z\cup\aubar Z$, such that the ideal classes of the places of $T$ generate the class group of
$E$ (possible, since every ideal class contains prime ideals outside any given finite set, by
Chebotarev in the Hilbert class field). Then every finite idele is the product of a principal idele
and of a \emph{$T$-idele}, an idele that is a unit at every place outside $T$. The group
$\mathcal O_{E,T}^\times$ of $T$-units is finitely generated (Dirichlet); let $u_1,\dots,u_k$
generate it.

\emph{Step 2: Chebotarev.} Let $F=E(\mu_{\ell^N},u_1^{1/\ell^N},\dots,u_k^{1/\ell^N})$, a finite
Galois extension, and $U_0=\operatorname{Gal}(\overline{\Q}/F)$. A place $w\in Z$ is prime to
$\ell$ and to every $u_i$ (since $w\notin T$), so $F/E$ is unramified at $w$ and every inertia group
$I_B$ over $w$ lies in the normal subgroup $U_0$; hence $z\in U_0$. Choose an open normal subgroup
$V\subseteq U_0$ of $\pgG$ with $zV\subseteq X$, and let $F'$ be its fixed field. By the Chebotarev
density theorem in $F'/E$, and since places of degree $>1$ over $\Q$ have density $0$, there are
infinitely many places $P$ of $E$ of prime norm, unramified in $F'$, with a place $B\mid P$ whose
Frobenius in $\operatorname{Gal}(F'/E)$ is the image of $z$. Exclude the finitely many $P$ in
$T\cup R\cup\aubar R$ and those ramified in $E/E^+$; a place of prime norm unramified over $E^+$ has
$P\ne\aubar P$ (otherwise $P\cap E^+$ is inert and $\auNm P$ is a square). Every Frobenius element
at $B$ lies in $zV$ (the inertia group $I_B$ lies in $V$, as $P$ is unramified in $F'$); this is
(iii), and every such element lies in $U_0$. Fix one, $\sigma$. As $\sigma$ fixes $\mu_{\ell^N}$,
$q=\auNm P\equiv1\pmod{\ell^N}$; as $\sigma$ fixes $u_i^{1/\ell^N}$ and acts as $x\mapsto x^q$ on the
residue field at $B$, $u_i^{(q-1)/\ell^N}\equiv1\pmod P$: the $T$-units are $\ell^N$-th powers in
$\auF_q^\times$.

\emph{Step 3: the character.} Let $\delta$ be a normalised tame generator for $\ell^N$ at $B$. By
local class field theory $\delta$ restricted to the maximal abelian extension of $E_P$ is the local
Artin symbol of a unit $u_0\in\auO_{E_P}^\times$; by the tame symbol (proof of
Lemma~\ref{pse:pairing}) $\bar u_0^{(q-1)/\ell^N}$ is a primitive $\ell^N$-th root of unity, so $\bar
u_0$ generates $\auF_q^\times/(\auF_q^\times)^{\ell^N}\cong\Z/\ell^N$. Let $\mathrm{lg}\colon
E_P^\times\to\Z/\ell^N$ send $x$ to the class of the residue of $x\varpi^{-v(x)}$ in this quotient,
identified with $\Z/\ell^N$ so that $\bar u_0\mapsto1$; it is continuous, $\mathrm{lg}
(u_0)=1$, and it kills the units $\equiv1\pmod P$ and the $T$-units (which are $\ell^N$-th powers
modulo $P$). Define $\theta$ on finite ideles by $\theta(xa)=\mathrm{lg}(a_P)$ for $x\in E^\times$ and
$a$ a $T$-idele; this is well defined because $E^\times\cap\{T\text{-ideles}\}$ is the group of
$T$-units, and it is continuous, trivial on $E^\times$, on $E_w^\times$ for $w\in T$, and on the
units at places $\ne P$. The character $y\mapsto\theta(y/\aubar y)$ is continuous, of finite order
and trivial on $E^\times$; by the reciprocity law of global class field theory there is a continuous
$\lambda_1\colon\pgG\to\Z/\ell^N$ with $\lambda_1(\pseArt(y))=\theta(y/\aubar y)$. For $y\in
E_w^\times$, $w\in S$, the idele $y/\aubar y$ is supported at $w,\aubar w\in T$, so $\lambda_1$ kills
$D_{B'}$ for $B'$ over $S$. It is unramified outside $P,\aubar P$ because $\theta$ kills the units at
all other places. Since $\pseArt(\aubar y)=c\,\pseArt(y)c$, $\lambda_1$ is anticyclotomic. For the idele
$y=u_0$ at $P$ (and $1$ elsewhere), $y/\aubar y$ has component $u_0$ at $P$ and a unit at $\aubar P\ne
P$, so $\lambda_1(\delta)=\mathrm{lg}(u_0)=1$.

\emph{Step 4: the Frobenius element.} Put $\phi=\sigma\delta^{-\lambda_1(\sigma)}$ (exponent read in
$[0,\ell^N)$). It is a Frobenius element at $B$, so $\phi\in X$ by (iii), and $\lambda_1(\phi)=0$.
Take $\lambda=\lambda_1$.
\end{proof}

\subsubsection*{Admissible places}
We now fix the meaning of ``a place that can be added to $J$'' in Propositions~\ref{pg:mirror},
\ref{pg:step}. Let $T_1,\dots,T_s$ be Selmer data (Definition~\ref{pg:def-datum}) and $M<N$. Let
$U_i\subseteq\pgG$ be the kernel of the action on $T_i/\ell^M$ and on $\mu_{\ell^N}$; it is open and
normal. Let $J$ be a set of tied primes (as in \S\ref{pg:ss-twisting}), with first prime $P_0$; for
$v\in J$ the data are a place $B_v\mid v$, the tame generator $\gamma_v$ (a normalised tame
generator for $\ell^N$ at $B_v$), the Frobenius element $\phi_v$ and $\lambda_v$, the $\lambda_v$
being as in Proposition~\ref{pse:A}(i) for a set $S$ containing the places above $\ell$ and the
places of $J\setminus\{P_0\}$ and their conjugates. A place $P$ with data $(B,\gamma_P,\phi_P,
\lambda_P)$ is \emph{admissible} for $J$ if: $P$, $\aubar P$ are outside $S_0\cup J\cup\aubar J$;
$P$ has prime norm and $P\ne\aubar P$; the data satisfy the conclusions of Proposition~\ref{pse:A}
for $N$ and the set $S=\{w\mid\ell\}\cup(J\setminus\{P_0\})\cup(\aubar J\setminus\{\aubar P_0\})$
(with $\delta=\gamma_P$); and $\lambda_{P_0}(\phi_P)=1$, $\lambda_v(\phi_P)=0$ for $v\in J\setminus
\{P_0\}$. It is \emph{admissible for $T_i$} if moreover $\phi_P,\phi_P^{c}\in U_i$. Then $J\cup
\{P\}$, with slope $0$ for the data for which $P$ is not admissible, and any slopes for the others,
is again a set of tied primes (the condition $\lambda_P(\phi_v)=0$ for $v\in J\setminus\{P_0\}$ holds
because $\lambda_P$ kills $D_{B_v}$).

\begin{lemma}[Admissible places near $\gamma_{P_0}H$]\label{pse:near}
Let $T$ be one of the data, $H\subseteq\pgG$ the subgroup generated by the commutators
$[g,\gamma_*]$ and $[g,\gamma_*^c]$, $g\in\pgG$, where $\gamma_*=\gamma_*(T)$. For every $h\in H$,
every finite set $R$ and every open $X\ni\gamma_{P_0}h$ there is a place $P\notin R\cup\aubar R$
admissible for $T$ with $\phi_P\in X$.
\end{lemma}
\begin{proof}
The element $\gamma_*$ lies in an inertia group at a place $v_*\in S_0$, $v_*\nmid\ell$; so
$[g,\gamma_*]=(g\gamma_*g^{-1})\gamma_*^{-1}$ lies in the normal subgroup generated by inertia at
$v_*$, and $[g,\gamma_*^c]$ in that generated by inertia at $\aubar v_*$. Hence $\gamma_{P_0}h$ lies
in the normal subgroup generated by inertia at the places of $Z=\{P_0,v_*,\aubar v_*\}$, which are
prime to $\ell$ and outside $S=\{w\mid\ell\}\cup(J\setminus\{P_0\})\cup(\aubar J\setminus\{\aubar
P_0\})$ (because $J$, $\aubar J$, $S_0$ are pairwise disjoint). The element $\gamma_{P_0}h$ lies in
the open group $U\cap U^c$, $U=U(T)$: $\gamma_{P_0}$ and $\gamma_{P_0}^c$ are inertial at places
outside $S_0$, where $T$ is unramified and $\mu_{\ell^N}$ too; and $h$, $h^c$ act trivially on $T$
and on $\mu_{\ell^N}$ because both $\gamma_*$ and $\gamma_*^{c}$ act on $T$ by scalars and trivially on
$\mu_{\ell^N}$ (so all commutators with them act trivially). The characters $\lambda_v$, $v\in J$, are
locally constant and kill $H$ (an abelian quotient kills commutators), so on a neighbourhood of
$\gamma_{P_0}h$ they take the values $\lambda_v(\gamma_{P_0})$, which is $1$ for $v=P_0$ and $0$ for
$v\ne P_0$ ($\lambda_v$ is unramified at $P_0$). Apply Proposition~\ref{pse:A} with this $S$ and
$Z$, with $R\cup S_0\cup J$ in place of $R$, and with $X$ replaced by the intersection of $X$, of
$U\cap U^c$ and of that neighbourhood. The place it gives is admissible for $T$.
\end{proof}

\subsection{The dual of a tied Selmer group}\label{pse:ss-mirror}

\subsubsection*{Which dual}
In Proposition~\ref{pg:mirror} and Lemma~\ref{pg:count}, $\pgSel^{\mathrm{tie},*}_i(J;a)$ is the
Selmer group in $H^1(E,(T_i/a)^*)$ with the local conditions orthogonal, under local Tate duality,
to those of the tied group. One may also work with a \emph{module of relations}: equivariant
homomorphisms $y$ on $U_i$ with values in $(T_i/a)^*$, unramified outside a finite set, satisfying
``reciprocity at split primes'' $\sum_jy(\phi_j)(z(\delta_j))=0$ for cocycles $z$ ramified at a
family of split primes. We keep the Tate dual of Lemma~\ref{pg:count}, for which the
Greenberg--Wiles formula is a published theorem, and show (Lemma~\ref{pse:reciprocity}) that
restriction to $U_i$ maps it injectively into the module of relations; the argument below is then
carried out in the module of relations.

Fix $i$, write $T=T_i$, $u=u_i$, $n=n_i$, $\gamma_*$, $\zeta_*$, $U=U_i$, $a=\varpi^k\mid\ell^M$,
$M<N$, $X=T/a$, and $H$ as in Lemma~\ref{pse:near}.

\begin{lemma}[Restriction to $U$]\label{pse:restrict}
Let $y$ be a cocycle representing a class of $\pgSel^{\mathrm{tie},*}(J;a)$. Then $y|_U$ is a
continuous homomorphism $U\to X^*$, independent of the representative, satisfying
$y(gwg^{-1})=g\cdot y(w)$ for $g\in\pgG$, $w\in U$. The class is zero if $y|_U=0$, and $y|_U=0$ if
$y|_H=0$. Moreover $y(\gamma_{P_0})=y(\gamma_{P_0}^c)=0$.
\end{lemma}
\begin{proof}
$U$ acts trivially on $X^*$ (it fixes $X$ and $\mu_{\ell^N}\supseteq\mu_a$), so $y|_U$ is a
homomorphism, coboundaries vanish on $U$, and the cocycle relation gives the equivariance. The image
of $\gamma_*$ in $\pgG/U\subseteq\operatorname{Aut}(T/\ell^M)\times(\Z/\ell^N)^\times$ is central (a
scalar, and $\pgcyc(\gamma_*)=1$ because $\gamma_*$ is inertial at a place prime to $\ell$); it acts on
$X^*$ by $\zeta_*^{-1}$, and $\zeta_*^{-1}-1$ is a unit. By Sah's lemma $H^1(\pgG/U,X^*)=0$, so by
inflation--restriction the class is determined by $y|_U$. If $y|_H=0$, then for $w\in U$,
$y(w)=y([w,\gamma_*])+y(\gamma_*w\gamma_*^{-1})=\gamma_*\cdot y(w)=\zeta_*^{-1}y(w)$ (as $[w,
\gamma_*]\gamma_*w\gamma_*^{-1}=w$), so $(\zeta_*^{-1}-1)y(w)=0$ and $y(w)=0$. Finally $P_0$ has
slope $0$, so the local condition of the tied group at $P_0$ and $\aubar P_0$ is the unramified one;
its orthogonal complement is the unramified condition (local duality at a place prime to $\ell$
where $X$ is unramified); a cocycle unramified at $B_{P_0}$ takes values in $(\gamma-1)X^*=0$ on
$\gamma\in I_{B_{P_0}}$, and likewise at $cB_{P_0}$, which contains $\gamma_{P_0}^c$.
\end{proof}

\begin{lemma}[Reciprocity at an admissible place]\label{pse:reciprocity}
Let $y$ be as in Lemma~\ref{pse:restrict}, and let $P$ be admissible for $T$, with $P,\aubar P$
outside the finite set where $y$ is ramified. Let $z'$ be a cocycle $\pgG\to X$ with the local
conditions and ties of $\pgSel^{\mathrm{tie}}(J;a)$ that may ramify at $P,\aubar P$ (no condition
there). Then
\[
\bigl\langle z'(\gamma_P),y(\phi_P)\bigr\rangle+\bigl\langle z'(\gamma_P^{-c}),y(\phi_P^{c})\bigr
\rangle=0 .
\]
\end{lemma}
\begin{proof}
The cup product $z'\cup y\in H^2(E,\Q/\Z(1))$ has $\sum_w\pseinv_w(z'\cup y)=0$ (reciprocity law of
the Brauer group; \cite[Theorem~8.1.17]{NSW}). Here $S\subseteq S_0$ denotes the ramification set of the datum (Definition~\ref{pg:def-datum}; it contains the places above $\ell$), not the admissibility set above. The places are: those of $S\cup J\cup\aubar J$, where
$z'$ satisfies the condition $L_w$ and $y$ lies in $L_w^\perp$, so the term vanishes; the
archimedean places, which are complex; the places outside $S\cup J\cup\aubar J\cup\{P,\aubar P\}$
where both are unramified and $w\nmid\ell$, where the term vanishes; and $P$, $\aubar P$. At $B$
over $P$ the group $D_B$ acts trivially on $X$ (as $\phi_P\in U$ and $I_B\subseteq U$, $P\notin
S_0$) and $\auNm P\equiv1\pmod{\ell^N}$; by Lemma~\ref{pse:pairing}, $\pseinv_P(z'\cup y)=\epsilon
\langle z'(\delta),y(\phi_P)\rangle$ for a normalised tame generator $\delta$ for $\ell^N$ at $B$ (it is also
normalised for $\ell^M$, by the compatibility of the $\xi_m$). At
$cB$, over $\aubar P$, the element $\phi_P^c$ is a Frobenius element and $\delta^{-c}$ a normalised
tame generator (Lemma~\ref{pse:conjgen}), so $\pseinv_{\aubar P}(z'\cup y)=\epsilon\langle z'(
\delta^{-c}),y(\phi_P^c)\rangle$. Hence $\langle z'(\delta),y(\phi_P)\rangle+\langle z'(\delta^{-c}),
y(\phi_P^c)\rangle=0$. The restrictions of $z'$ to $I_B$ and $I_{cB}$ are homomorphisms ($T$ is
unramified at $P,\aubar P$); by Lemma~\ref{pse:kpower}, applied to $\gamma=\gamma_P$ and $\delta$, $z'(\gamma_P)=k z'(\delta)$ and $z'(\gamma_P^{-c})=kz'(\delta^{-c})$
for one integer $k$. Multiplying by $k$ gives the claim.
\end{proof}

\begin{lemma}[Sah's lemma for the tied group]\label{pse:span}
If $z\in\pgSel^{\mathrm{tie}}(J;a)$ has a value with a unit coordinate, then $z(H)$ spans $X$
over $\auO$.
\end{lemma}
\begin{proof}
As $z(\gamma_*)=0$ and $u(\gamma_*)=\zeta_*$, the cocycle rule gives $z([g,\gamma_*])=(1-\zeta_*)
z(g)$ for all $g$. Since $1-\zeta_*$ is a unit, $z(\pgG)\subseteq\auO z(H)$. The span $Y$ of
$z(\pgG)$ is $\pgG$-stable ($u(g)z(h)=z(gh)-z(g)$), hence stable under the $\auO$-span $M_n(\auO)$ of
the $u(g)$. It contains a vector $x$ not in $\varpi X$; then $M_n(\auO)x=X$. So $\auO z(H)=X$.
\end{proof}

\begin{proposition}[The dual of a tied Selmer group]
\label{pse:mirror}
Let $J$ be a set of tied primes, $a=\varpi^k\mid\ell^M$ with $k\ge1$, $M<N$, and suppose that
$\pgSel^{\mathrm{tie}}_i(J;a)$ contains a cocycle $z$ one of whose values has a unit coordinate.
Suppose that $z$ is \emph{balanced} at every place $P$ admissible for $T_i$: there is a cocycle
$z'$ with the local conditions and ties of $\pgSel^{\mathrm{tie}}_i(J;a)$, allowed to ramify at
$P,\aubar P$, with $z'(\gamma_P)=z(\phi_P)$ and $z'(\gamma_P^{-c})=z(\phi_P^{-c})$. Then
$\pglen\pgSel^{\mathrm{tie},*}_i(J;a)\le\mathrm{gl}(T_i)$.
\end{proposition}

This is the precise form of Proposition~\ref{pg:mirror}: the dual is the Tate dual,
identified through Lemmas~\ref{pse:restrict}, \ref{pse:reciprocity}; ``a place that can be added''
means a place admissible for $T_i$; balance is needed for the one element $z$ only; and $k\ge1$.

\begin{proof}
Let $y$ represent a class of the dual group, and let $R$ be the finite set of places where $y$ is
ramified, enlarged by $S_0\cup J\cup\aubar J$. For $t\in\auO$ and $h,h'\in H$ put
\[
\Psi_t(h,h')=\bigl\langle tz(h),y(h')\bigr\rangle-\bigl\langle tz(h'^c),y(h^c)\bigr\rangle\in\Q/\Z .
\]
Recall $H\subseteq U\cap U^c$ (proof of Lemma~\ref{pse:near}) and $H^c=H$; on $U$, $z$ and $y$ are
homomorphisms. So $\Psi_t$ is biadditive on $H\times H$ and additive in $t$.

\emph{Step 1: $\Psi_t(h,h)=0$.} The tie with slope $0$ at $P_0$ gives $z(\gamma_{P_0})=
z(\gamma_{P_0}^{-c})=0$, hence $z(\gamma_{P_0}^c)=0$; and $y(\gamma_{P_0})=y(\gamma_{P_0}^c)=0$
(Lemma~\ref{pse:restrict}). As $\gamma_{P_0},\gamma_{P_0}^c\in U$, for $g=\gamma_{P_0}h$ we get
$z(g)=z(h)$, $z(g^c)=z(h^c)$, $y(g)=y(h)$, $y(g^c)=y(h^c)$, so $\Psi_t(h,h)=\Psi_t(g,g)$, where we
extend the formula for $\Psi_t(g,g)$ to all $g\in U\cap U^c$. The map $g\mapsto\Psi_t(g,g)$ is
locally constant. By Lemma~\ref{pse:near} there is an admissible $P\notin R\cup\aubar R$ with
$\phi=\phi_P$ so close to $\gamma_{P_0}h$ that $\Psi_t(\phi,\phi)=\Psi_t(\gamma_{P_0}h,
\gamma_{P_0}h)$ and $\phi,\phi^c\in U$. Let $z'$ be the cocycle given by the balance of $z$ at $P$.
Then $z'(\gamma_P)=z(\phi)$ and $z'(\gamma_P^{-c})=z(\phi^{-c})=-u(\phi^{-c})z(\phi^c)=-z(\phi^c)$. So
\[
\Psi_t(\phi,\phi)=\bigl\langle tz'(\gamma_P),y(\phi)\bigr\rangle+\bigl\langle tz'(\gamma_P^{-c}),
y(\phi^c)\bigr\rangle=0
\]
by Lemma~\ref{pse:reciprocity} applied to $tz'$.

\emph{Step 2: $\Psi=0$.} By biadditivity and Step 1 (applied to $hh'$, $h$, $h'$),
$\Psi_t(h,h')+\Psi_t(h',h)=0$. For $h\in H$, $\gamma_*h\gamma_*^{-1}\in H$ ($H$ is normal) and, since
$z(\gamma_*)=0$ and $h$ acts trivially, $z(\gamma_*h\gamma_*^{-1})=\zeta_*z(h)$; by equivariance and
$u(\gamma_*^{-c})=\zeta_*$, $\pgcyc(\gamma_*^c)=1$, $y((\gamma_*h\gamma_*^{-1})^c)=y(h^c)\circ\zeta_*$.
Hence $\Psi_t(\gamma_*h\gamma_*^{-1},h')=\Psi_{\zeta_*t}(h,h')$ and, in the same way,
$\Psi_{\zeta_*t}(h',\gamma_*h\gamma_*^{-1})=\Psi_t(h',h)$. Antisymmetry for the pair
$(\gamma_*h\gamma_*^{-1},h')$ with parameter $\zeta_*s$ gives $\Psi_{\zeta_*^2s}(h,h')+\Psi_s(h',h)=0$;
subtracting antisymmetry for $(h,h')$ with parameter $s$ gives $\Psi_{(\zeta_*^2-1)s}(h,h')=0$. As
$\zeta_*^2-1$ is a unit, $\Psi_t=0$ for all $t$.

\emph{Step 3: the map $\eta$.} Suppose $\sum_ja_jz(h_j)=0$ with $a_j\in\auO$, $h_j\in H$. For every
$h'\in H$ and $t\in\auO$, by $\Psi=0$,
$\sum_ja_j\,y(h_j^c)(tz(h'^c))=\sum_j\langle ta_jz(h_j),y(h')\rangle=0$; as the $z(h'^c)$, $h'\in H$,
span $X$ (Lemma~\ref{pse:span} and $H^c=H$), $\sum_ja_jy(h_j^c)=0$. Hence there is a unique
$\auO$-linear $\eta_y\colon X\to X^*$ with $\eta_y(z(h))=y(h^c)$ for $h\in H$. It is
$\pgG$-equivariant for the action on $X^*$ through $g\mapsto g^c$: for $h\in H$ and $g\in\pgG$,
$u(g)z(h)=z(ghg^{-1})$ (cocycle rule, $h$ and $ghg^{-1}$ in $U$), and $\eta_y(z(ghg^{-1}))=
y(g^ch^cg^{-c})=g^c\cdot y(h^c)$. So $\eta_y\in\pgHom_{\pgG}(X,X^{*c})$.

\emph{Step 4.} The map $y\mapsto\eta_y$ is $\auO$-linear, and injective: if $\eta_y=0$ then $y$
vanishes on $H^c=H$, hence the class of $y$ is $0$ by Lemma~\ref{pse:restrict}. So
$\pglen\pgSel^{\mathrm{tie},*}(J;a)\le\pglen\pgHom_{\pgG}(X,X^{*c})\le\mathrm{gl}(T)$.
\end{proof}

\subsection{The step of the descent}\label{pse:ss-step}

\begin{lemma}[Tilting]\label{pse:tilting}
Let $\auO$ be a discrete valuation ring with residue field $k$, $C$ an $\auO$-module of finite
length, $F=(\auO/I)^m$ with $I\ne0$, and $e_\gamma,e_\phi\colon C\to F$ linear. There is a set
$\mathrm{Bad}\subseteq k$ with at most $m$ elements such that for every $\alpha\in\auO$ with residue
outside $\mathrm{Bad}$,
\[
\pglen\bigl(\ker(\alpha e_\gamma-e_\phi)\bigr)\le\pglen\bigl(\ker e_\gamma\cap e_\phi^{-1}(e_\gamma
C)\bigr).
\]
\end{lemma}
\begin{proof}
Let $C'=\sum_{\alpha\in\auO}\ker(\alpha e_\gamma-e_\phi)\subseteq C$. If $\alpha e_\gamma(x)=
e_\phi(x)$ then $e_\phi(x)=e_\gamma(\alpha x)$ with $\alpha x\in C'$; so $e_\phi(C')\subseteq
e_\gamma(C')$. Let $\Lambda\subseteq\auO^m$ be the preimage of $e_\gamma(C')$; it contains $I\auO^m$,
so it is free of rank $m$; let $\pi\colon\Lambda\to e_\gamma(C')$ be the reduction. As $\Lambda$ is
free there are $\sigma\colon\Lambda\to C'$ with $e_\gamma\sigma=\pi$, and then (since
$e_\phi\sigma(\Lambda)\subseteq e_\gamma(C')=\pi(\Lambda)$) $\mathcal U\colon\Lambda\to\Lambda$ with
$\pi\,\mathcal U=e_\phi\sigma$. Let $\mathrm{Bad}$ be the set of roots in $k$ of the reduction of the
characteristic polynomial of $\mathcal U$. For $\alpha$ with residue outside $\mathrm{Bad}$,
$\det(\alpha-\mathcal U)$ is a unit, $\alpha-\mathcal U$ is bijective, and $(\alpha e_\gamma-e_\phi)
\sigma=\pi(\alpha-\mathcal U)$ maps onto $e_\gamma(C')$. Hence the image of $\alpha e_\gamma-e_\phi$
on $C'$ is at least as long as that of $e_\gamma$, and since
$\pglen C'=\pglen f(C')+\pglen(C'\cap\ker f)$ for every $f$, $\pglen(C'\cap\ker(\alpha e_\gamma-
e_\phi))\le\pglen(C'\cap\ker e_\gamma)$. Finally $\ker(\alpha e_\gamma-e_\phi)\subseteq C'$, and
$C'\cap\ker e_\gamma\subseteq\ker e_\gamma\cap e_\phi^{-1}(e_\gamma C)$ because $e_\phi(C')\subseteq
e_\gamma(C')$.
\end{proof}

\begin{proposition}[One step of the descent; Proposition~\ref{pg:step}]\label{pse:step}
Let $J$ be a set of tied primes for data $T_1,\dots,T_s$ over $\auO$ with uniformiser $\varpi$,
$\varpi^e\mid\ell^M$, $M<N$, and suppose that the residue field of $\auO$ has more than $2n_ie+1$
elements for all $i$. If for some $i$, some $1\le k\le e$ and some place $P$ admissible for $T_i$, an
element of $\pgSel^{\mathrm{tie}}_i(J;\varpi^k)$ is not balanced at $P$, then $P$ can be added to $J$
with a unit slope for $T_i$ and slope $0$ for the other data so that
$\Sigma(J)=\sum_{i'}\sum_{k'=1}^e\pglen\pgSel_{i'}^{\mathrm{tie}}(J;\varpi^{k'})$ decreases.
\end{proposition}
\begin{proof}
\emph{The other data.} With slope $0$ at $P$, a cocycle of the new group for $T_{i'}$, $i'\ne i$,
vanishes at $\gamma_P$ and $\gamma_P^{-c}$; its restrictions to $I_{B}$ and $I_{cB}$ are
homomorphisms killing the generators of their maximal $\ell$-power quotients (Lemma~\ref{pse:kpower}
with $\delta=\gamma_P$, and Lemma~\ref{pse:conjgen}), so it is unramified at $P,\aubar P$ and lies in
the old group. Hence these terms do not increase.

\emph{The datum $T_i$.} For $1\le k'\le e$ let $C_{k'}$ be the module of cocycles modulo
$\varpi^{k'}$ with the local conditions and ties of $\pgSel^{\mathrm{tie}}_i(J;\varpi^{k'})$, allowed
to ramify at $P,\aubar P$; it is finite. Let $e_\gamma,e_\phi\colon C_{k'}\to(T_i/\varpi^{k'})^2$ be
evaluation at $(\gamma_P,\gamma_P^{-c})$ and at $(\phi_P,\phi_P^{-c})$. As just shown, $C_{k'}\cap\ker
e_\gamma=\pgSel_i^{\mathrm{tie}}(J;\varpi^{k'})$, and $\ker e_\gamma\cap e_\phi^{-1}(e_\gamma C_{k'})$
is its submodule $B_{k'}$ of elements balanced at $P$. With slope $s=\alpha^{-1}$, $\alpha$ a unit,
the new group is $C_{k'}\cap\ker(\alpha e_\gamma-e_\phi)$ (the tie $z(\gamma_P)=sz(\phi_P)$, and the
same at $\aubar P$). By Lemma~\ref{pse:tilting} ($m=2n_i$), for $\alpha$ with residue outside a set
of at most $2n_i$ elements its length is at most $\pglen B_{k'}$. The union of these sets over
$k'\le e$, together with $0$, has at most $2n_ie+1$ elements, so some unit $\alpha$ avoids all. For
it, the new group at level $k'$ has length at most $\pglen B_{k'}\le\pglen\pgSel_i^{\mathrm{tie}}
(J;\varpi^{k'})$, and at level $k$ strictly less, since $B_k$ is a proper submodule of a module of
finite length. The enlarged set is a set of tied primes (\S\ref{pse:ss-aux}, admissible places),
the condition $u_i(\phi_P)\equiv u_i(\phi_P^{-c})\equiv1\pmod{\ell^M}$ for the nonzero slope holding
because $\phi_P,\phi_P^c\in U_i$.
\end{proof}

\begin{remark}\label{pse:rem-poly}
In the proof of Proposition~\ref{pg:poly}, ``balanced at every admissible $P$'' is to be read with
the present meaning, and the application of Proposition~\ref{pg:mirror} is replaced by
Proposition~\ref{pse:mirror}: at the end of the descent every element of every tied group at every
level $\varpi^{k'}$, $k'\le e$, is balanced at every place admissible for its datum, in particular the
element with a unit coordinate. The first prime $P_0$ and its character are given by
Proposition~\ref{pse:A} with $S=\{w\mid\ell\}$, $Z=\{w_0\}$, $z=g_0$, $R=S_0$, and $X$ the open set
of $g$ acting as $g_0$ on all $T_i/\ell^M$, with $g^{-c}$ acting as $g_0^{-c}$, and (for
\S\ref{pse:ss-stick}) with $\vartheta_Y(g)\equiv\vartheta_Y(g_0)\pmod{\ell^M}$ for the finitely many
systems $Y$.
\end{remark}

\subsection{The systems of the \texorpdfstring{$L$}{L}-values}\label{pse:ss-stick}

\subsubsection*{The input from theta lifts}
A \emph{tame character of modulus $Q$} is a character $\vartheta$ of $\auU(1)(E^+)\backslash
\auU(1)(\auA_{E^+})$ of weight $0$, trivial at the infinite places, tamely ramified exactly on a
subset of $Q$ and unramified elsewhere, where $Q$ is a finite set of places of $E$ prime to $\ell p_0p_1$ and
to the level, i.e.\ outside $S_0$, with $Q\cap\aubar Q=\emptyset$; it is a twist of modulus $Q$ when it is ramified at every
place of $Q$. For a ramified critical pair $(\nu,z)$ let $r(\vartheta)=|L(z,(\nu\vartheta)^{-1})|/
|L(z,\nu^{-1})|$ for twists $\vartheta$, with the conventions on omitted Euler factors of
Theorem~\ref{pg:main}(b).

\noindent\emph{Statement (IM) (integral measures).} For every ramified critical pair $(\nu,z)$ there
are a finite set $I\subseteq\Z_{\ge0}$ with at most two elements ($I=\{z\}$ if $z\ge2$, $I=\{0,1\}$ if
$z=1$), an integer $m\in\{0,1\}$ ($m=0$ if $z\ge2$, $m=1$ if $z=1$) and a function $F$ on tame
characters such that: (a) $F(1)\ne0$; (b) for every modulus $Q$ there are finitely many
$x_a\in\auU(1)(\auA_{E^+,f})$ and $c_a\in\C$, with $c_a$ integral at every prime not dividing
$\auNm Q$ (under every field embedding into $\C_p$), such that $F(\vartheta)=\sum_ac_a\vartheta(x_a)$
for every tame $\vartheta$ of modulus $Q$; (c) if $v,\aubar v\notin\mathsf M\cup Q\cup\aubar Q$ ($\mathsf M$ the set of chosen places) and
$\vartheta'$ is $\vartheta$ regarded as of modulus $Q\cup\{v\}$, then $F(\vartheta')=F(\vartheta)
\prod_{e\in I}\prod_\pm\bigl(1-\alpha_v^{\pm1}(\auNm v)^{-e}\bigr)$, $\alpha_v$ the value at a
uniformiser of the base change to $E$ of $\nu\vartheta$ at $v$; (d) for every twist $\vartheta$ of
modulus $Q$, $F(\vartheta)=F(1)\,(\auNm Q)^m\,r(\vartheta)^{m+1}$. For $z\ge2$, $F$ is an integral
multiple of the squared norm of the theta lift of $\nu\vartheta$ from $\auU(1)$ to $\auU(2z)$; for
$z=1$, a theta period on $\auU(2)$ times its conjugate.

\smallskip\noindent (IM) is proved in {\cite[Proposition~3.8]{P54}} for $z\ge2$ (at every prime $\ell\nmid
\auNm Q$; $\ell=d$ by {\cite[Lemma~3.6]{P54}}) and in {\cite[Proposition~3.12]{P54}} for $z=1$ (from the
exchange identity, {\cite[Theorem~3.2]{P53}}). The ratio $\rho_p(\vartheta)$ there is $r(\vartheta)$, with
the conventions on omitted Euler factors of Theorem~\ref{pg:main}(b).\par\smallskip

\subsubsection*{From (IM) to a system}
Let $\pseCl_Q$ be the group of finite ideles $x$ with $x\aubar x=1$, modulo $\auU(1)(E^+)$ and the
units $\equiv1$ modulo the places of $Q$. By Hilbert's Theorem~90 for ideles every class is
$y/\aubar y$ for a finite idele $y$. A continuous character $\psi\colon\pgG\to\overline{\Q}_\ell^\times$
is a \emph{character for $Q$} if $\psi(\pseArt(y))=\eta(y/\aubar y)$ for a character $\eta$ of
$\pseCl_Q$ trivial at the places above $\ell$; then $\psi$ is read on $\pseCl_Q$ by
$\psi(y/\aubar y):=\psi(\pseArt(y))$, and extended linearly to the group ring $\pgZb[\pseCl_Q]$. Every
character for $Q$ of finite order is $\psi_\vartheta$ for a tame character $\vartheta$ of modulus $Q$
(\S\ref{pg:ss-reduction}), with $\psi_\vartheta(y/\aubar y)=
\iota\vartheta^{-1}(y/\aubar y)$.

Fix a ramified critical pair $(\nu,z)$, a sign $\delta=\pm1$, and $F$, $I$, $m$ from (IM). Let
$\vartheta_Y\colon\pgG\to\pgZb^\times$ be the continuous character attached through $\iota$ and class
field theory to the algebraic Hecke character $y\mapsto\nu^{-\delta}(y/\aubar y)$ of $E$. It is
unramified outside $S_0$, and its value at a Frobenius element $\phi$ at a place $v\notin S_0$ times
$\psi_\vartheta(\phi)$ is $\iota$ of $\alpha_v^{\pm\delta}$ for the pair $\nu\vartheta^{\delta}$ (one sign for
all $v$), by the definitions, because $\nu^{-\delta}\vartheta^{-1}=(\nu\vartheta^{\delta})^{-\delta}$. At $g_0=\gamma_{w_0}$, $\vartheta_Y(g_0)$ is the value of the tame part
of $\nu^{\mp\delta}$ at a generator of $\auF_{w_0}^\times$, a root of unity $\ne1$ of order dividing
$\auNm w_0-1$, because the pair is ramified at the chosen place $w_0$. Hence
\begin{equation}\label{pse:eq-g0}
\pgv(\vartheta_Y(g_0)-1)\le\pgv(\auNm w_0-1)
\end{equation}
($\prod_{\xi^n=1,\xi\ne1}(1-\xi)=n$ and each factor is integral).

For $Q$ auxiliary (places of prime norm, $Q$, $\aubar Q$, $S_0$ pairwise disjoint) take $x_a$,
$c_a$ from (IM)(b); $\iota(c_a)\in\pgZb$ since $\ell\nmid\auNm Q$. Put $\Theta_Q=\sum_a\iota(c_a)
[x_a^{-\delta}]\in\pgZb[\pseCl_Q]$, and $\Theta_\emptyset$ likewise. Then, for every character
$\psi=\psi_\vartheta$ for $Q$ of finite order,
\[
\psi(\Theta_Q)=\sum_a\iota(c_a)\,\iota\vartheta^{\delta}(x_a)=\iota F(\vartheta^\delta),
\]
and, by (IM)(c), for $v\notin Q$ with $Q\cup\{v\}$ auxiliary (so that $v,\aubar v\notin\mathsf M\cup
Q\cup\aubar Q$, the set $\mathsf M$ of chosen places being contained in $S_0$) and $\phi$ a Frobenius element at a
place over $v$,
\begin{equation}\label{pse:eq-norm}
\psi(\Theta_{Q\cup\{v\}})=\psi(\Theta_Q)\cdot\prod_{e\in I}\prod_\pm\bigl(1-(\vartheta_Y(\phi)
\psi(\phi))^{\pm1}(\auNm v)^{-e}\bigr),
\end{equation}
the product being symmetric under inversion of $\vartheta_Y(\phi)\psi(\phi)$, so that the sign
conventions do not matter. Let $v_0=\pgv(\iota F(1))=\pgv(1(\Theta_\emptyset))$, finite by (IM)(a),
and $\beta_Y=v_0+2\#I\,\pgv(\vartheta_Y(g_0)-1)\le v_0+4\pgv(\auNm w_0-1)$. By (IM)(d), for a twist
$\vartheta$ of auxiliary modulus $Q$ (so that $\auNm Q$ is an $\ell$-adic unit),
\begin{equation}\label{pse:eq-ratio}
\pgv\bigl(\psi_\vartheta(\Theta_Q)\bigr)=v_0+(m+1)\,\pgv\bigl(\iota r(\vartheta^\delta)\bigr),
\quad\text{hence}\quad \pgv\bigl(\iota r(\vartheta^{\delta})\bigr)\le\pgv\bigl(\psi_\vartheta(\Theta_Q)
\bigr)-v_0\ \text{ if }\pgv(\iota r(\vartheta^\delta))\ge0 .
\end{equation}

\subsubsection*{The polynomial and its line}
Let $J$ be a set of tied primes with first prime $P_0$ chosen as in Remark~\ref{pse:rem-poly}, so
that $\vartheta_Y(\phi_{P_0})\equiv\vartheta_Y(g_0)\pmod{\ell^M}$, and let $\lambda_v$, $v\in J$, be
its characters. Put
\[
P_Y=\sum_a\iota(c_a)\prod_{v\in J}(1+X_v)^{\lambda_v(x_a^{-\delta})}\in\pgZb[X_v:v\in J],
\]
the exponents read in $[0,\ell^N)$ and $\lambda_v$ read on $\pseCl_J$ as above (each $\lambda_v$ is
of the form of Proposition~\ref{pse:A}(i), hence a character for $J$). For $\ell^N$-th roots of unity
$\zeta_v$ and $\psi=\prod_v\zeta_v^{\lambda_v}$ one has $P_Y((\zeta_v-1)_v)=\psi(\Theta_J)$.

\begin{lemma}[Euler factors as polynomials]\label{pse:eulerpoly}
Let $\alpha$ and $q$ be $\ell$-adic units, $I$ a finite set, $\rho\ge0$. There is $R\in\pgZb[X]$
with Gauss valuation $\pgv_\rho(R)\le2\#I\rho$ and $R(\zeta-1)=\zeta^{\#I}\prod_{e\in I}\prod_\pm
(1-(\alpha\zeta)^{\pm1}q^{-e})$ for every $\zeta\in\pgZb^\times$.
\end{lemma}
\begin{proof}
For one $e$, with $b=q^{-e}$: $\zeta(1-\alpha b\zeta)(1-\alpha^{-1}b\zeta^{-1})=(1-\alpha b(1+X))
(1+X-\alpha^{-1}b)$ at $X=\zeta-1$. Each factor is linear with unit coefficient of $X$, so has
$\pgv_\rho\le\rho$. Take the product over $e$; $\pgv_\rho$ is additive.
\end{proof}

\begin{lemma}\label{pse:line}
Let $\rho\le1$ be rational and suppose $\beta_Y+1<M\le N\rho$, $s_0\rho<1$ and $4\#J\le s_0$. Then there is a
point $x=(x_v)$ with $x_{P_0}=y$, $\pgv(y)=\rho$, and $x_v=0$ for $v\ne P_0$, such that
$\pgv(P_Y(x))\le\beta_Y+s_0\rho$.
\end{lemma}
\begin{proof}
Write $h=\lambda_{P_0}$, $L(X)=P_Y(X,0,\dots,0)=\sum_a\iota(c_a)(1+X)^{h(x_a^{-\delta})}$ and
$R_0(X)=\sum_a\iota(c_a)(1+X)^{h(x_a^{-\delta})}$ computed with the elements of $\Theta_{\{P_0\}}$
instead of $\Theta_J$. For an $\ell^N$-th root of unity $\zeta$ the character $\psi=\zeta^{h}$ is a
character for $\{P_0\}$ (hence for $J$), with $\psi(\phi_v)=\zeta$ for $v\in J\setminus\{P_0\}$
(tied primes: $\lambda_{P_0}(\phi_v)=1$). By \eqref{pse:eq-norm} applied to the places of
$J\setminus\{P_0\}$ one after another, $L(\zeta-1)=R_0(\zeta-1)\prod_{v\ne P_0}\prod_{e,\pm}(1-
(\vartheta_Y(\phi_v)\zeta)^{\pm1}(\auNm v)^{-e})$. Let $R_v$ be the polynomial of
Lemma~\ref{pse:eulerpoly} for $\alpha=\vartheta_Y(\phi_v)$, $q=\auNm v$, and $R_1=\prod_{v\ne P_0}
R_v$, so $\pgv_\rho(R_1)\le2\#I\#J\rho\le s_0\rho$. Then $(1+X)^{j}L-R_0R_1$, $j=\#I(\#J-1)$,
vanishes at the $\ell^N$ distinct points $\zeta-1$; as $(1+X)^{\ell^N}-1$ is monic with these simple
roots, it divides $(1+X)^jL-R_0R_1$ in $\pgZb[X]$.

$\pgv_\rho(R_0)\le\pgv(R_0(0))=\pgv(1(\Theta_{\{P_0\}}))$, and by \eqref{pse:eq-norm} with $Q=
\emptyset$, $\psi=1$, $1(\Theta_{\{P_0\}})=1(\Theta_\emptyset)\prod_{e,\pm}(1-\vartheta_Y(\phi_{P_0})^{
\pm1}(\auNm P_0)^{-e})$. Since $\auNm P_0\equiv1\pmod{\ell^N}$, $N\ge M$, and $\vartheta_Y(
\phi_{P_0})\equiv\vartheta_Y(g_0)\pmod{\ell^M}$, each factor is $\equiv1-\vartheta_Y(g_0)^{\pm1}
\pmod{\ell^M}$, whose valuation $\pgv(\vartheta_Y(g_0)-1)\le\beta_Y<M$; so each factor has that
valuation and $\pgv_\rho(R_0)\le\beta_Y$. Thus $\pgv_\rho(R_0R_1)\le\beta_Y+s_0\rho<\beta_Y+1<
M\le N\rho$.

Choose $y\in\pgZb$ with $\pgv(y)=\rho$ and $\pgv((R_0R_1)(y))=\pgv_\rho(R_0R_1)$: write $y=\pi w$
with $\pgv(\pi)=\rho$; then $\pgv((R_0R_1)(\pi w))=\pgv_\rho(R_0R_1)$ unless the residue of the unit
$w$ is a root of a nonzero polynomial over the (infinite) residue field. Now $\pgv((1+y)^{\ell^N}-1)
\ge N\rho$ (each term $\binom{\ell^N}{i}y^i$, $i=\ell^bi'$, has valuation $\ge N-b+\ell^b\rho\ge
N\rho$ as $\rho\le1$), which exceeds $\pgv((R_0R_1)(y))$; by the divisibility,
$\pgv((1+y)^jL(y))=\pgv((R_0R_1)(y))$, and $1+y$ is a unit. So $\pgv(P_Y(x))=\pgv(L(y))\le
\beta_Y+s_0\rho$.
\end{proof}

\begin{proposition}[Valuations of the theta polynomials]\label{pse:stick}
Assume (IM). In the situation of Proposition~\ref{pg:poly}, let finitely many pairs of a ramified
critical pair $(\nu,z)$ and a sign $\delta$ be given, with their systems $Y$, and suppose moreover
$M>\beta_Y+1$ for each. Then, with $s_0$ replaced by $\max(s_0,4(1+b_0))$, for every rational
$0<\rho<1/s_0$ and every large $N$ the set $J$ of Proposition~\ref{pg:poly} can be chosen so that
moreover, for each $Y$, the polynomial $P_Y\in\pgZb[X_v:v\in J]$ and a point $x_Y$ with
coordinates of valuation at least $\rho$ satisfy
\[
\pgv(P_Y(x_Y))\le v_0(Y)+4\pgv(\auNm w_0-1)+s_0\rho ,
\]
and, for all $\ell^N$-th roots of unity $\zeta_v\ne1$ with $\psi=\prod_v\zeta_v^{\lambda_v}=
\psi_\vartheta$,
\[
\pgv\bigl(\iota\,r(\vartheta^{\delta})\bigr)\le\max\bigl(0,\ \pgv(P_Y((\zeta_v-1)_v))-v_0(Y)\bigr).
\]
\end{proposition}
\begin{proof}
Choose $P_0$ as in Remark~\ref{pse:rem-poly}; the descent of Proposition~\ref{pg:poly} does not
change $P_0$, and adds at most $b_0$ primes, so $\#J\le1+b_0$ and $4\#J\le s_0$. The first bound is
Lemma~\ref{pse:line} with $\beta_Y\le v_0+4\pgv(\auNm w_0-1)$ (by \eqref{pse:eq-g0}; $M\le N\rho$
for large $N$). For the second: $\vartheta$ is a twist of modulus $J$ (each $\zeta_v\ne1$ and
$\psi(\gamma_v)=\zeta_v$), $J$ is auxiliary, and $P_Y((\zeta_v-1))=\psi(\Theta_J)$; apply
\eqref{pse:eq-ratio}.
\end{proof}

\subsection{Constants in the proof of the twisting theorem}\label{pse:ss-end}

\begin{remark}\label{pse:rem-constants}
The proof of Theorem~\ref{pg:II} in \S\ref{pg:ss-twisting} uses the results of this section with the following choices. (1) $M$ is chosen with
$M>\max(c,\max_Y\beta_Y)+1$, where $Y$ runs over the systems of the finitely many pairs and signs;
the residue field of $\auO$ is enlarged accordingly. (2) $s_0$ is as in Proposition~\ref{pse:stick}.
(3) In Proposition~\ref{pg:roots} the bound is $\beta=s(c+1)+\sum_Y(v_0(Y)+4\pgv(\auNm w_0-1)+1)$,
which is finite; $h_0$ depends on it, which is allowed. (4) In (c): by Proposition~\ref{pg:roots},
$\pgv(P_Y((\zeta_v^j-1)_v))\le\pgv(P_Y(x_Y))+\varepsilon/2\le v_0(Y)+4\pgv(\auNm w_0-1)+\varepsilon$,
so by Proposition~\ref{pse:stick} the ratio at $\vartheta=\tau^j$ has valuation at most
$4\pgv(\auNm w_0-1)+\varepsilon$. As $\auNm w_0-1$ divides $\#\mathrm{GL}_{rt}(\auF_{w_0})$, which
divides $\pgNch$, this is at most $4\pgv(\pgNch)+\varepsilon\le(4f+3)r\pgv(\pgNch)+\varepsilon=
e'_\ell$. (5) Proposition~\ref{pg:mirror} is replaced by Proposition~\ref{pse:mirror} and
Proposition~\ref{pg:step} by Proposition~\ref{pse:step} (Remark~\ref{pse:rem-poly}); statement (A)
is Proposition~\ref{pse:A}.
\end{remark}





\providecommand{\auO}{\mathcal O}
\providecommand{\auQlb}{\overline{\Q}_\ell}
\providecommand{\auGL}{\mathrm{GL}}
\providecommand{\auF}{\mathbb F}
\providecommand{\auU}{\mathrm U}
\providecommand{\auNm}{\operatorname{Nm}}
\providecommand{\auf}[1]{\texttt{#1}}
\providecommand{\panot}[1]{\par\smallskip{\raggedright\noindent\textbf{[Not proved here:} #1\textbf{]}\par}\smallskip}
\providecommand{\pgZb}{\overline{\Z}_\ell}
\providecommand{\pgv}{v}
\providecommand{\pgNch}{N_{\mathrm{ch}}}
\providecommand{\pgcodd}{c_{\mathrm{odd}}}
\providecommand{\pgloc}{\mathrm{loc}}
\providecommand{\pgG}{G_E}
\providecommand{\pgSel}{\operatorname{Sel}}
\providecommand{\pglen}{\operatorname{len}}
\providecommand{\pgtr}{\operatorname{tr}}
\providecommand{\pgFrob}{\operatorname{Frob}}
\providecommand{\pgpr}{\operatorname{pr}}
\providecommand{\pgcyc}{\varepsilon_\ell}
\providecommand{\pgS}{S_{\mathrm{bad}}}
\providecommand{\priFitt}{\operatorname{Fitt}}
\providecommand{\prilen}{\operatorname{len}}
\providecommand{\prirl}{\operatorname{rl}}
\providecommand{\prival}{\operatorname{val}}
\providecommand{\priAug}{\mathfrak I}
\providecommand{\priT}{\mathbb T}
\providecommand{\priX}{\mathcal X}
\providecommand{\priJ}{\mathcal J}
\providecommand{\prid}{\mathcal D}
\providecommand{\priN}{\mathcal N}
\providecommand{\priZ}{\mathcal Z}
\providecommand{\prics}{\mathrm{cs}}
\providecommand{\prity}{\mathrm{ty}}

\section{From Selmer groups to denominators}\label{pri:sec}

This section proves Theorem~\ref{pg:III} of \S\ref{pg:ss-ribet}, the passage from bounds for Selmer
groups to bounds for the denominators of the projection to a space of lifts, \emph{except} for the
local loss above $\ell$ (Proposition~\ref{pg:loss}), which is used in the exact form stated in
Proposition~\ref{pri:loss} below. The inputs listed in \S\ref{pri:ss-inputs} are proved in
\S\S\ref{pod:sec}, \ref{pku:sec}, \ref{pax:sec}, \ref{pay:sec}, \ref{prm:sec}. The
four propositions of \S\ref{pg:ss-ribet} that are assigned here are restated in their exact form:
Ribet's method to first order (Proposition~\ref{pri:ribet}, replacing~\ref{pg:ribet}), the
Taylor--Wiles places in a triangle (Proposition~\ref{pri:triangle}, replacing~\ref{pg:triangle}),
the odd classes (Proposition~\ref{pri:odd}, replacing~\ref{pg:odd}) and the operators at the
auxiliary places (Proposition~\ref{pri:norm}, replacing~\ref{pg:norm}). The inputs of
Theorem~\ref{pg:III} are listed in \S\ref{pri:ss-summary}.

\emph{Conventions.} Those of \S\ref{pg:sec}. In addition: $\auO$ is a discrete valuation ring with
uniformiser $\varpi$, residue field $k$, fraction field $K$ and additive valuation $\prival=
\prival_\auO$ with $\prival(\varpi)=1$; for an $\auO$-module $M$ of finite length, $\prilen M$ is
its (unnormalised) length, so that $\prilen\auO/a=\prival(a)$. When $\auO$ is a coefficient ring
in the sense of \S\ref{pg:sec}, with ramification index $e_\auO$ over $\Z_\ell$, the normalised
length of \S\ref{pg:sec} is $\prilen/e_\auO$ and $\pgv=\prival/e_\auO$. For submodules $S,T$ of a
module, $\prirl(S,T)=\prilen\bigl(S/(S\cap T)\bigr)$ (possibly $\infty$). For a group $G$ and a
ring $A$, $A[G]$ is the group algebra and $\priAug_A\subseteq A[G]$ its augmentation ideal; for a
finite set $J$, $A[X_J]=A[X_j:j\in J]$, $X_S=\prod_{j\in S}X_j$ for $S\subseteq J$, and the
\emph{top coefficient} of $f\in A[X_J]$ is its coefficient of $X_J$. For $g\in\pgG$,
$g^c=cgc$ and $g^{-c}=(g^c)^{-1}$, as in \S\ref{pg:sec}.

\subsection{The exact statement}\label{pri:ss-statement}

Fix $\ell$, $\iota$ and a twist $\vartheta$ with $\psi=\psi_\vartheta$ and $Q=Q_\vartheta$. Put
$N=\pgNch$ (\S\ref{pg:ss-statement}), $h=r^8$, and let $e'_\ell$, $e_\ell$, $\pgloc_\ell$,
$\beta_\ell$ be as in~\eqref{pg:eq-beta}.

\begin{definition}[Auxiliary twists]\label{pri:def-aux}
The twist $\vartheta$ is \emph{auxiliary with ramification at most $\mathrm{ram}$} if the places
of $Q$ have prime norm, $Q$, $\overline Q$ and $\pgS(\ell)$ are pairwise disjoint, $\psi$ is
anticyclotomic, unramified outside $Q\cup\overline Q$ and trivial on $D_v$ for $v\mid\ell$, and
$\sum_{w\in Q}\pgv(\psi(\gamma_w)-1)\le\mathrm{ram}$ for some choice of tame generators
$\gamma_w$. This is Theorem~\ref{pg:II}(a) with $\mathrm{ram}=e'_\ell$.
\end{definition}

Let $O$ be a pair (\S\ref{pg:ss-statement}), $r'=n_W$, $n=n_V$, $t'=t(O)=n-r'$. A family $c$ of
labels at the chosen places is \emph{in lift setting for $\vartheta$} if: the tame part of the
character of $O$ is $\chi_M^{\,j}$ times that of $\vartheta$ for some $1\le j\le m$; $r'$ and $t'$
are odd, $r'\le r$, $t'\le t$; each $c_w$ is a label of degree $r'$; above each of $p_0,p_1$ there
is a chosen $w$ with $r'=r$ and $c_w$ primitive, or $r'=1$ and $c_w=X-1$; and no rescaled label
has the root $1$. Let $e_c$ be the idempotent of $c$.

\begin{definition}[Lift generators]\label{pri:def-gen}
Let $\mathcal Z\subseteq\mathcal F(O)$ and $e$ a central idempotent. A form $y$ is a \emph{lift
generator of $\mathcal Z$ with Selmer bound $\mathrm{sel}\in[0,\infty]$} if $y\in\mathcal Z$; $y\in
e_c\mathcal F(O)\subseteq e\,\mathcal F(O)$ for a family $c$ in lift setting for $\vartheta$; $y$
has the Hecke polynomials of a lift of a system $q$; and there are a main summand $R$ of rank
$r'$ and string length $t'$ (Definition~\ref{pg:def-summand}) satisfying~\eqref{pg:eq-trace} for
every eigenform with the Hecke polynomials of a lift of $q$, an element $\gamma\in\pgG$ and
$\zeta\in\pgZb$, $\zeta-1$ a unit, such that for all large $\auO$ some lattice $u$ in $R$ has
$\psi(\gamma)u(\gamma)=\zeta$ and $\pglen\pgSel_Q(\psi u,\gamma;K/\auO)\le\mathrm{sel}$.
\end{definition}

\begin{theorem}[From Selmer groups to denominators, exact form]\label{pri:III}
Let $\vartheta$ be auxiliary with ramification at most $e'_\ell$, let $O$ be a pair with central
idempotent $e$, and let $\mathcal Z\subseteq\mathcal F(O)$ be spanned by lift generators of
$\mathcal Z$ with Selmer bound $e'_\ell$. Then there is $\Lambda\in\operatorname{End}(\mathcal Z)$
with
\[
\pgv(\det\Lambda)\le\dim\mathcal Z\cdot r'\cdot\beta_\ell
\]
such that for every $y\in\mathcal F(O)$ integral at $\iota$ and orthogonal to the false lifts
(Definition~\ref{pg:def-den}(a), for $e$ and $\mathcal Z$), $\Lambda\pgpr_{\mathcal Z}(y)$ is
integral at $\iota$ and $\langle y',\Lambda\pgpr_{\mathcal Z}(y)\rangle\in\pgZb$ for every $y'$
integral at every isomorphism $\C\to\auQlb$.
\end{theorem}

\begin{remark}[Comparison with Theorem~\ref{pg:III}]\label{pri:rem-compare}
(1) The conclusion is Definition~\ref{pg:def-den}(b) with $\beta=\beta_\ell$ and $n_W=r'$, i.e.\
the conclusion of Theorem~\ref{pg:III}. (2) The hypotheses of Theorem~\ref{pg:III} give those of
Theorem~\ref{pri:III}: Theorem~\ref{pg:II}(a) is Definition~\ref{pri:def-aux} with
$\mathrm{ram}=e'_\ell$, and every spanning eigenform $y$ of Theorem~\ref{pg:III} is a lift
generator with Selmer bound $e'_\ell$ (Theorem~\ref{pg:II}(b)), \emph{provided} its family $c$ is in
lift setting. Theorem~\ref{pg:III} does not mention two clauses of the lift setting: the tame part
of the character, which holds for $O_j(\vartheta)$ and $O_{\nu,z}(\vartheta)$ by construction, and
the parity of $r'$ and $t'$. For $O_{\nu,z}$, $r'=1$ and $t'=2z-1$ are odd; for $O_j$, $r'=r$ and
$t'=t=5r$ are odd if and only if $r$ is odd, and $r$ is odd by (S7) of {\cite[Lemma~13.4]{P46}}.
Theorem~\ref{pg:III} states the parity clause explicitly. With these additions the deduction of
Theorem~\ref{pg:main} in \S\ref{pg:ss-reduction} is unchanged.
\end{remark}

\subsection{What is used from outside this section}\label{pri:ss-inputs}

Besides published results named where they are used, the proof uses the following statements. They
are not proved in this section; each is followed by a pointer to the place where it is proved.

\begin{proposition}[The local loss, form used]\label{pri:loss}
Let $\auO$ be a large coefficient ring, $u_x\colon\pgG\to\auGL_n(\auO)$ ($x\in Y$, finite)
continuous, $\rho\colon\pgG\to\auGL_{d}(\auO)$, $M\ge1$. Let $Z^{\mathrm R}_M(\rho)$ be the module of
cocycles $z\colon\pgG\to(\auO/\ell^M)^d$ for $\rho$ satisfying $\sum_ia_iz(g_i)=0$ whenever
$a_i\in\auO$ and $\sum_ia_i(u_x(g_i)-1)=0$ for all $x$ (\emph{Ribet cocycles}), and
$Z^{\mathrm{cr}}_M(\rho)$ the module of functions obeying the crystalline relations of $\rho$
(Definition~\ref{pg:def-sel}). Put $\prics_M(\rho)=\prirl\bigl(Z^{\mathrm R}_M(\rho),
Z^{\mathrm{cr}}_M(\rho)\bigr)$.
\begin{enumerate}
\item[(a)] If, at each $v\mid\ell$, the restrictions to $D_v$ of all $u_x$ and of $\rho$ lie in a
class $\mathcal C_v$ of local lattices (as in Proposition~\ref{pg:loss}), then
$\prics_M(\rho)\le\sum_{v\mid\ell}\delta_v(\rho)$, where the \emph{defect} $\delta_v(\rho)$ is the
length of the crystalline cocycles of the Tate dual of $\rho|_{D_v}$ modulo those whose cup product
with every crystalline cocycle of every member of $\mathcal C_v$, under every pairing to
$\Z/\ell^m(1)$, is a coboundary.
\item[(b)] Let $c$ be in lift setting, $R$ a main summand of rank $r'$ and string length $t'$,
and $\rho_x$ ($x\in Y$) representations of rank $n$ satisfying the conclusion (c) of
Proposition~\ref{pri:BL} for $c$. For all large $\auO$, all lattices $u_x$ in $\rho_x$, all
lattices $T_0$ in $R$ and $T$ equal to $T_0$ on the $D_v$, $v\mid\ell$, and one $\pi\in Y$ with
$u_\pi=T\oplus\bigoplus_{e<t'}\auO(-e)$, there are classes $\mathcal C_v$ as in (a) containing all
$u_x|_{D_v}$ and the restrictions of the $t'+1$ summands $L$ of $u_\pi$, with
\[
\sum_{L}\sum_{v\mid\ell}\delta_v(L)/e_\auO\le r'\,\pgloc_\ell .
\]
\end{enumerate}
\end{proposition}

This is Proposition~\ref{pg:loss} in the form in which it is used (the passage from the local
statements of \S\ref{pll:sec} to (b), a sum over the local degrees $\le2f$ and over the summands with
weights $1$ and $1/t$, is elementary). Proposition~\ref{pg:loss} is stated in the form (a) with Ribet
cocycles. Its proof is given in \S\ref{pll:sec}: for (a), the restrictions to $D_v$ of the Ribet
cocycles come from $\mathcal C_v$ by Lemma~\ref{pll:ribet} (with $\Gamma=D_v$ inside $\pgG$, the
restrictions of the $u_x$ lying in $\mathcal C_v$ by hypothesis), and Proposition~\ref{pll:dual}
with $T=\rho|_{D_v}$ bounds each local contribution; a cocycle obeys the crystalline relations of
$\rho$ as soon as its restriction to every $D_v$, $v\mid\ell$, does. The explicit bounds of (b) come with the explicit exponents of
Theorem~\ref{pll:exponent} proved in {\cite[\S2.3]{P56}}.

\begin{proposition}[Base level of a label]\label{pri:BL}
Let $c$ be in lift setting for an auxiliary $\vartheta$. There are Hodge--Tate data $J'$ and a
space of forms of ``base level'' (the space $\mathcal F(O)$ of Definition~\ref{prm:def-level}: level
$J_O$ and character $\chi_O$; these forms are of principal level at the chosen places) with: (a) $e_c\mathcal F(O)$ consists of forms of base level; (b) if
$y$ is integral at $\iota$, so is $N\,e_cy$; (c) for every finite set $U$ of places off
$\pgS(\ell)\cup Q\cup\overline Q$ and every system $x$ of Hecke eigenvalues occurring in base level
with principal level at $U$, the normalised Galois representation $\rho_x=S_x\otimes\eta_O$ is
conjugate self-dual with string length $t'$, has Hodge--Tate data $J'$, is unramified outside $\ell
p_0p_1$, $Q\cup\overline Q$, $U\cup\overline U$, inertia at a chosen $w\nmid\ell$ acts through the roots of
the rescaled label times $(X-1)^{t'}$ and semisimply (the monodromy of $\operatorname{rec}(\Pi_w)$
is $0$), so that $o_w(\rho_x(s))(\rho_x(s)-1)=0$ for $s\in I_w$, $o_w$ the monic polynomial whose
roots are the (images of the) roots of the rescaled label, and $(\rho_x(s)-\psi(s))(\rho_x(s)-1)=0$
for $s$ in inertia above $Q\cup\overline Q$.
\end{proposition}
\noindent Proved in \S\ref{pax:ss-BL}, from Propositions~\ref{pga:G123} and \ref{pga:G45} and the
cuspidal-support part of Lemma~\ref{pga:depth0}, with Lemma~\ref{pay:eta}; the semisimplicity at a
chosen $w$ is the paragraph ``No monodromy at a chosen $w$'' of that proof (it uses the clause ``no
rescaled label has the root $1$'' of the lift setting); the base level and the
product labels at the chosen places are Definition~\ref{prm:def-level}, Lemmas~\ref{prm:level},
\ref{prm:label} and Corollary~\ref{prm:inertia}.\par\smallskip

\begin{lemma}[Block-diagonal lattices]\label{pri:block}
Let $\rho_x$ ($x\in Y$ finite) be continuous semisimple representations of rank $n$ and,
for $x\in\priX\subseteq Y$, $\pgtr\rho_x=\psi\pgtr R_x+\sum_{e<t'}\chi^e$ with $R_x$ main summands,
and let $u'_x$ be lattices in $R_x$ over large $\auO$ whose values span $M_{r'}(\auO)$ (M4). Then for
all large $\auO$ there are lattices $u_x$ in the $\rho_x$ with $u_x=\psi u'_x\oplus\bigoplus_{e<t'}
\auO(-e)$ for $x\in\priX$, and the $\psi u'_x(g)$ span $M_{r'}(\auO)$.
\end{lemma}
\noindent Proved in Lemma~\ref{pax:block} (Brauer--Nesbitt in characteristic $0$), where the two
hypotheses used (semisimplicity of the $\rho_x$, spanning (M4) of the $u'_x$) are stated.\par\smallskip

\begin{lemma}[Inputs on forms]\label{pri:forms}
(a) (Higman operators; freeness of the level.) For a finite set of places $w_j$ off the base level,
with $\Gamma=\prod_j\auGL_n(\auF_{w_j})$ acting on the forms of principal level at the $w_j$, and for
every subgroup $\Gamma'$ of a factor, there is an integral operator $\Psi$ with
$\sum_{k\in\Gamma'}k^{-1}\Psi k=c_0$, where $c_0$ is a fixed common multiple of the orders of the
stabilisers of the base level, a divisor of $N$; for $\Gamma'=\Gamma$ this says that $c_0y$ is a norm
$\sum_{k\in\Gamma}k\Psi y$ for $y$ of base level. (b) (Pairings.) If $y'$ is integral at every
isomorphism and $y''$ is integral at $\iota$ then $N\langle y',y''\rangle\in\pgZb$.
\end{lemma}
\noindent Proved in \S\ref{pax:ss-forms} (with $c_0=\#(J_O/J_0)$, $J_0$ the free sublevel of the
base level $J_O$, a divisor of $N$) and \S\ref{pax:ss-pairing}.\par\smallskip

\begin{lemma}[Chain lemma]\label{pri:chain}
Let $y$ be an eigenform in base level whose normalised Galois representation has the eigenvalue
$1$ at every element of $\pgG$ (a \emph{weak lift}). At every split place $w\nmid\ell$ off the level
the rescaled Hecke polynomial $H_w(y;X_0X)$, $X_0$ the value at a geometric Frobenius of
$\iota^{-1}\eta_O^{-1}$, is divisible by $\prod_{i<t'}(X-(\auNm w)^i)$.
\end{lemma}
\noindent Proved in \S\ref{pax:ss-chain}, from (G1) and the spherical chain lemma
(Lemma~\ref{pax:tadic}, proved directly in Corollary~\ref{pay:roots}). The statement holds at every
split $w\nmid\ell$ off the level at which $y$ is spherical.\par\smallskip

\begin{proposition}[Odd classes]\label{pri:odd}
Recall that $E$ is a CM field contained in $\Q(\zeta_{q_1d})$. There is an integer
$0<\pgcodd\le2^{4(q_1d+t+3)^{27}}h_E$ such that for every prime $\ell$, every
coefficient ring $\auO$ containing the values of the Dirichlet characters modulo $q_1d$, all $\nu\ge0$
and $0\le e<t$, the module $\mathrm{Sha}^-_{e,\nu}$ has length at most $\prival_\auO(\pgcodd)$. Here,
for $e\ne0$, $\mathrm{Sha}^-_{e,\nu}$ consists of the continuous cocycles $z\colon\pgG\to
(\auO/\ell^\nu)(-e)$ unramified outside $\ell$, vanishing on the elements of trivial cyclotomic
character in the decomposition groups above $\ell$, with $z(g^c)=-z(g)$; for $e=0$ it consists of
the everywhere unramified homomorphisms $\pgG\to\auO/\ell^\nu$.
\end{proposition}
\noindent Proved in Proposition~\ref{pod:odd}, for every prime $\ell$ including $\ell=2$
(Remark~\ref{pod:rem-two}). The argument is an Euler-system argument with Jacobi-sum units, in the
manner of Kolyvagin and of Rubin's proof of the Main Conjecture for $\Q(\mu_{p^\infty})$ (\cite[Appendix by K.~Rubin, p.~397 ff.]{LangCyclo},
whose \S1 works over $\Q(\mu_m)^+$ for general $m$), with explicit constants.\par\smallskip

\begin{proposition}[Kummer maps and detecting points]\label{pri:kummer}
Let $F=E^+$, $\nu>0$, $G_\nu=\ker(\pgcyc\bmod\ell^\nu)\subseteq\pgG$.
\begin{enumerate}
\item[(a)] Let $x=(w_i)_{i\in I}$ be a finite family of degree-one places $w_i\nmid\ell$ of $E$ above
distinct primes of $F$ with Frobenius elements in $G_\nu$, let $S_0(x)$ be the set of primes of $F$
above $\ell$ or below some $w_i$, let $e\in\Z$, and let $(f_i)_{i\in I}$ be maps $G_\nu\to\auF_\ell$
indexed by the same set. There is a Kummer map $\kappa\colon\pgG\to\Z/\ell^\nu$ of exponent $e$ for
$x$ off $T=S_0(x)$ (Definition~\ref{pku:def-kummer}) whose reduction on $G_\nu$ is not in the span
of the $f_i$. Here a Kummer map is given by a locally constant function $y$ on $G_F$, zero outside
$G_\nu$, which on $G_\nu$ is a homomorphism to $\operatorname{Hom}((\Z/\ell^\nu)((\pgcyc^+)^{-e}),\Q/\Z)$,
equivariant for the action of $G_F$ through $(\pgcyc^+)^{e}\pgcyc$, unramified outside $T$,
written $y(g)(a)=a\kappa(g)/\ell^\nu$, which satisfies the reciprocity law:
$\sum_p y(\phi_p)(z(\gamma_p))=0$ for every finite family of distinct primes $p\notin T$ of $F$ with
Frobenius elements $\phi_p\in G_\nu$ and normalised tame generators $\gamma_p$, and every cocycle
$z$ of $G_F$ in $(\Z/\ell^\nu)((\pgcyc^+)^{-e})$ unramified outside $S_0(x)$ and the $p$, vanishing at
complex conjugations, on the decomposition groups of the places $w_i$, and on the part in
$\ker\pgcyc^+$ of the decomposition groups above $\ell$. ($\pgcyc^+$ extends $\pgcyc$ by $1$ at $c$.)
\item[(b)] Let $g_0$ be a topological generator of $\pgG/(\pgG\cap\ker\pgcyc)$. There is a finite set
$s$ of elements of trivial cyclotomic character such that a locally constant cocycle
$\pgG\to(\auO/\ell^\nu)(-e)$ unramified outside $\ell$ that vanishes at $g_0$ and on $s$ is zero,
for all $\auO$, $\nu$, $e$. If $\ell\le h^2$, $s$ can be chosen with $(t+1)(\#s+1)\le h^{21}$.
\end{enumerate}
\end{proposition}
\noindent Proved in \S\ref{pku:sec}: (a) is Proposition~\ref{pku:main} (from the Greenberg--Wiles
formula, local and global duality and Chebotarev's theorem; with $T=S_0(x)$ the relations are
restrictions of dual Selmer classes), and (b) is Proposition~\ref{pku:b} (from
Lemma~\ref{pod:detect}). In (a) the number of maps $f_i$ is at most the number of places of $x$, and
in (b) the size bound is claimed only for $\ell\le h^2$; these are the forms used in
Proposition~\ref{pri:triangle} (Proposition~\ref{pku:triangle}).\par\smallskip

\begin{proposition}[Dictionary at an auxiliary place]\label{pri:dict}
Let $w$ be a degree-one place off the base level with $\auNm w\equiv1\pmod{\ell^\nu}$, $\gamma$ the
tame generator and $\phi$ a Frobenius element at $w$, and $\Gamma_w=\auGL_n(\auF_w)$ acting on the
forms of principal level at $w$. Fix a uniformiser $\varpi$ and a logarithm $\mathrm{ind}\colon
\auF_w^\times\to\Z/\ell^\nu$ with $\mathrm{ind}(\bar z)=1$ for the unit $z$ whose Artin symbol is
$\gamma$. Write $\langle\gamma\rangle=\mathrm{diag}(\bar z,1,\dots,1)\in\Gamma_w$, $U$ for the
Hecke operator of $\mathrm{diag}(\varpi,1,\dots,1)$ at pro-$p$ Iwahori level, $\tilde U=
\eta_O(\phi)U$, and $\operatorname{tr}[\,\cdot\,]$ for the traces of the normalised Galois
representations as elements of the Hecke algebra $\priT$ at principal level at $w$.
\begin{enumerate}
\item[(a)] On a form on which $\Gamma_w$ acts by $\chi\circ\mathrm{ind}\circ\det$,
$\operatorname{tr}[(\gamma-1)g\,p(\phi)]=(\langle\gamma\rangle-1)\operatorname{tr}[g\,p(\phi)]$ for
all $g$ and $p\in\pgZb[X]$.
\item[(b)] On a form on which the stabiliser of a line in $\Gamma_w$ acts by $\chi(\mathrm{ind}(
k_{11}))$ with $\chi(n!)\ne1$, $\operatorname{tr}[(\gamma-1)g\,p(\phi)]=p(\tilde U)
\operatorname{tr}[(\gamma-1)g]$.
\item[(c)] $\sum_{k\in\Gamma_w}k$ kills $\operatorname{tr}[(\gamma-1)g]$.
\item[(d)] If $\delta\in\priT$ and every base system on which $\delta\ne0$ has, at $w$, a rescaled
Hecke polynomial equal to $\prod_{i<t'}(X-(\auNm w)^i)$ times a divisor of $H$, then $\delta
\sum_{k\in\Gamma}k$ kills the value at the weight $1$ of the operator $B$ of~\eqref{pri:eq-B}.
\end{enumerate}
\end{proposition}
\noindent Proved in \S\ref{pax:ss-dict} (parts (a)--(d)); part (d) uses Lemma~\ref{pax:toproot},
proved in \S\ref{pay:ss-cons}.\par\smallskip

\subsection{Fitting ideals}\label{pri:ss-fitting}

Let $R=\auO[X_J]$, $J$ finite, and $s$ a finite set. For an $R$-submodule $N\subseteq R^s$ let
$\priFitt(N)$ be the ideal generated by the determinants of the $s\times s$ matrices all of whose
rows lie in $N$ (the zeroth Fitting ideal of $R^s/N$). Let $\Phi\colon\auO^n\to R^s$ be
$\auO$-linear, $p\in\auO\smallsetminus\{0\}$ with $\Phi(p\auO^n)\subseteq N$, and $z\colon G\to\auO^n$
a function on a set $G$. The \emph{forms} of $N$ are
\[
\mathrm{F}(N)=\{\Psi\circ z:\ \Psi\colon\auO^n\to\auO/p\ \text{linear},\ \Psi(\Phi^{-1}N)=0\}
\subseteq\operatorname{Map}(G,\auO/p).
\]

\begin{lemma}\label{pri:fitt-basic}
(a) If $N\subseteq N'$ then $\priFitt(N)\subseteq\priFitt(N')$. (b) If $a\in R$, $v\in R^s$ and
$av\in N$, then $a\priFitt(N+Rv)\subseteq\priFitt(N)$. (c) If $w\in\auO^n$ and $\Psi(w)=0$ for every
form-defining $\Psi$ (that is, every $\Psi$ with $\Psi(\Phi^{-1}N)=0$), then $\Phi(w)\in N$.
(d) Let $\mathfrak m=\ker(R\to\auO,\ X_j\mapsto0)$. Then $\priFitt(N+\mathfrak mR^s)\subseteq
\priFitt(N)+\mathfrak m$.
\end{lemma}
\begin{proof}
(a) is clear. (b) A row of $N+Rv$ is $n_i+r_iv$; expand the determinant multilinearly: the terms
with two rows proportional to $v$ vanish, those with no such row lie in $\priFitt(N)$, and in a term
with one such row multiplication by $a$ turns that row into $r_i(av)\in N$. (c) $M=\auO^n/\Phi^{-1}N$
is killed by $p$, hence a finitely generated $\auO/p$-module; $\auO/p$ is self-injective (a
quotient of a discrete valuation ring by a nonzero ideal), so linear maps $M\to\auO/p$ separate
points; thus $w\in\Phi^{-1}N$. (d) Reduce modulo $\mathfrak m$ entrywise: the determinant of a matrix
with rows $n_i+m_i$ is congruent modulo $\mathfrak m$ to that with rows $n_i$.
\end{proof}

\begin{lemma}[Cost of a submodule of forms]\label{pri:fitt-cost}
Let $W\subseteq\operatorname{Map}(G,\auO/p)$ be a submodule. There are $e\in\auO$ with $\prival(e)\le
\prirl(\mathrm F(N),W)$ and $K\supseteq N$ such that $\Phi(w)\in K$ for every $w$ with $\Psi(w)=0$
for all form-defining $\Psi$ with $\Psi\circ z\in W$, and $e\priFitt(K)\subseteq\priFitt(N)$.
\end{lemma}
\begin{proof}
Let $M=\auO^n/\Phi^{-1}N$ (finite length, killed by $p$) and $M^*=\operatorname{Hom}(M,\auO/p)$; the
map $\Psi\mapsto\Psi\circ z$ maps $M^*$ onto $\mathrm F(N)$. Let $M^*_W$ be the preimage of $W$
and $M'=\{m\in M:\Psi(m)=0\ \forall\Psi\in M^*_W\}$. By Matlis duality over the artinian
self-injective ring $\auO/p$, $\prilen(M')=\prilen(M^*/M^*_W)$; and $M^*/M^*_W$ maps onto
$\mathrm F(N)/(\mathrm F(N)\cap W)$ with kernel the image of those $\Psi$ with $\Psi\circ z\in W$
that already lie in $M^*_W$, i.e.\ isomorphically; so $\prilen M'=\prirl(\mathrm F(N),W)$. Choose a
composition series $0=M'_0\subset\dots\subset M'_k=M'$, $k=\prilen M'$, and lifts $w_i\in\auO^n$ of
generators of $M'_i/M'_{i-1}$; then $\varpi w_i\in\auO w_1+\dots+\auO w_{i-1}+\Phi^{-1}N$. Put
$K_i=N+R\Phi(w_1)+\dots+R\Phi(w_i)$, $K=K_k$, $e=\varpi^k$. Then $\varpi\Phi(w_i)\in K_{i-1}$ and
Lemma~\ref{pri:fitt-basic}(b) gives $\varpi\priFitt(K_i)\subseteq\priFitt(K_{i-1})$. Every $w$ as in
the statement has image in $M'$, so $\Phi(w)\in K$.
\end{proof}

\begin{definition}[Triangle cost]\label{pri:def-cost}
Let $\priZ\subseteq\operatorname{Map}(G,B)$ be a submodule, $B$ an $\auO$-module, $\tau,\tau',
\phi,\phi'\colon J\to G$ with $J$ totally ordered, $I\subseteq J$, $C\in\auO$, $b\in[0,\infty]$.
$\priZ$ has \emph{triangle cost $b$ along $I$ with constant $C$} if there is $A\subseteq G$ such
that, with $\priZ_0=\{c\in\priZ:c|_A=0\}$: (a) for $j\in I$ and $c\in\priZ_0$ with $c(\tau_i)=
c(\tau'_i)$ for all $i<j$ and $c(\phi_i)=c(\phi'_i)$ for $i<j$ in $I$, one has $C(c(\tau_j)-
c(\tau'_j))=0$; (b) for every submodule $\priZ_1\subseteq\priZ_0$ whose members vanish at all
$\tau_j,\tau'_j$ and satisfy $c(\phi_j)=c(\phi'_j)$ for $j\in I$, $\prirl(\priZ,\priZ_0)+\prirl(
\priZ_1,\priZ_1[C])\le b$, where $\priZ_1[C]$ is the $C$-torsion.
\end{definition}

\begin{proposition}[Fitting ideal of a triangle]\label{pri:fitt-triangle}
Keep $\Phi$, $p$, $z$, and put $y_g=\Phi(z_g)$. Suppose: modulo $N$, the basis vector $\epsilon_g$
of $R^s$ is $y_g$ for $g\in s\subseteq G$; $y_{\tau_j}-X_jy_{\phi_j}\in N$ and $y_{\tau'_j}-X_j
y_{\phi'_j}\in N$ for all $j$; every $i$ below an element of $I$ lies in $I$ or satisfies $X_iR^s
\subseteq N$; and $\mathrm F(N)$ has triangle cost $b$ along $I$ with constant $C$. Then there is
$F\in R$ with $F\cdot X_I\in\priFitt(N)$ and $\prival(F(0))\le b+(\#I+\#s)\prival(C)$.
\end{proposition}
\begin{proof}
Let $A$ be as in Definition~\ref{pri:def-cost}. \emph{Stage 1.} Lemma~\ref{pri:fitt-cost} with
$W=\{c:c|_A=0\}$ gives $e$, $K\supseteq N$, with $\prival(e)\le\prirl(\mathrm F(N),\mathrm F(N)_0)$
and $e\priFitt(K)\subseteq\priFitt(N)$; every $K'\supseteq K$ has $\mathrm F(K')\subseteq\mathrm
F(N)_0$. \emph{Stage 2.} Put $v_j=\Phi(z_{\phi_j}-z_{\phi'_j})$, so $\Phi(z_{\tau_j}-z_{\tau'_j})
\equiv X_jv_j\pmod K$. Going up $I$ in increasing order, let $T_j=T_{j^-}+Rv_j$ ($T_{\min^-}=K$). We
claim $CX_jv_j\in T_{j^-}$. For $i<j$, $X_iv_i\in T_{j^-}$ (if $i\in I$, $v_i\in T_{j^-}$; otherwise
$X_iR^s\subseteq N$). So every form-defining $\Psi$ of $T_{j^-}$ has $\Psi(z_{\tau_i})=\Psi(
z_{\tau'_i})$ for $i<j$ and $\Psi(z_{\phi_i})=\Psi(z_{\phi'_i})$ for $i<j$ in $I$; as $\Psi\circ z\in
\mathrm F(N)_0$, Definition~\ref{pri:def-cost}(a) gives $C\Psi(z_{\tau_j}-z_{\tau'_j})=0$. By
Lemma~\ref{pri:fitt-basic}(c), $C\Phi(z_{\tau_j}-z_{\tau'_j})\in T_{j^-}$, hence $CX_jv_j\in
T_{j^-}$. Lemma~\ref{pri:fitt-basic}(b) gives $CX_j\priFitt(T_j)\subseteq\priFitt(T_{j^-})$; with
$T=T_{\max}$, $\priFitt(T)\prod_{j\in I}CX_j\subseteq\priFitt(K)$. \emph{Stage 3.} Let $T_0=T+
\mathfrak mR^s$. Lemma~\ref{pri:fitt-cost} for $T_0$ and $W=\{C\text{-torsion}\}$ gives $f$ and
$K'\supseteq T_0$ with $\prival(f)\le\prirl(\mathrm F(T_0),\mathrm F(T_0)[C])$ and $f\priFitt(K')
\subseteq\priFitt(T_0)$. Since $\Psi(Cz_g)=0$ whenever $\Psi\circ z$ is $C$-torsion, $\Phi(Cz_g)\in
K'$, and $C\epsilon_g=\Phi(Cz_g)-C(y_g-\epsilon_g)\in K'$; so $C^{\#s}\in\priFitt(K')$ and $fC^{\#s}\in
\priFitt(T_0)\subseteq\priFitt(T)+\mathfrak m$ (Lemma~\ref{pri:fitt-basic}(d)). Thus some $d\in
\priFitt(T)$ has $d(0)=fC^{\#s}$. \emph{Costs.} The members of $\mathrm F(T_0)$ vanish at all
$\tau_j,\tau'_j$ (as $y_{\tau_j}\in X_jR^s+N\subseteq T_0$) and satisfy $c(\phi_j)=c(\phi'_j)$ for
$j\in I$ (as $v_j\in T$), and lie in $\mathrm F(N)_0$; so by (b) of the definition $\prival(e)+
\prival(f)\le b$. Then $F=eC^{\#I}d$ satisfies $FX_I\in\priFitt(N)$ and $\prival(F(0))\le b+
(\#I+\#s)\prival(C)$.
\end{proof}

\subsection{Ribet's method over an algebra}\label{pri:ss-ribet-alg}

Let $\mathsf R$ be a commutative ring, $\Lambda$ an $\mathsf R$-algebra and $\rho\colon\Lambda\to
M_\sigma(\mathsf L)$ an $\mathsf R$-algebra map to matrices over a commutative $\mathsf
R$-algebra $\mathsf L$, such that the characteristic polynomial of every $\rho(\lambda)$ has
coefficients in the image of $\mathsf R$. For a matrix $P\in M_m(\Lambda)$ let $\rho(P)\in
M_{m\sigma}(\mathsf L)$ be the block matrix $(\rho(P_{ij}))$.

\begin{lemma}[Amitsur]\label{pri:amitsur}
For every $P\in M_m(\Lambda)$ the characteristic polynomial of $\rho(P)$ has coefficients in the
image of $\mathsf R$.
\end{lemma}
\begin{proof}
Write $\rho(P)=\sum_qE_q$ with $E_q$ the block matrix whose only nonzero block row is the $q$-th
row of $\rho(P)$. By Amitsur's formula (Proposition~\ref{pax:amitsur}; \cite{Amitsur}) the coefficients of the characteristic
polynomial of a sum of matrices are integral polynomials in the coefficients of the characteristic
polynomials of the products $E_{q_1}\cdots E_{q_k}$. Such a product is
$\epsilon_{q_1}\,\rho(P_{q_1q_2}P_{q_2q_3}\cdots P_{q_{k-1}q_k})\,\pi_{q_k}$, where $\epsilon_q$ is
the $q$-th block column of $1$ and $\pi_q$ the $q$-th block row of $\rho(P)$; its characteristic
polynomial agrees, up to a power of $X$, with that of $\pi_{q_k}\epsilon_{q_1}\rho(\cdots)=
\rho(P_{q_kq_1}P_{q_1q_2}\cdots P_{q_{k-1}q_k})$, which has coefficients in $\mathsf R$.
(Details in Remark~\ref{pax:rem-amitsur}.)
\end{proof}

\begin{lemma}[Ribet's polynomial]\label{pri:ribet-poly}
Let $\mathfrak a\subseteq\Lambda$ be a two-sided ideal containing $\ker\rho$, $y\in\Lambda^m$,
$I=\sum_k\Lambda y_k$, $\rho_0\in\Lambda$, and $B\in M_m(\Lambda)$ with the entries of $By$ in
$\mathfrak a+I\rho_0$. Then there is $P\in\mathsf R[X]$ with $P(0)=\det\rho(B)$ (up to the image of
$\mathsf R$) and $I\,P(\rho_0)\subseteq\mathfrak a$.
\end{lemma}
\begin{proof}
Write $By\equiv B_1y\rho_0\pmod{\mathfrak a^m}$ with $B_1\in M_m(\Lambda)$. Let $\Lambda[X]$ act on
$\Lambda/\mathfrak a$ with $\Lambda$ acting on the left and $X$ by right multiplication by
$\rho_0$ (well defined as $\mathfrak a$ is two-sided). Then $B_X=B-XB_1$ kills the image of $y$.
Extend $\rho$ to $\rho_X\colon\Lambda[X]\to M_\sigma(\mathsf L[X])$ coefficientwise; as
$\Lambda[X]=\Lambda\otimes_{\mathsf R}\mathsf R[X]$, Lemma~\ref{pri:amitsur} (for $\mathsf R[X]$)
puts the characteristic polynomial $\chi(T)$ of $\rho_X(B_X)$ over $\mathsf R[X]$. By
Cayley--Hamilton, $\chi(\rho_X(B_X))=0$; writing $\chi(T)=TQ(T)+\chi(0)$ and $B'=-Q(B_X)\in
M_m(\Lambda[X])$ (a polynomial in $B_X$ with coefficients in $\mathsf R[X]$), we get $\rho_X(B'B_X)=
P\cdot1$ with $P=\chi(0)=\pm\det\rho_X(B_X)$, so $P(0)=\pm\det\rho(B)$. The entries of $B'B_X-P$
have all coefficients in $\ker\rho\subseteq\mathfrak a$, so $P$ kills the image of $y$ too:
$y_kP(\rho_0)\in\mathfrak a$ for all $k$.
\end{proof}

\begin{proposition}[Ribet element of a matrix of tied rows]\label{pri:ribet-elt}
Let $u_x\colon\auO[G]\to M_\sigma(\auO)$ ($x\in Y$ finite) be representations, $\mathsf L$ an
$\auO$-algebra, $S$ an $\auO$-subalgebra of $\operatorname{Map}(Y,\mathsf L)$ containing the
constants and, as functions of $x$, the coefficients of the characteristic polynomials of the
$u_x(g)$, and $\priJ$ an $S$-submodule of $\operatorname{Map}(Y,\mathsf L)$. Let $\mathfrak a_0$ be
the left ideal of $\auO[G]$ spanned by the $r$ with $(x\mapsto\pgtr u_x(rg))\in\priJ$ for all $g$, let
$y\in\auO[G]^m$ span a left ideal $I\supseteq\priAug$ modulo $\bigcap_x\ker u_x$, and $B\in
M_m(\auO[G])$ with the entries of $By$ in $\mathfrak a_0+I\rho_0$. Suppose that for $x$ in a subset
$\priX$, $u_x$ is (up to a fixed change of basis) $\psi_x\oplus\bigoplus_{e<t'}\eta_e$ with characters
$\eta_e$, that $\rho_1$ acts as $\mu_0$ on $\eta_0$ and as $0$ on the other summands, and that
$\rho_0$ acts as $0$ on $\eta_0$. Then there is $f\in S$ with $f\cdot\det u_x(w)\in\priJ$ (as functions
of $x$) for every $w\in\priAug$, and $f(x)=\mu_0\det\psi_x(B)\prod_e\det\eta_e(B)$ for $x\in\priX$.
\end{proposition}
\begin{proof}
By Lemma~\ref{pri:amitsur} applied to the family, the map $\rho=(u_x)_x$ extends to an
$S$-algebra map $\rho\colon S[G]\to M_\sigma(\operatorname{Map}(Y,\mathsf L))$ all of whose
characteristic polynomials have coefficients in $S$. Let $\mathfrak a$ be the set of
$\lambda\in S[G]$ with $\pgtr\rho(\lambda\mu)\in\priJ$ for all $\mu\in S[G]$: a two-sided ideal
(cyclicity of the trace) containing $\ker\rho$ and the image of $\mathfrak a_0$ (expand $\mu$ in
group elements). Lemma~\ref{pri:ribet-poly} gives $P\in S[X]$ with $P(0)=\det\rho(B)$ and $I\,
P(\rho_0)\subseteq\mathfrak a$. Put $f=\pgtr\rho(\rho_1P(\rho_0))\in S$. For $x\in\priX$,
$\rho_1\rho_0$ acts as $0$, so $\rho_1P(\rho_0)$ acts as $P(0)\rho_1$ and $f(x)=P(0)(x)\mu_0$; and
$\det\rho(B)$ at $x$ is the product of the determinants on the summands. For $w\in\priAug$, Cayley--
Hamilton gives $w''\in S[w]$ with $\rho(w''w)=\det\rho(w)$. Since $\priAug$ is two-sided, $w\rho_1\in
\priAug\subseteq I+\ker\rho$, so $w''w\rho_1P(\rho_0)\in\mathfrak a$, and $\det\rho(w)\,f=\pgtr\rho(
w''w\rho_1P(\rho_0))\in\priJ$.
\end{proof}

\subsection{Ribet's method to first order}\label{pri:ss-ribet}

Fix the following data. $G$ is a group, $\auO$ a discrete valuation ring with fraction field $K$,
$u_x\colon G\to\auGL_\sigma(\auO)$ ($x\in Y$, finite), $S$, $\priJ$ as in
Proposition~\ref{pri:ribet-elt} but with values in $\mathsf L[X_J]$, $J$ a finite totally ordered
set; $p\in\auO\smallsetminus\{0\}$; elements $\tau_k,\phi_k\in G$ for $k\in J\times\{\pm\}$;
polynomials $H_j\in\auO[X]$. \emph{The tied relations} are $\tau_k-1-X_j(\phi_k-1)H_j(\phi_k)$,
$k=(j,\pm)$. We say $\priJ$ is a \emph{trace module} if $S$ contains the $X_j$ and the coefficients
of the characteristic polynomials of the $u_x(g)$, and $\priJ$ contains the families $x\mapsto
p\,\pgtr u_x(g)$ and $x\mapsto\pgtr u_x\bigl((\tau_k-1-X_j(\phi_k-1)H_j(\phi_k))g\bigr)$ for all $g$
and $k$. A \emph{common Ribet element} at $\pi\in\priX\subseteq Y$ with value $d$ is $f\in S$ such
that $f\cdot\det(1-u_x(g))\in\priJ$ for all $g$, no monomial of $f(x)$, $x\in\priX$, properly divides
$X_J$, and the top coefficient of $f(\pi)$ is $d$.

For a character $\xi\colon G\to\auO^\times$, a \emph{character cocycle for $\xi$} (relative to
$\delta_v\in G$, $\gamma$, $a\in G$) is a cocycle $c\colon G\to\auO/p$ for $\xi$ that obeys all
linear relations among the families $(u_x(g)-1)_x$, vanishes at the $\delta_v$, and, if $c(\gamma)
=0$, satisfies $(\xi(g)-1)c(a)=0$ for all $g$.

\begin{proposition}[Ribet's method to first order; exact form of Proposition~\ref{pg:ribet}]
\label{pri:ribet}
Assume:
\begin{enumerate}
\item[(1)] for $x\in\priX$, $u_x=\rho_x\oplus\bigoplus_{e<t'}\xi_e$ (fixed change of basis) with
$\xi_0=1$, and the $\rho_\pi(g)$ span $M_{r'}(\auO)$;
\item[(2)] $\gamma_x\in G$ acts on $\rho_x$ by a scalar $\zeta_x$ with $\zeta_x-1\in\auO^\times$, and
trivially on the $\xi_e$;
\item[(3)] $\priJ$ is a trace module, and $H_j(\xi_e(\phi_k))\equiv1\pmod p$ for all $e,k$;
\item[(4)] $a\in G$ and $C\in\auO$ with $\prod_{e'\ne e}(\xi_e(a)-\xi_{e'}(a))\mid C$ for every $e$;
and $\prity\colon J\to\{0,\dots,t'-1\}$ is monotone;
\item[(5)] the cocycles $G\to(K/\auO)^{r'}$ for $\rho_\pi$ obeying all linear relations among the
$(u_x(g)-1)_x$ and vanishing at the $\tau_k$ and at $\gamma_\pi$ form a module of length $\le b_\rho$;
\item[(6)] $\delta_v\in G$ ($v$ in a finite set) are trivial on the $\xi_e$, and $o_v(u_x(\delta_v))
(u_x(\delta_v)-1)=0$ for all $x$, with $o_v\in\auO[X]$;
\item[(7)] for each $e$, every module of character cocycles for $\xi_e$ has triangle cost $b_e$
along $\{j:\prity(j)=e\}$ with constant $C$ (for $\tau_j=\tau_{(j,+)}$, $\tau'_j=\tau_{(j,-)}$ and
likewise $\phi$), and a finite set $s\subseteq G$ detects it (a member vanishing on $s$ is $0$).
\end{enumerate}
Then there is a common Ribet element at $\pi$ whose value $d$ lies in $\auO$ with
\[
\prival(d)\le b_\rho+\sum_{e<t'}b_e+t'\Bigl(\sum_v\prival(o_v(1))+2\#s\,\prival(C)\Bigr)+(\#J+1)
\prival(C).
\]
\end{proposition}

\begin{proof}
\emph{Generators.} As $Y$ is finite and $\auO$ noetherian, the $\auO$-module spanned by the families
$(u_x(g)-1)_x$ is finitely generated; choose $s_0\subseteq G$ finite so that, with the generators
$\mathrm{gen}=(\delta_v)_v\sqcup s_0\sqcup s$, there are vectors $z_g\in\auO^{\mathrm{gen}}$ with
$u_x(g)-1=\sum_iz_{g,i}(u_x(\mathrm{gen}_i)-1)$ for all $x,g$. A \emph{row} is a vector $r\in
\auO[X_J][G]^{\mathrm{gen}}$, with \emph{value} $\sum_ir_i(\mathrm{gen}_i-1)$. Let $c_\zeta=(\zeta_\pi-1)
^{-1}$, $\nu_0=c_\zeta(\gamma_\pi-1)$ and $\rho_0=\nu_0+(1-\nu_0)(a-1)$. The \emph{tie ideal} is the left
ideal spanned by $\bigcap_x\ker u_x$, $p\priAug$, the tied relations, and $\priAug\rho_0$; a row is
\emph{tied} if its value lies in the tie ideal. Since $\priJ$ is a trace module, $p$ and the tied
relations lie in $\mathfrak a_0$ (Proposition~\ref{pri:ribet-elt}); and $\priAug=\sum_i\auO[X_J][G]
(\mathrm{gen}_i-1)$ modulo $\bigcap\ker u_x$. Thus the values of tied rows lie in $\mathfrak a_0+
I\rho_0$ with $I$ the left ideal of the $\mathrm{gen}_i-1$, and Proposition~\ref{pri:ribet-elt}
applies to every square matrix $B$ of tied rows (with $y_i=\mathrm{gen}_i-1$). The following rows are
tied: the cocycle rows $z_{gh}-z_g-g z_h$; $z_{\mathrm{gen}_i}-\epsilon_i$; $\sum\theta_gz_g$ whenever
$\sum\theta_g(u_x(g)-1)=0$ for all $x$; $p\cdot r$ for any row $r$; the ties $z_{\tau_k}-X_jH_j(
\phi_k)z_{\phi_k}$; $(g-1)(c_\zeta z_{\gamma_\pi}+(1-\nu_0)z_a)$; and the rows $o_v(\delta_v)\epsilon_v$.
(Each value acts as $0$ under every $u_x$, or is one of the generators of the tie ideal.)

\emph{The block.} Read rows through $\rho_\pi$ at $X_J=0$: $g\mapsto\rho_\pi(g)$, so a row gives a
matrix in $M_{r'}(\auO)^{\mathrm{gen}}$, i.e.\ a row of $r'\times r'$ blocks. Let $N_\rho$ be the
$\auO$-span of the images of tied rows. A vector $w\in((K/\auO)^{r'})^{\mathrm{gen}}$ killed by
$N_\rho$ defines $\Psi\colon\auO^{\mathrm{gen}}\to(K/\auO)^{r'}$, and $c=\Psi\circ z$ is a cocycle for
$\rho_\pi$ (cocycle rows), obeys the linear relations, vanishes at the $\tau_k$ (at $X_J=0$ the tie
is the row $z_{\tau_k}$) and at $\gamma_\pi$ (the normaliser row at $X_J=0$, through $\rho_\pi$, is
$c_\zeta z_{\gamma_\pi}$ times $(\rho_\pi(g)-1)$, and the $\rho_\pi(g)-1$ span; with $c_\zeta(\zeta_\pi-1)=
1$). Conversely $c$ determines $w$ (generator rows). By (5) the module of such $w$ has length $\le
b_\rho$, so (structure of finitely generated modules over $\auO$, and $(K/\auO)$-duality) the cokernel of
$N_\rho\to M_{r'}(\auO)^{\mathrm{gen}}$, read as $\auO$-modules of rank $r'\#\mathrm{gen}$, has length
$\le b_\rho$, and there is a square matrix $B_0$ of tied rows whose block determinant through $\rho_\pi$
at $X_J=0$ has valuation $\le b_\rho$.

\emph{The characters.} Fix $e$. Read rows through $\xi_e$ (each $g\mapsto\xi_e(g)\in\auO$),
forgetting the $\delta$-coordinates, and let $N_e\subseteq\auO[X_J]^{s_0\sqcup s}$ be the image of
the tied rows. The forms of $N_e$ (with $\Phi$ the inclusion of $\auO^{s_0\sqcup s}$ and $z$ as
above) are character cocycles for $\xi_e$: cocycle and linearity from the corresponding rows, the
vanishing at $\delta_v$ as explained below, and $(\xi_e(g)-1)(c_\zeta c(\gamma)+c(a))$ from the normaliser rows, as
$\xi_e(\gamma_\pi)=1$. On $\xi_e$ the tie reads $z_{\tau_k}-X_jH_j(\xi_e(\phi_k))z_{\phi_k}$, and
$H_j(\xi_e(\phi_k))=1+pc$ by (3); subtracting the tied row $pX_jcz_{\phi_k}$ gives the relation
$y_{\tau_k}\equiv X_jy_{\phi_k}$ of Proposition~\ref{pri:fitt-triangle}. Since $s$ detects the forms,
the elements of $s_0$ can be eliminated: by Lemma~\ref{pri:fitt-basic}(c) there are tied rows whose
$s_0$-part, on $\xi_e$, is the unit vector $\epsilon_i$; one row can be made to serve all $e$ at once
by correcting it successively on each character with the elements $\pi_e\in\auO[G]$ that are $C$ on
$\xi_e$ and $0$ on the other $\xi_{e'}$ (they exist by (4): $C\prod_{e'\ne e}(a-\xi_{e'}(a))/
\prod_{e'\ne e}(\xi_e(a)-\xi_{e'}(a))$). After elimination, Proposition~\ref{pri:fitt-triangle}
applies with $I=\{\prity=e\}$ to $N'_e=N_e+\mathfrak n'R^{s}$ in place of $N_e$, where $\mathfrak n$ is
the ideal of monomials not dividing $X_J$ and $\mathfrak n'=(\mathfrak n:X_{\{\prity<e\}})$, which
contains every $X_i$ with $\prity(i)<e$: thus $X_iR^s\subseteq N'_e$ for these $i$, the $i$ below
$I$ ($\prity$ is monotone), so the hypothesis of Proposition~\ref{pri:fitt-triangle} holds. Its
other hypotheses pass from $N_e$ to $N'_e$: $\Phi^{-1}N'_e\supseteq\Phi^{-1}N_e$, so the forms of
$N'_e$ are forms of $N_e$, hence character cocycles, with triangle cost $b_e$ by (7). Finally
$\priFitt(N'_e)\subseteq\priFitt(N_e)+\mathfrak n'$, by multilinear expansion of a determinant with
rows $n_i+m_i$, $n_i\in N_e$, $m_i\in\mathfrak n'R^s$. (The vanishing of the forms at $\delta_v$ is
automatic: the tied generator row $z_{\delta_v}-\epsilon_v$ shows that, modulo the image of the tied
rows, $z_{\delta_v}$ has zero component in $\auO^{s_0\sqcup s}$; the price is that the rows $o_v(\delta_v)\epsilon_v$, which read as
$o_v(1)\epsilon_v$ on $\xi_e$, must be used to complete the matrices below, which is the factor
$\prod_vo_v(1)$ in $c_\chi$.) The result: matrices $B_{e,i}$ over $\auO[X_J]$ and
weights $w_{e,i}$ with $\sum_iw_{e,i}\det B_{e,i}=\Gamma_eX_{\{\prity=e\}}+\zeta_e$, $\zeta_e\in(
\mathfrak n:X_{\{\prity<e\}})$, $\prival(\Gamma_e(0))\le b_e+(\#\{\prity=e\}+\#s)\prival(C)$; and
for each choice of one $B_{e,i(e)}$ for every $e$ a single matrix $M_i$ of tied rows with
$\det\xi_e(M_i)=c_\chi\det B_{e,i(e)}$ for all $e$, $c_\chi=C^{\#s}\prod_vo_v(1)$, obtained by the
common-unit-row construction just described (and completing with the rows $o_v(\delta_v)\epsilon_v$,
which makes the matrix block triangular on each $\xi_e$).

\emph{Gluing.} Let $k\in\auO[G]$ be $1$ on the $\xi_e$ and $0$ on every $\rho_x$, $x\in\priX$: take
$k=\prod_x(\zeta_x-1)^{-1}(\zeta_x-\gamma_x)$, each factor being $0$ on $\rho_x$ and $1$ on the
characters. Apply Proposition~\ref{pri:ribet-elt} to the matrices
$(1-k)B_0+kM_i$ (rows tied, as tied rows form a left $\auO[X_J][G]$-module), with $\rho_1=k\pi_0$ ($C$
on $\xi_0$, $0$ on the other $\xi_e$ and on every $\rho_x$; $\rho_0$ is $0$ on $\xi_0=1$), and sum with weights $\prod_ew_{e,i(e)}$.
For $x\in\priX$ the value is $C\det\rho_x(B_0)\cdot c_\chi^{t'}\prod_e(\Gamma_eX_{\{\prity=e\}}+
\zeta_e)$. Since each factor $\Gamma_eX_{\{\prity=e\}}+\zeta_e$ is preceded by the earlier factors,
which are multiples of $X_{\{\prity<e\}}$, the product has no monomial properly dividing $X_J$ and
top coefficient $\prod_e\Gamma_e(0)$ times the constant term of the rest. The top coefficient at
$\pi$ is $d=C\,c_\chi^{t'}\det\rho_\pi(B_0)|_{X=0}\prod_e\Gamma_e(0)$ (the remaining factors are taken
at $X_J=0$). Hence $\prival(d)\le\prival(C)+t'(\#s\prival(C)+\sum_v\prival(o_v(1)))+b_\rho+\sum_eb_e+
(\#J+t'\#s)\prival(C)$, which is the bound stated. Finally $\det(1-u_x(g))=\det u_x(w)$ for
$w=-(g-1)\in\priAug$ (computed through the coordinates $z$), so $f$ is a common Ribet element.
\end{proof}

\begin{remark}
The proof is written in two places, in this order. The paragraphs ``Generators'' and ``The block'' are given above in full. The paragraphs ``The characters'' and ``Gluing'' are given above in outline and are written in full in \S\ref{prm:ss-chars}
(Lemmas~\ref{prm:tied}--\ref{prm:coeff}, Propositions~\ref{prm:block}, \ref{prm:chars}); that text takes from the present proof the notation of ``Generators'', the matrix $B_0$ of ``The block'' and the identity $\det(1-u_x(g))=\det u_x(w)$ of the last sentence, and nothing from the two paragraphs that it replaces. The passage
from coefficients modulo $\ell^M$ to $K/\auO$ in (5) is Lemma~\ref{pri:divisible}.
\end{remark}

\begin{lemma}[Divisible coefficients]\label{pri:divisible}
Let $\mathcal V_B$ be the module of cocycles $G\to B^{r'}$ for $\rho$ obeying the linear relations
among the $(u_x(g)-1)_x$ and vanishing on a set $T$ and at $\gamma$. If $\prilen\mathcal V_{\auO/
p^M}\le b$ for all $M$, with $p$ a nonzero non-unit, then $\prilen\mathcal V_{K/\auO}\le b$.
\end{lemma}
\begin{proof}
Every element of $\mathcal V_{K/\auO}$ is torsion (its values are combinations, with coefficients
$z_g$, of its finitely many values at the generators). With $a=p^{b+1}$, multiplication by $a^{-1}$
identifies $\mathcal V_{\auO/a}$ with the $a$-torsion of $\mathcal V_{K/\auO}$. A torsion module all of
whose $a$-torsion has length $\le b<\prival(a)$ is killed by $\varpi^b$, hence equals its
$a$-torsion.
\end{proof}

\subsection{Even classes, the triangle and the odd classes}\label{pri:ss-triangle}

For a set $S$ of places, an \emph{even class} of exponent $e$ for $S$ is a locally constant cocycle
$z\colon\pgG\to(\auO/\ell^\nu)(-e)$ unramified outside the places above $\ell$, $S$ and $\overline
S$, with $z(g^c)=z(g)$, vanishing on the elements of trivial cyclotomic character in the
decomposition groups above $\ell$.

\begin{definition}[Triangle of primes]\label{pri:def-tri}
Let $J$ be finite totally ordered, $\prity\colon J\to\{1,\dots,t'-1\}$ monotone, $a\in\pgG$, $h_0,
c_{\mathrm{tri}}\in\N$. Degree-one places $w_j\nmid\ell$ with tame generators $\tau_j$ and Frobenius elements
$\phi_j\in G_\nu$ form a \emph{triangle} for $(a,h_0,c_{\mathrm{tri}})$ if: (i) there is a finite set $s$ of
elements of trivial cyclotomic character, with $(\#J+\#s+1)h_0\le c_{\mathrm{tri}}$, detecting (with $a$) the
cocycles unramified outside $\ell$, such that for every $e$ and $g\in s$ some $j$ with $\prity(j)=e$
has $z(\phi_j)=2z(g)$ for all even classes $z$ of exponent $e$ for $\emptyset$; (ii) $\prod_{k\ne e}
(\pgcyc(a)^{-e}-\pgcyc(a)^{-k})\mid\ell^{h_0}$ ($k,e<t'$) and $4(\pgcyc(a)^{-e}-1)^2\mid\ell^{h_0}$
for $0<e<t'$; (iii) if $z$ is an even class of exponent $\prity(j)$ for $\{w_i\}$ with $z(\tau_i)=0$
for $i<j$ and $z(\phi_i)=0$ for $i<j$ with $\prity(i)=\prity(j)$, then $\ell^{h_0}z(\tau_j)=0$.
\end{definition}

\begin{proposition}[Taylor--Wiles places in a triangle; exact form of
Proposition~\ref{pg:triangle}]\label{pri:triangle}
Let $t'\le t$, $h_0\le\pgv(N)$ with $\ell\le h^2$ or $h_0=0$, and $a$ a topological generator of
$\pgG/\ker\pgcyc$ with $\pgcyc(a)^e-1\mid\ell^{h_0}$ for $0<e<t$ (such $a$, $h_0$ exist: below).
Let $\mathcal W$ be a set of Frobenius data such that for every $\nu$ and $g$ there is $\delta$ of
trivial cyclotomic character with $\mathcal W$ containing almost all Frobenius data with Frobenius in
some neighbourhood of every point of a nonempty open subset of a neighbourhood of $\delta^{\ell^\nu}
g$. Then there are $J$, $\prity$, independent of $\nu$, and for every $\nu>0$ places $w_j$ in
$\mathcal W$, above distinct primes of $E^+$, with $\auNm w_j\equiv1\pmod{\ell^\nu}$, forming a
triangle for $(a^{-1},h\,h_0,h^{22}h_0)$.
\end{proposition}

\begin{proof}
\emph{Choice of $a$ and $h_0$}. The image of $\pgcyc$ on $\pgG$ is a closed subgroup of
$\Z_\ell^\times$, which is pro-cyclic for odd $\ell$; for $\ell=2$ it lies in $1+4\Z_2$ (as $E$
contains $\mu_{5040}\supseteq\mu_4$), which is pro-cyclic. So $\pgG/\ker\pgcyc$ is pro-cyclic and has a
topological generator $a$. Every $\pgcyc(g)^e-1$ is a multiple of $\pgcyc(a)^e-1$: write
$\pgcyc(g)=\pgcyc(a)^x$ with $x\in\Z_\ell$ and put $y=\pgcyc(a)^e$. If $\ell$ is odd and
$y\not\equiv1\bmod\ell$, then $y-1$ is a unit. Otherwise ($y\equiv1\bmod\ell$ for odd $\ell$, and
$y\equiv1\bmod4$ for $\ell=2$) one has $\prival(y^x-1)=\prival(y-1)+\prival(x)\ge\prival(y-1)$. If $\ell>h^2$,
then $\ell\nmid q_1d$ (as $q_1d\le h^2$), so an inertia element at $\ell$ whose cyclotomic character
reduces to a generator of $\auF_\ell^\times$ fixes $E\subseteq\Q(\mu_{q_1d})$; since $t<\ell-1$,
$\pgcyc^e-1$ is a unit at it for $0<e<t$, and $h_0=0$. Otherwise a Frobenius element at a chosen place
above $p_i\ne\ell$ has $\pgcyc=p_i$, and $p_i^e-1$ divides $N$ for $0<e\le t\le rt$ (the order of
$\auGL_{rt}(\auF_{p_i})$); take $h_0=\pgv(N)$.

\emph{Separation.} With $u=\pgcyc(a^{-1})$: $u^e-u^k=u^k(u^{e-k}-1)$, a unit times a divisor of
$\ell^{h_0}$, so the product over $k\ne e$, $k<t'$, divides $\ell^{(t'-1)h_0}$; and $4\mid\ell^{h_0}$
unless $h_0=0$, in which case $u^e-1$ is a unit and $\ell$ is odd. As $t+2\le h$ these divide
$\ell^{hh_0}$. This is (ii).

\emph{Test points.} Proposition~\ref{pri:kummer}(b) for $g_0=a^{-1}$ gives $s$, independent of
$\nu$. $J$ is the set of pairs $(e,g)$, $0<e<t'$, $g\in s$, ordered lexicographically with $e$
first, and $\prity(e,g)=e$; so $\#J+\#s+1\le(t+1)(\#s+1)$. If $h_0=0$, $(\#J+\#s+1)hh_0\le h^{22}h_0$
is trivial; if $h_0\ne0$, then $\ell\le h^2$ by hypothesis and the size bound of
Proposition~\ref{pri:kummer}(b) gives $(\#J+\#s+1)hh_0\le h^{22}h_0$.

\emph{Induction}.
Suppose places $w_i$ ($i<j$) with Kummer maps $\kappa_i$ of exponent $\prity(i)$ have been chosen,
the reductions of the $\kappa_i$ with $\prity(i)=e:=\prity(j)$ being linearly independent on $G_\nu$.
By Proposition~\ref{pri:kummer}(a), with $x$ the earlier places of exponent $e$ and $(f_i)$ the
reductions of their Kummer maps, there is a Kummer map $\kappa_0$ of exponent $e$ for $x$, off
$T_j=S_0(x)$, whose reduction is outside the span of the earlier ones. Since the earlier maps of exponent $e$ together with $\kappa_0$ are
additive on $G_\nu$ with independent reductions, they are jointly surjective onto $(\Z/\ell^\nu)^m$
(Nakayama); and conjugation by $a$ followed by a power multiplies a map of exponent $d$ on $G_\nu$
by $\pgcyc(a)^{d}$ (the equivariance), so composing with $\prod_{k\ne e}(\mathrm{conj}_a-\pgcyc(a)^k)$
kills the maps of the other exponents. This gives $u\in G_\nu$ with $\kappa_i(u)=0$ for $i<j$ and
$\kappa_0(u)=\prod_{k\ne e}(\pgcyc(a)^e-\pgcyc(a)^k)$, which divides $\ell^{hh_0}$ by (ii). Let
$g=g_j$ be the test point of $j$ and $\sigma=\delta^{\ell^\nu}\,g\,g^c\,u\,(u^c)^{-1}$. A Kummer map
$\kappa$ of exponent $e$ takes near $\sigma$ the value $2\kappa(u)$ (it is additive on $G_\nu$,
negated by $c$, and $\ell^\nu\kappa(\delta)=0$), and every even class $z$ of exponent $e$ unramified
outside $\ell$ takes near $\sigma$ the value $2z(g)$ ($z$ even: the values at $g,g^c$ add up and
those at $u,u^c$ cancel, as $u\in G_\nu$ acts trivially). By Chebotarev (the hypothesis on $\mathcal
W$) choose a degree-one place $w_j\in\mathcal W$ with Frobenius $\phi_j$ near $\sigma$, above a
prime of $F$ not below the earlier places and not in the $T_i$, $i\le j$. Then $\kappa_i(\phi_j)=0$
for $i<j$, $2\kappa_j(\phi_j)=4\kappa_0(u)=4\prod_{k\ne e}(\pgcyc(a)^e-\pgcyc(a)^k)$, which divides
$\ell^{h_0}\ell^{(t'-1)h_0}\mid\ell^{hh_0}$ in $\Z/\ell^\nu$ (as $4\mid\pgcyc(a)^e-1\mid\ell^{h_0}$),
and $z(\phi_j)=2z(g)$. This gives (i) for $(e,g)$, and $\phi_j\in G_\nu$ gives $\auNm w_j\equiv1\pmod{\ell^\nu}$.

\emph{Property (iii)} . Let $z$ be an even class of exponent $e=\prity(j)$ for the $w_i$,
with $z(\tau_i)=0$ for $i<j$ and $z(\phi_i)=0$ for $i<j$ of exponent $e$. Vanishing at $\tau_i$ and
(by evenness) at $\tau_i^{-c}$ makes $z$ unramified at $w_i,\overline w_i$ for $i<j$ (a tame
generator generates the image of inertia on a module on which inertia acts trivially, and the
cyclotomic character is unramified there). Extend $z$ to $z^+$ on $G_F$ by $z^+(gc)=z(g)$; then
$2z^+$ is a cocycle for $(\pgcyc^+)^{-e}$ with the local conditions of
Proposition~\ref{pri:kummer}(a) at $S_0(x)=T_j$ (at $\ell$ because $g^2\in\pgG$ for $g\in\ker\pgcyc^+$
in a decomposition group, at the earlier places of exponent $e$ because $z$ is unramified there and
vanishes at their Frobenius elements), vanishing at complex conjugations, and unramified outside
$T_j$ and the primes below $w_i$, $i\ge j$ (for $\auO\ne\Z_\ell$ apply this to $f\circ z$ for the
additive maps $f\colon\auO/\ell^\nu\to\Z/\ell^\nu$, which separate points). The reciprocity law of $\kappa_j$ gives $\sum_{i\ge j}\kappa_j(\phi_i)\,
2z(\tau_i)=0$; since $\kappa_j(\phi_i)=0$ for $i>j$ (the later choices), $2\kappa_j(\phi_j)z(
\tau_j)=0$. Values of $z$ at inertia elements at $w_j$ are multiples of $z(\tau_j)$; hence every
multiple of $2\kappa_j(\phi_j)$ in $\Z/\ell^\nu$ kills $z(\tau_j)$, in particular $\ell^{hh_0}$.
\end{proof}

\begin{remark}
The details of the four steps above (choice of $a$ and $h_0$, test points, induction with
$T_j=S_0(x)$, property (iii)) are checked in Proposition~\ref{pku:triangle}. Note that ``even'' here
means invariant under $g\mapsto g^c$ with the local condition at $\ell$ of $\mathrm{Sha}$; this
differs in form from the description in Proposition~\ref{pg:triangle}.
\end{remark}

\begin{lemma}[Cost of a character]\label{pri:cost}
Let the $w_j$ form a triangle for $(a^{-1},h_1,c_1)$ as in Definition~\ref{pri:def-tri}, $0\le e<t'$,
and $\priZ$ a module of cocycles $\pgG\to(\auO/\ell^\nu)(-e)$ unramified outside $\ell$, the $w_j$
and the $\overline w_j$, such that $\pgcyc(a)^{-e}-1$ kills $z(a^{-1})$ for every member unramified
outside $\ell$. Let $b$ be $\prirl(\priZ,Z^{\mathrm{cr}})$ plus the length of $\mathrm{Sha}^-_{e,\nu}$.
Then $\priZ$ has triangle cost $b$ along $\{\prity=e\}$ with constant $\ell^{h_1}$, for $\tau'_j=
\tau_j^{-c}$, $\phi'_j=\phi_j^{-c}$.
\end{lemma}
\begin{proof}
Take $A$ the set on which the local condition at $\ell$ is imposed (elements of trivial cyclotomic
character of the decomposition groups above $\ell$, or inertia above $\ell$ if $e=0$); members of
$\priZ$ obeying the crystalline relations vanish on $A$ (they are combinations of crystalline cocycles
there, which vanish on $A$), so
$\prirl(\priZ,\priZ_0)\le\prirl(\priZ,Z^{\mathrm{cr}})$. Since $\tau_j,\phi_j$ act trivially on
$(\auO/\ell^\nu)(-e)$ (trivial resp.\ $\equiv1\bmod\ell^\nu$ cyclotomic character), $z(g^{-c})=
-z(g^c)$ there, so $z(\tau_j)=z(\tau'_j)$ iff the even part $z^+=z+z\circ c$ vanishes at $\tau_j$,
and likewise for $\phi$. For $e\ne0$, $z^+$ is an even class;
for $e=0$ the condition (a) is vacuous as $\prity>0$. (a) is then (iii) of the triangle applied to
$z^+$. (b): let $\priZ_1$ be as in Definition~\ref{pri:def-cost}. Its members are unramified outside
$\ell$ (vanishing at $\tau_j,\tau'_j$). The map $z\mapsto z\circ c-z$ sends $\priZ_1$ into
$\mathrm{Sha}^-_{e,\nu}$. Let $e\ne0$. If $z\circ c=z$, then $z$ is even, hence an even class of
exponent $e$ for $\emptyset$; as $\phi_j$ acts trivially, $z(\phi'_j)=z(\phi_j^{-c})=-z(\phi_j^c)=
-z(\phi_j)$, so $z(\phi_j)=z(\phi'_j)$ gives $2z(\phi_j)=0$ for $\prity(j)=e$, and by (i) of the
triangle $4z(g)=2z(\phi_j)=0$ for $g\in s$. As $d=\pgcyc(a)^{-e}-1$ kills $z(a^{-1})$, $4dz$ vanishes
at $a^{-1}$ and on $s$, hence $4dz=0$; and $4d\mid4d^2\mid\ell^{h_1}$ by (ii). So the kernel is
$\ell^{h_1}$-torsion and $\prirl(\priZ_1,\priZ_1[\ell^{h_1}])\le\prilen\mathrm{Sha}^-_{e,\nu}$. Let
$e=0$. The coefficients are trivial, so the members of $\priZ_1$ are homomorphisms; they are
unramified outside $\ell$ and vanish on $A$, the inertia groups above $\ell$; so
$\priZ_1\subseteq\mathrm{Sha}^-_{0,\nu}$ (the everywhere unramified homomorphisms,
Proposition~\ref{pri:odd}) and $\prirl(\priZ_1,\priZ_1[\ell^{h_1}])\le\prilen\priZ_1\le\prilen
\mathrm{Sha}^-_{0,\nu}$.
\end{proof}

\subsection{Ribet's element for the Galois representations}\label{pri:ss-columns}

\begin{proposition}[Ribet's element]\label{pri:columns}
Let $c$ be in lift setting for $\vartheta$, auxiliary with ramification at most $\mathrm{ram}$;
let the $w_j$ ($j\in J$) form a triangle for $(a,hh_0,h^{22}h_0)$ with $h_0\le\pgv(N)$ (so $a$ here
is the element written $a^{-1}$ in Proposition~\ref{pri:triangle} and Lemma~\ref{pri:cost}); let
$\rho_x$, $x\in Y$ finite, satisfy the conclusion (c) of Proposition~\ref{pri:BL} with $U=\{w_j\}$;
for $x$ in $\priX\subseteq Y$ let $\pgtr\rho_x=\psi\pgtr R_x+\sum_{e<t'}\chi^e$ with $R_x$ main summands
and $\psi R_x$ having a lattice on which some $\gamma_x$ acts by a scalar $\zeta_x$, $\zeta_x-1$ a
unit; let $\priJ$ be a trace module for the $\rho_x$ with $p=\ell^\nu$, $\tau_{(j,+)}=\tau_j$,
$\tau_{(j,-)}=\tau_j^{-c}$, likewise $\phi$, and $H_j(1)=1$; and at $\pi\in\priX$ let $\pglen
\pgSel_Q(\psi R_\pi,\gamma_\pi;K/\auO)\le\mathrm{sel}<\infty$ (normalised). Then there is a common Ribet
element at $\pi$ with value $d$ satisfying
\[
\pgv(d)+2\pgv(N)\le\mathrm{sel}+2t'\,\mathrm{ram}+r'\pgloc_\ell+e_\ell .
\]
\end{proposition}

\begin{proof}
Choose a large coefficient ring $\auO$ over which: Proposition~\ref{pri:loss}(b) holds; the lengths
of $\mathrm{Sha}^-_{e,\nu}$ are $\le\prival(\pgcodd)$ (Proposition~\ref{pri:odd}); the $H_j$ are
defined; and there are lattices $u_x$ in the $\rho_x$ which, for $x\in\priX$, are block diagonal $T_x
\oplus\bigoplus_{e<t'}\auO(-e)$ with $T_x=\psi u'_x$, $u'_x$ a lattice in $R_x$, $T_x$ spanning
$M_{r'}(\auO)$ (Lemma~\ref{pri:block}, whose hypotheses hold: the $\rho_x$ are semisimple and the $R_x$ satisfy (M4), Remark~\ref{pax:rem-block}), and at $\pi$ the lattice of the Selmer hypothesis. We check the hypotheses of Proposition~\ref{pri:ribet},
with $G=\pgG$, $\xi_e=\chi^e$, $\rho_x=T_x$:

(1), (2): by construction, $\chi(\gamma_x)=1$ since $\gamma_x$ is inertial at a place prime to
$\ell$. (3): by the trace module hypothesis, after doubling the relations by $g\mapsto g^{-c}$ (the
conjugate self-duality turns a relation at $(j,+)$ into one at $(j,-)$ up to multiples of $\ell^\nu
\pgtr$); $\chi^e(\phi_k)\equiv1\bmod\ell^\nu$, so
$H_j(\chi^e(\phi_k))\equiv H_j(1)=1$. (4): $C=\ell^{hh_0}$ by (ii) of the triangle. (6): from (c) of
Proposition~\ref{pri:BL}: there are places covering $Q$ and the places above $p_0p_1$ prime to
$\ell$, with tame generators $\delta_v$ and polynomials $o_v$ with $o_v(\delta_v)(\delta_v-1)=0$ on
every $u_x$ and $\sum_v\pgv(o_v(1))\le2\,\mathrm{ram}+8f\pgv(N)$ (
at $w\in Q$, $o=X-\psi(\delta)$ with $\pgv(o(1))=\pgv(1-\psi(\delta))$, counted for $w$ and $\overline
w$; at a chosen place, $o=o_w=\prod(X-\text{roots of the rescaled label})$ with $o(1)$ dividing $N$ by
the lift setting; at most $8f$ such places). At a chosen place the identity $o_w(\delta_v)(\delta_v-1)=0$
on $u_x$, for every $x\in Y$, is the semisimplicity clause of Proposition~\ref{pri:BL}(c): $\rho_x|_{I_w}$
is semisimple with eigenvalues among the roots of $o_w(X)(X-1)$, the monodromy of
$\operatorname{rec}(\Pi_w)$ being $0$ (paragraph ``No monodromy at a chosen $w$'' of the proof in
\S\ref{pax:ss-BL}); the characteristic polynomial alone would not suffice. (5): Ribet cocycles of $T_\pi$ with values in $K/\auO$
vanishing at the $\tau_k$ and $\gamma_\pi$ are unramified outside $\ell p_0p_1Q\overline Q$ (vanishing
at both $\tau_j$ and $\tau_j^{-c}$ kills inertia at $w_j$ and $\overline w_j$), so modulo $\ell^M$
their length is at most $\prics_M(T_\pi)$ plus the length of the Selmer group; by
Proposition~\ref{pri:loss}(a) and Lemma~\ref{pri:divisible}, $b_\rho=\mathrm{loss}_T+\pglen\pgSel$ with
$\mathrm{loss}_T=\sum_v\delta_v(T_\pi)$. (7): with $s'=\{a\}\cup s\cup\{\tau_k\}$, $\#s'\le2(\#J+\#s+1)$,
Lemma~\ref{pri:cost} gives triangle cost $b_e=\prics_\nu(\auO(-e))+\prilen\mathrm{Sha}^-_{e,\nu}$; the
hypothesis on $z(a^{-1})$ holds because a character cocycle with $c(\gamma_\pi)=0$ (which holds for
those unramified outside $\ell$, as $\gamma_\pi$ is inertial off $\ell$) has $(\chi^e(g)-1)c(a)=0$;
and $s'$ detects (Proposition~\ref{pri:kummer}(b), plus the $\tau_k$ to remove ramification at the
$w_j$).

Proposition~\ref{pri:ribet} gives $d$ with (normalising lengths by $e_\auO$)
\[
\pgv(d)\le\mathrm{loss}_T+\mathrm{sel}+\sum_{e<t'}\bigl(\prics_\nu(\auO(-e))+\pgv(\pgcodd)\bigr)+
t'\bigl(2\mathrm{ram}+8f\pgv(N)+4(\#J+\#s+1)hh_0\bigr)+(\#J+1)hh_0 .
\]
By Proposition~\ref{pri:loss}, $\mathrm{loss}_T+\sum_e\prics_\nu(\auO(-e))\le r'\pgloc_\ell$. With
$(\#J+\#s+1)hh_0\le h^{22}h_0\le h^{22}\pgv(N)$, $t'\le t$ and $h_0\le\pgv(N)$, the remaining terms plus
$2\pgv(N)$ are at most
\[
2t'\mathrm{ram}+t\pgv(\pgcodd)+8ft\pgv(N)+2\pgv(N)+\bigl((4t+1)h^{22}+h\bigr)\pgv(N)=2t'\mathrm{ram}+e_\ell,
\]
which is the definition of $e_\ell$ in~\eqref{pg:eq-beta} .
\end{proof}

\subsection{The operators at the auxiliary places}\label{pri:ss-norm}

Fix places $w_j$ ($j\in J$) off the base level with $\auNm w_j\equiv1\pmod{\ell^\nu}$, the forms of
principal level at the $w_j$, $\Gamma=\prod_j\Gamma_j$, $\Gamma_j=\auGL_n(\auF_{w_j})$,
$\mathrm{Nm}=\sum_{k\in\Gamma}k$, $\mathrm{Nm}_j=\sum_{k\in\Gamma_j}k$, and let $\priT$ be the Hecke
algebra away from the level, acting by integral operators commuting with $\Gamma$. Let $c_0$ be as
in Lemma~\ref{pri:forms}(a). A ring map $\sigma$ from the integral operators is a \emph{reduction
modulo $\ell^\nu/k$} if it kills every $T$ with $kT\in\ell^\nu\cdot(\text{integral operators})$. The
\emph{tied ideal} $\mathfrak a\subseteq\priT[X_J]$ is generated by the $\ell^\nu\pgtr[g]$ and the
$\pgtr[(\gamma_j-1-X_j(\phi_j-1)H_j(\phi_j))g]$, $\gamma_j$ the tame generator and $\phi_j$ a
Frobenius at $w_j$.

\begin{lemma}[One place]
\label{pri:oneplace}
Let $\delta'\in\priT$, $t'\ne0$, and suppose every base system on which $\delta'\ne0$ has at $w=w_j$
a rescaled Hecke polynomial $\prod_{i<t'}(X-(\auNm w)^i)$ times a divisor of $H=H_j$. Let $\sigma$ be a
reduction modulo $\ell^\nu/(2n!)$. There is an integral operator $\partial$ made of translates at $w$
such that for all $g$: $\mathrm{Nm}_j\pgtr[(\gamma-1)g]=0$; $\sigma(\partial\pgtr[(\gamma-1)g])=
\sigma(\mathrm{Nm}_j\pgtr[(\phi-1)H(\phi)g])$; and for each $k\in\Gamma_j$ there is an integral $Z$ with
$\sigma(c_0\partial(k-1))=\sigma(\mathrm{Nm}_jZ)$ and $\delta'\,\mathrm{Nm}\,Z=0$.
\end{lemma}
\begin{proof}
Let $K_1\subseteq\Gamma_j$ be the stabiliser of a line, $\mathrm{ind}_1(k)=\mathrm{ind}(k_{11})$ on
$K_1$, and $\mathrm{Tr}=\sum_hh$, $\mathrm{Tr}'=\sum_hh^{-1}$ over fixed representatives $h$ of
$\Gamma_j/K_1$. For a weight $f\colon\Z/\ell^\nu\to\Z$ put $L(f)=\sum_{k\in K_1}f(\mathrm{ind}_1(k))k$,
and with $p(X)=(X-(\auNm w)^{t'-1})H(X)$ let
\begin{equation}\label{pri:eq-B}
B(f)=p(\tilde U)\,L(f),\qquad\partial=-\mathrm{Tr}\,B(\lambda)\,\mathrm{Tr}',
\end{equation}
$\lambda$ the lift $\Z/\ell^\nu\to\Z$ of least absolute value; $B(f)$ is made of translates at $w$. The first claim is Proposition~\ref{pri:dict}(c).
For the second, modulo $\sigma$: (1) by Proposition~\ref{pri:dict}(b) (decompose under the
characters of $K_1$ and note $\chi(n!)\ne1$ is needed only where the weight is nonzero modulo the
reduction), $B(\lambda)\pgtr[(\gamma-1)g]\equiv L(\lambda)\pgtr[(\gamma-1)g\,p(\phi)]$; (2) the
transfer from $K_1$ to $\Gamma_j$ of $k\mapsto k_{11}$ is $\det$ (an explicit computation on diagonal
matrices and transvections), and for any homomorphism $\alpha$ to
$\Z/\ell^\nu$, $\mathrm{Tr}\,L_\alpha(\lambda)\,\mathrm{Tr}'\equiv L_{\alpha^{\mathrm{tr}}}(\lambda)$
modulo $\ell^\nu$ (the coefficient of $g$ on the left is a sum over representatives of the logarithms
of the components of $g$, congruent to the logarithm of the transfer); so $\mathrm{Tr}\,L(\lambda)\,
\mathrm{Tr}'\equiv L_{\det}(\lambda):=\sum_{k\in\Gamma_j}\lambda(\mathrm{ind}\det k)k$; (3) by
Proposition~\ref{pri:dict}(a), $L_{\det}(\lambda)\pgtr[(\gamma-1)g\,p(\phi)]=(\langle\gamma\rangle-1)
L_{\det}(\lambda)\pgtr[g\,p(\phi)]$; (4) $(\langle\gamma\rangle-1)L_{\det}(\lambda)\equiv-\mathrm{Nm}_j$
modulo $\ell^\nu$, as $\mathrm{ind}\det\langle\gamma\rangle=1$ shifts the weight by $1$ and
$\lambda(x-1)-\lambda(x)\equiv-1$; (5) $p(\phi)=(\phi-(\auNm w)^{t'-1})H(\phi)\equiv(\phi-1)H(\phi)$ as
$\auNm w\equiv1$. Together, $\partial\pgtr[(\gamma-1)g]\equiv\mathrm{Nm}_j\pgtr[(\phi-1)H(\phi)g]$.
For the third claim we use four facts.

(i) \emph{$B(1)$ is fixed by $K_1$ on the left}: $\kappa B(1)=B(1)$ for $\kappa\in K_1$. Indeed the image
of $L(1)$ consists of $K_1$-invariant forms; these are fixed by the pro-$p$ Iwahori subgroup $I_1$
(whose image in $\Gamma_j$ is the upper unipotent group, contained in $K_1$); and $\tilde U$ preserves
the $K_1$-invariant vectors in the $I_1$-invariants, by the argument in the proof of
Lemma~\ref{pax:Ueig} showing that $U$ preserves the vectors of type $(\mathcal P,\tilde\chi(k_{11}))$
(that argument uses only the decomposition of $\mathcal P$ and of $I_1dI_1$, and applies verbatim to
the trivial character). Hence $p(\tilde U)L(1)$ has $K_1$-invariant image. In particular
$L(1)B(1)=\#K_1\cdot B(1)$.

(ii) \emph{Right translation shifts the weight.} As $K_1$ stabilises a line, $k\mapsto k_{11}$ and
hence $\mathrm{ind}_1$ are homomorphisms on $K_1$. For $\kappa\in K_1$ with $a=\mathrm{ind}_1(\kappa)$,
substituting $k'=k\kappa$ gives $L(f)\kappa=\sum_{k'}f(\mathrm{ind}_1(k')-a)k'=L(f(\cdot-a))$.
Since $\lambda(x-a)\equiv\lambda(x)-\lambda(a)\pmod{\ell^\nu}$, the weight $g_a=\ell^{-\nu}(\lambda(\cdot
-a)-\lambda+\lambda(a))$ is $\Z$-valued, and $B(\lambda)(\kappa-1)=-\lambda(a)B(1)+\ell^\nu p(\tilde
U)L(g_a)$, the last term being $\ell^\nu$ times an integral operator, hence killed by $\sigma$.

(iii) \emph{Rewriting $\mathrm{Tr}'(k-1)$.} For each representative $h$, $k^{-1}h$ lies in a unique
coset $h'K_1$; write $h^{-1}k=\kappa_hh'^{-1}$ with $\kappa_h\in K_1$. As $h\mapsto h'$ permutes the
representatives, $\mathrm{Tr}'(k-1)=\sum_h(\kappa_h-1)h'^{-1}$. With $a_h=\mathrm{ind}_1(\kappa_h)$,
(ii) gives
\[
\sigma\bigl(\partial(k-1)\bigr)=\sigma\Bigl(\sum_h\lambda(a_h)\,\mathrm{Tr}\,B(1)\,h'^{-1}\Bigr).
\]

(iv) \emph{Inserting $c_0$.} Lemma~\ref{pri:forms}(a) with $\Gamma'=K_1$ gives an integral $\Psi$ with
$\sum_{\kappa\in K_1}\kappa^{-1}\Psi\kappa=c_0$. By (i), $c_0B(1)=\sum_\kappa\kappa^{-1}\Psi B(1)=
L(1)\Psi B(1)$, and $\mathrm{Tr}\,L(1)=\mathrm{Nm}_j$. Put
\[
Z=\sum_h\lambda(a_h)\,\Psi\,B(1)\,h'^{-1},
\]
an integral operator. Then $\sigma(c_0\partial(k-1))=\sigma(\mathrm{Nm}_jZ)$ by (iii).

It remains to show $\delta'\mathrm{Nm}Z=0$. Write $\mathrm{Nm}=\mathrm{Nm}^{(j)}\mathrm{Nm}_j$, with
$\mathrm{Nm}^{(j)}$ the product of the norms at the other places; it commutes with the operators at
$w$ and with $\delta'$. By the computation in (iv), $\mathrm{Nm}_j\Psi B(1)=\mathrm{Tr}\,L(1)\Psi B(1)=
c_0\,\mathrm{Tr}\,B(1)$. By (i), $\mathrm{Nm}\,B(1)=\mathrm{Nm}^{(j)}\mathrm{Tr}\,L(1)B(1)=\#K_1\cdot
\mathrm{Nm}^{(j)}\mathrm{Tr}\,B(1)$, and $\delta'\mathrm{Nm}B(1)=0$ by Proposition~\ref{pri:dict}(d).
As the operators act on a vector space over a field of characteristic $0$,
$\delta'\mathrm{Nm}^{(j)}\mathrm{Tr}\,B(1)=0$. Hence
$\delta'\mathrm{Nm}Z=c_0\sum_h\lambda(a_h)\,\delta'\mathrm{Nm}^{(j)}\mathrm{Tr}\,B(1)\,h'^{-1}=0$.
(The factor $2n!$ in $\sigma$ is needed only for the second claim.)
\end{proof}

\begin{proposition}[The deformed norm; exact form of Proposition~\ref{pg:norm}]\label{pri:norm}
Let $\delta'\in\priT$, $\delta=c_0^2\delta'$, $t'\ne0$, and suppose the hypothesis of
Lemma~\ref{pri:oneplace} at every $w_j$ (with $H_j$). Let $\sigma$ be a reduction modulo
$\ell^\nu/(2n!)$. There are integral operators $\priN_S$ ($S\subseteq J$), commuting with $\priT$,
and $Y_S$, with $\priN_\emptyset=\mathrm{Nm}$, $Y_\emptyset=\delta^{\#J}$, such that
\begin{enumerate}
\item[(a)] for every $G\in\mathfrak a$, $\sigma\bigl(\sum_S\priN_{J\smallsetminus S}\,G_S\bigr)=0$,
where $G_S$ is the coefficient of $X_S$ in $G$;
\item[(b)] $\sigma(\delta^{\#J}\priN_S)=\sigma(Y_S\,\mathrm{Nm})$ for every $S$.
\end{enumerate}
Consequently: let $\prid\in\priT[X_J]$ with $(\delta_0-n_0)\prid\in\mathfrak a$ for some $\delta_0\in
\priT$ and $n_0\in\pgZb$, let $A\in\priT$ with $A^2=A$, and suppose the hypothesis above holds for
$\delta'=A\delta_0$. Let $l$ be a form of base level with $d_1l$ integral, and $y$ a form with
$A\delta_0l=n_0Ay$, $\delta_0y=n_0y$ and $\prid_Sy=0$ for $S\ne J$. If $2n!(n_0c_0^2)^{\#J+1}$ divides
$\ell^\nu$, then $d_1c_0A\prid_Jy$ is integral.
\end{proposition}

\begin{proof}
For each $j$ let $\partial_j$ be the operator of Lemma~\ref{pri:oneplace}. Operators made of
translates at different places commute, and commute with $\priT$. Let $\priN_S$ be the coefficient of
$X_S$ in $\prod_j(\mathrm{Nm}_j+X_j\partial_j)$, i.e.\ $\priN_S=\prod_{j\in S}\partial_j\prod_{j\notin
S}\mathrm{Nm}_j$; then $\priN_\emptyset=\mathrm{Nm}$, and for $j\notin S$, $\priN_S=\priN'\mathrm{Nm}_j$,
$\priN_{S\cup\{j\}}=\priN'\partial_j$ with the same $\priN'$.

(a) It suffices to treat $G=G'G_0$ with $G_0$ a generator of $\mathfrak a$ and $G'$ arbitrary
(the expression is additive in $G$, and $\mathfrak a$ is spanned by such products). For $G_0=
\ell^\nu\pgtr[g]$, $\sigma$ kills it. For $G_0=t_0-X_jr_0$ with $t_0=\pgtr[(\gamma_j-1)g]$, $r_0=
\pgtr[(\phi_j-1)H_j(\phi_j)g]$: the coefficient of $X_S$ in $G'(t_0-X_jr_0)$ is $G'_St_0$ if $j\notin S$
and $G'_{S}t_0-G'_{S\smallsetminus j}r_0$ if $j\in S$. Pair $S\not\ni j$ with $S\cup\{j\}$: the
complements are $S^c\ni j$ and $S^c\smallsetminus j$, so
\[
\priN_{S^c}G'_St_0+\priN_{S^c\smallsetminus j}(G'_{S\cup j}t_0-G'_Sr_0)=\priN''\bigl(\partial_j
G'_St_0+\mathrm{Nm}_jG'_{S\cup j}t_0-\mathrm{Nm}_jG'_Sr_0\bigr),
\]
and $\mathrm{Nm}_jt_0=0$, $\sigma(\partial_jt_0)=\sigma(\mathrm{Nm}_jr_0)$ (Lemma~\ref{pri:oneplace};
the $G'_S\in\priT$ commute with $\partial_j$). So the sum is $\equiv0$.

(b) is Proposition~\ref{prm:normb}. (The induction must be run for all $S$ simultaneously and at
every place of $S$; the operators $Y_S$ are built from the operator $\Psi$ of
Lemma~\ref{pri:forms}(a), which is not a Hecke operator, so they need not commute with $\priT$, and
this is not used below; see Remark~\ref{prm:rem-norm}.)

\emph{Consequence}. Put $n_1=n_0c_0^2$,
$i_0=\#J+1$ and let $G$ be the top coefficient of $\priN\prid$, i.e.\ $G=\sum_S\priN_{J\smallsetminus
S}\prid_S$. Since $(\delta_0-n_0)\prid\in\mathfrak a$, (a) gives $\sigma(G\delta_0)=\sigma(n_0G)$
after multiplying by the commuting $A$; more precisely $\sigma(\delta G)=\sigma(n_1AG)$ with
$\delta=c_0^2A\delta_0$. Then
\[
\sigma(n_1^{i_0}A G)=\sigma(\delta^{i_0}G)=\sigma\Bigl(\sum_S\delta^{\#J}\priN_{J\smallsetminus
S}\,\delta\prid_S\Bigr)=\sigma\Bigl(\sum_SY_{J\smallsetminus S}\,\mathrm{Nm}\,\delta\prid_S\Bigr).
\]
Apply both sides to the integral form $u=\Psi(d_1l)$, whose norm is $c_0d_1l$
(Lemma~\ref{pri:forms}(a)). Since $\mathrm{Nm}\,\delta\prid_S u=\delta\prid_S(c_0d_1l)=c_0^3d_1
\prid_SA\delta_0l=c_0^3d_1n_0\prid_SAy$, which vanishes for $S\ne J$, only $S=J$ survives, and
$Y_\emptyset=\delta^{\#J}$ acts on $Ay$ as $n_1^{\#J}A$ ($\delta_0y=n_0y$). Hence
$n_1^{i_0}AGu\equiv n_1^{i_0}d_1c_0A\prid_Jy$ modulo $\ell^\nu/(2n!)$ times integral forms; the left
side is $n_1^{i_0}$ times an integral form. As $2n!\,n_1^{i_0}\mid\ell^\nu$, cancelling $n_1^{i_0}$
shows that $d_1c_0A\prid_Jy$ is integral.
\end{proof}


\subsection{Proof of Theorem~\ref{pri:III}}\label{pri:ss-proof}

\begin{lemma}[Orthogonality to the false lifts]\label{pri:weak}
Let $y$ be orthogonal to the false lifts, $c$ in lift setting with $e_c\le e$, and suppose
$\mathcal Z$ is stable under $e_c$ and the Hecke operators at split places off the level. Then
$e_c(y-\pgpr_{\mathcal Z}y)$ lies in the span of the Hecke eigenforms in $e_c\mathcal F(O)$ of base
level that are not weak lifts.
\end{lemma}
\begin{proof}
By Lemma~\ref{pri:chain}, applied to the finitely many weak-lift systems of base level, at all but
finitely many split $w\nmid\ell$ the polynomial $H_w(\cdot;X)$ of each of them vanishes at $\lambda_w$
(the lift root, the largest root of the rescaled chain); fix one such $w$. The Hecke operators are
normal and commute with $e_c$ (which is self-adjoint, being a central idempotent of a unitary group
action); $\mathcal Z^\perp$ is stable under them (the adjoint of a Hecke operator at $w$ is one at
$\overline w$). Let $y'=e_c(y-
\pgpr_{\mathcal Z}y)\in e_c\mathcal F\cap\mathcal Z^\perp$, and $y'_x$ its orthogonal component in a
weak-lift system $x$. Then $y'_x\in e\mathcal F\cap\mathcal Z^\perp\cap\ker H_w(\lambda_w)$ is a false
lift, so $\langle y,y'_x\rangle=0$; also $\langle\pgpr_{\mathcal Z}y,y'_x\rangle=0$; hence
$\|y'_x\|^2=\langle y',y'_x\rangle=\langle y-\pgpr_{\mathcal Z}y,e_cy'_x\rangle=0$.
\end{proof}

\begin{lemma}[Taylor--Wiles places with root chains]\label{pri:TW}
In the setting of Theorem~\ref{pri:III} let $x_0$ be the system of a lift generator, with main
summand $R$. There are $J$, $\prity$, and for every $\nu>0$ places $w_j$ forming a triangle as in
Proposition~\ref{pri:triangle}, off $\pgS(\ell)\cup Q\cup\overline Q$ and the base level, and
polynomials $e_j\in\pgZb[X]$ with $e_j(1)=1$, such that every weak-lift base system $x$ whose
normalised characteristic polynomials are congruent modulo the maximal ideal to those of $x_0$ has,
at $w_j$, rescaled Hecke polynomial $\prod_{i<t'}(X-(\auNm w_j)^i)$ times a divisor of $e_j$.
\end{lemma}
\begin{proof}
Let $\mathcal W$ be the set of Frobenius data $(w,\phi)$ with $w$ off the bad places and the base
level, $\det(1-\psi(\phi)R(\phi))$ a unit, and such that for every weak-lift base system $x$ and every
$S$ with $S(1)$ a unit and $\det(X-\rho_x(\phi))\equiv S\cdot\prod_{i<t'}(X-\chi(\phi)^i)$, the rescaled
Hecke polynomial at $w$ is the chain times a divisor of an integral $e$ with $e(1)=1$. \emph{$\mathcal
W$ satisfies the hypothesis of Proposition~\ref{pri:triangle}}: by (M5) an inertial $\gamma$ above $p_0p_1$
acts on $R$ by a scalar $y$ with $y^{\ell^aj}-1$ a unit for $j\le4r'+2$, and is killed by $\psi$ and
$\pgcyc$; among $\gamma^i\delta^{\ell^\nu}g$, $i\le4r'+2$, one has $\det(1-\psi R)$ a unit on a
neighbourhood (for each eigenvalue $\mu$ of $\psi R(\delta^{\ell^\nu}g)$ at most one $i$ has $y^{i
\ell^\nu}\mu\equiv1$); the other conditions hold for almost all $w$ with Frobenius in a suitable open
set ((G1) in the ``most Frobenius'' form, and Lemma~\ref{pri:chain}: the chain divides the rescaled
Hecke polynomial $\iota^{-1}$-wise, with quotient $T$; then $\det(X-\rho_x(\phi))=\text{chain}\cdot
\iota T\equiv S\cdot\text{chain}$, monic factors cancel modulo the maximal ideal, so $\iota T\equiv S$,
$\iota T(1)$ is a unit and $e=T/T(1)$ works). Apply Proposition~\ref{pri:triangle}. For $x$ congruent
to $x_0$: $\det(X-\rho_{x_0}(\phi))=\det(X-\psi R(\phi))\cdot\prod_{i<t'}(X-\chi(\phi)^i)$ (equality of
traces of all powers, \eqref{pg:eq-trace}); so the condition in $\mathcal W$ applies with
$S=\det(X-\psi R(\phi_j))$. Take $e_j$ the product of the finitely many $e$ obtained.
\end{proof}

\begin{proof}[Proof of Theorem~\ref{pri:III}]
\emph{Step 1: one generator.} Let $v_0\ne0$ be a lift generator with Selmer bound $\mathrm{sel}=
e'_\ell$, of type $c$ and system $x_0$; let $\priX$ be the set of systems containing a nonzero lift
generator (with any Selmer bound), $Y$ the base systems. By Proposition~\ref{pri:BL} the $\rho_x$,
$x\in Y$, have integral characteristic polynomials and satisfy (c). Choose (Proposition~\ref{pri:triangle}
via Lemma~\ref{pri:TW}) $J$, and finitely many $g_k\in\pgG$ and nonzero integral $\mathrm{val}_k$ with
$\prod_k(\mathrm{val}_k-\det(1-\rho_x(g_k)))=0$ at every base system $x$ that is not a weak lift (for
each such $x$ pick $g$ with $\det(1-\rho_x(g))\ne0$). Let $\delta_0=\prod_k(\mathrm{val}_k-\det(1-
\rho(g_k)))\in\priT$, $n_0=\prod_k\mathrm{val}_k$; on $\priX$ each $\det(1-\rho_x(g))=0$ (the trivial
character is a summand), so $\delta_0=n_0$ there. Choose $\nu$ with $2n!(n_0N^2)^{\#J+1}\mid\ell^\nu$,
and places $w_j$ for this $\nu$ by Lemma~\ref{pri:TW}; let $A\in\priT$ be the idempotent of the
systems congruent to $x_0$ (it exists in the Hecke algebra: the product of differences of separating
coefficients). Then every base system with $A\delta_0\ne0$ is a
weak lift congruent to $x_0$, so has the root chain at each $w_j$ with $H_j=e_j$.

Now that the $w_j$ are fixed, replace $Y$ by the finite set of systems of eigenvalues of $\priT$
occurring in base level with principal level at the $w_j$; it contains $\priX$ (forms of base level
are of principal level at the $w_j$), and Proposition~\ref{pri:BL}(c) with $U=\{w_j\}$ applies to it,
as Proposition~\ref{pri:columns} requires. Each $x\in Y$ is a ring map $\priT\to\overline{\Q}_\ell$
(through $\iota$). Let $S$ be the image of $\priT[X_J]$ under $\epsilon\colon\priT[X_J]\to
\operatorname{Map}(Y,\overline{\Q}_\ell[X_J])$, $\epsilon(t)(x)=x(t)$ coefficientwise, and $\priJ=
\epsilon(\mathfrak a)$. $S$ contains the $X_j$ and the characteristic polynomial coefficients (they
lie in $\priT$, being limits of Hecke-polynomial coefficients), and $\priJ$ is an $S$-submodule
containing the families of the definition (the generators of $\mathfrak a$); so $\priJ$ is a trace
module. The map $\epsilon$ is injective: $\priT$ is a commutative algebra generated by normal
operators on a finite-dimensional space, hence simultaneously diagonalisable, so an element acting
by $0$ on every system is $0$. By Proposition~\ref{pri:columns} (with $\mathrm{ram}=e'_\ell$,
$\mathrm{sel}=e'_\ell$) there is a common Ribet element $f\in S$ at $x_0$ with value $d$,
$\pgv(N^2d)\le e'_\ell+2t'e'_\ell+r'\pgloc_\ell+e_\ell$. Write $f=\epsilon(\prid)$ with $\prid\in
\priT[X_J]$, unique by injectivity. For each $g$, $\det(1-\rho(g))\in\priT$ and $f\cdot\det(1-\rho(g))
=\epsilon(G)$ for some $G\in\mathfrak a$; injectivity gives $\prid\det(1-\rho(g))=G\in\mathfrak a$.
Hence
$(\delta_0-n_0)\prid\in\mathfrak a$ ($\delta_0-n_0$ lies in the ideal generated by the
$\det(1-\rho(g_k))$), $\prid_S=0$ on $\priX$ for $S\ne J$, and $\prid_J(x_0)=d$.

Let $y$ be integral and orthogonal to the false lifts, $l=e_cy$ and $y_1=e_c\pgpr_{\mathcal Z}y$. By
Proposition~\ref{pri:BL}(b), $Nl$ is integral; $y_1$ is supported on $\priX$, so $\delta_0y_1=n_0y_1$
and $\prid_Sy_1=0$ for $S\ne J$; and by Lemma~\ref{pri:weak}, $l-y_1$ is supported on non-weak-lift
systems, where $A\delta_0=0$, so $A\delta_0l=A\delta_0y_1=n_0Ay_1$. Proposition~\ref{pri:norm} (with
$c_0\mid N$, so that $2n!(n_0c_0^2)^{\#J+1}$ divides $2n!(n_0N^2)^{\#J+1}$ and hence $\ell^\nu$, and
$d_1=N$) shows that $Nc_0A\prid_Jy_1$ is integral, hence so is $N^2A\prid_Jy_1=(N/c_0)\cdot Nc_0A
\prid_Jy_1$, $N/c_0$ being an integer. Define $G(z)=N^2A\prid_J e_cz$; it acts
on each lift generator $v$ of type $c$ in system $x$ by the scalar $N^2A(x)\prid_J(x)$, on lift
generators killed by $e_c$ by $0$, and $G(\pgpr_{\mathcal Z}y)$ is integral; its eigenvalue at $v_0$
has $\pgv\le\mathrm{bound}:=e'_\ell+2t'e'_\ell+r'\pgloc_\ell+e_\ell$.

\emph{Step 2: all generators at once}. Choose a basis $b_1,\dots,b_m$ of $\mathcal Z$ among lift
generators. Step~1 for $v_0=b_p$ gives an operator $G_p$, diagonal in this basis, with eigenvalues
$l_{pk}$, $\pgv(l_{pp})\le\mathrm{bound}$, and with $G_p\pgpr_{\mathcal Z}y$ integral. The set of
operators with this integrality property is closed under integral linear combinations. For a square
array $(l_{pk})$ there are integral $d_p$ with $|\sum_pl_{pk}d_p|_\ell\ge|l_{kk}|_\ell$ for every $k$
(choose $d$ outside the finitely many hypersurfaces where some linear form $\sum_pl_{pk}X_p$ is smaller
than its largest coefficient; the residue field of $\pgZb$ is infinite). So $\Lambda_0=\sum_pd_pG_p$
is diagonal with eigenvalues $\mu_k$, $\pgv(\mu_k)\le\mathrm{bound}$, and preserves the Hecke
eigenspaces.

\emph{Step 3: the determinant.} Put $\Lambda=N\Lambda_0$. Its eigenvalues have $\pgv\le\mathrm{bound}+
\pgv(N)\le r'\beta_\ell$: indeed, with $r'\ge1$, $t'\le t$, $\mathrm{sel},\mathrm{ram}\le e'_\ell$,
\[
e'_\ell+2t'e'_\ell+r'\pgloc_\ell+e_\ell+\pgv(N)\le r'\bigl(\pgloc_\ell+(1+2t)e'_\ell+e_\ell+\pgv(N)
\bigr)=r'\beta_\ell .
\]
So $\pgv(\det\Lambda)\le\dim\mathcal Z\cdot r'\beta_\ell$. $\Lambda\pgpr_{\mathcal Z}y$ is integral at
$\iota$, and by Lemma~\ref{pri:forms}(b), $\langle y',\Lambda\pgpr_{\mathcal Z}y\rangle=N\langle y',
\Lambda_0\pgpr_{\mathcal Z}y\rangle\in\pgZb$ for $y'$ integral at every isomorphism. If $\mathcal Z=0$
there is nothing to prove.
\end{proof}

\subsection{The inputs of Theorem~\ref{pg:III}}\label{pri:ss-summary}

Theorem~\ref{pri:III}, hence Theorem~\ref{pg:III} (with the parity clause of Remark~\ref{pri:rem-compare}),
is proved from: (i) Proposition~\ref{pri:loss} (the local loss, Proposition~\ref{pg:loss}), proved in
\S\ref{pll:sec} with Theorem~\ref{pll:exponent} ({\cite[\S2.3]{P56}}); (ii) statements (G1)--(G5) of
\S\ref{pg:ss-galois}, proved in \S\ref{pga:sec} from published theorems; (iii) Chebotarev's theorem, the
Greenberg--Wiles formula and the other published results quoted in \S\S\ref{pod:sec}, \ref{pku:sec},
\ref{pax:sec}. The other inputs are proved as follows: Proposition~\ref{pri:ribet} in \S\ref{prm:ss-chars},
Proposition~\ref{pri:odd} as Proposition~\ref{pod:odd}, Proposition~\ref{pri:kummer} in \S\ref{pku:sec},
Proposition~\ref{pri:BL} in \S\ref{pax:ss-BL} and \S\ref{prm:ss-level}, Lemmas~\ref{pri:block},
\ref{pri:forms}, \ref{pri:chain} in \S\ref{pax:sec} and \S\ref{pay:sec}, Proposition~\ref{pri:dict} in
\S\ref{pax:ss-dict} and \S\ref{pay:ss-cons}, Proposition~\ref{pri:norm}(b) with Proposition~\ref{prm:normb},
and Lemma~\ref{pri:amitsur} as Proposition~\ref{pax:amitsur}.




\providecommand{\auO}{\mathcal O}
\providecommand{\auQlb}{\overline{\Q}_\ell}
\providecommand{\panot}[1]{\par\smallskip{\raggedright\noindent\textbf{[Not proved here:} #1\textbf{]}\par}\smallskip}
\providecommand{\pgv}{v}
\providecommand{\pgloc}{\mathrm{loc}}
\providecommand{\pgcyc}{\varepsilon_\ell}
\providecommand{\pglen}{\operatorname{len}}
\providecommand{\pgHom}{\operatorname{Hom}}
\providecommand{\pllZf}{Z_f}
\providecommand{\pllperp}{\perp}
\providecommand{\pllrk}{\operatorname{rk}}
\providecommand{\plldef}{\delta^{\vee}}
\providecommand{\pllB}{B_{\mathrm{cris}}}
\providecommand{\pllD}{\mathrm{dB}}

\section{The local loss}\label{pll:sec}

This section concerns Proposition~\ref{pg:loss}. We first fix exactly what is asserted
(\S\ref{pll:ss-statement}), then prove: the reduction, by local Tate duality, of the defect to a
\emph{dual defect} defined by cup products (\S\ref{pll:ss-duality}); the passage from an exponent
killing cup products to a bound for the dual defect (\S\ref{pll:ss-exponent}); a remark on another route for large
primes, through the theory of Fontaine and Laffaille (\S\ref{pll:ss-FL}); and the sum over places and
summands (\S\ref{pll:ss-sum}); the explicit exponents for the other primes
(\S\ref{pll:ss-gap}) are proved in {\cite[\S\S2--9]{P56}}.

The logical structure is as follows. Proposition~\ref{pll:dual} (which uses local Tate duality and
the results of Bloch and Kato), Lemma~\ref{pll:exp} and Lemma~\ref{pll:sum} reduce
Proposition~\ref{pg:loss}, for every prime $\ell$, to the explicit exponent bounds of
Theorem~\ref{pll:exponent}, which are proved in {\cite[\S2]{P56}} from
{\cite[\S\S3.1--9]{P56}}; the deduction is carried out in Remark~\ref{pll:status}.

\emph{Conventions.} As in \S\ref{pg:sec}. Fix a prime $\ell$, a place $v\mid\ell$ of $E$, and write
$\Gamma=D_v$, identified with $\operatorname{Gal}(\auQlb/E_v)$, an open subgroup of
$G_{\Q_\ell}=\operatorname{Gal}(\auQlb/\Q_\ell)$ of index $[E_v:\Q_\ell]$. $\auO$ is a coefficient
ring with uniformiser $\varpi$; lengths are normalised ($\pglen\auO/\ell=1$), so that the length of
a finite $\auO$-module $X$ is $\log_\ell\#X/[\auO:\Z_\ell]$. A \emph{lattice} is a finite free $\auO$-module $T$
with a continuous $\auO$-linear action of $\Gamma$; $V=T[1/\ell]$; $T^*=\pgHom_\auO(T,\auO(1))$ is
its \emph{Tate dual}. The Hodge--Tate weights (``jumps'') of a potentially crystalline $V$ are
indexed by the embeddings $\gamma\colon\auO\to\C_\ell$, as in \S\ref{pg:sec}.

\subsection{The statement}\label{pll:ss-statement}

\begin{definition}\label{pll:def}
Let $T$ be a lattice and $\pllB$ Fontaine's crystalline period ring.
\begin{enumerate}
\item The \emph{crystalline cocycles} of $T$ are the continuous maps $c\colon\Gamma\to T$ for which some
$y\in\pllB\otimes_{\Z_\ell}T$ satisfies $1\otimes c(g)=gy-y$ for all $g\in\Gamma$; they form an
$\auO$-module $\pllZf(T)$, which contains the coboundaries $g\mapsto gx-x$, $x\in T$. (These are
the continuous cocycles whose class lies in $H^1_f(\Gamma,T)$, the preimage of the Bloch--Kato
subspace $H^1_f(\Gamma,V)=\ker(H^1(\Gamma,V)\to H^1(\Gamma,\pllB\otimes V))$.)
\item A map $s\colon\Gamma\to T/\ell^M$ \emph{obeys the crystalline relations} of $T$ if every
finite $\auO$-linear relation $\sum_g\lambda_g\,c(g)=0$ ($\lambda_g\in\operatorname{End}_\auO T$)
that holds for all $c\in\pllZf(T)$ also holds for $s$. The reductions modulo $\ell^M$ of the
elements of $\pllZf(T)$ obey them; we write $\overline{\pllZf}(T)_M$ for the module of these
reductions.
\item Let $\mathcal C$ be a class of lattices. $s\colon\Gamma\to T/\ell^M$ \emph{comes from
$\mathcal C$} if $s=\mathsf t\circ a$ with $Y\in\mathcal C$, $a\in\pllZf(Y)$ and $\mathsf t\colon
Y\to T/\ell^M$ $\auO$-linear and $\Gamma$-equivariant. For a submodule $N$ of maps
$\Gamma\to T/\ell^M$ all of whose elements come from $\mathcal C$, its \emph{defect} is the length
of $N/(N\cap\{\text{maps obeying the crystalline relations of }T\})$.
\item A \emph{pairing modulo $\ell^m$} of lattices $Y$, $L$ is a pair $(\psi,c_\psi)$ with
$\psi\in(Y\otimes_\auO L)^*$ and $c_\psi\in\pllZf((Y\otimes L)^*)$ such that $g\psi-\psi=
\ell^mc_\psi(g)$ for all $g$; with a fixed $\Z_\ell$-linear $\mathrm{f}\colon\auO\to\Z_\ell$ for
which $(x,y)\mapsto\mathrm f(xy)$ is a perfect pairing, it gives the $\Gamma$-equivariant pairing
$b_\psi(y,l)=\mathrm f(\psi(y\otimes l))\bmod\ell^m$ with values in $\Z/\ell^m(1)$.
\item The \emph{cup-orthogonal} $L^{\pllperp\mathcal C}$ of $\mathcal C$ in $L$ is the module of
maps $c\colon\Gamma\to L$ such that for every $Y\in\mathcal C$, every $m$, every pairing
$(\psi,c_\psi)$ modulo $\ell^m$ of $Y$ and $L$ and every $a\in\pllZf(Y)$, the $2$-cochain
$(g,h)\mapsto b_\psi(a(g),g\,c(h))$ is a coboundary in $C^2(\Gamma,\Z/\ell^m(1))$. The \emph{dual
defect} is
\[
\plldef_{\mathcal C}(L)=\pglen\bigl(\pllZf(L)/(\pllZf(L)\cap L^{\pllperp\mathcal C})\bigr),
\]
and $\mathcal C$ has \emph{cup exponent} $\kappa$ on $L$ if $\ell^\kappa\pllZf(L)\subseteq L^{\pllperp
\mathcal C}$.
\end{enumerate}
\end{definition}

The classes $\mathcal C_v$ of Proposition~\ref{pg:loss} are classes $\mathcal C(U)$ or
$\mathcal C(U,\mathrm{tame})$: for finite sets $U(\gamma)\subset\Z$, one for each embedding
$\gamma\colon\auO\to\C_\ell$, $\mathcal C(U)$ consists of the crystalline lattices whose jumps at
each $\gamma$ lie in $U(\gamma)$; if $\ell=p_i$, $\mathcal C(U,\mathrm{tame})$ consists of the
lattices that become crystalline over the tamely ramified extension of $E_v$ of degree
$e'=\operatorname{lcm}_{1\le a\le r}(\ell^a-1)$, have jumps in $U$, and such that for every inertial
$s$ both $X^{e'}-1$ and $q_s$ kill $s$ on $D_{\mathrm{pcris}}$ ({\cite[Definition~9.2]{P56}}). Here
$q_s=\prod_{y\in M_s}m_y\in\Z_\ell[X]$, with $m_y$ the minimal polynomial of $y$ over $\Q_\ell$ and
$M_s$ the multiset formed by the eigenvalues of $s$ on the inertial type that (G4)
(Proposition~\ref{pga:G45}) attaches to the label at the place above $\ell$, rescaled by the tame part
of the twist, together with eigenvalues $1$ up to the rank of the representation; $q_s$ is monic of degree at most
$\max(r+t,2r)\,r$ $(=(r+t)r$, as $t=5r)$ and depends only on $s$, not on the lattice
({\cite[Proposition~9.23]{P56}} and the discussion after it). (Becoming crystalline over the tame extension of
degree $e'$ then follows from the condition on $X^{e'}-1$, {\cite[Theorem~9.9]{P56}}.) In Proposition~\ref{pg:loss}, $U(\gamma)$ is the union of the
jumps at $\gamma$ of the $t'+1$ summands $L_c$, $c\in\{\mathrm{main}\}\cup\{0,\dots,t'-1\}$: the
main summand $L_{\mathrm{main}}$ of rank $r'$ with the jumps $J_\sigma$ of (M1)--(M4) at the complex
embedding $\sigma=\sigma_\gamma$ attached to $\gamma$ through $\iota$ and $v$ ({\cite[Convention~7.1]{P56}},
reading (R$^-$): $\sigma_\gamma=\iota^{-1}\circ g^{-1}\circ e_v$ for any $g$ extending $\gamma$), and the characters $L_e=\chi^e$ with the single jump $e$.

\begin{proposition}[Exact form of Proposition~\ref{pg:loss}]\label{pll:main}
In the situation of Proposition~\ref{pg:loss}, with $\mathcal C_v$ as above and $\auO$ large, let
$N_v(L_c)$ be the module of maps $\Gamma\to L_c/\ell^M$ that come from $\mathcal C_v$. Then its
defect $\delta_v(L_c)$ satisfies (here $N_v(L_c)$ is indeed a module: the classes
$\mathcal C(U)$ and $\mathcal C(U,\mathrm{tame})$ are stable under finite direct sums, since the
jump condition and the conditions on $X^{e'}-1$ and $q_s$ are checked summand by summand, and
$\mathsf t_1\circ a_1+\mathsf t_2\circ a_2=(\mathsf t_1,\mathsf t_2)\circ(a_1,a_2)$ with
$(a_1,a_2)\in\pllZf(Y_1\oplus Y_2)$)
\begin{equation}\label{pll:eq-place}
\delta_v(L_c)\le\plldef_{\mathcal C_v}(L_c^*)\le[E_v:\Q_\ell]\cdot\frac{\pllrk_\auO L_c}{x_c}
\cdot\pllD_\ell,\qquad
\pllD_\ell=10^6\,\frac mk\bigl(b\,\pgv(M)+\lambda_\ell\bigr),
\end{equation}
with $x_{\mathrm{main}}=1$, $x_e=t$ for the characters, and $\lambda_\ell$ as in
\S\ref{pg:ss-statement}. Consequently $\sum_{v\mid\ell}\sum_c\delta_v(L_c)\le r'\pgloc_\ell$.
\end{proposition}

The first inequality in \eqref{pll:eq-place} is Proposition~\ref{pll:dual}; the second is
Theorem~\ref{pll:exponent} with Lemma~\ref{pll:exp} (Remark~\ref{pll:status}); when $\ell\ge4b+2$ and
$\ell\nmid q_1dM$ see also Remark~\ref{pll:FL} for another route; the last assertion is Lemma~\ref{pll:sum}. Note that the statement concerns the
module of \emph{all} such cocycles, not only those that Ribet's method produces. The cocycles that
Ribet's method produces come from $\mathcal C_v$ by Lemma~\ref{pll:ribet} below, so they form a
submodule of $N_v(L_c)$, whose defect is no larger.

\subsection{Reduction to the dual defect}\label{pll:ss-duality}

We use three published facts, for a lattice $T$ over $\Gamma$ and $m\ge1$.
\begin{enumerate}
\item[(D1)] (Local Tate duality \cite[Theorem~7.2.6]{NSW}.) For a finite $\Z_\ell[\Gamma]$-module
$A$, cup product and the invariant map give a perfect pairing $H^1(\Gamma,A)\times
H^1(\Gamma,\pgHom(A,\mu_{\ell^\infty}))\to\Q_\ell/\Z_\ell$. We apply it to $A=T/\ell^m$, whose
dual is identified with $T^*/\ell^m$ through $\mathrm f$; the resulting pairing $H^1(\Gamma,
T/\ell^m)\times H^1(\Gamma,T^*/\ell^m)\to\Z/\ell^m$ is $\langle[s],[c]\rangle=
\operatorname{inv}[(g,h)\mapsto\mathrm f(c(h)\text{ evaluated at }s(g))]$, up to the sign
convention for $g$ acting on the second variable.
\item[(D2)] (Bloch--Kato \cite[Proposition~3.8]{BlochKato}, integral form.) The images of
$H^1_f(\Gamma,T)$ in $H^1(\Gamma,T/\ell^m)$ and of $H^1_f(\Gamma,T^*)$ in $H^1(\Gamma,T^*/\ell^m)$
are exact annihilators of each other under (D1), when $V$ is de Rham. (Bloch and Kato prove that
$H^1_f(V)$ and $H^1_f(V^*)$ are exact annihilators; the integral statement is deduced in
Lemma~\ref{pll:D2} below, after (D3). Compare \cite[Chapter~I, Proposition~4.3]{RubinES}, where the finite-level groups are defined by
propagation from $H^1_f(\Gamma,V)$.)
\item[(D3)] (Bloch--Kato \cite[Corollary~3.8.4]{BlochKato}.)
$\dim_{\Q_\ell}H^1_f(\Gamma,V)=\dim_{\Q_\ell}H^0(\Gamma,V)+\dim_{\Q_\ell}D_{\mathrm{dR}}(V)/
\mathrm{Fil}^0D_{\mathrm{dR}}(V)$ for de Rham $V$.
\end{enumerate}

\begin{lemma}\label{pll:D2}
Statement (D2) holds for every lattice $T$ with $V$ de Rham and every $m\ge1$.
\end{lemma}

\begin{proof}
All lengths in this proof are $\Z_\ell$-lengths; write $d=[E_v:\Q_\ell]$, $n=\pllrk_{\Z_\ell}T$,
$A\subseteq H^1(\Gamma,\bar T)$ and $A^*\subseteq H^1(\Gamma,\bar T^*)$ for the two images, where
$\bar T=T/\ell^m$, $\bar T^*=T^*/\ell^m$.

\emph{Orthogonality.} The pairing (D1) of the reductions of $x\in H^1(\Gamma,T)$ and $y\in
H^1(\Gamma,T^*)$ is the reduction modulo $\ell^m$ of the cup product $x\cup y\in H^2(\Gamma,
\Z_\ell(1))\cong\Z_\ell$, and $x\cup y$ is the image of $(x\cup y)\otimes1\in H^2(\Gamma,\Q_\ell(1))$
in which the rational classes are orthogonal by \cite[Proposition~3.8]{BlochKato}. As
$H^2(\Gamma,\Z_\ell(1))\cong\Z_\ell$ is torsion free, $x\cup y=0$ for $x\in H^1_f(T)$, $y\in
H^1_f(T^*)$. So $A^*\subseteq A^{\pllperp}$, and by (D1) it suffices to prove
$\pglen A+\pglen A^*=\pglen H^1(\Gamma,\bar T)$.

\emph{Length of $A$.} $H^1_f(T)$ is the preimage in the finitely generated module $H^1(\Gamma,T)$
of a subspace of $H^1(\Gamma,V)$; it contains the torsion $H^1(\Gamma,T)_{\rm tors}$, and it is
saturated, so $H^1_f(T)\cap\ell^mH^1(\Gamma,T)=\ell^mH^1_f(T)$. The sequence
$0\to H^1(\Gamma,T)/\ell^m\to H^1(\Gamma,\bar T)$ (from $0\to T\xrightarrow{\ell^m}T\to\bar T\to0$)
therefore gives $A\cong H^1_f(T)/\ell^m$, of length $m\,h_f+\pglen H^1(\Gamma,T)[\ell^m]$, where
$h_f=\dim_{\Q_\ell}H^1_f(\Gamma,V)$ (for a finite group $X$, $X/\ell^m$ and $X[\ell^m]$ have the same
order). The same sequence gives $0\to H^0(\Gamma,T)/\ell^m\to H^0(\Gamma,\bar T)\to
H^1(\Gamma,T)[\ell^m]\to0$, and $H^0(\Gamma,T)$ is free of rank $h_0=\dim H^0(\Gamma,V)$. Hence
$\pglen A=m(h_f-h_0)+\pglen H^0(\Gamma,\bar T)=m\dim_{\Q_\ell}D_{\rm dR}(V)/\mathrm{Fil}^0+
\pglen H^0(\Gamma,\bar T)$ by (D3). Likewise $\pglen A^*=m\dim_{\Q_\ell}D_{\rm dR}(V^*)/\mathrm{Fil}^0
+\pglen H^0(\Gamma,\bar T^*)$.

\emph{Comparison.} Since $V^*=\pgHom(V,\Q_\ell(1))$ (over $\auO$, equivalently over $\Q_\ell$ after
the trace form $\mathrm f$), $\mathrm{Fil}^0D_{\rm dR}(V^*)$ is the annihilator of
$\mathrm{Fil}^0D_{\rm dR}(V)$ under the perfect pairing $D_{\rm dR}(V)\times D_{\rm dR}(V^*)\to
D_{\rm dR}(\Q_\ell(1))$ (whose filtration jump is at $-1$), so the two quotient dimensions add up to
$\dim_{\Q_\ell}D_{\rm dR}(V)=dn$. By the local Euler characteristic formula of Tate \cite[Theorem~7.3.1]{NSW},
$\pglen H^1(\Gamma,\bar T)=\pglen H^0(\Gamma,\bar T)+\pglen H^2(\Gamma,\bar T)+dmn$, and by
\cite[Theorem~7.2.6]{NSW} $H^2(\Gamma,\bar T)$ is dual to $H^0(\Gamma,\bar T^*)$. Adding the two
lengths gives $\pglen A+\pglen A^*=dmn+\pglen H^0(\Gamma,\bar T)+\pglen H^0(\Gamma,\bar T^*)=
\pglen H^1(\Gamma,\bar T)$.
\end{proof}

\begin{proposition}\label{pll:dual}
Let $\mathcal C$ be a class of lattices containing $T$, all of whose members are potentially
crystalline, and let $N$ be a module of maps $\Gamma\to T/\ell^M$ that come from $\mathcal C$. Then
\[
\pglen\bigl(N/(N\cap\overline{\pllZf}(T)_M)\bigr)\le\plldef_{\mathcal C}(T^*).
\]
In particular the defect of $N$ is at most $\plldef_{\mathcal C}(T^*)$.
\end{proposition}

\begin{proof}
Write $m=M$, $\bar T=T/\ell^m$, $\bar T^*=T^*/\ell^m$.

\emph{Step 1: the elements of $N$ are cocycles, and $\overline{\pllZf}(T)_m$ is the module of
cocycles whose class lies in the image $A$ of $H^1_f(T)$.} If $s=\mathsf t\circ a$ then $s$ is a continuous cocycle, as $a$ is one and $\mathsf t$ is equivariant. A coboundary
$g\mapsto g\bar x-\bar x$ of $\bar T$ is the reduction of the coboundary of any lift $x\in T$,
which is crystalline. Hence $\overline{\pllZf}(T)_m$ is exactly the module of cocycles whose class
lies in $A$.

\emph{Step 2: $[N]$ is orthogonal to the image of $H:=\pllZf(T^*)\cap T^{*\pllperp\mathcal C}$.}
Let $s=\mathsf t\circ a\in N$, with $a\in\pllZf(Y)$, $Y\in\mathcal C$. As $Y$ is free, $\mathsf t$
lifts to an $\auO$-linear $q\colon Y\to T$, equivariant modulo $\ell^m$. Put
$\psi(y\otimes\phi)=\phi(q(y))$ on $Y\otimes T^*$. For $g\in\Gamma$,
$(g\psi-\psi)(y\otimes\phi)=\phi(g\,q(g^{-1}y)-q(y))\in\ell^m\auO$, so $g\psi-\psi=\ell^m
c_\psi(g)$ for a unique $c_\psi(g)\in(Y\otimes T^*)^*$; $\ell^mc_\psi$ is the coboundary of $\psi$,
hence crystalline, and $\pllZf$ is saturated in the maps $\Gamma\to(Y\otimes T^*)^*$ (if $\ell^mc$
is crystalline with period $y$, then $c$ is crystalline with period $\ell^{-m}y$ since
$\pllB\otimes$ is torsion-free over $\Z_\ell$ and $\ell$ is invertible in $\pllB$). Thus
$(\psi,c_\psi)$ is a pairing modulo $\ell^m$ of $Y$ and $T^*$, and
$b_\psi(a(g),g\,c(h))=\mathrm f\bigl((g\,c(h))(\mathsf t\,a(g))\bigr)=\mathrm f\bigl((g\,c(h))
(s(g))\bigr)\bmod\ell^m$. For $c\in H$ this cochain is a coboundary by definition of
$T^{*\pllperp\mathcal C}$; it is a cocycle representing the cup product of $[s]$ and $[\bar c]$
under the pairing of (D1). Hence $\langle[s],[\bar c]\rangle=0$.

\emph{Step 3: counting.} Let $A^*\subseteq H^1(\Gamma,\bar T^*)$ be the image of $H^1_f(T^*)$,
i.e.\ of $\pllZf(T^*)$, and $\bar H\subseteq A^*$ the image of $H$. By Step 2,
$[N]\subseteq\bar H^{\pllperp}$; by (D2), $A=A^{*\pllperp}$. By (D1) and the duality of finite
abelian groups, $\bar H^{\pllperp}/A^{*\pllperp}$ is dual to $A^*/\bar H$, hence has the same
order. By Step 1, $N/(N\cap\overline{\pllZf}(T)_m)$ injects into $\bar H^\pllperp/A$. Finally
$\pllZf(T^*)/H$ surjects onto $A^*/\bar H$. Taking normalised lengths (that is, $\log_\ell$ of orders divided by $[\auO:\Z_\ell]$, the same factor on both sides) gives the claim. The last
assertion holds because the elements of $\overline{\pllZf}(T)_m$ obey the crystalline relations.
\end{proof}

The next lemma shows that the cocycles of Ribet's method (cocycles obeying every $\auO$-linear
relation among the families $(u_x(g)-1)_x$) come from $\mathcal C_v$. Both classes
$\mathcal C(U)$ and $\mathcal C(U,\mathrm{tame})$ satisfy its hypotheses: they are stable under
finite direct sums (proof of Proposition~\ref{pll:main}); all members become crystalline over one
finite Galois extension $E'$ of $E_v$ ($E'=E_v$ for $\mathcal C(U)$; for $\mathcal C(U,\mathrm{tame})$,
the Galois closure of the tame extension of degree $e'$, since a representation crystalline over a
finite extension is crystalline over every larger finite extension); membership depends only on $V$; and for
a $\Gamma$-stable sublattice $S\subseteq Y$, $S[1/\ell]$ is a subrepresentation of $Y[1/\ell]$, so
it is crystalline (resp.\ crystalline over the same tame extension), its jumps form a
sub-multiset of those of $Y$, and $D_{\mathrm{pcris}}(S[1/\ell])$ is an $s$-stable subspace of
$D_{\mathrm{pcris}}(Y[1/\ell])$, on which $X^{e'}-1$ and $q_s$ still kill $s$.

\begin{lemma}\label{pll:ribet}
Let $\mathcal C$ be a class of lattices over $\Gamma$, stable under finite direct sums and under
passage to $\Gamma$-stable sublattices, all of whose members become crystalline over a fixed finite
Galois extension $E'/E_v$. Let $u_x$ ($x$ in a finite set $X$) be members of $\mathcal C$, of ranks
$n_x$, and let $z\colon\Gamma\to T/\ell^M$ be a continuous cocycle such that $\sum_i\lambda_i\,
z(g_i)=0$ whenever $\lambda_i\in\auO$, $g_i\in\Gamma$ and $\sum_i\lambda_i(u_x(g_i)-1)=0$ for every
$x$. Then $z$ comes from $\mathcal C$. The same holds if $\Gamma$ is replaced by a group $G$ in which
$\Gamma$ is a subgroup, the $u_x$ and $z$ are defined on $G$ and only the restrictions to $\Gamma$
are asserted to lie in $\mathcal C$ and to come from $\mathcal C$.
\end{lemma}

\begin{proof}
The last sentence follows from the first: relations among elements of $\Gamma$ are relations among
elements of $G$, so the restriction of $z$ to $\Gamma$ satisfies the hypothesis.
Let $Y=\bigoplus_x\operatorname{End}_\auO(u_x)$, with $g\in\Gamma$ acting on the $x$-th summand
by left multiplication by $u_x(g)$; then $Y\cong\bigoplus_xu_x^{\oplus n_x}\in\mathcal C$. Let
$1\in Y$ be the identity in each summand and $a(g)=(u_x(g)-1)_x=g\cdot1-1$, the coboundary of $1$,
so $a\in\pllZf(Y)$. Let $S\subseteq Y$ be the $\auO$-span of the values of $a$. Since
$g\,a(h)=a(gh)-a(g)$, $S$ is $\Gamma$-stable; it is finite free (a submodule of $Y$), so
$S\in\mathcal C$. By hypothesis $\Phi(\sum_i\lambda_ia(g_i)):=\sum_i\lambda_iz(g_i)$ is a
well-defined $\auO$-linear map $\Phi\colon S\to T/\ell^M$, and it is equivariant:
$\Phi(g\,a(h))=z(gh)-z(g)=g\,z(h)=g\,\Phi(a(h))$ by the cocycle relation. Thus $z=\Phi\circ a$,
and it remains to show $a\in\pllZf(S)$, i.e.\ to find $y\in\pllB\otimes_{\Z_\ell}S$ with
$gy-y=a(g)$.

Write $W_\Q=W\otimes_{\Z_\ell}\Q_\ell$, and recall $\pllB\otimes_{\Z_\ell}W=\pllB\otimes_{\Q_\ell}
W_\Q$ for a finite free $W$. The image $\bar1$ of $1$ in $(Y/S)_\Q$ is $\Gamma$-invariant, since
$g1-1\in S$. Put $\Gamma'=\operatorname{Gal}(\auQlb/E')$, a normal subgroup of finite index. The
sequence $0\to S_\Q\to Y_\Q\to(Y/S)_\Q\to0$ consists of representations crystalline over $E'$, and
$D_{\mathrm{cris},E'}=(\pllB\otimes-)^{\Gamma'}$ is exact on such representations (it is left exact
and its value has dimension $\dim_{\Q_\ell}$ of the representation over
$E'_0$), so $(\pllB\otimes Y_\Q)^{\Gamma'}\to(\pllB\otimes(Y/S)_\Q)^{\Gamma'}$ is onto. These are
$\Q_\ell$-vector spaces with an action of the finite group $\Gamma/\Gamma'$, so taking invariants
keeps the map onto. Hence some $y'\in(\pllB\otimes Y_\Q)^\Gamma$ maps to $1\otimes\bar1$. Then
$y=1\otimes1-y'$ maps to $0$ in $\pllB\otimes(Y/S)_\Q$, so $y\in\pllB\otimes S_\Q=\pllB\otimes S$,
and $gy-y=g\cdot1-1=a(g)$.
\end{proof}

\subsection{From an exponent to the dual defect}\label{pll:ss-exponent}

\begin{lemma}\label{pll:exp}
Let $L$ be a potentially crystalline lattice. Then $\pllZf(L)$ is a free $\auO$-module of rank at
most $2[E_v:\Q_\ell]\pllrk_\auO L$. Consequently, if $\mathcal C$ has cup exponent $\kappa$ on $L$, then
$\plldef_{\mathcal C}(L)\le2\kappa\,[E_v:\Q_\ell]\,\pllrk_\auO L$.
\end{lemma}

\begin{proof}
The coboundary map $L\to\pllZf(L)$ has kernel $L^\Gamma$, and $\pllZf(L)/\{\text{coboundaries}\}
=H^1_f(\Gamma,L)\subseteq H^1(\Gamma,L)$, a finitely generated $\Z_\ell$-module
(\cite[Theorem~7.1.8(iii) and Corollary~2.7.10]{NSW}). So $\pllZf(L)$ is finitely generated; it
is torsion free, being a module of $L$-valued maps; hence free. Over $\Z_\ell$ its rank is
$\pllrk L-\pllrk L^\Gamma+\dim H^1_f(\Gamma,V)$, and by (D3) this equals
$\pllrk_{\Z_\ell}L+\dim_{\Q_\ell}D_{\mathrm{dR}}(V)/\mathrm{Fil}^0\le\pllrk_{\Z_\ell}L+
[E_v:\Q_\ell]\pllrk_{\Z_\ell}L$, because $D_{\mathrm{dR}}(V)$ is free of rank $\dim_{\Q_\ell}V$
over $E_v$ (potentially crystalline implies de Rham). Dividing by $[\auO:\Z_\ell]$ gives the rank
bound. If $\ell^\kappa\pllZf(L)\subseteq L^{\pllperp\mathcal C}$, then $\pllZf(L)/(\pllZf(L)\cap
L^{\pllperp\mathcal C})$ is a quotient of $\pllZf(L)/\ell^\kappa\cong(\auO/\ell^\kappa)^{\pllrk}$, of
normalised length $\kappa\cdot\pllrk_\auO\pllZf(L)$.
\end{proof}

\subsection{Large primes: another route}\label{pll:ss-FL}

\begin{remark}\label{pll:FL}
Suppose $\ell\nmid q_1dM$ and that all jumps of all members of $\mathcal C_v$ and of the summands
$L_c$ lie in an interval $[w_0,w_0+\ell-2]$ containing $0$ (this holds if $\ell\ge4b+2$, since all
these jumps lie in $[-2b,2b]$). Then $E_v/\Q_\ell$ is unramified, $\mathcal C_v$ is a crystalline
class $\mathcal C(U)$, and after a twist by $\pgcyc^{-w_0}$ all jumps lie in $[0,\ell-2]$. In this
situation the conclusion $\delta_v(L_c)=0$ is given by Theorem~\ref{pll:exponent}(d) with
Proposition~\ref{pll:dual} (Remark~\ref{pll:status}). One could also try to argue through the
theory of Fontaine and Laffaille \cite{FontaineLaffaille} and \cite[\S4]{BlochKato}; this would
require integral statements on lattices and on subquotients of torsion objects that we do not find
in these sources in the needed form, and the paper does not use this route.
\end{remark}

\subsection{The remaining primes}\label{pll:ss-gap}

For the remaining primes we prove explicit cup exponents by an argument in integral $\ell$-adic
Hodge theory, written in {\cite[\S\S2--9]{P56}}. We state precisely what is needed.

\begin{theorem}\label{pll:exponent}
Let $\ell$ be any prime, $v\mid\ell$, $\auO$ large, $\mathcal C_v$ and the summands $L_c$ as in
Proposition~\ref{pll:main}, with $L_c\in\mathcal C_v$. Then $\mathcal C_v$ has, on $L_c^*$, a cup
exponent $k_c$ with $2k_c\le\pllD_\ell/x_c$; or, in the sub-case below where no exponent is
claimed, directly $\plldef_{\mathcal C_v}(L_c^*)\le[E_v:\Q_\ell]\pllrk(L_c)\pllD_\ell/x_c$. In
detail, with $h=r^8$, $Q=kt$ and $\lambda=\lambda_\ell$:
\begin{enumerate}
\item[(a)] (tame, $\ell=p_i$) exponent $3r^2(N+1)^2+150(h^2+4t+2)$, $N=\max(r+t,2r)r$;
\item[(b)] (crystalline, $\ell\le Q$) exponent $10^4\lfloor2\lambda/Q\rfloor(\lfloor\log_\ell
h^3\rfloor+1)$;
\item[(c)] (crystalline, $\ell>Q$, $\lambda\ge h^3$) exponent $10^4\lfloor2\lambda/Q\rfloor m$;
\item[(d)] (crystalline, $\ell>Q$, $\lambda<h^3$) $\plldef=0$ if $\lambda<\lfloor h/5\rfloor$ or
$L_c$ is a character, and in all cases $\plldef_{\mathcal C_v}(L_c^*)\le
[E_v:\Q_\ell]\cdot10^5\bigl(\pllrk(L_c)\lfloor2\lambda/Q\rfloor m+\mu\lambda\bigr)$ with
$\mu=\lfloor\pllrk(L_c)m/(kx_c)\rfloor$.
\end{enumerate}
\end{theorem}

Theorem~\ref{pll:exponent} is proved in {\cite[\S2.3]{P56}} (proof following
{\cite[Theorem~2.13]{P56}}) from {\cite[Propositions~2.5--2.10]{P56}}, which are proved in
{\cite[\S\S7, 8.1 and~8.2]{P56}} from {\cite[\S\S3.1--9]{P56}}; the
dependency chain is recorded at the end of {\cite[\S2.2]{P56}}. What that proof takes from \S\ref{pg:sec} and from the present section (the definitions, Lemma~\ref{pll:exp}, and the description of the summands and of the classes $\mathcal C_v$) does not depend on the theorem. No published theorem is used in
place of this argument ({\cite[Remark~2.1]{P56}}).

\begin{remark}[Proof of \eqref{pll:eq-place} from Theorem~\ref{pll:exponent}]\label{pll:status}
Given Theorem~\ref{pll:exponent}, \eqref{pll:eq-place} follows: in cases (a)--(c) by
Lemma~\ref{pll:exp} applied to $L=L_c^*$ ($\pllrk L_c^*=\pllrk L_c$) and the inequality
$2k_c\le\pllD_\ell/x_c$; in case (a) this inequality uses $b\,\pgv(M)\ge b$ and the sizes
{\cite[(38)]{P46}} ($r^2(N+1)^2\le49r^6$ and $h^2$ are polynomial in $r$, far below $b\,m/(kt)$); in
cases (b)--(c), $\lfloor\log_\ell h^3\rfloor+1\le m$ because $2\log(4b)<m$, so $2k_c\le2\cdot10^4\cdot(2\lambda/(kt))\,m\le10^6(m/k)\lambda/x_c$ as $x_c\le t$. In case (d), $10^5(\rho\lfloor2\lambda/(kt)\rfloor m+\lfloor\rho
m/(kx)\rfloor\lambda)\le3\cdot10^5\rho m\lambda/(kx)\le\rho\,\pllD_\ell/x$ for $0<x\le t$.
For $\ell>2b$ one has $\ell\nmid(2b)!$ and $\ell>b+30h^3$ (as $30h^3<b$), hence $\lambda_\ell=0$, so if
moreover $\ell\nmid M$, case (d) gives $\plldef=0$: this is the ``no loss'' clause of
Proposition~\ref{pg:loss} (for $\ell>Q$, {\cite[Theorem~7.9]{P56}} gives letters that read
thickness, and {\cite[Definition~2.4]{P56}} then forces ramification index $e=1$, so
$E_v/\Q_\ell$ is unramified and in particular $\ell\nmid q_1d$). For $\ell\ge4b+2$ prime to $q_1d$
see also Remark~\ref{pll:FL} for another route.
\end{remark}

\subsection{The sum over places and summands}\label{pll:ss-sum}

\begin{lemma}\label{pll:sum}
If \eqref{pll:eq-place} holds for all $v\mid\ell$ and all $c$, then $\sum_{v\mid\ell}\sum_c
\delta_v(L_c)\le r'\pgloc_\ell$.
\end{lemma}

\begin{proof}
$\sum_{v\mid\ell}[E_v:\Q_\ell]=[E:\Q]=2f$. Over the summands, $\sum_c\pllrk(L_c)/x_c=r'+t'/t\le
r'+1\le2r'$. Hence the double sum is at most $2f\cdot2r'\cdot\pllD_\ell=r'\cdot4f\cdot10^6
\frac mk(b\,\pgv(M)+\lambda_\ell)=r'\pgloc_\ell$.
\end{proof}




\section{Galois representations attached to the lifts}\label{C3:g5}
\providecommand{\auO}{\mathcal O}
\providecommand{\auQlb}{\overline{\Q}_\ell}
\providecommand{\auG}{\overline G}
\providecommand{\auNm}{\operatorname{Nm}}
\providecommand{\auU}{\mathrm U}
\providecommand{\auGL}{\mathrm{GL}}
\providecommand{\auA}{\mathbb A}
\providecommand{\auF}{\mathbb F}
\providecommand{\aubar}[1]{\overline{#1}}
\providecommand{\auf}[1]{\texttt{#1}}
\providecommand{\pgG}{G_E}
\providecommand{\pgtr}{\operatorname{tr}}
\providecommand{\pgFrob}{\operatorname{Frob}}
\providecommand{\pgcyc}{\varepsilon_\ell}
\providecommand{\pgNch}{N_{\mathrm{ch}}}
\providecommand{\pgarec}{\operatorname{rec}}
\providecommand{\pgaWD}{\operatorname{WD}}
\providecommand{\pgaInd}{\operatorname{Ind}}
\providecommand{\pgaHT}{\operatorname{HT}}
\providecommand{\pgaord}{\operatorname{ord}}

\subsection{The inputs of Theorem~\ref{pg:I}}\label{pga:sec}

This section supplies the inputs of Theorem~\ref{pg:I} used in \S\ref{pg:ss-galois}:
(i) the properties (G1)--(G5) of Galois representations, deduced from published theorems whose
hypotheses we check (\S\S\ref{pga:ss-need}--\ref{pga:ss-local}); (ii) statement (L), the Hecke
polynomial of a theta lift, and statement (C), the central action on the lift from $\auU(1)$,
both proved here (\S\S\ref{pga:ss-L}, \ref{pga:ss-C}); (iii) condition (M5) and the two
normalisations of the proof of Theorem~\ref{pg:I} (\S\ref{pga:ss-M}). The final paragraph
(\S\ref{pga:ss-status}) lists the published results on which Theorem~\ref{pg:I} rests.

Notation is that of \S\ref{pg:sec}. Throughout, $\pgFrob_w$ is a \emph{geometric} Frobenius,
$\pgcyc$ the cyclotomic character and $\chi=\pgcyc^{-1}$, so that $\chi(\pgFrob_w)=\auNm w$. For an
algebraic Hecke character $\phi$ of $E$ we write $\phi_\ell$ for its $\ell$-adic Galois character
(through $\iota$), normalised so that $\phi_\ell(\pgFrob_w)=\iota\phi(\varpi_w)$ for $w\nmid\ell$
where $\phi$ is unramified, $\varpi_w$ a uniformiser at $w$ placed in the idele at $w$. Local
class field theory $\mathrm{Art}_w$ sends uniformisers to geometric Frobenius elements.

\subsubsection{What is needed, and in which form}\label{pga:ss-need}

Recall (G1)--(G5) of \S\ref{pg:ss-galois}. We make two of them precise.

\emph{(G1) and (G3), exact form.} Let $V$ be a positive definite standard hermitian space of rank
$n$ and $y$ an algebraic modular form on $\auU(V)$ of weight $\lambda$ that is an eigenform of the
$T_{w,i}$ at all split $w$ off a finite set. We use the Harish-Chandra parameter
$\varsigma=\lambda+\rho$ at each $\tau$, $\rho=(\frac{n-1}2,\frac{n-3}2,\dots,\frac{1-n}2)$, in the
normalisation in which the parameter of a character of $\auU(1)$ of weight $c$ is $c$. (G1) asks
for $\det(X-S_y(\pgFrob_w))=\iota H_w(y;X)$; (G3) asks that at $\sigma$ in the CM type the
Hodge--Tate weights be $\frac{n-1}2-\varsigma_{\sigma,i}$ and at $\aubar\sigma$ be
$\frac{n-1}2+\varsigma_{\sigma,i}$ ($1\le i\le n$), with the sign convention in which $\pgcyc$ has
weight $-1$, that is, $\chi$ has weight $1$. (With $\varsigma=\lambda+\rho$ these are the numbers
$i-1-\lambda_{\sigma,i}$ and $\lambda_{\sigma,i}+n-i$ of (G3).)

\emph{Where each property is used.} (G1): in the trace identity \eqref{pg:eq-trace} and in (M1).
(G2): in (M2) (unramified outside $\ell p_0p_1$, crystalline above $\ell\nmid M$) and in the
Selmer argument of Theorem~\ref{pg:III}. (G3): in (M2) (weights) and in the balance of the
Bloch--Kato conditions. (G4): in (M2) for $\ell=p_i$, (M3), (M4), (M5). (G5): only in
Theorem~\ref{pg:III}, at the places of $Q_\vartheta$.

\emph{Weak forms suffice.} Every use of (G1)--(G5) in Theorems~\ref{pg:I} and \ref{pg:III} is
through one of the following \emph{weak} statements: (G1$'$) $\pgtr S_y(\pgFrob_w)=\iota\,
(\text{coefficient})$ for a set of split $w$ of Dirichlet density one; (G2$'$), (G4$'$), (G5$'$)
``a given polynomial kills $S_y(g)$'' for $g$ a Frobenius, respectively inertial, element (for
instance $(S_y(s)-\psi_w(s))(S_y(s)-1)=0$); plus the crystalline and Hodge--Tate statements above
$\ell$. Indeed (G1$'$) gives \eqref{pg:eq-trace} by Chebotarev and continuity of traces, and the
characteristic polynomial of a semisimple representation is determined by its traces in
characteristic $0$; (M3)--(M5) use only that $R(\gamma_w)$ is killed by the polynomial whose roots
are the images of the roots of the rescaled label (which, $R(\gamma_w)$ having finite order, gives
that $R(\gamma_w)$ is semisimple with eigenvalues among these roots, and the trace identity then
gives the multiplicities). There is no short self-contained proof of the weak forms, since
already (G1$'$) contains the existence of the Galois representations. We therefore derive the
strong forms from the literature and note that they imply the weak ones.

\subsubsection{The global input}\label{pga:ss-global}

\begin{theorem}[Published results]\label{pga:lit}
Let $F$ be a CM field with maximal totally real subfield $F^+$, $\ell$ a prime and
$\iota\colon\auQlb\simeq\C$.
\begin{enumerate}
\item[(BC)] \emph{(Labesse~\cite{Labesse}, with M\oe glin--Waldspurger~\cite[Introduction, Th\'eor\`eme, p.~606]{MW89} for the
form of the discrete constituents.)} Assume $[F^+:\Q]\ge2$ and $F/F^+$ unramified at all finite
places. Let $G=\auU(V)$ for a positive definite
hermitian space $V$ of rank $n$ over $F$, quasi-split at all finite places, and $\pi$ an
automorphic representation of $G(\auA_{F^+})$ whose archimedean components are irreducible
algebraic representations (that is, $\pi_\infty$ is cohomological). Then there are a partition
$n=\sum_i n_im_i$ and cuspidal automorphic representations $\sigma_i$ of $\auGL_{n_i}(\auA_F)$ such
that $\Pi=\boxplus_i\operatorname{Speh}(\sigma_i,m_i)$ is an automorphic representation of
$\auGL_n(\auA_F)$ which is conjugate self-dual, cohomological with the infinitesimal character
determined by $\pi_\infty$, and is the local base change of $\pi_v$ at every place $v$ of $F^+$
split in $F$ and at every finite $v$ at which $G$ and $\pi_v$ are unramified; each
$\operatorname{Speh}(\sigma_i,m_i)$ is conjugate self-dual.
\item[(GR)] \emph{(Harris--Taylor~\cite{HarrisTaylor}, Shin~\cite{Shin}, Chenevier--Harris
\cite{ChenevierHarris}; collected in \cite[Thm.~2.1.1]{BLGGT}.)} Let $\Sigma$ be a regular
algebraic cuspidal automorphic representation of $\auGL_N(\auA_F)$ with $\Sigma^c\simeq\Sigma^\vee
\otimes(\phi\circ\auNm_{F/F^+}\circ\det)$ for an algebraic character $\phi$ of
$\auA_{F^+}^\times/F^{+\times}$ with $\phi_v(-1)$ independent of $v\mid\infty$. Then there is a
continuous semisimple $r_\iota(\Sigma)\colon G_F\to\auGL_N(\auQlb)$ with $r_\iota(\Sigma)^c\simeq
r_\iota(\Sigma)^\vee\pgcyc^{1-N}\phi_\ell$, de Rham above $\ell$ with Hodge--Tate weights given by
the infinitesimal character of $\Sigma_\infty$, and such that for every finite $v\nmid\ell$
\[
\pgaWD\bigl(r_\iota(\Sigma)|_{G_{F_v}}\bigr)^{\mathrm{F\text{-}ss}}\simeq\pgarec\bigl(\Sigma_v\otimes
|\det|_v^{(1-N)/2}\bigr)
\]
(after $\iota$). The equality \emph{including the monodromy operator} is \cite[Thm.~2.1.1(2)]{BLGGT}, stated
there for every regular algebraic cuspidal polarised pair, which rests on
Caraiani~\cite{CaraianiDuke}. For $v\mid\ell$, $r_\iota(\Sigma)|_{G_{F_v}}$ is de Rham, hence
potentially semistable, and for the Weil--Deligne representation of Fontaine's $D_{\mathrm{pst}}$
one has $\iota\pgaWD(r_\iota(\Sigma)|_{G_{F_v}})^{\mathrm{F\text{-}ss}}\prec\pgarec(\Sigma_v\otimes
|\det|_v^{(1-N)/2})$ in the sense of \cite[\S1]{BLGGTlp2}: the two sides have isomorphic
semisimplifications, and for each irreducible $W_{F_v}$-type the partial sums of the sizes of
the Jordan blocks of the monodromy on the left are bounded by those on the right
(Barnet-Lamb--Gee--Geraghty--Taylor \cite[Theorem~A]{BLGGTlp2}, stated there for pairs
$(\Sigma,\phi)$ as here, with no further hypothesis); with equality (``$\simeq$'') when the
weight is Shin-regular (loc.\ cit.), and in general when $\phi=1$ by
Caraiani~\cite[Theorem~1.1]{CaraianiANT}. In particular $r_\iota(\Sigma)$ is crystalline at
$v\mid\ell$ if $\Sigma_v$ is unramified \cite[Thm.~2.1.1(4)]{BLGGT}.
\end{enumerate}
\end{theorem}
These statements are quoted, not proved. In \cite{Labesse} the base change is stated
for a representation $\pi$ of the discrete spectrum with $\pi_\infty$ in the discrete series and
satisfying a sign condition on the archimedean traces of its near-equivalence class, which
holds automatically when $G(F^+\otimes\R)$ is compact (loc.\ cit.); the conclusion there is
$\Pi=\boxplus_i\Pi_i$ with $\Pi_i$ discrete automorphic and conjugate self-dual, and by
M\oe glin--Waldspurger~\cite[Introduction, Th\'eor\`eme, p.~606]{MW89} a discrete automorphic representation of $\auGL_N$ is
$\operatorname{Speh}(\sigma,m)$ with $\sigma$ cuspidal, $\sigma$ being conjugate self-dual when
$\operatorname{Speh}(\sigma,m)$ is. The assumption $[F^+:\Q]\ge2$ is made throughout
\cite{Labesse}. The sign condition is that of \cite[\S2.1]{BLGGT}
(polarized automorphic representations), and \cite[Theorem~2.1.1]{BLGGT} gives the identity at
$v\mid\ell$ only when $\Sigma_v$ has an Iwahori-fixed vector. Existence, the polarization,
de Rham-ness with the Hodge--Tate weights, the identity at $v\nmid\ell$ and crystallinity at
unramified $v\mid\ell$ are all contained in \cite[Theorem~2.1.1(1)--(4)]{BLGGT}; at ramified
$v\mid\ell$ the paper uses only the relation $\prec$ of \cite[Theorem~A]{BLGGTlp2}, see the proof
of Proposition~\ref{pga:G45}.

\begin{proposition}[(G1)--(G3), and no monodromy where the level is tame]\label{pga:G123}
Let $V$, $y$, $\lambda$ be as in \S\ref{pga:ss-need}, let $\pi$ be an irreducible constituent of
the automorphic representation generated by $y$ (on which the $T_{w,i}$ act by the eigenvalues
of $y$), and $\Pi=\boxplus_i\operatorname{Speh}(\sigma_i,m_i)$ as in (BC). Let $s_i\in\frac12\Z$
with $s_i\equiv(n_i-1-n+m_i)/2\pmod1$; then $\sigma_i|\det|^{s_i}$ is regular algebraic and satisfies
the hypothesis of (GR) (Proposition~\ref{pax:sign}). Put
\[
S_y=\bigoplus_i\bigoplus_{a=0}^{m_i-1}r_\iota\bigl(\sigma_i|\det|^{s_i}\bigr)\otimes\pgcyc^{\,e_{i,a}},
\quad e_{i,a}=\tfrac{m_i-1}2-a-s_i+\tfrac{n_i-n}2.
\]
Then $S_y$ satisfies (G1) at every split $w$ at which $y$ is spherical, (G2) and (G3); and for
every finite $v\nmid\ell$ the monodromy of $\pgaWD(S_y|_{G_{E_v}})$ vanishes as soon as the
monodromy of $\pgarec(\Pi_v)$ vanishes.
\end{proposition}
\begin{proof}
$\pi_\infty$ is the irreducible algebraic representation of weight $\lambda$ (definite group), so
(BC) applies; the hypothesis ``quasi-split at all finite places'' holds because $E/E^+$ is
unramified at all finite places and $V$ has a lattice unimodular at every finite place: the
stabiliser of such a lattice is hyperspecial, so $\auU(V)$ is unramified, in particular
quasi-split, at every finite place of $E^+$. Moreover $[E^+:\Q]\ge2$: every quadratic
extension of $\Q$ is ramified at some prime (Minkowski), while $E/E^+$ is unramified at all
finite places. No hypothesis at the places of $E^+$ ramified over $\Q$ is needed in
\cite{Labesse}. Since $\Pi$ is
regular (its infinitesimal character is $\varsigma$, regular), the constituents
$\operatorname{Speh}(\sigma_i,m_i)$ have pairwise disjoint infinitesimal characters, and each
$\sigma_i|\det|^{s_i}$ is regular algebraic (Proposition~\ref{pax:sign}). The Galois representation
of $\operatorname{Speh}(\sigma,m)$ is $\bigoplus_ar_\iota(\sigma|\det|^{s})\pgcyc^{e_a}$, $e_a=
\frac{m-1}2-a-s+\frac{N-n}2$ ($N$ the rank of $\sigma$), because the Satake parameters, respectively
the Langlands parameters, of $\operatorname{Speh}(\sigma,m)_v$ are those of $\sigma_v$ times
$|\cdot|^{(m-1)/2-a}$, $0\le a<m$ (Proposition~\ref{pax:sign}). Hence:

(G1). At a split $w$ at which $y$ is spherical, $\pi_v$ is unramified and $\Pi_w=\pi_v$ under
$\auU(V)(E^+_v)\simeq\auGL_n(E_w)$; the Satake parameters of $\Pi_w$ are the roots of
$H_w(y;X)$ rescaled by $(\auNm w)^{-(n-1)/2}$ (the normalisation of $H_w$: $\sum_i(-1)^i(\auNm
w)^{i(i-1)/2}T_{w,i}X^{n-i}=\prod_k(X-(\auNm w)^{(n-1)/2}\alpha_k)$ for the Satake parameters
$\alpha_k$ of $\Pi_w$, Satake's formula for $\auGL_n$), and $\pgarec(\Pi_w|\det|^{(1-n)/2})$ sends
$\pgFrob_w$ to $\operatorname{diag}((\auNm w)^{(n-1)/2}\alpha_k)$. By (GR) for each constituent
and the remark on Speh representations, $\det(X-S_y(\pgFrob_w))=\iota H_w(y;X)$ exactly, with no further twist: write
$q=\auNm w$, $\sigma=\sigma_i$, $m=m_i$, $s=s_i$, and let $\alpha_\sigma$ be a Satake parameter of
$\sigma_w$. Since $\pgarec(|\cdot|^u)(\pgFrob_w)=q^{-u}$ and $\pgcyc(\pgFrob_w)=q^{-1}$, the summand
$r_\iota(\sigma|\det|^{s})\otimes\pgcyc^{e_{i,a}}$ has the Frobenius eigenvalue
$\alpha_\sigma q^{-s}q^{(n_i-1)/2}q^{-e_{i,a}}=q^{(n-1)/2}\alpha_\sigma q^{a-(m-1)/2}$ by the
value of $e_{i,a}$ in the statement; as $a$ runs over $0,\dots,m-1$ these are $q^{(n-1)/2}$
times the Satake parameters $\alpha_\sigma q^{(m-1)/2-a}$ of $\operatorname{Speh}(\sigma,m)_w$.

(G2). If $y$ is fixed by a hyperspecial subgroup at $v$, $\pi_v$ and hence $\Pi_v$ is unramified
(compatibility of (BC) at unramified places), so each $\sigma_{i,v}$ is unramified (the cuspidal
support of an unramified representation is unramified); by (GR) the Frobenius-semisimplified
Weil--Deligne representation of $S_y$ at $v\nmid\ell$ is unramified with vanishing monodromy, so
$S_y$ is unramified at $v$; at $v\mid\ell$ each summand is crystalline.

(G3). By (GR) the Hodge--Tate weights of $r_\iota(\Sigma)$ at $\sigma$ are the entries of the
infinitesimal character of $\Sigma_\sigma$ shifted by $\frac{N-1}2$; the union over the summands
of $S_y$ gives, at $\sigma$ in the CM type, $\{\frac{n-1}2-\varsigma_{\sigma,i}\}$ and at
$\aubar\sigma$ the numbers $\frac{n-1}2+\varsigma_{\sigma,i}$ (the sign being fixed by the
convention of \S\ref{pga:ss-need}; the infinitesimal character of $\Pi_{\aubar\sigma}$ is the
negative of that of $\Pi_\sigma$ by conjugate self-duality).

Monodromy. $\pgarec(\operatorname{Speh}(\sigma,m)_v)=\bigoplus_a\pgarec(\sigma_v)|\cdot|^{(m-1)/2-a}$
with the monodromy of $\pgarec(\sigma_v)$ on each summand (the local component of a Speh
representation is the Langlands quotient, whose parameter carries no monodromy beyond that of
$\sigma_v$). By (GR), \cite[Thm.~2.1.1(2)]{BLGGT}, the monodromy of $S_y$ at $v$ is the direct sum of those of
the $\pgarec(\sigma_{i,v})$, hence it vanishes if that of $\pgarec(\Pi_v)$ does.
\end{proof}

\begin{remark}[The sign hypothesis]\label{pga:sign}
For $m_i>1$ the twist $|\det|^{(m_i-1)/2}$ need not make $\sigma_i$ algebraic; the correct twist is
$|\det|^{s_i}$ with $s_i\equiv(n_i-1-n+m_i)/2\pmod1$, and the polarisation character is then
$|\cdot|^{2s_i}$, whose sign is $+1$ at every infinite place. So the hypothesis of (GR) always holds;
see Proposition~\ref{pax:sign}.
\end{remark}

\subsubsection{Tame level: (G4) and (G5)}\label{pga:ss-local}

Let $K$ be a finite extension of $\Q_p$ with residue field $\auF_q$ and ring of integers $\auO_K$,
$\ell\ne p$ or $\ell=p$, and fix an injective character $\omega\colon\overline{\auF}_q^\times
\to\auQlb^\times$. For $d\ge1$ let $K_d$ be the unramified extension of degree $d$; the tame
inertia character $\theta_d\colon I_K\to\auF_{q^d}^\times$ is $s\mapsto s(\varpi^{1/(q^d-1)})/
\varpi^{1/(q^d-1)}\bmod\varpi$. A \emph{label} $c$ of degree $n$ over $\auF_q$ (monic, squarefree,
$c(0)\ne0$) factors as $c=\prod_kf_k$ with $f_k$ irreducible of degree $d_k$, roots $\xi_k^{q^a}$;
it names (Green) the irreducible character $\sigma_c$ of $\auGL_n(\auF_q)$ which is, up to sign,
the Deligne--Lusztig character $R_T(\theta)$ for the torus $T=\prod_k\operatorname{Res}\auF_{q^{
d_k}}^\times$ and the character $\theta$ of $T$ named by the $\xi_k$ through $\omega$; it is
Harish-Chandra induced from the cuspidal representation $\otimes_k\sigma_{f_k}$ of the Levi
$\prod_k\auGL_{d_k}(\auF_q)$.

\begin{lemma}[Depth zero]\label{pga:depth0}
Let $\Pi$ be an irreducible smooth representation of $\auGL_n(K)$ whose vectors fixed by the
principal congruence subgroup $K(1)=1+\varpi M_n(\auO_K)$ contain $\sigma_c$, inflated from
$\auGL_n(\auO_K)\to\auGL_n(\auF_q)$, for a label $c$. Then $\Pi\simeq\pi_1\times\dots\times\pi_s$
(irreducible normalised parabolic induction), where $\pi_k$ is a depth-zero supercuspidal
representation of $\auGL_{d_k}(K)$ containing $\sigma_{f_k}$; $\pgarec(\Pi)$ has monodromy $0$;
and $\pgarec(\Pi)|_{I_K}$ is semisimple, factors through the tame quotient, and, on an element
$s\in I_K$ with $\theta_d(s)$ a generator of $\auF_{q^d}^\times$, $d=\operatorname{lcm}(d_k)$, it
has the eigenvalues $\omega(\xi_k^{\,q^a})^{\epsilon}$, $0\le a<d_k$, where the roots $\xi_k$ are
normalised by $\xi_k=\theta_{d}(s)^{e_k}$ through the naming of $\sigma_c$, for a sign $\epsilon\in\{\pm1\}$ depending only on the normalisations of
$\pgarec$, $\mathrm{Art}$ and of Green's parametrisation, not on $c$ or $\Pi$.
\end{lemma}
\begin{proof}
\emph{Cuspidal support.} We prove the following statement, used also elsewhere with labels that
are not squarefree: \emph{if $\Pi^{K(1)}$ contains an irreducible representation $W$ of
$\auGL_n(\auF_q)$ with cuspidal support $(L,\tau)$, $L=\prod_{k=1}^s\auGL_{d_k}$,
$\tau=\otimes_k\sigma_{f_k}$ ($f_k$ irreducible, repetitions allowed), then $\Pi$ has depth zero and
is a subquotient of $\pi_1\times\dots\times\pi_s$, with $\pi_k$ a depth-zero supercuspidal
representation of $\auGL_{d_k}(K)$ containing $\sigma_{f_k}$; so its supercuspidal support is
$\{\pi_1,\dots,\pi_s\}$.} We use two papers of Moy and Prasad: \cite{MoyPrasad94} and
\cite{MoyPrasad}.

\emph{Step 1: $\Pi$ contains a depth-zero minimal $K$-type.} Let $Q=LU$ be the block upper
triangular parabolic subgroup of type $(d_1,\dots,d_s)$. Since $W$ has cuspidal support $(L,\tau)$, it
is a constituent of the Harish-Chandra induced representation $\operatorname{Ind}_{Q(\auF_q)}^{
\auGL_n(\auF_q)}\tau$ (inflated to $Q(\auF_q)$); by Harish-Chandra's theory this does not depend on
the parabolic subgroup with Levi factor $L$, \cite[Proposition~6.1(1)]{MoyPrasad}. Frobenius
reciprocity for the finite group gives $0\ne\operatorname{Hom}_{\auGL_n(\auF_q)}(W,\operatorname{Ind}
\tau)=\operatorname{Hom}_{Q(\auF_q)}(W,\tau)$. As $U(\auF_q)$ acts trivially on $\tau$, a non-zero
$Q(\auF_q)$-map $W\to\tau$ factors through the $U(\auF_q)$-coinvariants of $W$, which, the group being
finite and the coefficients of characteristic $0$, are identified $L(\auF_q)$-equivariantly with the
invariants $W^{U(\auF_q)}$. So $W^{U(\auF_q)}$ contains $\tau$ as an $L(\auF_q)$-subrepresentation
(complete reducibility). Let $\mathcal Q\subseteq\auGL_n(\auO_K)$ be the inverse image of
$Q(\auF_q)$ and $\mathcal Q^+$ the inverse image of $U(\auF_q)$. Then $\mathcal Q$ is a parahoric
subgroup $\mathcal G_x$ (for a point $x$ of the apartment of the diagonal torus), $\mathcal Q^+=
\mathcal G_x^+$ is its pro-unipotent radical, $\mathcal Q/\mathcal Q^+=L(\auF_q)$ is the group of
points of its reductive quotient, and $\Pi^{\mathcal Q^+}=(\Pi^{K(1)})^{U(\auF_q)}\supseteq
W^{U(\auF_q)}$. Hence the $(\mathcal Q,\tau)$-isotypic subspace of $\Pi$ is non-zero. Since $\tau$ is
cuspidal, $(\mathcal Q,\tau)$ is an unrefined minimal $K$-type of depth zero in the sense of
\cite[\S3.4, (i)]{MoyPrasad} (= \cite[Definition~5.1(i)]{MoyPrasad94}). Moreover $\Pi$ has depth zero:
$K(1)=P_{x_0,0+}$ for the hyperspecial vertex $x_0$ fixed by $\auGL_n(\auO_K)$, and $\Pi^{K(1)}\ne0$, so
the depth $\rho(\Pi)$, the smallest number $r\ge0$ with $\Pi^{P_{x,r+}}\neq0$ for some $x$
\cite[Theorem~5.2(1)]{MoyPrasad94} (\cite[Theorem~3.5(1)]{MoyPrasad} in any characteristic), is $0$.

\emph{Step 2: the supercuspidal support.} By \cite[Theorem~6.11(2)]{MoyPrasad}, an irreducible
admissible representation whose $(\mathcal G_x,\sigma)$-isotypic subspace is non-zero, for a minimal
$K$-type $(\mathcal G_x,\sigma)$ of depth zero, is a subquotient of the (unnormalised) induced
representation $I(\mathcal P,\theta)=\operatorname{Ind}_{\mathcal P}^{\mathcal G}\theta$, where
$\mathcal M$ is the Levi subgroup attached to $\mathcal G_x$ in \cite[6.3]{MoyPrasad}, $\mathcal P$ is
a parabolic subgroup with Levi factor $\mathcal M$, $\theta=\operatorname{c-Ind}_{\tilde{\mathcal
M}_x}^{\mathcal M}t$, $\tilde{\mathcal M}_x$ is the normaliser in $\mathcal M$ of the parahoric
$\mathcal M_x=\mathcal M\cap\mathcal G_x$, and $t$ is an irreducible representation of $\tilde{\mathcal
M}_x$ containing $\sigma$ (the set $\mathcal E(\sigma)$ of \cite[6.5]{MoyPrasad}); when $\mathcal G_x$ is
maximal, $\mathcal M=\mathcal G$ and $I(\mathcal P,\theta)=\theta$. We apply it to $(\mathcal
Q,\tau)$ (Step 1; $\Pi$ is admissible). The maximal $\auF_q$-split torus in the centre of $L$ is the
torus of block scalars; its lift $C$ to the diagonal torus has centraliser the block diagonal group,
so $\mathcal M=L(K)=\prod_k\auGL_{d_k}(K)$, $\mathcal M_x=\prod_k\auGL_{d_k}(\auO_K)$ (with
$\mathcal M_x/\mathcal M_x^+\cong\mathcal Q/\mathcal Q^+$, \cite[Proposition~6.4]{MoyPrasad}), and
$\tilde{\mathcal M}_x=\prod_kK^\times\auGL_{d_k}(\auO_K)$. An irreducible representation of this
product is $t=\otimes_kt_k$ with $t_k$ irreducible; $t$ contains $\tau=\otimes_k\sigma_{f_k}$, so each
$t_k$ contains $\sigma_{f_k}$, and $\theta=\otimes_k\theta_k$ with $\theta_k=\operatorname{c-Ind}_{
K^\times\auGL_{d_k}(\auO_K)}^{\auGL_{d_k}(K)}t_k$, an irreducible supercuspidal representation of
$\auGL_{d_k}(K)$ \cite[Proposition~6.6]{MoyPrasad} (applied to $\auGL_{d_k}$ and its maximal parahoric
$\auGL_{d_k}(\auO_K)$) which contains $\sigma_{f_k}$. Normalised and unnormalised induction differ by
$\delta_{\mathcal P}^{1/2}$, whose restriction to $\mathcal M$ is $\otimes_k|\det|^{a_k}$ for real
$a_k$, an unramified character; put $\pi_k=\theta_k\otimes|\det|^{a_k}=\operatorname{c-Ind}(t_k\otimes
|\det|^{a_k})$ (last assertion of \cite[Proposition~6.6]{MoyPrasad}), again a supercuspidal containing
$\sigma_{f_k}$, of depth zero by Step 1's criterion. Then $\Pi$ is a subquotient of
$\pi_1\times\dots\times\pi_s$ (normalised induction, in some order of the blocks; the order does not
change the supercuspidal support), as claimed.

\emph{Depth-zero supercuspidals with a given $\sigma_f$.} By the second assertion of
\cite[Proposition~6.6]{MoyPrasad}, an irreducible representation $\pi$ of $\auGL_d(K)$ containing the
(cuspidal) $\sigma_f$ is $\operatorname{c-Ind}_{K^\times\auGL_d(\auO_K)}^{\auGL_d(K)}\tilde\sigma_f$
for an irreducible $\tilde\sigma_f$ containing $\sigma_f$; as $\varpi$ is central, $\tilde\sigma_f$
restricts irreducibly to $\auGL_d(\auO_K)$, so it is an extension of $\sigma_f$, determined by the
scalar $\tilde\sigma_f(\varpi)$. Hence two such, $\pi$ and $\pi'$ with the same $\sigma_f$, differ by
twisting $\tilde\sigma_f$ by a character $\chi$ of $K^\times\auGL_d(\auO_K)/\auGL_d(\auO_K)\cong
\varpi^{\mathbb Z}$; choosing an unramified character $\mu$ of $K^\times$ with $\mu(\varpi)^d=
\chi(\varpi)$, $\chi$ is the restriction of $\mu\circ\det$, and $\pi'\simeq\pi\otimes(\mu\circ\det)$
by the last assertion of \cite[Proposition~6.6]{MoyPrasad}. Conversely, if $\pi$ contains $\sigma_f$
and $\pi\otimes(\mu\circ\det)$, $\mu$ unramified, contains $\sigma_{f'}$, then $\pi$ contains both
$(\auGL_d(\auO_K),\sigma_f)$ and $(\auGL_d(\auO_K),\sigma_{f'})$ ($\mu\circ\det$ is trivial on
$\auGL_d(\auO_K)$), which are therefore associates \cite[Theorem~5.2, last assertion]{MoyPrasad94}:
there is $g$ with $\auGL_d(\auO_K)\cap g\auGL_d(\auO_K)g^{-1}$ mapping onto both reductive quotients
and $\sigma_f\simeq\sigma_{f'}^g$ \cite[\S3.4, definition of associates, (i)]{MoyPrasad}. The image of
$\auGL_d(\auO_K)\cap g\auGL_d(\auO_K)g^{-1}$ in $\auGL_d(\auF_q)$ is the group of points of a
parabolic subgroup, proper unless $g$ fixes the vertex of $\auGL_d(\auO_K)$ (as in the proof of
\cite[Proposition~6.6]{MoyPrasad}); so $g\in K^\times\auGL_d(\auO_K)$, conjugation by $g$ induces an
inner automorphism of $\auGL_d(\auF_q)$, and $\sigma_f\simeq\sigma_{f'}$, that is $f=f'$. So two
depth-zero supercuspidals containing $\sigma_f$, $\sigma_{f'}$ are inertially equivalent (differ by an
unramified twist) if and only if $f=f'$.

\emph{Irreducibility.} Two supercuspidals $\pi_k$, $\pi_{k'}$ can be linked in Zelevinsky's sense
only if $\pi_{k'}\simeq\pi_k|\det|^{u}$, in particular only if they are inertially equivalent, so
only if $f_k=f_{k'}$, which does not happen because $c$ is squarefree. By Zelevinsky's
irreducibility criterion~\cite[Theorem~4.2]{Zelevinsky} (segments of length one, pairwise unlinked) the
induced representation $\pi_1\times\dots\times\pi_s$ is irreducible, and $\Pi$, which is a
subquotient of it, equals it. Its Langlands parameter is $\bigoplus_k\pgarec(\pi_k)$, with
monodromy $0$.

\emph{Inertia.} For a depth-zero supercuspidal $\pi$ of $\auGL_d(K)$ containing $\sigma_f$, $f$
irreducible with root $\xi$ generating $\auF_{q^d}$ over $\auF_q$ (that is, $\xi$ regular),
$\pgarec(\pi)\simeq\pgaInd_{W_{K_d}}^{W_K}\Lambda$ for a tamely ramified character $\Lambda$ of
$K_d^\times$ whose restriction to $\auO_{K_d}^\times$ is the inflation of the character
$\omega^{\pm1}$ of $\auF_{q^d}^\times$ named by $\xi$, the sign being fixed by the conventions
(the level zero case of the essentially tame correspondence: by
\cite[Corollary~3.3]{BHtame} the naive bijection is corrected by a tamely ramified character of
$K_d^\times$, and in level zero this character is the \emph{unramified} character
$\chi_2^{d-1}$, $\chi_2$ unramified of order $2$, \cite[\S7.5, (7.5.1)]{BHeffective}, which
does not change the restriction to inertia). Restricted to
$I_K=I_{K_d}$, the induced representation is $\bigoplus_{a<d}\Lambda^{q^a}\circ\mathrm{Art}^{-1}$
(Mackey), and $\Lambda\circ\mathrm{Art}^{-1}$ on $I_{K_d}$ factors through $\theta_d$; so the
eigenvalues at $s$ are the images under $\omega$ of the conjugates $\xi^{q^a}$, or uniformly of their
inverses. Summing over $k$ gives the
claim.
\end{proof}

\begin{lemma}[Parahoric level with a tame character]\label{pga:parahoric}
Let $n=n_1+n_2$, $\mathcal P\subseteq\auGL_n(\auO_K)$ the parahoric subgroup of the matrices that
are block lower triangular modulo $\varpi$ for the blocks $(n_1,n_2)$, $\psi$ a non-trivial
character of $\auF_q^\times$, and $\psi_{\mathcal P}(k)=\psi(\det k_{11}\bmod\varpi)$. If an
irreducible smooth $\Pi$ has a non-zero vector on which $\mathcal P$ acts by $\psi_{\mathcal P}$,
then $\Pi\simeq\Pi_1\times\Pi_2$ with $\Pi_1\otimes(\tilde\psi^{-1}\circ\det)$ unramified for an
extension $\tilde\psi$ of $\psi$ to a tame character of $K^\times$ and $\Pi_2$ unramified; the
monodromy of $\pgarec(\Pi)$ vanishes; and every $s\in I_K$ acts on $\pgarec(\Pi)$ by a matrix
$A$ with $(A-\tilde\psi(\mathrm{Art}^{-1}s)^{\epsilon})(A-1)=0$.
\end{lemma}
\begin{proof}
Let $P=MN$ be the upper block parabolic, $\bar N$ the opposite unipotent radical. $\mathcal P$ has
an Iwahori decomposition $\mathcal P=(\mathcal P\cap\bar N)(\mathcal P\cap M)(\mathcal P\cap N)$
with $\mathcal P\cap M=\auGL_{n_1}(\auO_K)\times\auGL_{n_2}(\auO_K)$ and $\psi_{\mathcal P}$ trivial on
the two unipotent parts. Let $\mathcal P_1$ be the pro-unipotent radical of $\mathcal P$ (the
kernel of reduction to $\auGL_{n_1}(\auF_q)\times\auGL_{n_2}(\auF_q)$) and $M_1=\mathcal P_1\cap M$.
We use the following statement (J), which is Moy--Prasad~\cite[Proposition~6.7, p.~116]{MoyPrasad}, and for
the Iwahori subgroup and trivial character is Borel's theorem:
\begin{enumerate}
\item[(J)] for every admissible $\Pi$ (in particular every irreducible smooth $\Pi$, irreducible smooth
representations of $\auGL_n(K)$ being admissible), the projection $\Pi\to\Pi_N$ restricts to an isomorphism
$\Pi^{\mathcal P_1}\to(\Pi_N)^{M_1}$, equivariant for $\mathcal P/\mathcal P_1=\mathcal P\cap M/M_1$.
\end{enumerate}
(The hypotheses of~\cite[6.3, Proposition~6.4, Proposition~6.7]{MoyPrasad}: $\mathcal P=\mathcal G_x$ is a
non-maximal parahoric ($n_1,n_2\ge1$) whose point $x$ lies in the apartment of the diagonal torus $S$; the
maximal $\auF_q$-split torus in the centre of its reductive quotient $\auGL_{n_1}\times\auGL_{n_2}$ is the
torus of block scalars, its lift $C\subset S$ has centraliser $M$, the block diagonal Levi subgroup, and
$P=MN$ may be either of the two opposite parabolics with Levi $M$; $\mathcal M_x=M\cap\mathcal G_x=
\mathcal P\cap M$, and $\mathcal M_x/\mathcal M_x^+\cong\mathcal G_x/\mathcal G_x^+$ by Proposition~6.4(2), so
$\mathcal M_x^+=M_1$ and $\mathcal G_x^+=\mathcal P_1$. Equivariance holds because the projection is
$M$-equivariant.)
(Casselman's notes \cite{CasselmanNotes} give the \emph{surjectivity} of $\Pi^{K_0}\to(\Pi_N)^{K_0\cap M}$
for open $K_0$ with an Iwahori factorisation (Theorem~3.3.3 there) and an isomorphism on the subspace
$\pi(K_0aK_0)\Pi^{K_0}$ (Proposition~4.1.4 there); injectivity on all of $\Pi^{\mathcal P_1}$ is the
additional content of (J).) Taking $(\mathcal P\cap M,\psi\circ\det\otimes1)$-isotypic parts in (J),
the isotypic vector of $\Pi$ gives a non-zero vector of $\Pi_N$ on which
$\auGL_{n_1}(\auO_K)\times\auGL_{n_2}(\auO_K)$ acts by $\psi\circ\det\otimes1$. As taking isotypic
parts for a compact open subgroup is exact and $\Pi_N$ has finite length, some irreducible
subquotient $\tau_1\otimes\tau_2$ of $\Pi_N$ has such a vector, so that $\tau_1\otimes(\tilde\psi^{-1}
\circ\det)$ and $\tau_2$ are spherical. The supercuspidal support of $\Pi$ is that of
$\tau_1\otimes\tau_2$ (every irreducible subquotient of the Jacquet module of an irreducible
representation has the supercuspidal support of that representation, Bernstein--Zelevinsky): $n_1$
characters with restriction $\psi$ to $\auO_K^\times$ and $n_2$ unramified characters. As $\psi\ne1$, no
segment of the first group is linked to one of the second, so by Zelevinsky's classification
$\Pi\simeq\Pi_1\times\Pi_2$ with $\Pi_1$, $\Pi_2$ irreducible with these two supports. By the
geometric lemma of Bernstein--Zelevinsky, the composition factors of the Jacquet module of
$\Pi_1\times\Pi_2$ for the parabolic of type $(n_1,n_2)$ are the representations
$(\alpha_1\times\alpha_2)\otimes(\beta_1\times\beta_2)$, with $\alpha_i\otimes\beta_i$ running over the
composition factors of Jacquet modules of $\Pi_i$; the first factor has supercuspidal support
consisting only of characters with restriction $\psi$ to $\auO_K^\times$ only when it is
$\Pi_1\otimes\Pi_2$ itself. Hence $\tau_1\otimes\tau_2\simeq\Pi_1\otimes\Pi_2$, so that
$\Pi_1\otimes(\tilde\psi^{-1}\circ\det)$ and $\Pi_2$ are spherical. Spherical
representations have parameters with monodromy $0$ and unramified; so $\pgarec(\Pi)=
\pgarec(\Pi_1)\oplus\pgarec(\Pi_2)$ has monodromy $0$, $I_K$ acts trivially on the second summand
and by $\tilde\psi\circ\mathrm{Art}^{-1}$ (to the power $\epsilon$) on the first.
\end{proof}

\begin{proposition}[(G4), (G5)]\label{pga:G45}
Let $y$, $S_y$ be as in Proposition~\ref{pga:G123} and $w$ a split
place, $w\mid p$, $v$ the place of $E^+$ below.
\begin{enumerate}
\item[(G4)] If $y$ is fixed by the principal congruence subgroup at $w$ and lies in the
$\sigma_c$-isotypic part for a label $c$, then $S_y|_{I_w}$ is tame and, if $w\nmid\ell$, of finite
order, semisimple, with the eigenvalues $\omega(\xi^{\epsilon})$, $\xi$ running over the roots of
$c$, on a generator $s$ of tame inertia with $\theta_n(s)$ a generator; if $w\mid\ell$, $S_y$ is
potentially crystalline at $w$, crystalline over $E_w(\varpi^{1/(q^d-1)})$, $d=\operatorname{lcm}_k(d_k)$,
with this inertial type. (Lemma~\ref{pga:depth0} is phrased with $\theta_d(s)$; since the norm
$\auF_{q^n}^\times\to\auF_{q^d}^\times$ is surjective, $\theta_n(s)$ a generator implies
$\theta_d(s)$ a generator.)
\item[(G5)] If $w\nmid\ell$ and $y$ is an eigenvector of the parahoric $\mathcal P$ of type
$(n_W,t)$ with character $\psi_{\mathcal P}$, $\psi=\psi_w\ne1$, then $(S_y(s)-\psi_w(s)^{
\epsilon})(S_y(s)-1)=0$ for $s\in I_w$, $\psi_w$ being viewed on $I_w$ through $\mathrm{Art}_w$.
\end{enumerate}
The sign $\epsilon$ is the same in (G4) and (G5) and for all $w$.
\end{proposition}
\begin{proof}
$\Pi_w=\pi_v$ ((BC) at split places). (G4): Lemma~\ref{pga:depth0} gives the parameter of
$\Pi_w\otimes|\det|^{(1-n)/2}$ (the twist is unramified); its monodromy is $0$, so by
Proposition~\ref{pga:G123} the monodromy of $S_y$ at $w$ is $0$ and $S_y|_{I_w}$ is the
restriction of the parameter, of finite order and semisimple, with the stated eigenvalues. For
$w\mid\ell$ we use only the relation $\prec$ of (GR) at $w$ (\cite[Theorem~A]{BLGGTlp2}, valid for
the polarized pairs $(\sigma_i|\det|^{s_i},|\cdot|^{2s_i})$ of Proposition~\ref{pga:G123}, so no
reduction to the case $\phi=1$ is needed). Let $\Sigma=\sigma_i|\det|^{s_i}$, $\rho=\iota\pgaWD
(r_\iota(\Sigma)|_{G_{E_w}})^{\mathrm{F\text{-}ss}}$ and $\rho'=\pgarec(\Sigma_w\otimes|\det|^{(1-
N)/2})$. Since $\Pi_w$ has a parameter with monodromy $0$, so has $\rho'$ (the parameter of $\Pi_w$
is the sum of those of the $\operatorname{Speh}(\sigma_{i,w},m_i)$, each containing that of
$\sigma_{i,w}$ up to an unramified twist, as in Proposition~\ref{pga:G123}). Thus in every
isotypic part $\rho'[\tau]$ all Jordan blocks have size $1$, and $\rho\prec\rho'$ gives first
$\rho^{\mathrm{ss}}\simeq\rho'^{\mathrm{ss}}$ (so $\rho[\tau]\ne0$ only if $\rho'[\tau]\ne0$) and then
that the largest block of $\rho[\tau]$ has size $\le1$; hence the monodromy of $\rho$ is $0$ and
$\rho\simeq\rho'$. Now $r_\iota(\Sigma)|_{G_{E_w}}$ is de Rham, hence potentially semistable
(Berger), and the monodromy of $D_{\mathrm{pst}}$ is that of $\pgaWD$, which is $0$; a potentially
semistable representation whose $D_{\mathrm{pst}}$ has $N=0$ is potentially crystalline, and it is
crystalline over a finite extension $L/E_w$ exactly when the Weil--Deligne representation is
unramified on $I_L$. The cyclotomic twists in $S_y$ are crystalline. So $S_y|_{G_{E_w}}$ is
potentially crystalline, and its inertial type is the restriction to $I_{E_w}$ of the parameter of
Lemma~\ref{pga:depth0}, which becomes trivial on $I_{E_w(\varpi^{1/(q^{d}-1)})}$. (For $n_i=1$,
$\Sigma$ is an algebraic Hecke character and the monodromy is $0$ for dimension reasons.) (G5):
Lemma~\ref{pga:parahoric} and Proposition~\ref{pga:G123} in the same way.
\end{proof}

\emph{The sign $\epsilon$ in Theorem~\ref{pg:I}.} Replacing every eigenvalue by its inverse does
not affect (M3)--(M5) nor Lemmas~\ref{pg:cyclic}, \ref{pg:burnside}: the inverse of a root of a
primitive label is a root of a primitive label, $\xi\mapsto\xi^{-1}$ preserves orders and the
property ``$\ne1$'', and $z\mapsto z^{-1}$ preserves $\auF_p^\times$. In Theorem~\ref{pg:III} the
sign must agree with the one in the definition of $\psi_\vartheta$ restricted to $I_w$; since
$\psi_\vartheta$ is itself defined through class field theory with the same $\mathrm{Art}_w$, and
(G5) involves the tame part of the same character $\mu_j^0\vartheta$ that defines the parahoric
character, the two signs coincide once $\epsilon$ is fixed, which we do by \emph{defining}
$\omega$ so that $\epsilon=1$.

\subsubsection{Statement (L): the Hecke polynomial of a lift}\label{pga:ss-L}

We state (L) for the Weil representation $\omega_\mu$ of the pair $(\auU(W),\auU(V))$ of
{\cite[\S2.1]{P54}}: $(g,h)$ acts by
$\Psi(x)\mapsto a(\det h)\,b(\det g)^{-1}\cdots\Psi(h^{-1}xg)$ at a split place, where
$a=\chi_W=\upsilon^{r_-}$ is a power of the splitting character $\upsilon$ ($r_-$ the number of
negative signs of $W$) and $b=\chi_V=\mu\,\upsilon^{\lfloor n/2\rfloor}$ with $\mu=\mu^0_j\vartheta$.
Thus $\mu$ enters only through $\det g$, $g\in\auU(W)$, and $\auU(V)$ acts through a splitting
given by a power of $\upsilon$ alone; this is formula {\cite[(1)]{P54}}. For an
automorphic character $c$ of $\auU(1)$ unramified at a split $w$ and $N\ge1$ put
\[
c^{(N)}_w=\bigl(c(\varpi_w/\aubar\varpi_w)\,\upsilon(\varpi_w)^{-N}\bigr)(\auNm w)^{N/2}
.
\]
Let $r=\operatorname{rk}W<n=\operatorname{rk}V$, $t=n-r$, $q=\auNm w$, $\alpha_w=\chi_{W,w}^{(r)}$,
$\beta_w=\chi_{V,w}^{(n)}$, and $P_w(X)=\sum_i(-X)^iq^{\binom i2}T_{w,i}=X^{\operatorname{rk}}
H_w(1/X)$ the reversed Hecke polynomial.

\begin{proposition}[(L)]\label{pga:L}
Let $w$ be split with $w,\aubar w$ off the level of the pair, let $y_0$ be a form on $\auU(W)$
spherical at $w$ with $P^W_w(X)y_0=p(X)y_0$, $p\in\C[X]$, and let $y=\Theta(\Phi)y_0$ be its theta
lift for a datum $\Phi$ which is the characteristic function of the integral matrices at $w$. Then
$y$ is spherical at $w$ and
\begin{equation}\label{pga:eq-L}
P_w^V(X)\,y=\prod_{l=0}^{t-1}(1-q^l\alpha_wX)\cdot p\bigl(q^t\alpha_w\beta_w^{-1}X\bigr)\,y .
\end{equation}
Equivalently, if $y\ne0$ and $\gamma_1,\dots,\gamma_r$ are the roots of $H^W_w(y_0;X)$, the roots
of $H^V_w(y;X)$ are $\alpha_wq^l$ ($0\le l<t$) and $q^t\alpha_w\beta_w^{-1}\gamma_k$; this is
\eqref{pg:eq-liftpoly} with $\lambda_w=\alpha_w$ and rescaling factor $q^t\alpha_w/\beta_w$.
Moreover the lift is $\auG_O$-equivariant as in {\cite[Proposition~14.27]{P46}}.
\end{proposition}
\begin{proof}
\emph{Coset representatives.} At the split place $w$, $\auU(X)(E^+_v)\simeq\auGL_N(E_w)$, $N=
\operatorname{rk}X$, and on vectors fixed by $K_w=\auGL_N(\auO_w)$ one has
\[
P_w(X)=\sum_{U}(-X)^{\dim U}q^{\binom{\dim U}2}\,R\bigl({}^t\gamma_U\bigr),\qquad
P_{\aubar w}(X)=\sum_U(-X)^{\dim U}q^{\binom{\dim U}2}\,R\bigl(\gamma_U^{-1}\bigr),
\]
the sums over the subspaces $U\subseteq\auF_q^N$, where $\gamma_U\in M_N(\auO_w)$ has
$\det\gamma_U$ of valuation $\dim U$ and the kernel of $\gamma_U\bmod\varpi$ is $U$, and
$\{\gamma_U:\dim U=i\}$ is a set of representatives of $K_w\backslash K_w
\operatorname{diag}(\varpi1_i,1_{N-i})K_w$ (the classical description of the Hecke operator
$T_i$ by lattices $\Lambda'$ with $\varpi\Lambda\subset\Lambda'\subset\Lambda$; the second formula is
the first transported by $g\mapsto{}^t\aubar g^{-1}$, which exchanges $w$ and $\aubar w$).

\emph{Translates of the vacuum.} Let $\mathbf 1$ be the characteristic function of
$M_{n\times r}(\auO_w)$. In the Schr\"odinger model at $w$ an element $h\in\auGL_n(E_w)$ acts by
$\mathbf 1\mapsto\chi_W(\det h)|\det h|^{-r/2}\cdots\mathbf 1(h^{-1}\,\cdot)$ and $g\in
\auGL_r(E_w)$ by $\mathbf 1\mapsto\chi_V(\det g)^{-1}|\det g|^{n/2}\cdots\mathbf 1(\cdot\,g)$ (the
character factors are those of the splitting: $a=\chi_W$ on the $V$ side, $b=\chi_V$ on the $W$
side). For $h={}^t\gamma_U$ this gives
\[
\omega({}^t\gamma_U)\mathbf 1(x)=\alpha_w^{\dim U}\,\mathbf 1\bigl({}^tx\,\gamma_U^{-1}\bigr)
\]
and for $g=\gamma_U^{-1}$ on the $W$ side, $\omega(\gamma_U^{-1})\mathbf 1(x)=\beta_w^{\dim U}
\mathbf 1(x\gamma_U^{-1})$; here $\alpha_w$, $\beta_w$ arise as $\chi(\varpi)q^{N/2}$ from the
character value at $\det\gamma_U$, of valuation $\dim U$, and from $|\det\gamma_U|^{-N/2}$.

\emph{The $q$-binomial identity.} For a matrix $z$ over $E_w$ with $J$ columns, $z\gamma_U^{-1}$ is
integral if and only if $z$ is integral and $U\subseteq\ker(z\bmod\varpi)$. By the $q$-binomial
theorem $\sum_{U\subseteq K}(-Y)^{\dim U}q^{\binom{\dim U}2}=\prod_{l<\dim K}(1-q^lY)$ for a
subspace $K$ of an $\auF_q$-space. Hence
\[
P^V_w(X/\alpha_w)\mathbf 1(x)=\mathbf 1(x)\prod_{l<n-\rho(x)}(1-q^lX),\qquad
P^W_{\aubar w}(Y/\beta_w)\mathbf 1(x)=\mathbf 1(x)\prod_{l<r-\rho(x)}(1-q^lY),
\]
$\rho(x)$ the rank of $x\bmod\varpi$ (for $P^V$ the kernel of ${}^tx\bmod\varpi$ has dimension
$n-\rho$, for $P^W$ that of $x\bmod\varpi$ has dimension $r-\rho$). As $n-\rho=t+(r-\rho)$,
$\prod_{l<n-\rho}(1-q^lX)=\prod_{l<t}(1-q^lX)\prod_{l<r-\rho}(1-q^{l+t}X)$, and therefore
\begin{equation}\label{pga:eq-vac}
P^V_w(X/\alpha_w)\,\mathbf 1=\prod_{l<t}(1-q^lX)\cdot P^W_{\aubar w}(q^tX/\beta_w)\,\mathbf 1 .
\end{equation}
(Both sides are also $K_w\times K_w$-fixed, the characters being unramified.)

\emph{Transfer to the lift.} By {\cite[(5)]{P54}}, $R_h\Theta(\Phi)y_0=\Theta(\omega(1,h)\Phi)
y_0$ and $\Theta(\omega(g,1)\Phi)y_0=\Theta(\Phi)(R_{g^{-1}}y_0)$. The first gives that $y$ is
$K_w$-fixed and $P^V_w(Z)y=\Theta(P^V_w(Z)\Phi)y_0$, with $P^V_w$ acting on the datum at $w$. The
second, summed with the coefficients of $P^W_{\aubar w}$, gives $\Theta(P^W_{\aubar w}(Z)\Phi)y_0
=\Theta(\Phi)(P^{W,*}_{\aubar w}(Z)y_0)$ where $P^{W,*}_{\aubar w}(Z)=\sum_U(\cdots)R_{\gamma_U}$
is the adjoint of $P^W_{\aubar w}(\bar Z)$ for the invariant pairing; on $K_w$-fixed vectors it is
$P^W_w(Z)$, because the adjoint of $T_{\aubar w,i}$ is $T_{w,i}$ (the double coset of
$\operatorname{diag}(\varpi1_i,1)^{-1}$ at $\aubar w$ is that of $\operatorname{diag}(\varpi 1_i,
1)$ at $w$ under the identification, and $\sum_{K\backslash KdK}R_{\gamma^{-1}}$ and
$\sum_{KdK/K}R_{\gamma}$ coincide on $K$-fixed vectors). Applying $\Theta(\cdot)y_0$ to \eqref{pga:eq-vac} with
$X=\alpha_wZ$,
\[
P^V_w(Z)y=\prod_{l<t}(1-q^l\alpha_wZ)\,\Theta(\Phi)\bigl(P^W_w(q^t\alpha_w\beta_w^{-1}Z)y_0\bigr)
=\prod_{l<t}(1-q^l\alpha_wZ)\,p(q^t\alpha_w\beta_w^{-1}Z)\,y .
\]
The statement about roots follows from $P(X)=X^{\deg}H(1/X)$. The $\auG_O$-equivariance is
(L2) of {\cite[Lemma~2.2]{P54}}.
\end{proof}

\emph{Comparison with {\cite[Lemma~2.2]{P54}}.} {\cite[Lemma~2.2(L3)]{P54}} proves that the single
factor $1-q^{t-1}\alpha_wX$ kills the lift, which is enough for the false lifts; it does not give
the factor $p(q^t\alpha_w\beta_w^{-1}X)$, which is what the trace identity needs.
{\cite[Lemma~2.2(L2)]{P54}} is the equivariance part of (L). Both statements, and the trace identity
below, are for the same representation $\omega_\mu$. (The lifts of the upper chain in
{\cite[\S5.1.1]{P49}} are normalised differently: there the character twists both groups. That
convention is not used for the single lift $W_j\to V_j$.)

\subsubsection{Statement (C): the central action on the lift from $\auU(1)$}\label{pga:ss-C}

Let $O=O_{\nu,z}(\vartheta)$, $W$ of rank one, $V$ of rank $n=2z$, $J_1$ the level of $\auU(W)=
\auU(1)$ at the chosen places, free by {\cite[Lemma~14.5]{P46}}. Write $k(g,h)$ for the theta kernel
($g\in\auU(1)(\auA_f)$, $h\in\auU(V)(\auA_f)$, with values in the weight space), so that the lift
of the constant function is $L(h)=\sum_{c}k(s_c,h)$ over representatives $s_c$ of
$\auU(1)(E^+)\backslash\auU(1)(\auA_{E^+,f})/J_1$ (the level is free, so no stabiliser weights),
and the descent of $\Xi$ is $D\Xi(g)=\langle\Xi,k(g,\cdot)\rangle$ up to conjugation, as in
{\cite[(6)]{P54}}. For $z\in\auU(1)(\auA_{E^+,f})$, viewed as the central element $z\cdot1_n$
of $\auU(V)(\auA_f)$, put $\varrho(z)\Xi=\chi_V(z)\,R_z\Xi$, where $\chi_V(z)$ is the value of the
character by which the Weil representation of the scalar pair $(z,z\cdot1_n)$ acts on the datum,
inverted.

\begin{proposition}[(C)]\label{pga:C}
\begin{enumerate}
\item[(a)] $\varrho(z)^{-1}k(g,\cdot)=k(gz,\cdot)$; hence $\varrho(z)L=L$.
\item[(b)] $\langle\varrho(z)\Xi,\Xi'\rangle=\langle\Xi,\varrho(z)^{-1}\Xi'\rangle$, and
$\varrho(z)$ commutes with every Hecke operator at a split place off the level and with the action
of $\auG_O$.
\item[(c)] If $z$ is integral above $\ell$ (and hence a unit there), $\varrho(z)^{\pm1}$ preserves
the forms integral at every $j\colon\C\to\C_\ell$.
\item[(d)] Every class in $\auU(1)(E^+)\backslash\auU(1)(\auA_{E^+,f})/J_1$ has a representative
integral above $\ell$.
\item[(e)] If $\Xi'$ is orthogonal to the false lifts in $\ker D$ and $s_1,\dots,s_{h_1}$ represent
the classes, then $\sum_c\varrho(s_c)^{-1}\Xi'$ is orthogonal to the false lifts in $(\C L)^\perp$.
\end{enumerate}
\end{proposition}
\begin{proof}
(a) The pair $(z,z1_n)\in\auU(W)\times\auU(V)$ acts on $\operatorname{Hom}(W,V)$ by $x\mapsto
z^{-1}xz=x$, so in the Weil representation it acts on every datum by the scalar given by the
splitting characters, which by definition is $\chi_V(z)^{-1}$: $\omega(z,z1_n)\Phi=\chi_V(z)^{-1}
\Phi$. Using the equivariance of the kernel, $k(gz,h)=$ (the kernel for the datum $\omega(z,1)\Phi$)
at $(g,h)$ $=\chi_V(z)^{-1}\cdot$(the kernel for $\omega(1,z^{-1}1_n)\Phi$) $=\chi_V(z)^{-1}k(g,hz^{-1})$,
i.e.\ $k(g z,\cdot)=\varrho(z^{-1})k(g,\cdot)=\varrho(z)^{-1}k(g,\cdot)$ . Summing over $g=s_c$: $\varrho(z)^{-1}L=\sum_ck(s_cz,\cdot)$,
and $c\mapsto$ class of $s_cz$ is a bijection of the class set (right translation by $z$ in the
abelian group $\auU(1)(\auA_f)$ permutes the double cosets; $k(\gamma g u,\cdot)=k(g,\cdot)$ for
$\gamma$ rational and $u\in J_1$ since the datum is $J_1$-invariant up to the character, which is
trivial on $J_1$ for the constant function); so $\varrho(z)^{-1}L=L$.

(b) $R_z$ preserves the inner product (the measure on $\Gamma\backslash G_f/J$ and the stabilisers
are translation invariant, and $z$ is central so $R_z$ preserves the level and the character), and
$|\chi_V(z)|=1$ ($\chi_V$ is unitary on $\auU(1)(\auA_f)$); so $\varrho(z)^*=\overline{\chi_V(z)}
R_{z^{-1}}=\varrho(z^{-1})$. Central translations commute with all translations, hence with Hecke
operators and with $\auG_O$, which acts by translations times characters of $\det$.

(c) $(\varrho(z^{-1})\Xi)(g)=\chi_V(z^{-1})\Xi(gz^{-1})$. If $\Xi$ is integral at $j$,
$j(\Xi(gz^{-1}))\in g_\ell z^{-1}_\ell W_\lambda(\auO_{\C_\ell})=g_\ell W_\lambda(\auO_{\C_\ell})$,
since $z_\ell^{\pm1}\cdot1_n\in\auGL_n(\auO_{E\otimes\Z_\ell})$ preserves $W_\lambda(\auO)$ (the
Schur matrices of integral matrices are integral). And $|j\chi_V(z)|=1$: $\chi_V$ is an algebraic
character of $\auU(1)$, so its value at a finite idele factors as $j\chi_V(z)=\chi_{V,\ell}(z)\cdot
\prod_{v\mid\ell}\theta_v(z_v)$, where $\chi_{V,\ell}$ is the $\ell$-adic avatar of $\chi_V$, a continuous
character of the compact group $\auU(1)(E^+)\backslash\auU(1)(\auA_{E^+,f})$ and hence unit-valued, and
each $\theta_v$ is an algebraic character of $\auU(1)(E^+_v)$, which takes unit values on the units
$z_v$. The same for $z^{-1}$.

(d) $\auU(1)(E^+)$ is dense in $\prod_{v\mid\ell}\auU(1)(E^+_v)$ (weak approximation for the torus
$\auU(1)$, which holds because $\auU(1)$ is a rational variety, being split by the quadratic
extension $E$; Cayley parametrisation). Given $g$, choose $\gamma\in\auU(1)(E^+)$ with $\gamma
g_v\in\auU(1)(\auO_{E^+_v})$ for all $v\mid\ell$ (an open subgroup); $\gamma g$ represents the same
class and is integral above $\ell$.

(e) Let $y$ be a false lift in $(\C L)^\perp$ at $w$: $y\in e\mathcal F(O)$, $H_w(\lambda_w)y=0$,
$y\perp L$. By (b), $\varrho(z)y\in e\mathcal F(O)\cap\ker H_w(\lambda_w)$ for every $z$. Put
$Y=\sum_c\varrho(s_c)y$. Then $D Y=0$: for every $g$, by (a) and (b),
\[
\langle Y,k(g,\cdot)\rangle=\Bigl\langle y,\sum_c\varrho(s_c)^{-1}k(g,\cdot)\Bigr\rangle=
\Bigl\langle y,\sum_ck(gs_c,\cdot)\Bigr\rangle=\langle y,\varrho(g)^{-1}L\rangle=\langle y,L\rangle=0,
\]
because $\sum_ck(gs_c,\cdot)=\varrho(g)^{-1}\sum_ck(s_c,\cdot)=\varrho(g)^{-1}L=L$ by (a) (the
group is abelian), and $y\perp L$. So $Y$ is a false lift in $\ker D$, and $\langle Y,\Xi'\rangle=0$ by
hypothesis; by (b) this is $\langle y,\sum_c\varrho(s_c)^{-1}\Xi'\rangle=0$. As $y$ was arbitrary,
the claim follows.
\end{proof}

This is statement (C) of Lemma~\ref{pg:line} exactly.

\subsubsection{Condition (M5) and the two normalisations}\label{pga:ss-M}

Recall from {\cite[\S14.1]{P46}} (data (D3), (D7)): $q_2=(q_0d)^{s+6}$ with $18480\mid q_0$,
$2^s=m^{22}$, $p_0,p_1\equiv1\pmod{q_2^2}$, $r\le3x^5$, $\sqrt x\le m\le2\sqrt x$; and $\chi_M$ has
order $q_2$ at every chosen place.

\begin{lemma}\label{pga:q2}
For $1\le j\le m$, $1\le u\le4r+2$, $b\ge0$ and every prime $\ell$: $q_2\nmid j\,u\,\ell^b$.
\end{lemma}
\begin{proof}
Let $p\in\{2,5\}$, $p\ne\ell$; $p\mid18480\mid q_0$, so $p^{s}\mid q_2$. If $q_2\mid ju\ell^b$ then
$p^s\mid ju$. But $ju\le m(4r+2)\le2\sqrt x(12x^5+2)<x^{11}\le m^{22}=2^s\le p^s$ (as $x\ge2^{20}$
by (S1)), a contradiction.
\end{proof}

\begin{proposition}[(M5)]\label{pga:M5}
In the proof of Theorem~\ref{pg:I}, (M5) holds for $R_{j,q}$ and for $R_{\nu,z}$.
\end{proposition}
\begin{proof}
Let $i$ with $p_i\ne\ell$, $p=p_i$, $w\mid p$ a chosen place whose label is primitive (case (a))
or $X-1$ (case (b)). By (G4) (Proposition~\ref{pga:G45}, $\epsilon=1$) and the twist by $\eta_j$,
$R(\gamma_w)$ is semisimple with eigenvalues $\omega(\xi')$, $\xi'$ running over the roots of the
rescaled label, that is, the roots $\xi$ of $c$ multiplied by the element $z\in\auF_p^\times$
naming the tame part of $\eta_j$ at $w$; this tame part is a power of $\chi_M$ (the twist being
trivial at the chosen places), so $z^{q_2}=1$.

(a) Let $\xi$ be a root of $c$; its order is $p^d-1$, $d$ the degree of the label. Put $e=(p^d-1)/(p-1)$ and $y_0=
\auNm_{\auF_{p^d}/\auF_p}(\xi z)=\xi^ez^d\in\auF_p^\times$. Every root $\xi'=(\xi z)^{p^a}$
satisfies $\xi'^e=y_0^{p^a}=y_0$. So $s=\gamma_w^e$ acts by the scalar $y=\omega(y_0)$. \emph{The
order of $y_0$ is divisible by $q_2^2$}: for a prime $p'\mid q_2$ let $p'^{\alpha}\|p-1$; then
$\alpha\ge2v_{p'}(q_2)$. The $p'$-primary component of $\xi^e$ (which has order $p-1$) has order
$p'^\alpha$, that of $z^d$ has order at most $p'^{v_{p'}(q_2)}<p'^\alpha$; in a cyclic $p'$-group
the product of an element of order $p'^\alpha$ and one of smaller order has order $p'^\alpha$. So
$q_2^2\mid\pgaord(y_0)$.

(b) The only root is $\xi'=z_w$, the element naming $\chi_M^{-j}$ at $w$ (the twist is trivial at
the chosen places); $s=\gamma_w$, $y=\omega(z_w)$, and $\pgaord(z_w)=q_2/\gcd(q_2,j)$.

In both cases: if $y_0^x=1$ then $q_2\mid jx$ (in (a) even $q_2^2\mid x$). Now
$y^{\ell^au}-1$ is a non-unit if and only if $y^{\ell^au}$ has $\ell$-power order, i.e.\ if and only
if $y_0^{\ell^{a+b}u}=1$ for some $b\ge0$; that would give $q_2\mid ju\ell^{a+b}$, which
Lemma~\ref{pga:q2} excludes for $1\le u\le4r+2$ ($\ge4n+2$ as $n\le r$). So $y^{\ell^au}-1$ is a
unit for all $a\ge0$, $1\le u\le4n+2$.
\end{proof}

Thus (M5) needs no further property of admissible data; it follows from $p_i\equiv1\pmod{q_2^2}$, $q_2=(q_0d)^{s+6}$
and $2^s=m^{22}$.

\begin{proposition}[First normalisation: the trace identity]\label{pga:trace}
With $\alpha_w,\beta_w$ as in Proposition~\ref{pga:L} for the pair $O_j(\vartheta)$, $n=r+t$:
\[
q^t\beta_w^{-1}=\psi_\vartheta(\pgFrob_w)\,\eta_j(\pgFrob_w),\qquad
\eta_j(y)=(\mu_j^0)^{-1}(y/\aubar y)\,|y|^{(r-t)/2},
\]
and $w\mapsto\alpha_w^{-1}$ is the Frobenius value of the $\ell$-adic character $\eta_O$ of the
algebraic Hecke character $y\mapsto(\chi_W(y/\aubar y)\upsilon(y)^{-r}|y|^{-r/2})^{-1}$. Hence
\eqref{pg:eq-trace} holds with this $\eta_O$ and $R_{j,q}=S_q\otimes\eta_j$.
\end{proposition}
\begin{proof}
$n=r+t$ is even ($r$, $t=5r$ odd). For even $N$ the space character of rank $N$ of the
character $\mu$ (the character $\chi_V$) is the
base change $y\mapsto\mu(y/\aubar y)$, the powers of $\upsilon$ cancelling; so $\beta_w=\mu(\varpi_w/\aubar\varpi_w)q^{n/2}$ with
$\mu=\mu_j^0\vartheta$. Hence $q^t\beta_w^{-1}=(\mu_j^0)^{-1}(\varpi_w/\aubar\varpi_w)\,
\vartheta^{-1}(\varpi_w/\aubar\varpi_w)\,q^{(t-r)/2}$, and $|\varpi_w|^{(r-t)/2}=q^{(t-r)/2}$:
this is $\eta_j(\pgFrob_w)\psi_\vartheta(\pgFrob_w)$ by definition of $\psi_\vartheta$
(\S\ref{pg:ss-reduction}) and of $\eta_j$. By (G1) and
Proposition~\ref{pga:L}, for split $w$ off the level the eigenvalues of $S_y(\pgFrob_w)$ are
$\alpha_wq^l$ and $q^t\alpha_w\beta_w^{-1}\gamma_k$, those of $S_q(\pgFrob_w)$ being the
$\gamma_k$; multiplying by $\eta_O(\pgFrob_w)=\alpha_w^{-1}$ gives $\chi(\pgFrob_w)^l$ ($l<t$) and
the eigenvalues of $(\psi_\vartheta\eta_jS_q)(\pgFrob_w)$. Traces agree on a set of Frobenius
elements of density one, hence everywhere (Chebotarev, continuity). The only properties of
$\eta_O$ used later are that it is the character of an algebraic Hecke character depending on the
pair only; it is unramified at every finite place not above $\ell$ (Lemma~\ref{pay:eta}).
\end{proof}

\begin{proposition}[Second normalisation: the weights]\label{pga:weights}
Let $\epsilon_\sigma=1$ for $\sigma$ in the CM type and $-1$ otherwise. The Hodge--Tate weights of
$R_{j,q}$ at $\sigma$ are
\[
\tfrac{t-1}2-\epsilon_\sigma\bigl(\varsigma_{\sigma,i}+\kappa_j(\sigma)\bigr),\qquad1\le i\le r,
\]
and (M2) holds. For $R_{\nu,z}$ the weight at $\sigma$ is $z-1-\epsilon_\sigma c(\sigma)$, $-c$ the
weight of $\nu$, and (M2) holds.
\end{proposition}
\begin{proof}
By (G3) for $y_0$ on the definite $\auU(W_m)$ (rank $r$, parameter $\varsigma$), $S_q$ has weights
$\frac{r-1}2-\epsilon_\sigma\varsigma_{\sigma,i}$. By (G3) in rank one (same convention: the
parameter of a character of $\auU(1)$ is its weight), the character $y\mapsto(\mu_j^0)^{-1}(y/\aubar
y)$, of $\auU(1)$-weight $+\kappa_j$, has weight $-\epsilon_\sigma\kappa_j(\sigma)$; and $|y|^{(r-t)/
2}$ has weight $\frac{t-r}2$ at every $\sigma$ (it is $\chi^{(t-r)/2}$ times a character of finite
order, up to the unitary normalisation). Adding, the weights of $R=S_q\otimes\eta_j$ are as
stated. (The two signs come from one and the same convention (G3), so a change of convention
$\sigma\leftrightarrow\aubar\sigma$ changes both, and the relative sign is convention-free.
Consistency check: the weights of $R$ at $\sigma$ and $\aubar\sigma$ pair to $t-1$, as (M1)
requires.)

Now $\varsigma_{\sigma,i}+\kappa_j=\kappa_j+(\frac k2-i)$ for $i<k$ and $=(j-i'+\frac12)\Delta$ for
the $i'$-th entry of the second group ($1\le i'\le m$). With $|\frac{t-1}2|<t$, $k\le r$ and
$t=5r$: the first $k-1$ weights $u$ satisfy $|u+\epsilon_\sigma\kappa_j|\le\frac k2+\frac{t-1}2\le
t$ (high, with the sign of (M2)); the other $m$ satisfy $\bigl||u|-|j-i'+\frac12|\Delta\bigr|\le
\frac{t-1}2<3r$, and $\frac12\Delta\le|j-i'+\frac12|\Delta\le(m-\frac12)\Delta$ with $\Delta=2r^8$, so
$\Delta/4\le|u|\le2m\Delta$ (low). Weights with $|u|\le2h^2=2r^{16}$: not the high ones, since
$\kappa_j\ge\kappa\ge xb/x_1>2r^{16}+t$ by (S1); all low ones, since $2m\Delta\le2r^{16}$; so there
are exactly $m\le rm/k$. Negative weights: at $\sigma$ in the CM type the $k-1$ high ones and the
low ones with $i'\le j$ ($j$ of them); at $\aubar\sigma$ the low ones with $i'>j$ ($m-j$); in total
$k-1+m=r$. For $R_{\nu,z}$ the same computation in rank one, with the character
$\nu^{-1}(y/\aubar y)|y|^{1-z}$, gives $z-1-\epsilon_\sigma c(\sigma)$; it is high since
$|c-\kappa_j|\le t$ and $z\le r$, and exactly one of the weights at $\sigma,\aubar\sigma$ is
negative.
\end{proof}

\emph{The convention.} By Lemma~\ref{pax:conv}, the automorphic representation attached to a form
of $\mathcal A_\lambda$ has $\pi_\infty\simeq\rho_\lambda\circ\tau$ at $\tau$, of parameter
$\varsigma=\lambda+\rho$, uniformly in the rank, including rank one; so the parameters of
$\mathcal A(W_m)$ and the weight of $\mu^0_j$ are measured in the same convention, and the relative
sign above is justified.

\subsubsection{The published inputs of Theorem~\ref{pg:I}}\label{pga:ss-status}

Theorem~\ref{pg:I} follows from the proof of \S\ref{pg:ss-galois} supplemented by
Propositions~\ref{pga:L} (statement (L)), \ref{pga:M5} ((M5)), \ref{pga:trace} and
\ref{pga:weights} (the two normalisations), Proposition~\ref{pax:sign} (the twist for Speh
constituents) and Lemma~\ref{pax:conv} (the convention), and the following published results:
\begin{enumerate}
\item[(1)] the published theorems quoted in Theorem~\ref{pga:lit} (base change of Labesse;
Galois representations and local--global compatibility of Harris--Taylor, Shin, Chenevier--Harris,
Caraiani, Barnet-Lamb--Gee--Geraghty--Taylor);
\item[(2)] the local results quoted in Lemmas~\ref{pga:depth0} and \ref{pga:parahoric}: Moy--Prasad
on depth-zero types and Jacquet functors, Zelevinsky's irreducibility criterion, the level-zero
case of the tame local Langlands correspondence (Bushnell--Henniart), Jacquet's lemma;
\end{enumerate}
Statement (C) (Proposition~\ref{pga:C}) is proved in full; it is used in Lemma~\ref{pg:line}, not
in Theorem~\ref{pg:I}. {\cite[Proposition~14.36]{P46}} is used in the proof of Theorem~\ref{pg:main},
not in Theorem~\ref{pg:I}.



\section{Odd classes and Kummer maps}\label{C3:g6}
\providecommand{\auO}{\mathcal O}
\providecommand{\pgcodd}{c_{\mathrm{odd}}}
\providecommand{\pgG}{G_E}
\providecommand{\pgcyc}{\varepsilon_\ell}
\providecommand{\prilen}{\operatorname{len}}
\providecommand{\prival}{\operatorname{val}}
\providecommand{\anrd}{\mathrm{rd}}
\providecommand{\panotodd}[1]{\par\smallskip{\raggedright\noindent\textbf{[Not proved here:} #1\textbf{]}\par}\smallskip}
\providecommand{\podSha}{\mathrm{Sha}^-}
\providecommand{\podC}{\mathcal C}
\providecommand{\podm}{\mathfrak m}
\providecommand{\podord}{\operatorname{ord}}
\providecommand{\podfr}[1]{\langle #1\rangle}
\providecommand{\podth}{\theta}
\providecommand{\podGal}{\operatorname{Gal}}
\providecommand{\podLam}{\Lambda}

\subsection{The odd classes}\label{pod:sec}

This section proves Proposition~\ref{pri:odd} of \S\ref{pri:sec} (the odd classes), for every prime
$\ell$, the prime $2$ and the primes dividing $q_1d$ included. The method is that of Kolyvagin's Euler
system of Gauss sums, as presented by Rubin~\cite[\S\S1--4]{RubinTexel} (there for $\Q(\mu_p)$), made effective at finite level and
combined with a crude class-field-theoretic count that disposes of the small primes. No form of the
main conjecture of Iwasawa theory is used.

\subsubsection{Statement and plan}\label{pod:ss-statement}

\emph{Conventions.} $\Omega=\overline\Q\subset\C$, $G_\Q=\podGal(\Omega/\Q)$, and $c\in G_\Q$ is
the complex conjugation induced by the fixed embedding. For $M\ge1$, $\varepsilon_M\colon G_\Q\to
(\Z/M)^\times$ is the cyclotomic character modulo $M$ and $H_M=\ker\varepsilon_M=G_{\Q(\mu_M)}$. For a
prime $\ell$, $\pgcyc\colon G_\Q\to\Z_\ell^\times$ is the $\ell$-adic cyclotomic character. Places of
$\Omega$ are its valuation rings; for a place $A$ of $\Omega$, $D_A\supseteq I_A$ are its
decomposition and inertia groups in $G_\Q$ (in $G_K$ for a subfield $K$, by intersection), and
$\podord_A$ is the valuation of $A$ normalised by $\podord_A(r)=1$ for the rational prime $r$ below
$A$. A \emph{Frobenius element at $A$} is a $\phi\in D_A$ acting on the residue field of $A$ as
$t\mapsto t^r$. Throughout, $N=q_1d$, and $E\subseteq\Q(\mu_N)$ is the CM field of
{\cite[\S14.3]{P46}} (it is abelian over $\Q$, contained in $\Q(\zeta_{q_1d})$, and totally complex);
$h_E$ is its class number.

$\auO$ is the ring of integers of a finite extension of $\Q_\ell$, with maximal ideal $\podm$,
uniformiser $\varpi$, valuation $\prival=\prival_\auO$ ($\prival(\varpi)=1$) and \emph{unnormalised}
length $\prilen$ (so $\prilen\auO/a=\prival(a)$). We assume that $\auO$ contains the values of all
Dirichlet characters modulo $N$. Put $\podLam_\nu=\auO/\ell^\nu\auO$. A Dirichlet character $\psi$
modulo $N$ (not necessarily primitive, $\psi(a)=0$ for $(a,N)>1$) is viewed as a character of
$G_\Q$ through $\varepsilon_N$, and for $e\in\Z$ we write
\[
\Upsilon_{e,\psi}=\pgcyc^{\,e}\,\psi\colon G_\Q\to\auO^\times .
\]
On an $\auO$-module killed by $M'$, an element $g$ with $\varepsilon_{NM'}(g)=a$ acts through
$\Upsilon_{e,\psi}$ as the scalar $a^e\psi(a)$, because $\pgcyc(g)\equiv a\pmod{\ell^{v_\ell(M')}}$.
For $x\in\Q/\Z$ let $\podfr x\in[0,1)$ be its representative, and for an integer $c\ge1$
\[
m_c(x)=\lfloor c\podfr x\rfloor-\lfloor c\podfr{-x}\rfloor\in\Z .
\]
For $M\ge1$ the units $u\in(\Z/M)^\times$ are represented in $[0,M)$, and
\[
\podth_{c,M}(e,\psi)=\sum_{u\in(\Z/M)^\times}m_c(u/M)\,u^e\psi(u)\in\auO
\]
is the \emph{twisted Stickelberger sum}.

A \emph{cocycle for a character} $\chi\colon G\to\auO^\times$ with values in an $\auO$-module $V$
is a map $x\colon G\to V$ with $x(gh)=x(g)+\chi(g)x(h)$; all cocycles below are locally constant
(continuous for the discrete topology on $V$). The module of Proposition~\ref{pri:odd} is, for
$\nu\ge0$ and $0\le e$:
\begin{definition}\label{pod:def-sha}
For $e\ne0$, $\podSha_{e,\nu}$ is the $\auO$-module of locally constant cocycles $z\colon\pgG\to
\podLam_\nu$ for $\pgcyc^{-e}|_{\pgG}$ (that is, with values in $\podLam_\nu(-e)$) such that (i) $z$
vanishes on $I_A\cap\pgG$ for every place $A$ of $\Omega$ not above $\ell$; (ii) $z$ vanishes on
the elements $g\in D_A\cap\pgG$ with $\pgcyc(g)=1$, for every place $A$ above $\ell$; (iii) $z(cgc)=
-z(g)$ for all $g\in\pgG$. For $e=0$, $\podSha_{0,\nu}$ is the module of continuous homomorphisms
$\pgG\to\podLam_\nu$ that vanish on every inertia group.
\end{definition}
This is the module of Proposition~\ref{pri:odd}: ``unramified outside $\ell$'' is (i), and $g^c=cgc$.
No comparison with a Tate--Shafarevich group is needed or made.

\begin{proposition}[$=$ Proposition~\ref{pri:odd}]\label{pod:odd}
Let $t\ge0$ and $y=N+t+3$. There is an integer $0<\pgcodd\le2^{4y^{27}}h_E$ such that for every
prime $\ell$, every $\auO$ as above, all $\nu\ge0$ and all $0\le e<t$,
\[
\prilen_\auO\podSha_{e,\nu}\le\prival_\auO(\pgcodd).
\]
\end{proposition}

\emph{Plan.} For $e=0$ the class number serves (\S\ref{pod:ss-e0}). For $e\ne0$ the module is split,
over $\Q$, into components indexed by the characters $\psi$ modulo $N$ (\S\ref{pod:ss-split}); only
the components with $\Upsilon_{e,\psi}$ odd matter. An integer $X_e$ of size $\le2^{(N+e+3)^6}$ with
a ``Stickelberger property'' is produced from Bernoulli numbers and Voronoi's congruence
(\S\ref{pod:ss-voronoi}). With Gauss and Jacobi sums (\S\ref{pod:ss-gauss}) and Artin reciprocity,
$4X_e^2$ kills every odd component, at \emph{every} prime $\ell$ (\S\ref{pod:ss-annih}); so
$Y_e=16\varphi(N)X_e^2$ kills $\podSha_{e,\nu}$. To pass from an annihilator to a length one needs
a bounded number of elements of $\pgG$ at which a cocycle that vanishes is zero (``detecting
points''). For $\ell\le y^2$ a count of generators of a ray class group gives $\le y^{20}$ such points
(\S\ref{pod:ss-detect}); for $\ell>y^2$ Kolyvagin's derivative classes give $v_\ell(X_e)+1$ points for
each odd component (\S\S\ref{pod:ss-koly}--\ref{pod:ss-descent}). The two regimes are put together
in \S\ref{pod:ss-proof}.

\begin{remark}[The prime $2$ and the primes dividing $N$]\label{pod:rem-two}
For $\ell\le y^2$, in particular for $\ell=2$ and for every $\ell\mid N$, only two ingredients are
used: the annihilator $Y_e$ (Proposition~\ref{pod:annih} and Lemma~\ref{pod:split}), whose proof
makes no assumption on $\ell$ (no division by $2$, by $\varphi(N\ell)$ or by $\ell-1$ occurs; the
factor $4$ in $4X_e^2$ and in $4\varphi(N)$ is exactly what makes this possible), and the detecting
points of Lemma~\ref{pod:detect}, which come from class field theory and are valid for all $\ell$.
The Euler-system descent, which needs $2$, $\varphi(N\ell)$ and some $a^e-1$, $a^{e+1}-1$ to be units
of $\auO$, is only used for $\ell>y^2\ge N+e+3$. The price for small $\ell$ is the exponent $y^{20}$
in place of $1$, which is absorbed in the bound for $\pgcodd$. We do not know a published statement
covering $\ell=2$ in this form.
\end{remark}

\subsubsection{Generalities on cocycles}\label{pod:ss-gen}

\begin{lemma}\label{pod:coc}
Let $G$ be a profinite group, $\chi\colon G\to\auO^\times$ a continuous character, $V$ an
$\auO$-module and $x\colon G\to V$ a locally constant cocycle for $\chi$.
\begin{enumerate}
\item[(a)] The zeros of $x$ form an open subgroup $Z(x)$ of $G$; in particular $x$ vanishes on the
closed subgroup generated by any set on which it vanishes.
\item[(b)] If $g,s\in G$ and $\chi(g)$ fixes $x(s^{-1})$ \textup(e.g.\ $\chi(g)$ acts trivially on
$V$\textup), then $x(sgs^{-1})=\chi(s)x(g)$.
\item[(c)] \textup(Sah\textup) If $x$ vanishes at the commutator $k=\delta g\delta^{-1}g^{-1}$ and
$\chi(k)=1$, then $(\chi(\delta)-1)x(g)=(\chi(g)-1)x(\delta)$.
\item[(d)] Let $\gamma\in G$ with $x(\gamma)=0$ and $(\chi(\gamma)-1)x\ne0$. Then $x$ does not vanish
identically on any closed normal subgroup $H$ with $G/H$ abelian.
\item[(e)] If $g\in G$ satisfies $g^2=1$ and $\chi(g)$ acts trivially on $V$, then $2x(g)=0$.
\end{enumerate}
\end{lemma}
\begin{proof}
(a) $x(1)=0$; if $x(g)=x(h)=0$ then $x(gh)=0$ and $0=x(gg^{-1})=x(g)+\chi(g)x(g^{-1})$; the set is
open as $x$ is locally constant. (b) $x(sgs^{-1})=x(s)+\chi(s)x(g)+\chi(sg)x(s^{-1})$ and $x(s^{-1})
=-\chi(s)^{-1}x(s)$; as $\chi(g)x(s^{-1})=x(s^{-1})$ the first and last terms cancel. (c) Expand
$x(\delta g)=x(kg\delta)$: $x(\delta)+\chi(\delta)x(g)=x(k)+\chi(k)\bigl(x(g)+\chi(g)x(\delta)\bigr)$.
(d) If $x|_H=0$, then for $g\in G$ the element $h=g^{-1}\gamma^{-1}g\gamma$ lies in $H$ and
$g\gamma=\gamma gh$, so $x(g)=x(g)+\chi(g)x(\gamma)=x(g\gamma)=x(\gamma g)+\chi(\gamma g)x(h)=
\chi(\gamma)x(g)$; thus $(\chi(\gamma)-1)x=0$. (e) $x(g)=x(g^3)=x(g)+\chi(g)x(g)+\chi(g)^2x(g)=3x(g)$.
\end{proof}

\begin{lemma}\label{pod:length}
Let $M$ be an $\auO$-submodule of the maps $G\to\podLam_\nu$, killed by an integer $m\ge1$, and let
$p_1,\dots,p_k\in G$ be such that $x\in M$ and $x(p_i)=0$ for all $i$ imply $x=0$. Then
$\prilen M\le\prival(m^k)$.
\end{lemma}
\begin{proof}
$x\mapsto(x(p_i))_i$ embeds $M$ into $\podLam_\nu[m]^k$, and $\podLam_\nu[m]\cong\auO/(\ell^\nu,m)$
has length $\le\prival(m)$.
\end{proof}

\subsubsection{Detecting points from class field theory}\label{pod:ss-detect}

Let $K$ be a totally complex number field, $S$ a finite set of finite places of $K$ and $M\ge1$.
Call a map $x\colon G_K\to V$ to an abelian group an \emph{$(S,M)$-homomorphism} if it is a locally
constant homomorphism, $Mx=0$, and $x$ vanishes on the inertia groups at all finite places of $K$
outside $S$.

\begin{lemma}\label{pod:cft}
Let $U_S^{(M)}$ be the group of finite ideles of $K$ that are $M$-th powers at the places of $S$.
Then $[\mathbb A_{K,f}^\times:U_S^{(M)}]\le M^{2\#S+[K:\Q]}$, and there are finitely many $p_i\in
G_K$ and integers $n_i\ge2$ with $\prod_in_i$ dividing $h_K[\mathbb A_{K,f}^\times:U_S^{(M)}]$,
such that every $(S,M)$-homomorphism $x$ satisfies $n_ix(p_i)=0$ for all $i$, and $x=0$ if
$x(p_i)=0$ for all $i$.
\end{lemma}
\begin{proof}
The index is a product over $v\in S$ of $[K_v^\times:K_v^{\times M}]=M\,\#\mu_M(K_v)\,|M|_v^{-1}$ ($|\cdot|_v$ normalised, so that
$|M|_v^{-1}=\#\auO_v/M\auO_v$). (Proof: for an endomorphism $\phi$ of an abelian group $A$ with finite kernel
and cokernel put $h(\phi)=\#\operatorname{coker}\phi/\#\ker\phi$; $h$ is multiplicative in short exact
sequences preserved by $\phi$ and equals $1$ on finite groups. Apply this to $\phi=(\cdot)^M$ on
$K_v^\times=\varpi^\Z\times\auO_v^\times$: on $\varpi^\Z$, $h=M$; let $p$ be the residue characteristic, $e$
the absolute ramification index and $k>e/(p-1)$; the $p$-adic logarithm is an isomorphism of
$U=1+\varpi^k\auO_v$ onto the additive group $\varpi^k\auO_v\cong\auO_v$ (inverse: the exponential), so
$h((\cdot)^M|_U)=\#\auO_v/M\auO_v=|M|_v^{-1}$; as $\auO_v^\times/U$ is finite,
$h((\cdot)^M|_{\auO_v^\times})=|M|_v^{-1}$; and $\ker\phi=\mu_M(K_v)$, so
$\#\operatorname{coker}\phi=M\,\#\mu_M(K_v)\,|M|_v^{-1}$. See also \cite[Chapter~II, Corollary~(5.8)]{Neukirch}.) Further $\prod_{v\in S}|M|_v^{-1}\le M^{[K:\Q]}$ and
$\#\mu_M(K_v)\le M$. Let $\mathcal A$ be the subgroup of the finite ideles generated by $K^\times$
and by the ideles that are units at every finite place outside $S$ and lie in $U_S^{(M)}$. Since $K$
is totally complex, the Artin map $\mathrm{Art}\colon\mathbb A_{K,f}^\times/\overline{K^\times}\to
G_K^{\mathrm{ab}}$ is surjective (global class field theory: the archimedean components have trivial Artin image, $K$
being totally complex; for every finite abelian extension $L/K$ the norm residue symbol $\mathbb A_K^\times\to
\mathrm{Gal}(L/K)$ is surjective \cite[Chapter~VI, Theorem~(5.5)]{Neukirch}, so the image is dense; and it
is compact, the source being a quotient of the compact group $C_K^0$ of idele classes of absolute value $1$
\cite[Chapter~VI, Theorem~(1.6)]{Neukirch}). For an $(S,M)$-homomorphism $x$ and any
homomorphism $f\colon V\to\Q/\Z$, $f\circ x$ is a character of $G_K$ that kills
$\mathrm{Art}(\mathcal A)$: it kills $\mathrm{Art}(K^\times)$ by reciprocity, the local Artin images
of units outside $S$ (inertia), and the images of $M$-th powers. Hence $f\circ x$ factors through
the finite abelian group $Q=G_K^{\mathrm{ab}}/\mathrm{Art}(\mathcal A)$, and since the $f$ separate
points of $V$, so does $x$. Write $Q=\prod_i\Z/n_i$ with $n_i\ge2$ and let $p_i\in G_K$ lift the
generators. The order of $Q$ divides $[\mathbb A^\times_{K,f}:\mathcal A]$, which divides
$[\mathcal AU_\emptyset:\mathcal A]$ times $[\mathbb A^\times_{K,f}:\mathcal AU_\emptyset]$, where
$U_\emptyset$ are the ideles that are units at all finite places outside $S$ (no condition at $S$);
the second index divides $h_K$ (the quotient by $K^\times U_\emptyset$ is a quotient of the class
group), and the first divides $[U_\emptyset:U_\emptyset\cap U_S^{(M)}]$, which divides the index of
$U_S^{(M)}$.
\end{proof}

\begin{lemma}\label{pod:hK}
For $M\ge3$ and $K=\Q(\mu_M)$, $h_K\cdot[\mathbb A_{K,f}^\times:U_S^{(M)}]\le2^{M^3+5M^2+M}$ when
$S$ is the set of places above one rational prime.
\end{lemma}
\begin{proof}
$d=[K:\Q]=\varphi(M)\le M$, $|d_K|\le M^d\le(2^M)^M$, and $\#S\le d$. Every ideal class contains an
integral ideal of norm $\le|d_K|^{1/2}\le|d_K|$ (Minkowski), and there are at most $n^d$ integral
ideals of norm $n$ (at most $d$ prime ideals above each prime and $\le$ the number of factorisations,
crudely $\le n^d$), so $h_K\le|d_K|^{d+1}\le2^{M^2(M+1)}$. The index is $\le M^{3M}\le2^{3M^2}$ by
Lemma~\ref{pod:cft}. Hence the product is $\le2^{M^3+4M^2}$.
\end{proof}
Any bound $2^{O(M^3)}$ suffices below.

\begin{lemma}[Detecting points]\label{pod:detect}
Let $\ell$ be a prime and $y\ge N$ with $\ell\le y^2$. There is a set $s\subseteq\pgG$ with $\#s\le
y^{20}$ such that for all $\auO$, $\nu$ and $e$, every locally constant cocycle $z\colon\pgG\to
\podLam_\nu$ for $\pgcyc^{-e}$ that vanishes on $I_A\cap\pgG$ for all places $A\nmid\ell$ and
vanishes on $s$ is zero. Moreover, if $g_0\in\pgG$ topologically generates $\pgG/(\pgG\cap\ker
\pgcyc)$, there is such a set $s'$ with $\#s'\le y^{20}$, contained in $\ker\pgcyc$, such that
$z(g_0)=0$ and $z|_{s'}=0$ imply $z=0$.
\end{lemma}
\begin{proof}
Put $M=N\ell^2\ge3$ and $K=\Q(\mu_M)\supseteq E$, so $G_K=H_M\subseteq\pgG$; let $S$ be the places of
$K$ above $\ell$. Let $s$ consist of a set of representatives $\{\rho_t\}$ of $\pgG/G_K$ (at most
$\varphi(M)$ elements) and the points $p_i$ of Lemma~\ref{pod:cft} for $(K,S,M)$; by
Lemma~\ref{pod:hK} there are at most $\log_2\bigl(h_K[\mathbb A^\times_{K,f}:U_S^{(M)}]\bigr)\le
M^3+5M^2+M$ of the latter. Since $M\le N\ell^2\le y^5$, $\#s\le M+M^3+5M^2+M\le M^4\le y^{20}$.

Let $z$ vanish on $s$, and suppose $\ell^kz=0$; we show $z=0$ by induction on $k$, the case $k=0$
being trivial. The cocycle $Mz$ vanishes on $s$ and is killed by $\ell^{k-1}$, since
$\ell^{k-1}Mz=N\ell\cdot\ell^kz=0$; so $Mz=0$ by induction. For $h\in G_K$ one has
$\pgcyc(h)\equiv1\pmod{\ell^{v_\ell(M)}}$, and $z(g)$ is killed by $M$, so $\pgcyc(h)^{-e}z(g)=z(g)$. Thus $z(hg)=z(h)+z(g)$ for $h\in G_K$, $g\in
\pgG$; in particular $z|_{G_K}$ is an $(S,M)$-homomorphism (an inertia group of $K$ at a place not
above $\ell$ lies in one of $E$ at a place not above $\ell$). It vanishes at the $p_i$, so
$z|_{G_K}=0$ by Lemma~\ref{pod:cft}. For $g\in\pgG$ write $g=h\rho_t$ with $h\in G_K$; then
$z(g)=z(h)+z(\rho_t)=0$.

For the second statement, write each $p\in s$ as $p=h_pq_p$ with $h_p\in\ker\pgcyc$ and $q_p$ in the
closed subgroup $\overline{\langle g_0\rangle}$ (possible as $g_0$ generates modulo the kernel and
$\pgcyc$ is continuous with closed image), and let $s'=\{h_p\}$. If $z(g_0)=0$ and $z|_{s'}=0$, then
$z$ vanishes on $\overline{\langle g_0\rangle}$ by Lemma~\ref{pod:coc}(a), hence at $p=h_pq_p$, and
$z=0$ by the first part.
\end{proof}

\begin{corollary}[Proposition~\ref{pri:kummer}(b) for $\ell\le(N+t+2)^2$]\label{pod:kummer-b}
Let $h=r^8$ as in \S\ref{pg:ss-ribet}, $t=5r$, and let $\ell\le y_0^2$ with $y_0=N+t+2$. With $g_0$ a
topological generator of $\pgG$ modulo $\ker\pgcyc$, there is $s\subseteq\ker\pgcyc$ with
$(t+1)(\#s+1)\le y_0^{21}\le h^{21}$ such that a locally constant cocycle $\pgG\to\podLam_\nu(-e)$
unramified outside $\ell$ that vanishes at $g_0$ and on $s$ is zero, for all $\auO,\nu,e$.
\end{corollary}
\begin{proof}
The second statement of Lemma~\ref{pod:detect} with $y=y_0\ge N$ gives $s\subseteq\ker\pgcyc$ with $\#s\le y_0^{20}$. As $t+2\le y_0$ and
$y_0\le y_0^{20}$: $(t+1)(\#s+1)\le(y_0-1)y_0^{20}+(t+1)\le y_0^{21}-y_0^{20}+y_0\le y_0^{21}$.
Finally $y_0\le h$: by Remark~\ref{pod:rem-size}, $N\le x^{15}$, and $x^5\le r$, so $y_0\le x^{15}+5r+2
\le r^3+5r+2\le r^8$.
\end{proof}

\begin{remark}[The range of Proposition~\ref{pri:kummer}(b)]\label{pod:rem-kummer-b}
For $\ell>y_0^2$ the inequality $(t+1)(\#s+1)\le h^{21}$ is not obtained from
Corollary~\ref{pod:kummer-b}: the detecting set grows with $\ell$ (Lemma~\ref{pod:detect} gives
$\#s\le\max(N,\ell^{1/2})^{20}$). Applying Lemma~\ref{pod:detect} with $y=h$ instead of $y_0$ gives
the inequality for all $\ell\le h^2$ (Proposition~\ref{pku:b}), and Proposition~\ref{pri:triangle}
uses it only when $h_0\ne0$, which forces $\ell\le h^2$ (Proposition~\ref{pku:triangle}).
\end{remark}

\begin{remark}[Proposition~\ref{pri:kummer}(a)]\label{pod:rem-kummer-a}
Proposition~\ref{pri:kummer}(a) is proved in Proposition~\ref{pku:main}: with $T$ equal to the set of
primes of $F=E^+$ above $\ell$ or below the given places, the relations are the restrictions of dual
Selmer classes, and the Greenberg--Wiles formula, applied to $\Z/\ell^\nu((\pgcyc^+)^{-e})$ and to its
$\ell$-torsion, gives at least $\#x+[F:\Q]-4$ independent reductions.
\end{remark}

\subsubsection{The case \texorpdfstring{$e=0$}{e=0}}\label{pod:ss-e0}

\begin{lemma}\label{pod:e0}
$\prilen\podSha_{0,\nu}\le\prival(h_E)$.
\end{lemma}
\begin{proof}
An element of $\podSha_{0,\nu}$ is a continuous homomorphism $\pgG\to\podLam_\nu$ vanishing on all
inertia groups; as $E$ is totally complex it factors through $\podGal(H/E)\cong\mathrm{Cl}_E$, $H$
the Hilbert class field. Writing $\mathrm{Cl}_E\cong\prod_i\Z/n_i$, $\podSha_{0,\nu}\cong
\prod_i\podLam_\nu[n_i]$, of length $\le\sum_i\prival(n_i)=\prival(h_E)$.
\end{proof}

\subsubsection{Splitting by characters}\label{pod:ss-split}

\begin{definition}\label{pod:def-comp}
For $e\ge1$, a Dirichlet character $\psi$ modulo $N$ and $\nu\ge0$, the \emph{component}
$\podC_{e,\psi,\nu}$ is the $\auO$-module of locally constant cocycles $x\colon G_\Q\to\podLam_\nu$
for $\Upsilon_{e,\psi}^{-1}$ that vanish at $c$ and at every $g\in H_N$ which lies either in $I_A$ for
a place $A\nmid\ell$, or in $D_A$ for a place $A\mid\ell$ with $\pgcyc(g)=1$. We call $\psi$ (or
$\Upsilon_{e,\psi}$) \emph{odd} if $\Upsilon_{e,\psi}(c)=(-1)^e\psi(-1)=-1$.
\end{definition}

\begin{lemma}\label{pod:split}
Let $e\ge1$. There are $\sigma_t\in G_\Q$ \textup($t\in(\Z/N)^\times$\textup) and $b(u,g)\in\pgG$
\textup($u,g\in G_\Q$\textup), independent of $\auO,\nu$, with the following property. For $x\in
\podSha_{e,\nu}$ and each $\psi$ put
\[
z_\psi(g)=\sum_{t\in(\Z/N)^\times}\Upsilon_{e,\psi}(\sigma_t)\,x\bigl(b(\sigma_t,g)\bigr)\qquad(g\in G_\Q).
\]
Then (a) $z_\psi\in\podC_{e,\psi,\nu}$; (b) if $z_\psi=0$ for every odd $\psi$, then $4\varphi(N)x=0$.
\end{lemma}
\begin{proof}
Choose $S_t\in G_\Q$ with $\varepsilon_N(S_t)=t$ and put $\sigma_t=S_tS_1^{-1}$, so $\sigma_1=1$. For
$u,g\in G_\Q$ the element $u(cgc^{-1}g^{-1})u^{-1}$ lies in $H_N\subseteq\pgG$ (as $E\subseteq
\Q(\mu_N)$ and $G_\Q/H_N$ is abelian); let $b(u,g)$ be this element of $\pgG$. Since $b(u,g)$ is a conjugate of a commutator,
$\pgcyc(b(u,g))=1$, and therefore
\begin{equation}\label{pod:eq-add}
x(b(u,g)\,k)=x(b(u,g))+x(k)\qquad(k\in\pgG).
\end{equation}
\emph{Cocycle rule.} One has $b(u,gk)=b(u,g)\,b(ug,k)$, so by \eqref{pod:eq-add}
$x(b(u,gk))=x(b(u,g))+x(b(ug,k))$. Next, $T(u,k):=\Upsilon_{e,\psi}(u)x(b(u,k))$ depends only on
$uH_N$: if $a\in H_N$ then $b(au,k)=ab(u,k)a^{-1}$, and with the cocycle rule for $x$ and
$\pgcyc(b(u,k))=1$, $x(ab(u,k)a^{-1})=\pgcyc(a)^{-e}x(b(u,k))$, while $\Upsilon_{e,\psi}(au)=
\pgcyc(a)^e\Upsilon_{e,\psi}(u)$. Hence
\[
z_\psi(gk)=z_\psi(g)+\sum_t\Upsilon_{e,\psi}(\sigma_t)x(b(\sigma_tg,k))
=z_\psi(g)+\Upsilon_{e,\psi}(g)^{-1}\sum_tT(\sigma_tg,k)=z_\psi(g)+\Upsilon_{e,\psi}(g)^{-1}z_\psi(k),
\]
because $\sigma_tg\in\sigma_{t\varepsilon_N(g)}H_N$ and $t\mapsto t\varepsilon_N(g)$ permutes
$(\Z/N)^\times$. Local constancy is clear.

\emph{Local conditions.} $b(u,c)=1$, so $z_\psi(c)=0$. Let $g\in H_N$ lie in $I_A$ for some $A\nmid
\ell$, or in $D_A\cap\ker\pgcyc$ for some $A\mid\ell$. For $u\in G_\Q$ let $w(u)\in\pgG$ be the
element $ugu^{-1}\in H_N$; it lies in $I_{uA}$ (resp.\ in $D_{uA}\cap\ker\pgcyc$), so $x(w(u))=0$ by
(i), (ii) of Definition~\ref{pod:def-sha}. As $b(u,g)w(u)=w(uc)$, \eqref{pod:eq-add} gives
$x(b(u,g))=x(w(uc))-x(w(u))=0$, so $z_\psi(g)=0$. This proves (a).

(b) If $\psi$ is even, $\Upsilon_{e,\psi}(c)=1$ and $z_\psi(c)=0$ give $z_\psi(cgc)=z_\psi(g)$. If all
odd $z_\psi$ vanish, then $\Sigma(h):=\sum_\psi z_\psi(h)$ satisfies $\Sigma(chc)=\Sigma(h)$ for
$h\in\pgG$. By orthogonality ($\auO$ contains the values of the characters),
$\sum_\psi\Upsilon_{e,\psi}(\sigma_t)=\pgcyc(\sigma_t)^e\sum_\psi\psi(t)=\varphi(N)[t=1]$, so
$\Sigma(h)=\varphi(N)x(b(1,h))=\varphi(N)\bigl(x(chc)-x(h)\bigr)$, by \eqref{pod:eq-add} applied to
$b(1,h)h=chc$. Comparing $h$ and $chc$ and using $x(chc)=-x(h)$:
$-2\varphi(N)x(h)=2\varphi(N)x(h)$, i.e.\ $4\varphi(N)x(h)=0$.
\end{proof}

\subsubsection{A Stickelberger bound from Voronoi's congruence}\label{pod:ss-voronoi}

$B_n(T)$ are the Bernoulli polynomials and $B_n=B_n(0)$; recall $B_n(T+1)-B_n(T)=nT^{n-1}$,
$B_n'(T)=nB_{n-1}(T)$, and that $(j+1)!\,B_j\in\Z$ (induction on $j$ via $\sum_{k\le j}\binom{j+1}{k}
B_k=0$ for $j\ge1$).

\begin{definition}\label{pod:def-stick}
Let $N,e\ge1$. An integer $X\ge1$ is a \emph{Stickelberger bound for $(N,e)$} if
\begin{enumerate}
\item[(S1)] for every prime $\ell$ some integer $a\ge2$ with $a\equiv1\pmod N$ prime to $\ell$ has $(a^e-1)(a^{e+1}-1)\mid
X$;
\item[(S2)] for every $a\equiv1\pmod N$, every domain $R$ of characteristic $0$ and every Dirichlet
character $\psi$ modulo $N$ with values in $R$ and $\psi(-1)=-(-1)^e$, there is $w\in R$ such that
for all $m\ge1$ prime to $a$
\[
(a^{e+1}-1)X\equiv2w\sum_{0\le j<Nm}\Bigl\lfloor\frac{aj}{Nm}\Bigr\rfloor j^e\psi(j)\pmod{mR}.
\]
\end{enumerate}
\end{definition}

\begin{lemma}\label{pod:powersum}
For $a\in\Z_{\ge0}$ put $\beta_a=(e+2)!\,N^{e+1}B_{e+1}(a/N)$. Then $\beta_a\in\Z$ and, for every
$u\ge1$, $(e+2)!\sum_{i<u}(a+iN)^{e+1}\equiv\beta_au\pmod{u^2}$.
\end{lemma}
\begin{proof}
Let $P(T)=(e+1)!\,N^{e+2}\bigl(B_{e+2}(T/N)-B_{e+2}\bigr)$. The coefficient of $T^{k+1}$ is
$(e+1)!\binom{e+2}{k+1}N^jB_j$ with $j+k=e+1$, an integer since $(j+1)!$ divides $(e+1)!\binom{e+2}
{k+1}$ (for $k=0$ this is $(e+2)!$, otherwise $j+1\le e+1$). From $B_{e+2}(T+1)-B_{e+2}(T)=(e+2)T^{e+1}$
one gets $P(T+N)-P(T)=N(e+2)!\,T^{e+1}$, hence $P(a+uN)-P(a)=N(e+2)!\sum_{i<u}(a+iN)^{e+1}$. Taylor
expansion of the integer polynomial $P$ gives $P(a+uN)\equiv P(a)+P'(a)uN\pmod{(uN)^2}$, and
$P'(a)=(e+2)!N^{e+1}B_{e+1}(a/N)=\beta_a$ is an integer. Divide by $N$.
\end{proof}

\begin{lemma}[Voronoi's congruence]\label{pod:voronoi}
Let $R$ be a commutative ring without additive torsion, $\chi$ a character modulo $N$ with values in
$R$, $m\ge1$, $f=Nm$, and $a\equiv1\pmod N$ prime to $m$. Then
\[
(e+2)!(e+1)N a^e\sum_{0\le j<f}\Bigl\lfloor\frac{aj}{f}\Bigr\rfloor j^e\chi(j)\equiv
(a^{e+1}-1)\sum_{0\le j<N}\chi(j)\beta_j\pmod{mR}.
\]
\end{lemma}
\begin{proof}
As $(a,f)=1$, $j\mapsto aj-fk_j$ with $k_j=\lfloor aj/f\rfloor$ permutes $[0,f)$, and $\chi(aj-fk_j)=
\chi(j)$ since $N\mid f$ and $a\equiv1\pmod N$. Expanding $(aj-fk_j)^{e+1}$ to first order,
\[
\sum_{j<f}\chi(j)j^{e+1}\equiv a^{e+1}\sum_{j<f}\chi(j)j^{e+1}-(e+1)a^ef\sum_{j<f}k_j\chi(j)j^e
\pmod{f^2}.
\]
Write $j=b+iN$ ($b<N$, $i<m$) and use Lemma~\ref{pod:powersum}: $(e+2)!\sum_{j<f}\chi(j)j^{e+1}\equiv
m\sum_{b<N}\chi(b)\beta_b\pmod{m^2}$. Multiplying the display by $(e+2)!$ and substituting,
$(a^{e+1}-1)m\sum_b\chi(b)\beta_b\equiv(e+2)!(e+1)a^eNm\sum_jk_j\chi(j)j^e\pmod{m^2}$. Divide by $m$.
\end{proof}

\begin{lemma}\label{pod:L}
Let $\chi$ be a complex character modulo $N$ \textup(possibly imprimitive\textup) and $e\ge1$. Then
$\sum_{0\le j<N}\chi(j)\beta_j=-(e+2)!(e+1)N\,L(\chi,-e)$, and $L(\chi,-e)=0$ if and only if
$\chi(-1)=(-1)^e$.
\end{lemma}
\begin{proof}
$L(\chi,s)=N^{-s}\sum_{j=1}^N\chi(j)\zeta(s,j/N)$ and $\zeta(-e,x)=-B_{e+1}(x)/(e+1)$; the term $j=N$
equals the term $j=0$ as $B_{e+1}(1)=B_{e+1}(0)$ for $e\ge1$. This gives the formula. For the zeros:
write $\chi=\chi_0\chi_f$ with $\chi_f$ primitive of conductor $f$; then $L(\chi,s)=L(\chi_f,s)
\prod_{p\mid N,p\nmid f}(1-\chi_f(p)p^{-s})$, and the Euler factors do not vanish at $s=-e$ (their
zeros lie on $\Rel s=0$). For primitive $\chi_f$, the functional equation relates $L(\chi_f,-e)$ to
$L(\overline{\chi_f},1+e)\ne0$ times a ratio of Gamma factors, which vanishes exactly when $-e$ is a
pole of $\Gamma\bigl((s+\kappa)/2\bigr)$, $\kappa\in\{0,1\}$ the parity of $\chi_f$, i.e.\ when
$e\equiv\kappa\pmod2$ (see the functional equation in \cite[Chapter~4]{Washington}). For $\chi_f=1$ ($\zeta$) the same holds since $e\ge1$.
\end{proof}

\begin{lemma}\label{pod:X}
For $N,e\ge1$ there is a Stickelberger bound $X_e$ for $(N,e)$ with $1\le X_e\le2^{(N+e+3)^6}$.
\end{lemma}
\begin{proof}
Put $g(t)=2\beta_t+[t\equiv1](-1)^e+[t\equiv-1]$ for $t\in\Z/N$ (with $\beta_t=\beta_{t_0}$,
$t_0\in[0,N)$ the representative) and let $T$ be the integer matrix indexed by
$(\Z/N)^\times$ with $T_{s,t}=g(t/s)$. For every complex character $\chi$ modulo $N$, the vector
$(\chi(t))_t$ is an eigenvector with eigenvalue $\lambda_\chi=\sum_{t\in\Z/N}\chi(t)g(t)=
2\sum_t\chi(t)\beta_t+(-1)^e+\chi(-1)$. If $\chi(-1)=-(-1)^e$, $\lambda_\chi=2\sum\chi\beta\ne0$ by
Lemma~\ref{pod:L}; otherwise $\sum\chi\beta=0$ and $\lambda_\chi=2(-1)^e\ne0$. The characters form a
basis, so $\det T=\prod_\chi\lambda_\chi\ne0$. Put
\[
X_e=\bigl((1+N)^e-1\bigr)\bigl((1+N)^{e+1}-1\bigr)\bigl((1+2N)^e-1\bigr)\bigl((1+2N)^{e+1}-1\bigr)\,
|\det T|.
\]
(S1): $\ell$ cannot divide both $1+N$ and $1+2N$; take $a$ the one prime to $\ell$.
(S2): for $\psi$ with values in $R$ (we may enlarge $R$ to contain $\Z[\psi]$, a domain of
characteristic $0$), $(\psi(t))_t$ is an eigenvector of $T$ with eigenvalue $\lambda_\psi$, so applying
the adjugate, $\lambda_\psi\mid\det T$ in $R$; write $X_e=\lambda_\psi u$. By parity $\lambda_\psi=
2\sum_t\psi(t)\beta_t$. Lemma~\ref{pod:voronoi} multiplied by $2u$ gives (S2) with $w=u(e+2)!(e+1)Na^e$.
\emph{Size.} $|B_{e+1}(x)|\le2(e+1)!\zeta(e+1)(2\pi)^{-e-1}<17(e+1)!$ on $[0,1]$ (Fourier series), so
$|\beta_t|<17(e+2)!N^{e+1}(e+1)!$ and the entries of $T$ are $<64\,(e+2)!N^{e+1}(e+1)!=:64M_e$. Hence
$|\det T|\le\varphi(N)!\,(64M_e)^{\varphi(N)}$ (Leibniz expansion).
Now let $u=N+e+3\ge5$; then $\log_2N\le u$, $\log_2(e+2)\le u$, $\varphi(N)\le N\le u$. Each of the
four factors in front of $|\det T|$ is $<(1+2N)^{e+1}$, and $\log_2(1+2N)\le2N$, so their product is
$<2^{8N(e+1)}\le2^{8u^2}$. Next, $(e+2)!(e+1)!\le(e+2)^{2(e+2)}$, so $\log_2(64M_e)\le6+2(e+2)
\log_2(e+2)+(e+1)\log_2N\le6+3u^2$, and $\log_2\varphi(N)!\le N\log_2N\le u^2$. Hence
\[
\log_2X_e\le8u^2+u^2+u(6+3u^2)=3u^3+9u^2+6u\le7u^3\le u^6 ,
\]
using $u\ge5$ (so $9u^2+6u\le4u^3$ and $7\le u^3$). That is, $X_e\le2^{(N+e+3)^6}$.
\end{proof}

\subsubsection{Gauss sums and a family of Jacobi-type quotients}\label{pod:ss-gauss}

\begin{lemma}[Hermite's identity]\label{pod:hermite}
For $c,r\ge1$ coprime and $x\in\Q/\Z$, $\sum_{b\in\Z/r}m_c(x+b/r)=m_c(rx)$.
\end{lemma}
\begin{proof}
$\lfloor c\podfr z\rfloor=c\podfr z-\podfr{cz}$. The representatives of $x+b/r$ are
$(b'+\podfr{rx})/r$, $b'\in[0,r)$, so $\sum_b\podfr{x+b/r}=\frac{r-1}2+\podfr{rx}$; as $c$ is prime to
$r$, also $\sum_b\podfr{cx+cb/r}=\frac{r-1}2+\podfr{crx}$. Hence $\sum_b\lfloor c\podfr{x+b/r}\rfloor=
\frac{(c-1)(r-1)}2+\lfloor c\podfr{rx}\rfloor$; the same with $-x$, and subtract.
\end{proof}

Fix a prime $p$ and a place $B$ of $\Omega$ above $p$. Let $\psi_0$ be a nontrivial additive
character of $\mathbb F_p$ and $\chi_0$ a generator of the group of characters of $\mathbb F_p^\times$,
both with values in $\Omega$; for $t\in\Z/(p-1)$ let $G(t)=\sum_{x\in\mathbb F_p^\times}\chi_0^t(x)
\psi_0(x)$ be the Gauss sum. Recall: $G(0)=-1$; $G(t)G(-t)=\chi_0^t(-1)p$ for $t\ne0$; hence $G(t)\mid
p$ and $G(t)$ is a unit at every place not above $p$. For $\sigma\in G_\Q$ with $\varepsilon_p(\sigma)
=\alpha$ and $\varepsilon_{p-1}(\sigma)=a$ one has $\sigma G(t)=\chi_0^{at}(\alpha)^{-1}G(at)$
(substitute $x\mapsto\alpha^{-1}x$); write $K_\sigma(s)=\chi_0^s(\alpha)^{-1}$, a homomorphism in $s$.

\begin{lemma}[Weak Stickelberger]\label{pod:stick}
There is an injective homomorphism $S\colon\Z/(p-1)\to\Q/\Z$ with $\podfr{S(t)}=\podord_B G(t)$ for
all $t$.
\end{lemma}
\begin{proof}
If $\sigma\in I_B$ then $\varepsilon_{p-1}(\sigma)=1$, so $\sigma G(t)=K_\sigma(t)G(t)$ and
$\sigma$ fixes $G(a+b)/(G(a)G(b))$. An element fixed by $I_B$ lies in a field unramified at $B$ over
$\Q_p$ (after completion), so its $\podord_B$ is an integer; hence $t\mapsto\podord_BG(t)\bmod\Z$ is a
homomorphism $S$. For $t\ne0$: $\psi_0$ takes $p$-th roots of unity, $\equiv1$ modulo $B$, so
$G(t)\equiv\sum_x\chi_0^t(x)=0$ modulo $B$ and $\podord_BG(t)>0$; the same for $p/G(t)=\pm
\sum\chi_0^{-t}(x)\psi_0^{-1}(x)$, so $\podord_BG(t)<1$. With $\podord_BG(0)=0$ this gives
$\podfr{S(t)}=\podord_BG(t)$, and $S(t)=0$ forces $t=0$.
\end{proof}

\begin{lemma}[Jacobi family]\label{pod:J}
Let $c\ge1$. There is $J\colon\Q/\Z\to\Omega^\times$ such that for every $M$ with $p\equiv1\pmod M$
and all $x$ with $Mx=0$:
\begin{enumerate}
\item[(J1)] $\sigma J(x)=J(ax)$ whenever $\varepsilon_M(\sigma)=a$;
\item[(J2)] $J(x)$ is a unit at every place not above $p$;
\item[(J3)] $\podord_BJ(x)=m_c(x)$;
\item[(J4)] $|\iota J(x)|=1$ for every embedding $\iota\colon\Omega\to\C$;
\item[(J5)] if $r\mid M$ and $(r,c)=1$, then $\prod_{b\in\Z/r}J(x+b/r)\big/J(rx)$ is a root of unity;
\item[(J6)] if $q$ is a prime dividing $p-1$, $qb=0$, and $A$ is a place above $q$, then
$J(x+b)/J(x)\equiv1$ modulo $A$ \textup(and both are units at $A$\textup).
\end{enumerate}
\end{lemma}
\begin{proof}
Put $j(t)=G(t)^c/G(ct)$ and $J(S(t))=j(t)/j(-t)$ (and $J=1$ outside the $(p-1)$-torsion, which is
$S(\Z/(p-1))$). (J1): with $a=\varepsilon_{p-1}(\sigma)$, $\sigma j(t)=K_\sigma(at)^cG(at)^c/
(K_\sigma(cat)G(cat))=j(at)$, and $S(at)=aS(t)$; for $Mx=0$, $ax$ depends only on $a\bmod M$.
(J2): Gauss sums divide $p$. (J3): $\podord_Bj(t)=c\podfr{S(t)}-\podfr{cS(t)}=\lfloor c\podfr{S(t)}
\rfloor$, and $S(-t)=-S(t)$. (J4): if $\sigma\in G_\Q$ induces complex conjugation through $\iota$,
then $\varepsilon_{p-1}(\sigma)=-1$ and $\sigma j(t)=j(-t)$, so $\iota J=\iota j/\overline{\iota j}$.
(J5): Let $Q$ be the quotient. It is a unit off $p$ by (J2) and has absolute value $1$ everywhere by
(J4). A place $A$ above $p$ is $\sigma^{-1}B$ for some $\sigma$; with $a=\varepsilon_M(\sigma)$ (a unit
modulo $r$), (J1) and (J3) give $\podord_A Q=\podord_B\sigma Q=\sum_bm_c(ax+ab/r)-m_c(rax)=0$ by
Lemma~\ref{pod:hermite}. So $Q$ is an algebraic integer all of whose conjugates have absolute value
$1$; by Kronecker's theorem it is a root of unity. (J6): for $qu=0$ in $\Z/(p-1)$, $\chi_0^u$ takes
$q$-th roots of unity, $\equiv1$ modulo $A$, so $G(w+u)\equiv G(w)$ modulo $A$, and both are units at
$A$ ($q\ne p$). Hence $j(w+u)\equiv j(w)$ and $J(x+b)\equiv J(x)$ modulo $A$, using $b=S(u)$ with
$qu=0$.
\end{proof}

\subsubsection{Artin reciprocity for components and the annihilator}\label{pod:ss-annih}

\begin{lemma}\label{pod:locals}
Let $x\in\podC_{e,\psi,\nu}$ with $mx=0$ \textup($m\ge1$\textup), and let $a\equiv1\pmod N$ be
prime to $\ell$.
\begin{enumerate}
\item[(a)] If $\phi\in H_{Nm}$ and $\varepsilon_{Nm}(s)=u$, then $x(s^{-1}\phi s)=u^e\psi(u)x(\phi)$.
\item[(b)] If $g\in H_{Nm}\cap D_A$ for a place $A\mid\ell$, then $(a^e-1)x(g)=0$.
\item[(c)] If $g\in H_{Nm}$ and $g^2=1$, then $2x(g)=0$.
\item[(d)] If $g\in H_{Nm}\cap I_A$ for a place $A\nmid\ell$, then $x(g)=0$.
\end{enumerate}
\end{lemma}
\begin{proof}
Elements of $H_{Nm}$ act trivially through $\Upsilon_{e,\psi}^{\pm1}$ on $\podLam_\nu[m]$.
(a) is Lemma~\ref{pod:coc}(b) with $\chi=\Upsilon^{-1}_{e,\psi}$, $\chi(s^{-1})=\Upsilon_{e,\psi}(s)$.
(c) is Lemma~\ref{pod:coc}(e). (d) is the defining condition. (b): as $\Q_\ell(\mu_{\ell^\infty})/\Q_\ell$
is totally ramified with group $\Z_\ell^\times$, there is $\delta\in I_A$ with $\pgcyc(\delta)=a$;
$\delta$ fixes $\mu_N$ (the prime-to-$\ell$ part because $\delta$ is in inertia above $\ell$, the
$\ell$-part because $a\equiv1\pmod N$), so $\Upsilon_{e,\psi}(\delta)=a^e$. The commutator
$k=\delta g\delta^{-1}g^{-1}$ lies in $H_N\cap D_A\cap\ker\pgcyc$, so $x(k)=0$, and
Lemma~\ref{pod:coc}(c) gives $(a^{-e}-1)x(g)=(\Upsilon_{e,\psi}(g)^{-1}-1)x(\delta)=0$.
\end{proof}

\begin{lemma}[Artin reciprocity for components]\label{pod:artin}
Let $m\ge1$ be an integer, $F=\Q(\mu_{Nm})$, $q$ a prime
with $q\equiv1\pmod{Nm}$, $B$ a place above $q$, $\phi$ a Frobenius element at $B$. Let $x\in
\podC_{e,\psi,\nu}$ with $mx=0$, and $a\equiv1\pmod N$ prime to $\ell$. Let $y\in F^\times$ be such
that at every place $A$ of $\Omega$ above neither $\ell$ nor $q$, either $x$ vanishes on $D_A$, or
$y=\eta\,r^m$ with $\eta$ a unit at $A$ and $r\in\Omega^\times$ fixed by $I_A\cap H_{Nm}$. For $u\in
(\Z/Nm)^\times$ choose $s_u\in G_\Q$ with $\varepsilon_{Nm}(s_u)=u$ and let $k_u=\podord_B(s_uy)$.
Then
\[
\Bigl(\sum_{u\in(\Z/Nm)^\times}k_u\,u^e\psi(u)\Bigr)\cdot2(a^e-1)\,x(\phi)=0 .
\]
\end{lemma}
\begin{proof}
Put $d=2(a^e-1)$ and let $f\colon\podLam_\nu\to\Q/\Z$ be any homomorphism. Since $H_{Nm}=G_F$ acts
trivially on $\podLam_\nu[m]$, $\chi_f(h)=f(d\,x(h))$ is a continuous character of $G_F$ with $m\chi_f
=0$. By Lemma~\ref{pod:locals}, $\chi_f$ is unramified at all places not above $\ell$, kills the
decomposition groups above $\ell$ (by (b)), and kills every element of order $2$ (by (c)), hence the
archimedean local Artin images. Global reciprocity for $F$ gives $\sum_w\chi_f(\mathrm{Art}_w(y))=0$
over all places $w$ of $F$ ($\chi_f$ factors through a finite abelian extension of $F$, and the
product formula for the norm residue symbol of a principal idele is \cite[Chapter~VI, Corollary~(5.7)]{Neukirch}). The terms at places above $\ell$ vanish; at a place
$w\nmid\ell q$ of $F$ below $A$, $\chi_f$ is unramified, so $\chi_f(\mathrm{Art}_w(\alpha))=
\podord_w(\alpha)\chi_f(\mathrm{Frob}_w)$ for $\alpha\in F_w^\times$; either $\chi_f$ kills $D_A\cap
G_F$, or $y=\eta r^m$ as above, and then $\podord_w(\eta)=0$ while $\podord_w(r^m)=m\,\podord_w(r)$
with $\podord_w(r)\in\Z$ because $F_w(r)/F_w$ is unramified ($r$ is fixed by inertia); as $m\chi_f=0$,
the term vanishes. The prime $q$ splits
completely in $F$; the places of $F$ above $q$ are the $w_u$ below $s_u^{-1}B$, pairwise distinct
(if $w_u=w_{u'}$ then $s_{u'}gs_u^{-1}\in D_B\subseteq H_{Nm}$ for some $g\in G_F$, forcing
$u=u'$) and exhausting (any place above $q$ is $\sigma^{-1}B$, and $\sigma^{-1}B$ lies over the same
place of $F$ as $s_u^{-1}B$ for $u=\varepsilon_{Nm}(\sigma)$). At $w_u$, $\chi_f$ is unramified, a
Frobenius is $s_u^{-1}\phi s_u$, and $\podord_{w_u}(y)=\podord_{s_u^{-1}B}(y)=\podord_B(s_uy)=k_u$
($q$ is unramified in $F$). Hence $\sum_uk_u\chi_f(s_u^{-1}\phi s_u)=0$, i.e.\ by
Lemma~\ref{pod:locals}(a), $f\bigl(\sum_uk_uu^e\psi(u)\,d\,x(\phi)\bigr)=0$. As the $f$ separate points,
the claim follows.
\end{proof}

\begin{lemma}[Euler factor]\label{pod:thetadiv}
Let $X$ be a Stickelberger bound for $(N,e)$, $\psi$ odd \textup($\Upsilon_{e,\psi}(c)=-1$, i.e.\
$\psi(-1)=-(-1)^e$\textup), $a\equiv1\pmod N$ prime to $\ell$, $n\ge0$. Then in $\podLam_n$, the
Stickelberger sum $\podth_{a,N\ell^n}(e,\psi)$ divides $(a^{e+1}-1)X$.
\end{lemma}
\begin{proof}
For $n=0$ the ring is zero. Let $n\ge1$ and write $n=n'+1$. Put $T_m=\sum_{0\le j<Nm}\lfloor aj/(Nm)
\rfloor j^e\psi(j)$ and $U=\sum_{u\in(\Z/N\ell^n)^\times}\lfloor au/(N\ell^n)\rfloor u^e\psi(u)$. An
integer $j<N\ell^n$ with $\psi(j)\ne0$ is either prime to $\ell$ or $\ell j'$ with $j'<N\ell^{n'}$, and
$\lfloor a\ell j'/(N\ell^n)\rfloor=\lfloor aj'/(N\ell^{n'})\rfloor$; so $T_{\ell^n}=U+\ell^e\psi(\ell)
T_{\ell^{n'}}$. Modulo $\ell^n$ (indeed modulo $N\ell^n$), $\podth_{a,N\ell^n}=2U$: $m_a(u/M)=
\lfloor au/M\rfloor-\lfloor a(M-u)/M\rfloor$ and $u\mapsto M-u$ multiplies $u^e\psi(u)$ by
$(-1)^e\psi(-1)=-1$ modulo $M$. By (S2) at $m=\ell^n$ and $m=\ell^{n'}$ (both prime to $a$),
$(a^{e+1}-1)X\equiv2wT_{\ell^n}\pmod{\ell^n}$ and $\equiv2wT_{\ell^{n'}}\pmod{\ell^{n'}}$; multiplying
the second by $\ell^e\psi(\ell)$ ($e\ge1$) makes it a congruence modulo $\ell^n$. Subtracting,
$(1-\ell^e\psi(\ell))(a^{e+1}-1)X\equiv2wU=w\podth_{a,N\ell^n}\pmod{\ell^n}$, and $1-\ell^e\psi(\ell)$
is a unit in $\podLam_n$ ($\ell^e$ is nilpotent there).
\end{proof}

\begin{proposition}[Annihilation, every $\ell$]\label{pod:annih}
Let $e\ge1$, $X$ a Stickelberger bound for $(N,e)$ and $\psi$ odd. Then $4X^2$ kills
$\podC_{e,\psi,\nu}$, for every prime $\ell$ and every $\nu$.
\end{proposition}
\begin{proof}
Let $z\in\podC=\podC_{e,\psi,\nu}$ and suppose $4X^2z\ne0$. The cocycle $z'=2X^2z\in\podC$ vanishes
at $c$, and $(\Upsilon_{e,\psi}^{-1}(c)-1)z'=-4X^2z\ne0$; by Lemma~\ref{pod:coc}(d), $z'$ does not
vanish on $H_{N\ell^\nu}$. So there is $\sigma_0\in H_{N\ell^\nu}$ with $z'(\sigma_0)\ne0$. Let $U$ be the intersection of
$H_{N\ell^\nu}$ with the normal core of the open subgroup $Z(z')$; it is an open normal subgroup of
$G_\Q$ contained in $H_{N\ell^\nu}$. For $k\in U$, $z'(\sigma_0k)=z'(\sigma_0)+\sigma_0\cdot z'(k)
=z'(\sigma_0)$, so $z'\ne0$ on the coset $\sigma_0U\subseteq H_{N\ell^\nu}$. By Chebotarev's density
theorem, applied to the finite Galois extension cut out by $U$, the coset $\sigma_0U$ contains a
Frobenius element at a place above a prime unramified in that extension; since the coset lies in
$H_{N\ell^\nu}$, that prime splits in $\Q(\mu_{N\ell^\nu})$. Hence there are a prime $p\equiv1\pmod{N\ell^\nu}$, a place $B$ above it and a Frobenius element $\sigma\in
H_{N\ell^\nu}$ at $B$ with $z'(\sigma)\ne0$. By (S1) choose $a\equiv1\pmod N$ prime to $\ell$ with
$X=(a^e-1)(a^{e+1}-1)b$, $b\in\Z$. Let $J$ be the Jacobi family of Lemma~\ref{pod:J} for $(p,B,c=a)$
and $y=J(1/(N\ell^\nu))$. By (J1) with $M=N\ell^\nu$, $y$ is fixed by $H_{N\ell^\nu}$, so $y\in F=
\Q(\mu_{N\ell^\nu})$, and $s_uy=J(u/(N\ell^\nu))$; by (J2) $y$ is a unit at every place not above $p$;
by (J3) $\podord_B(s_uy)=m_a(u/(N\ell^\nu))$. Lemma~\ref{pod:artin} (with $m=\ell^\nu$, $q=p$,
$\eta=y$, $r=1$) gives $\podth_{a,N\ell^\nu}(e,\psi)\cdot2(a^e-1)z(\sigma)=0$ in $\podLam_\nu$. By
Lemma~\ref{pod:thetadiv} write $(a^{e+1}-1)X=\podth\kappa$ in $\podLam_\nu$; then
$2(a^e-1)(a^{e+1}-1)X\,z(\sigma)=0$, and multiplying by $b$, $2X^2z(\sigma)=0$, a contradiction.
\end{proof}

\begin{corollary}\label{pod:kill}
For $e\ge1$ and $X_e$ as in Lemma~\ref{pod:X}, $Y_e=4\varphi(N)\cdot4X_e^2$ kills $\podSha_{e,\nu}$
for every prime $\ell$.
\end{corollary}
\begin{proof}
For $x\in\podSha_{e,\nu}$, the odd components of $4X_e^2x$ are $4X_e^2z_\psi=0$
(Lemma~\ref{pod:split}(a) and Proposition~\ref{pod:annih}); apply Lemma~\ref{pod:split}(b) to
$4X_e^2x$.
\end{proof}

\subsubsection{Kolyvagin's derivative classes}\label{pod:ss-koly}

In this subsection $\ell$ is a prime, $n\ge1$, $e\ge1$, $\psi$ is a character modulo $N$ with values
in $\auO$, and $\Upsilon_1=\Upsilon_{e+1,\psi}=\pgcyc^{e+1}\psi$. We fix a primitive $\ell^n$-th root
of unity $\zeta\in\Omega$; for $\omega\in\mu_{\ell^n}$, $\log\omega\in\Z/\ell^n$ is defined by
$\zeta^{\log\omega}=\omega$. For a finite set $T$ of primes let $\mathcal I(T)$ be the set of
$g\in H_N$ lying in $I_A$ for some place $A$ above neither $\ell$ nor a prime of $T$.

\begin{definition}\label{pod:def-KS}
Let $\gamma\in G_\Q$, $p$ a prime and $P$ a finite set of primes. A \emph{Kolyvagin system} for
$(\Upsilon_1,\gamma,p,P)$ with values in an $\auO$-module $V$ is a family of maps $\kappa_S\colon
G_\Q\to V$, $S\subseteq P$, such that for all $S\subseteq P$: (KS1) $\kappa_S$ is a locally constant
cocycle for $\Upsilon_1$ vanishing on $\mathcal I(S\cup\{p\})$; (KS2) $\kappa_S(\gamma)=0$; (KS3) for
$q\in P\smallsetminus S$, a place $A$ above $q$ and a Frobenius element $\phi$ at $A$, there is
$\tau\in I_A$ with $\kappa_{S\cup\{q\}}(\tau)=\kappa_S(\phi)$.
\end{definition}

\begin{proposition}\label{pod:KS}
Let $c\ge1$, $\gamma\in G_\Q$, $P$ a finite set of primes $\equiv1\pmod{cN\ell^n}$, and $p$ a prime
with $p\equiv1\pmod{N\ell^n\prod_{q\in P}q}$, $B$ a place above $p$. There is a Kolyvagin system
$(\kappa_S)$ for $(\Upsilon_1,\gamma,p,P)$ with values in $\podLam_n$ such that: if $\kappa_\emptyset$
vanishes on $I_B$, then $(\Upsilon_1(\gamma)-1)\,\podth_{c,N\ell^n}(e,\psi)=0$ in $\podLam_n$.
\end{proposition}

\emph{Construction.} Let $J$ be the Jacobi family of Lemma~\ref{pod:J} for $(p,B,c)$; it satisfies
(J1)--(J6) for every $M\mid p-1$, in particular for $M_S=N\ell^n\prod_{q\in S}q$ and $qM_S$ ($S
\subseteq P$, $q\in P\smallsetminus S$). Choose $y(x)\in\Omega^\times$ with $y(x)^{\ell^n}=J(x)$. For
$q\in P$ choose a surjective homomorphism $L_q\colon(\Z/q)^\times\to\Z/\ell^n$ (possible as
$\ell^n\mid q-1$) and $v_q$ with $L_q(v_q)=1$. For $x$ with $M_Sx=0$, $\sigma\in G_\Q$ and
$a=\varepsilon_{M_S}(\sigma)$, (J1) gives $(\sigma y(x)/y(ax))^{\ell^n}=\sigma J(x)/J(ax)=1$, i.e.
\begin{equation}\label{pod:eq-root}
\sigma\,y(x)/y(ax)\in\mu_{\ell^n}.
\end{equation}
For $g\in G_\Q$ the commutator $\gamma^{-1}g\gamma g^{-1}$ lies in $H_{M}$ for every $M$, hence fixes
$J(x)$, hence moves $y(x)$ by an element of $\mu_{\ell^n}$; applying $\gamma$, $\gamma y(x)/(g\gamma
g^{-1})y(x)\in\mu_{\ell^n}$. Put
\[
F_g(x)=\log\frac{\gamma\,y(x)}{(g\gamma g^{-1})\,y(x)}\in\Z/\ell^n .
\]
For a function $G\colon\Q/\Z\to\Z/\ell^n$, shifts $d=(d_q)\in(\Z/\ell^n)^P$ and $S\subseteq P$, let
$\mathcal D_q^{d}G(x)=\sum_{u\in(\Z/q)^\times}(L_q(u)-d_q)\,G(x+u/q)$ (these operators commute) and
\[
\Lambda_S^d(G)=\sum_{a\in(\Z/N\ell^n)^\times}a^e\psi(a)\,\Bigl(\prod_{q\in S}\mathcal D_q^d\Bigr)G
\Bigl(\frac a{N\ell^n}\Bigr)\in\podLam_n,
\qquad \kappa_S(g)=\Lambda_S^0(F_g).
\]
$\Lambda_S^d(G)$ depends only on $G$ on the $M_S$-torsion and on the $d_q$, $q\in S$.

\begin{lemma}[Equivariance]\label{pod:equiv}
Let $M$ be a multiple of $M_S$, $g\in G_\Q$, and $a\in\Z$ prime to $M$ with $a\equiv
\varepsilon_{N\ell^n}(g)\pmod{N\ell^n}$. If $G(x)=G'(ax)$ for all $x$ with $Mx=0$, then $\Upsilon_{e,\psi}(g)\Lambda_S^d(G)=\Lambda_S^{f}(G')$
with $f_q=L_q(a\bmod q)+d_q$.
\end{lemma}
\begin{proof}
Induction on $\#S$; only the two hypotheses on $a$ are used. For $S=\emptyset$, as $a\equiv
\varepsilon_{N\ell^n}(g)\pmod{N\ell^n}$, $\Upsilon_{e,\psi}(g)$ acts on $\podLam_n$ as $a^e\psi(a)$ and the
substitution $a'\mapsto aa'$ in the sum over $(\Z/N\ell^n)^\times$ gives the claim. For $S\cup\{q\}$,
the substitution $u\mapsto(a\bmod q)u$ (a bijection of $(\Z/q)^\times$, as $q\mid M$ and $a$ is prime to $M$) in $\mathcal D_q$ turns $\mathcal D_q^dG$ into $\mathcal
D_q^fG'$ composed with $x\mapsto ax$, using $L_q(au)-f_q=L_q(u)-d_q$ and $a(x+u/q)=ax+au/q$.
\end{proof}

\begin{lemma}[Shifts do not matter]\label{pod:shift}
For $S\subseteq P$, $h\in G_\Q$ and all shifts $d$, $\Lambda_S^d(F_h)=\Lambda_S^0(F_h)$.
\end{lemma}
\begin{proof}
\emph{Distribution relation.} For $r\in P$ and $M_{S\cup\{r\}}x=0$: $\sum_{b\in\Z/r}F_h(x+b/r)=
F_h(rx)$. Indeed $D=\prod_by(x+b/r)/y(rx)$ has $D^{\ell^n}$ a root of unity by (J5) ($r\equiv1\pmod c$
is prime to $c$), so $D$ is a root of unity, and $\gamma$ and $h\gamma h^{-1}$ act on it in the same
way (through the same cyclotomic character); the claim is $\log(\gamma D/(h\gamma h^{-1})D)=0$.
Since $r\equiv1\pmod{N\ell^n}$, Lemma~\ref{pod:equiv} (with $g=1$, $a=r$, $M=M_{S'}$, which is prime to $r$, and $r\equiv1=\varepsilon_{N\ell^n}(1)\pmod{N\ell^n}$, $G=\sum_b F_h(\cdot+b/r)$,
$G'=F_h$) gives, for $r\notin S'$: $\Lambda_{S'}^d\bigl(\sum_bF_h(\cdot+b/r)\bigr)=\Lambda_{S'}^{f}(F_h)$
for some shifts $f$.

\emph{Induction on $\#S$.} It suffices to change $d$ at one $r\in S$. Put $S'=S\smallsetminus\{r\}$.
Then $\Lambda_S^d(F_h)=\Lambda_{S'}^d(\mathcal D_r^dF_h)$, and changing $d_r$ to $d_r'$ changes this by
$(d_r-d_r')\Lambda_{S'}^d\bigl(\sum_{u\in(\Z/r)^\times}F_h(\cdot+u/r)\bigr)=(d_r-d_r')\bigl(
\Lambda_{S'}^d(\textstyle\sum_{b\in\Z/r}F_h(\cdot+b/r))-\Lambda_{S'}^d(F_h)\bigr)$. By the
distribution relation the first term is $\Lambda_{S'}^f(F_h)$, which equals $\Lambda_{S'}^d(F_h)$ by
induction. So the change is $0$.
\end{proof}

\begin{proof}[Proof of Proposition~\ref{pod:KS}]
(KS1), cocycle rule. Let $g,g'\in G_\Q$, $M=M_S$, $a=\varepsilon_M(g)$. For $Mx=0$,
\[
F_{gg'}(ax)-F_g(ax)=a\,F_{g'}(x).
\]
Indeed, as $\ell^n\mid M$, $g$ raises $\ell^n$-th roots of unity to the power $a$, so
$\zeta^{aF_{g'}(x)}=g\bigl(\gamma y(x)/(g'\gamma g'^{-1})y(x)\bigr)=(g\gamma g^{-1})(gy(x))\big/
(gg'\gamma(gg')^{-1})(gy(x))$. By \eqref{pod:eq-root}, $gy(x)=y(ax)\omega$ with $\omega\in\mu_{\ell^n}$,
on which all conjugates of $\gamma$ act alike, so $\omega$ cancels and the right side is
$\frac{(g\gamma g^{-1})y(ax)}{\gamma y(ax)}\cdot\frac{\gamma y(ax)}{(gg'\gamma(gg')^{-1})y(ax)}=
\zeta^{F_{gg'}(ax)-F_g(ax)}$. Hence, by Lemma~\ref{pod:equiv} with $G=aF_{g'}$, $G'=F_{gg'}-F_g$, and
Lemma~\ref{pod:shift},
\[
\kappa_S(gg')-\kappa_S(g)=\Lambda_S^0(F_{gg'}-F_g)=\Upsilon_{e,\psi}(g)\,\Lambda_S^0(aF_{g'})=
\Upsilon_{e,\psi}(g)\,a\,\kappa_S(g')=\Upsilon_1(g)\kappa_S(g'),
\]
as $\pgcyc(g)\equiv a\pmod{\ell^n}$. $\kappa_S$ is locally constant since $g\mapsto(g\gamma g^{-1})y(x)$
is, for the finitely many $x$ involved.

(KS1), vanishing. Let $g\in H_N\cap I_A$, $A$ above $r'$ with $r'\notin\{\ell,p\}\cup S$. Then $g$
fixes $\mu_{M_S}$ (it fixes $\mu_N$, and the roots of unity of order prime to $r'$ by inertia).
For $M_Sx=0$, $y(x)$ is a unit at $A$ (by (J2)), so $gy(x)\equiv y(x)$ modulo $A$, and $gy(x)/y(x)
\in\mu_{\ell^n}$ by \eqref{pod:eq-root}; as $\ell^n$-th roots of unity are distinct modulo $A$
($r'\ne\ell$), $gy(x)=y(x)$; the same holds for $g^{-1}$ and for $y(ax)$. With $\gamma y(x)=
y(ax)\omega$ ($a=\varepsilon_{M_S}(\gamma)$, $\omega\in\mu_{\ell^n}$ fixed by $g$), $(g\gamma
g^{-1})y(x)=g\gamma y(x)=y(ax)\omega=\gamma y(x)$. So $F_g=0$ on the $M_S$-torsion and $\kappa_S(g)=0$.

(KS2): $F_\gamma=0$.

(KS3). Let $q\in P\smallsetminus S$, $M=M_S$ (prime to $q$), $A$ above $q$, $\phi$ a Frobenius at $A$.
Choose $\tau\in I_A$ with $\varepsilon_{qM}(\tau)\equiv1\pmod M$ and $\equiv v_q^{-1}\pmod q$ (inertia
at $A$ maps onto $(\Z/q)^\times$ and fixes $\mu_M$). We claim, for $Mx=0$,
\begin{equation}\label{pod:eq-FI}
\sum_{u\in(\Z/q)^\times}L_q(u)\,F_\tau(x+u/q)=F_\phi(qx).
\end{equation}
Granting this, $\kappa_{S\cup\{q\}}(\tau)=\Lambda_S^0(\mathcal D_q^0F_\tau)=\Lambda^0_S(x\mapsto
F_\phi(qx))=\Lambda_S^f(F_\phi)=\kappa_S(\phi)$ by Lemma~\ref{pod:equiv} ($g=1$, $M=M_S$, $a=q$, prime to $M_S$ and $\equiv1\pmod
{N\ell^n}$) and Lemma~\ref{pod:shift}.

Proof of \eqref{pod:eq-FI}. Put $Y_b=y(x+b/q)$ ($b\in\Z/q$), $y'=y(qx)$, and let $\lambda(u)\in\Z$ lift
$L_q(u)$, so $\lambda(v_qu)-\lambda(u)-1=\ell^nk_u$ with $k_u\in\Z$. By (J1) at level $qM$:
$\tau Y_{v_qu}/Y_u\in\mu_{\ell^n}$ and, for every conjugate $h\tau h^{-1}$, the same. By (J1) at level
$M$: $\phi Y_0/y'\in\mu_{\ell^n}$, and the same for $h\phi h^{-1}$. By (J5), $\bigl(\prod_bY_b/y'
\bigr)^{\ell^n}$ is a root of unity. For $h\in G_\Q$ put
\[
W_h=\frac{\prod_bY_b}{(h\phi h^{-1})Y_0}\prod_{u\in(\Z/q)^\times}\Bigl(\frac{(h\tau h^{-1})Y_{v_qu}}
{Y_u}\Bigr)^{\lambda(u)} .
\]
Each $W_h$ is a root of unity (product of roots of unity and of $\prod_bY_b/y'$). \emph{Claim:}
$W_h\equiv1$ modulo $hA$. Since $h\tau h^{-1}\in I_{hA}$ acts trivially on the residue field and
$h\phi h^{-1}$ as $t\mapsto t^q$, and the $Y_b$ are units at $hA$ ($q\ne p$; (J2)), modulo $hA$:
$W_h\equiv\prod_bY_b\cdot Y_0^{-q}\prod_u(Y_{v_qu}/Y_u)^{\lambda(u)}=\prod_bY_b\cdot Y_0^{-q}
\prod_uY_u^{\lambda(v_q^{-1}u)-\lambda(u)}$, and $\lambda(v_q^{-1}u)-\lambda(u)=-1-\ell^nk_{v_q^{-1}u}$.
By (J6) (with $q\mid p-1$), $Y_u^{\ell^n}=J(x+u/q)\equiv J(x)=Y_0^{\ell^n}$ modulo $hA$, and since
$\sum_u(\lambda(v_q^{-1}u)-\lambda(u))=0$, summing gives $\ell^n\sum_uk_u=-(q-1)$. Therefore $\prod_uY_u^{\lambda(v_q^{-1}u)-\lambda(u)}
\equiv\prod_uY_u^{-1}\cdot Y_0^{q-1}$, and $W_h\equiv Y_0\cdot Y_0^{-q}\cdot Y_0^{q-1}=1$.

Now compare $\gamma W_1$ and $\gamma W_{\gamma^{-1}}$. Write $\phi Y_0=y'\omega_0$ and $\tau Y_{v_qu}=
Y_u\omega_u$ with $\omega$'s in $\mu_{\ell^n}$, on which all conjugates of $\gamma$ act alike. Then
$\gamma\phi Y_0=\gamma y'\cdot\gamma\omega_0$, $\phi\gamma Y_0=(\phi\gamma\phi^{-1})(y'\omega_0)=
(\phi\gamma\phi^{-1})y'\cdot\gamma\omega_0$, and similarly for $\tau$; hence
\[
\frac{\gamma W_1}{\gamma W_{\gamma^{-1}}}=\frac{\phi\gamma Y_0}{\gamma\phi Y_0}\prod_u\Bigl(
\frac{\gamma\tau Y_{v_qu}}{\tau\gamma Y_{v_qu}}\Bigr)^{\lambda(u)}
=\frac{(\phi\gamma\phi^{-1})y'}{\gamma y'}\prod_u\Bigl(\frac{\gamma Y_u}{(\tau\gamma\tau^{-1})Y_u}
\Bigr)^{\lambda(u)}=\zeta^{\sum_uL_q(u)F_\tau(x+u/q)-F_\phi(qx)}.
\]
$W_1$ is a root of unity $\equiv1$ modulo $A$, so $\gamma W_1$, a power of $W_1$, is $\equiv1$ modulo
$A$; and $\gamma W_{\gamma^{-1}}\equiv1$ modulo $\gamma\gamma^{-1}A=A$. So the $\ell^n$-th root of
unity on the right is $\equiv1$ modulo $A$, hence equal to $1$ ($q\ne\ell$). This is \eqref{pod:eq-FI}.

\emph{The bottom class.} Let $\kappa_\emptyset$ vanish on $I_B$, and let $a\in\Z$ with $a\equiv
\varepsilon_{N\ell^n}(\gamma)$. Choose $r\in\Omega$ with $r^{\ell^n}=p$ and $\tau\in I_B$ with $\tau r=
\zeta r$ (tame inertia at $B$ maps onto $\mu_{\ell^n}$, $p\ne\ell$). Note $\tau\in H_{N\ell^n}$ as
$p\nmid N\ell$. For $N\ell^nx=0$, $(y(x)/r^{m_c(x)})^{\ell^n}=J(x)p^{-m_c(x)}$ is a unit at $B$ by (J3),
so $\tau$ moves $y(x)/r^{m_c(x)}$ by an $\ell^n$-th root of unity $\equiv1$ modulo $B$, i.e.\ fixes
it: $\tau y(x)=\zeta^{m_c(x)}y(x)$. With $\gamma y(x)=y(ax)\omega$:
$(\tau\gamma\tau^{-1})y(x)=\tau\gamma(\zeta^{-m_c(x)}y(x))=\zeta^{-am_c(x)}\tau(y(ax)\omega)=
\zeta^{m_c(ax)-am_c(x)}\gamma y(x)$. Hence $F_\tau=am_c-m_c(a\,\cdot)$ on the $N\ell^n$-torsion, and
$0=\kappa_\emptyset(\tau)=a\Lambda_\emptyset(m_c)-\Lambda_\emptyset(m_c(a\cdot))$. By
Lemma~\ref{pod:equiv} with $g=\gamma$, $\Lambda_\emptyset(m_c(a\,\cdot))=\Upsilon_{e,\psi}(\gamma)^{-1}
\Lambda_\emptyset(m_c)$, and $\Lambda_\emptyset(m_c)=\podth_{c,N\ell^n}(e,\psi)$ in $\podLam_n$. So
$(a\Upsilon_{e,\psi}(\gamma)-1)\podth=0$, i.e.\ $(\Upsilon_1(\gamma)-1)\podth=0$.
\end{proof}

\subsubsection{Descent and detecting points for large \texorpdfstring{$\ell$}{l}}\label{pod:ss-descent}

In this subsection $e\ge1$, $\psi$ is odd, $\ell>N+e+3$, and $\podC=\podC_{e,\psi,\nu}$. Then $2$ and
$\varphi(N\ell)=\varphi(N)(\ell-1)$ are units of $\auO$ ($0<\varphi(N),\ell-1<\ell$). Moreover there
is $a\equiv1\pmod N$ prime to $\ell$ with $a^e-1$ and $a^{e+1}-1$ units of $\auO$: take $a\equiv1
\pmod N$ and $a\equiv g\pmod\ell$ with $g$ a primitive root modulo $\ell$ (as $\ell>N$, $\ell\nmid N$);
then $a^k\not\equiv1\pmod\ell$ for $0<k<\ell-1$, and $e+1<\ell-1$. Fix such an $a$. For a prime $r$,
$x\colon G_\Q\to V$ is \emph{trivial above $r$} if $x$ vanishes on $D_A$ for every place $A$ above
$r$. $\mathcal Z^{(m)}(T)$ denotes the module of locally constant cocycles $G_\Q\to\podLam_m$ for
$\Upsilon_1=\pgcyc^{e+1}\psi$ vanishing on $\mathcal I(T)$.

\begin{lemma}\label{pod:trivabove}
Let $r\equiv1\pmod{N\ell}$ be prime, $A$ a place above $r$, $\sigma$ a Frobenius at $A$. If $x\in
\podC$, $\ell x=0$ and $x(\sigma)=0$, then $x$ is trivial above $r$.
\end{lemma}
\begin{proof}
Let $A'=\rho A$ be any place above $r$. $I_{A'}\subseteq H_N$ ($\mu_N$ is unramified at $r$) and
$A'\nmid\ell$, so $x|_{I_{A'}}=0$. $\sigma\in D_A\subseteq H_{N\ell}$ acts trivially on
$\podLam_\nu[\ell]$, so $x(\rho\sigma\rho^{-1})=\Upsilon_{e,\psi}^{-1}(\rho)x(\sigma)=0$
(Lemma~\ref{pod:coc}(b)). $D_{A'}$ is topologically generated by $I_{A'}$ and the Frobenius
$\rho\sigma\rho^{-1}$; apply Lemma~\ref{pod:coc}(a).
\end{proof}

\begin{lemma}[Reciprocity modulo $\ell$]\label{pod:recl}
Let $q\equiv1\pmod{N\ell}$ be prime, $B$ a place above $q$, $\phi$ a Frobenius at $B$, $\tau\in I_B$.
Let $x\in\podC$ with $\ell x=0$, let $T$ be a finite set of primes, and let $\kappa\in\mathcal
Z^{(m)}(T)$ with $\ell\kappa=0$. Assume that at every place $A$ not above $q$, either $A$ lies above
no prime of $T$, or $x$ vanishes on $D_A$; and that $2(a^e-1)x(\phi)\ne0$. Then $\kappa(\tau)=0$.
\end{lemma}
\begin{proof}
Let $F=\Q(\mu_{N\ell})$. Let $V=\auO\kappa(\tau)$, an $\mathbb F_\ell$-vector space with basis
$\beta_1,\dots,\beta_k$. For each $i$ extend $v\mapsto\zeta_\ell^{\beta_i^*(v)}$ ($\zeta_\ell$ a fixed
primitive $\ell$-th root of unity, $\beta_i^*$ the coordinate) to a homomorphism $\lambda_i\colon\podLam_m[\ell]\to\mu_\ell$
(extend the coordinates from $V$ to the $\mathbb F_\ell$-vector space $\podLam_m[\ell]$, which contains
all values of $\kappa$ since $\ell\kappa=0$). On $G_F$, $\kappa$ is a homomorphism ($\Upsilon_1\equiv1$ modulo $\ell$
there) and $\lambda_i\circ\kappa|_{G_F}$ is a continuous homomorphism to $\mu_\ell$, i.e.\ a cocycle for
the natural action; by Hilbert's Theorem~90 there is $r_i\in\Omega^\times$ with
$\lambda_i(\kappa(g))=gr_i/r_i$ for $g\in G_F$, and $y_i=r_i^\ell\in F^\times$. Choose $\sigma_u\in
G_\Q$ with $\varepsilon_{N\ell}(\sigma_u)=u$ for $u\in(\Z/N\ell)^\times$ and let $k_{i,u}=
\podord_B(\sigma_uy_i)$. Apply Lemma~\ref{pod:artin} with $m=\ell$, $y=y_i$: at a place $A$ not above
$\ell q$, either $x|_{D_A}=0$, or $A$ is above no prime of $T$, and then $\kappa$ vanishes on
$I_A\cap H_N\supseteq I_A\cap G_F$, so $r_i$ is fixed by $I_A\cap G_F$ and $y_i=1\cdot r_i^\ell$.
Hence $s_i:=\sum_uk_{i,u}u^e\psi(u)$ kills $2(a^e-1)x(\phi)\ne0$, so $s_i\in\podm$.

Let $\pi\in\Omega$ with $\pi^\ell=q$ and $j\in\Z/\ell$ with $\tau\pi=\zeta_\ell^j\pi$. Note $\tau\in
I_B\subseteq H_{N\ell}=G_F$. For each $u$: $\kappa(\sigma_u^{-1}\tau\sigma_u)=\Upsilon_1(\sigma_u)^{-1}
\kappa(\tau)=(u^{e+1}\psi(u))^{-1}\kappa(\tau)$ (Lemma~\ref{pod:coc}(b)). Apply $\sigma_u$ to
$\lambda_i(\kappa(\sigma_u^{-1}\tau\sigma_u))=\sigma_u^{-1}\tau\sigma_u r_i/r_i$: the left side becomes
its $u$-th power, the right side $\tau(\sigma_ur_i)/\sigma_ur_i$. The element $\sigma_ur_i/\pi^{k_{i,u}}$
has $\ell$-th power $\sigma_uy_i\,q^{-k_{i,u}}$, a unit at $B$, so $\tau$ moves it by an $\ell$-th root
of unity congruent to $1$ modulo $B$, i.e.\ fixes it; hence $\tau(\sigma_ur_i)/\sigma_ur_i=
\zeta_\ell^{jk_{i,u}}$. Therefore the $i$-th coordinate of $u\cdot(u^{e+1}\psi(u))^{-1}\kappa(\tau)$ is
$jk_{i,u}$, i.e.\ $\kappa(\tau)=u^e\psi(u)\sum_ijk_{i,u}\beta_i$. Summing over $u$:
$\varphi(N\ell)\kappa(\tau)=\sum_ijs_i\beta_i\in\podm V=\podm\kappa(\tau)$. So $(\varphi(N\ell)-t)
\kappa(\tau)=0$ for some $t\in\podm$; $\varphi(N\ell)-t$ is a unit, and $\kappa(\tau)=0$.
\end{proof}

\begin{definition}\label{pod:def-tight}
A finite set $T$ of primes $\equiv1\pmod{N\ell}$ is \emph{tight} (at level $m$) if every $\kappa\in
\mathcal Z^{(m)}(T)$ with $\ell\kappa=0$ lies in $\mathcal Z^{(m)}(\emptyset)$.
\end{definition}

\begin{lemma}\label{pod:tight}
Let $T$ be tight, $z\in\podC$ with $\ell z=0$, trivial above every prime of $T$ but not above a prime
$r\equiv1\pmod{N\ell}$. Then $T\cup\{r\}$ is tight.
\end{lemma}
\begin{proof}
Let $\kappa\in\mathcal Z^{(m)}(T\cup\{r\})$, $\ell\kappa=0$. It suffices to show that $\kappa$
vanishes on $I_A\cap H_N$ for places $A$ above $r$ (then $\kappa\in\mathcal Z^{(m)}(T)$ and tightness
applies). Let $\sigma$ be a Frobenius at $A$. By Lemma~\ref{pod:trivabove}, $z(\sigma)\ne0$, hence
$2(a^e-1)z(\sigma)\ne0$. At a place $A'$ not above $r$, either $A'$ lies above no prime of $T\cup\{r\}$,
or $A'$ lies above a prime of $T$ and $z|_{D_{A'}}=0$. Lemma~\ref{pod:recl} (with $q=r$, the set
$T\cup\{r\}$) gives $\kappa(\tau)=0$ for $\tau\in I_A$.
\end{proof}

\begin{lemma}\label{pod:finite}
$\podC$ and $\mathcal Z^{(m)}(\emptyset)$ are finite. If $\podC\ne0$, then for every $f\ge1$ there is a
finite set $Q$ of primes $\equiv1\pmod f$ such that for every $\kappa\in\mathcal Z^{(m)}(\emptyset)$
not vanishing on $H_f$ and every $0\ne z\in\podC$ there are $r\in Q$, a place $A$ above $r$ and a
Frobenius $\sigma$ at $A$ with $\kappa(\sigma)\ne0\ne z(\sigma)$. Every nonzero $z\in\podC$ is
nonzero on every $H_f$.
\end{lemma}
\begin{proof}
On $H_{N\ell^{\max(\nu,m)}}$ these cocycles are homomorphisms to a finite group, unramified outside
$N\ell$ (they vanish on inertia in $H_N$ at places prime to $\ell$, and inertia at places prime to
$N\ell$ lies in $H_N$); by Hermite--Minkowski there are finitely many such, and a cocycle is determined
by its restriction and its values on coset representatives. The last statement is
Lemma~\ref{pod:coc}(d) with $\gamma=c$, as $\Upsilon_{e,\psi}^{-1}(c)-1=-2$ is a unit. For a pair
$(\kappa,z)$ as stated, the open subgroups $Z(\kappa)$ and $Z(z)$ do not contain $H_f$; a group is not
the union of two proper subgroups, so some $h\in H_f$ lies in neither. Let $U$ be the intersection of $H_f$ with the normal cores of
$Z(\kappa)$ and $Z(z)$, an open normal subgroup of $G_\Q$ inside $H_f$. For $k\in U$ the cocycle rule
gives $\kappa(hk)=\kappa(h)$ and $z(hk)=z(h)$, so every element of $hU\subseteq H_f$ lies in neither
subgroup. By Chebotarev, applied to the extension cut out by $U$, $hU$ contains a Frobenius at a place
above a prime $r$ unramified there, and $r\equiv1\pmod f$ because $hU\subseteq H_f$. Choose one $r$ for each of the
finitely many pairs.
\end{proof}

\begin{proposition}[Kolyvagin--Rubin detection]\label{pod:koly}
Let $X$ be a nonzero Stickelberger bound for $(N,e)$, $\psi$ odd, $\ell>N+e+3$. There are
$\gamma_0,\dots,\gamma_{v_\ell(X)}\in G_\Q$ such that $x\in\podC$ and $x(\gamma_i)=0$ for all $i$ imply
$x=0$.
\end{proposition}
\begin{proof}
Put $m=v_\ell(X)+1$ and suppose the contrary: for every $\delta_1,\dots,\delta_m$ some $0\ne x\in\podC$
vanishes at all $\delta_j$. \emph{Step 1.} For every set $T$ of at most $m$ primes $\equiv1\pmod
{N\ell}$ there is $0\ne z\in\podC$ with $\ell z=0$, trivial above every prime of $T$: choose for each
$r\in T$ a Frobenius $\sigma_r$ as in Lemma~\ref{pod:trivabove}, take $x\ne0$ vanishing at all
$\sigma_r$, and let $z=\ell^kx$ with $\ell^kx\ne0=\ell^{k+1}x$; $z$ still vanishes at the $\sigma_r$.

\emph{Step 2 (set-up).} As in the proof of Lemma~\ref{pod:locals}(b) there is $\gamma\in I_A$
($A\mid\ell$) with $\pgcyc(\gamma)=a$ and $\gamma\in H_N$, so
$\Upsilon_1(\gamma)-1=a^{e+1}-1$ is a unit. Let $Q$ be as in Lemma~\ref{pod:finite} for
$f=aN\ell^m$. By Step~1 with $T=\emptyset$ take $0\ne z_0\in\podC$, $\ell z_0=0$; by
Lemma~\ref{pod:finite} and Chebotarev there are a prime $p\equiv1\pmod{N\ell^m\prod_{r\in Q}r}$, a place
$B$ above $p$ and a Frobenius $\tau_0$ at $B$ with $z_0(\tau_0)\ne0$. So $z_0$ is not trivial above
$p$, and $\{p\}$ is tight by Lemma~\ref{pod:tight} (the empty set is tight). Let $(\kappa_S)_{S
\subseteq Q}$ be the Kolyvagin system of Proposition~\ref{pod:KS} for $(n=m,c=a,\gamma,P=Q,p,B)$; the
primes of $Q$ are $\equiv1\pmod{aN\ell^m}$.

\emph{Step 3 (descent).} We show by descending induction on $i\in[0,m]$: for $S\subseteq Q$ with
$\#S\le i$ and $S\cup\{p\}$ tight, $\ell^i\kappa_S\in\mathcal Z^{(m)}(\emptyset)$. For $i=m$,
$\ell^m\kappa_S=0$. Let $i<m$ and $S$ as stated. It suffices to show $\ell^{i+1}\kappa_S=0$: then
$\ell\cdot\ell^i\kappa_S=0$, $\ell^i\kappa_S\in\mathcal Z^{(m)}(S\cup\{p\})$ by (KS1), and tightness
gives the claim. Suppose $\kappa':=\ell^{i+1}\kappa_S\ne0$. By induction ($\#S\le i+1$) $\kappa'\in
\mathcal Z^{(m)}(\emptyset)$; $\kappa'(\gamma)=0$ (KS2) and $(\Upsilon_1(\gamma)-1)\kappa'\ne0$, so by
Lemma~\ref{pod:coc}(d) $\kappa'$ does not vanish on $H_{aN\ell^m}$. By Step~1 (with $T=S\cup\{p\}$, of
size $\le i+1\le m$) there is $0\ne w\in\podC$, $\ell w=0$, trivial above $p$ and the primes of $S$. By
Lemma~\ref{pod:finite} there are $r\in Q$, $A\mid r$ and a Frobenius $\sigma$ at $A$ with
$\kappa'(\sigma)\ne0\ne w(\sigma)$. Then $w$ is not trivial above $r$, so $r\notin S$, and
$S\cup\{r\}\cup\{p\}$ is tight by Lemma~\ref{pod:tight}. By induction ($\#(S\cup\{r\})\le i+1$),
$\ell^{i+1}\kappa_{S\cup\{r\}}\in\mathcal Z^{(m)}(\emptyset)$, which vanishes on $I_A$ (as $I_A\subseteq
H_N$, $A\nmid\ell$, so $I_A\subseteq\mathcal I(\emptyset)$). By (KS3) there is $\iota\in I_A$ with
$\kappa'(\sigma)=\ell^{i+1}\kappa_S(\sigma)=\ell^{i+1}\kappa_{S\cup\{r\}}(\iota)=0$, a contradiction.

\emph{Step 4.} With $i=0$, $S=\emptyset$: $\kappa_\emptyset\in\mathcal Z^{(m)}(\emptyset)$, so it
vanishes on $I_B\subseteq\mathcal I(\emptyset)$ ($p\equiv1\pmod{N\ell}$). Proposition~\ref{pod:KS}
gives $(a^{e+1}-1)\podth_{a,N\ell^m}(e,\psi)=0$, hence $\podth_{a,N\ell^m}=0$ in $\podLam_m$, and by
Lemma~\ref{pod:thetadiv} $\ell^m\mid(a^{e+1}-1)X$, so $\ell^m\mid X$ in $\auO$. As
$X=\ell^{m-1}X'$ with $\ell\nmid X'$, this gives $\ell\mid X'$ in $\auO$ ($\auO$ is a domain); but $X'$ is
a unit of $\auO$, so $\ell$ would be a unit, which is absurd.
\end{proof}

\subsubsection{Proof of Proposition~\ref{pod:odd}}\label{pod:ss-proof}

For $0<e<t$ let $X_e$ be as in Lemma~\ref{pod:X} and $Y_e=16\varphi(N)X_e^2$, and put
\[
\pgcodd=\Bigl(\prod_{0<e<t}Y_e\Bigr)^{y^{20}}h_E,\qquad y=N+t+3 .
\]
\emph{Size.} For $0<e<t$, $N+e+3\le y$, so $X_e\le2^{y^6}$. Also $16\varphi(N)\le2^{4+\log_2N}
\le2^{2y^6}$ (as $y\ge4$). Hence $Y_e=16\varphi(N)X_e^2\le2^{2y^6+2y^6}=2^{4y^6}$. There are $t-1<y$
values of $e$, so $\prod_{0<e<t}Y_e\le2^{4y^7}$ and $\pgcodd\le2^{4y^7\cdot y^{20}}h_E=2^{4y^{27}}h_E$.

\emph{$e=0$:} Lemma~\ref{pod:e0}, as $h_E\mid\pgcodd$.

\emph{$0<e<t$, $\ell\le y^2$:} by Corollary~\ref{pod:kill}, $Y_e$ kills $\podSha_{e,\nu}$; by
Lemma~\ref{pod:detect} (with $y\ge N$), at most $y^{20}$ points detect it (elements of
$\podSha_{e,\nu}$ vanish on inertia off $\ell$). Lemma~\ref{pod:length} gives
$\prilen\podSha_{e,\nu}\le\prival(Y_e^{y^{20}})\le\prival(\pgcodd)$.

\emph{$0<e<t$, $\ell>y^2$:} then $\ell>N+e+3$. For each odd $\psi$ let $\gamma_{\psi,i}$ ($0\le i\le
v_\ell(X_e)$) be the points of Proposition~\ref{pod:koly}, and consider the points $b(\sigma_u,
\gamma_{\psi,i})\in\pgG$ of Lemma~\ref{pod:split}. Their number is at most $\varphi(N)^2(v_\ell(X_e)+1)
\le y^2(y^6+1)\le y^{20}$ (as $v_\ell(X_e)\le\log_2X_e\le y^6$). If $x\in\podSha_{e,\nu}$ vanishes at all
of them, then $z_\psi(\gamma_{\psi,i})=0$, so $z_\psi=0$ for every odd $\psi$, and $4\varphi(N)x=0$ by
Lemma~\ref{pod:split}(b); since $4\varphi(N)<\ell$ is prime to $\ell$ and $\ell^\nu x=0$, $x=0$. With
the annihilator $Y_e$ (Corollary~\ref{pod:kill}), Lemma~\ref{pod:length} gives $\prilen
\podSha_{e,\nu}\le\prival(Y_e^{y^{20}})\le\prival(\pgcodd)$. \qed

\begin{remark}[The bound $\log\pgcodd\le x^{800}$]\label{pod:rem-size}
Proposition~\ref{pg:odd} and Lemma~\ref{pg:sums} use $\log\pgcodd\le x^{800}$. We check it from the
data ({\cite[Lemma~13.4]{P46}} and {\cite[\S14.3]{P46}}). \emph{Inputs:} $f\le q_1\le x^{12}$ (S3);
$\sqrt d\,q_1\prod_{p\mid q_0}p^{-1/(p-1)}=\anrd_E\le x^{12}$ with $q_1\ge x^{11}$ (S3, S4), so
$\sqrt d\le x\,q_0$, and $q_0\le\log x$ (S1); hence $N=q_1d\le x^{14}(\log x)^2\le x^{15}$. Further
$x^5\le r\le3x^5$ and $t=5r\le15x^5$ ({\cite[\S14.3]{P46}}). \emph{The first factor:} $y=N+t+3\le
x^{15}+15x^5+3\le2x^{15}$, so $4y^{27}\log2\le4\cdot2^{27}x^{405}\log2\le x^{406}$. \emph{The class
number:} $[E:\Q]=2f\le2x^{12}$ and $|d_E|=\anrd_E^{\,2f}$, so $\log|d_E|\le2x^{12}\cdot12\log x$.
Every ideal class of $E$ contains an integral ideal of norm $\le|d_E|^{1/2}$ (Minkowski: the
constant $(4/\pi)^{r_2}n!/n^n$ is $\le1$), and there are at most $k^{[E:\Q]}$ integral ideals of norm
$k$ (the number is the $[E:\Q]$-fold divisor function at most, and $\tau_n(k)\le n^{\Omega(k)}\le
k^{\log_2n}\le k^n$), so $h_E\le\sum_{k\le|d_E|^{1/2}}k^{2f}\le|d_E|^{(2f+1)/2}$ and
\[
\log h_E\le\frac{2f+1}2\log|d_E|\le(x^{12}+1)\cdot24x^{12}\log x\le x^{26}.
\]
Therefore $\log\pgcodd\le x^{406}+x^{26}\le x^{407}\le x^{800}$. (If $N\le x^{13}$, the same computation gives $x^{352}$.)
\end{remark}



\providecommand{\auO}{\mathcal O}
\providecommand{\auNm}{\operatorname{Nm}}
\providecommand{\pgG}{G_E}
\providecommand{\pgcyc}{\varepsilon_\ell}
\providecommand{\prilen}{\operatorname{len}}
\providecommand{\prival}{\operatorname{val}}
\providecommand{\prity}{\mathrm{ty}}
\providecommand{\pkuinv}{\operatorname{inv}}
\providecommand{\pkuHom}{\operatorname{Hom}}
\providecommand{\pkuGal}{\operatorname{Gal}}
\providecommand{\pkucp}{\varepsilon^{+}}
\providecommand{\pkuSel}{\operatorname{Sel}^{*}}
\providecommand{\pkuur}{\mathrm{ur}}
\providecommand{\pkuAnn}{\operatorname{Ann}}
\providecommand{\pkured}{\operatorname{red}}

\subsection{Kummer maps}\label{pku:sec}

This section proves part~(a) of Proposition~\ref{pri:kummer} of \S\ref{pri:sec}
(Proposition~\ref{pku:main}) and part~(b) (Proposition~\ref{pku:b}), and gives the details of the
proof of Proposition~\ref{pri:triangle}, the only place where Proposition~\ref{pri:kummer} is used
(\S\ref{pku:ss-triangle}). The proof of (a) is a count of Greenberg--Wiles type over the totally
real field $E^+$; it uses only the Greenberg--Wiles formula, local and global duality for finite
Galois modules, the explicit form of the local pairing at a split prime, and Chebotarev's theorem.
The set $T$ of part~(a) can be taken to be the smallest one allowed, the primes above $\ell$ and
below the given places; with this choice the ``relations'' defining a Kummer map are exactly
restrictions of dual Selmer classes, and the difficulty that a larger $T$ would create (relations
ramified at primes of $T$ where the test cocycles are unramified but not zero) does not arise.

\subsubsection{Setting and statement}\label{pku:ss-statement}

\emph{Conventions.} $\Omega=\overline\Q\subset\C$; $E\subseteq\Omega$ is the CM field of
{\cite[\S14.3]{P46}}: it contains a primitive $5040$-th root of unity, and $E/E^+$ is unramified at
all finite places. Put $F=E^+$, $G_F=\pkuGal(\Omega/F)\supseteq\pgG=\pkuGal(\Omega/E)$, and let
$c\in G_F\smallsetminus\pgG$ be the fixed complex conjugation. A \emph{complex conjugation} is any
conjugate of $c$ in $G_{\Q}$ lying in $G_F$; none lies in $\pgG$ ($E$ is totally complex). For a
place $B$ of $\Omega$, $D_B\supseteq I_B$ are its decomposition and inertia groups in $G_F$; a
\emph{Frobenius element} at $B$ over a finite prime $v$ of $F$ is an element of $D_B$ acting on
the residue field as $t\mapsto t^{\auNm v}$. Fix a prime $\ell$ and $\nu\ge1$, put $n=\ell^\nu$,
and let $\pgcyc$ be the $\ell$-adic cyclotomic character of $G_\Q$. Fix a primitive $n$-th root of
unity $\zeta\in\Omega$ and let $\theta\colon\mu_n\to\frac1n\Z/\Z$ be the isomorphism
$\zeta^a\mapsto a/n$. A \emph{normalised tame generator} (for $n$) at a place $B$ over a prime $p\nmid
\ell$ of $F$ is an element $\gamma\in I_B$ with $\gamma(\varpi^{1/n})=\zeta\varpi^{1/n}$ for some
uniformiser $\varpi$ of $F_p$ and some $n$-th root $\varpi^{1/n}$; these exist, and the quotient
$\gamma(\varpi^{1/n})/\varpi^{1/n}$ does not depend on the choice of $\varpi$ and of the root (as in
\S\ref{pse:ss-tame}, over $F$ instead of $E$).

Let $G_\nu=\{g\in\pgG:\pgcyc(g)\equiv1\bmod n\}$, and define $\pkucp\colon G_F\to\Z_\ell^\times$ by
$\pkucp(g)=\pgcyc(g)$ and $\pkucp(gc)=\pgcyc(g)$ for $g\in\pgG$. For $e\in\Z$ let $M=M_e$ be $\Z/n$
with $g\in G_F$ acting by $\pkucp(g)^{-e}$, let $M^\vee=\pkuHom(M,\Q/\Z)$ (as an abelian group), and
let $M^*=\pkuHom(M,\mu_n)$ be the Tate dual, with $(gf)(a)=g\bigl(f(g^{-1}a)\bigr)$. For a finite
$G_F$-module $A$ we write $A[\ell]$ for its $\ell$-torsion and $H^i(F,A)$, $H^i(F_v,A)$ for
continuous cohomology of $G_F$ and of $D_{B}$, $B$ a fixed place over $v$.

\begin{lemma}[The characters]\label{pku:basic}
\begin{enumerate}
\item[(i)] $\pkucp$ is a continuous character, equal to $\pgcyc\eta$ with $\eta$ the quadratic
character of $E/F$. It is $1$ at every complex conjugation, unramified at every prime of $F$ not
above $\ell$, and if $\ell=2$ it takes values in $1+4\Z_2$.
\item[(ii)] $G_\nu$ is open and normal in $G_F$, contains no complex conjugation, and acts trivially
on $M$, $M^*$ and $\mu_n$. The group $G_F/G_\nu$ is abelian, and its $\ell$-Sylow subgroup is
generated by at most one element if $\ell$ is odd and by at most two if $\ell=2$.
\item[(iii)] $F$ is totally real with $r_F=[F:\Q]\ge576$ real places.
\item[(iv)] Let $w$ be a place of $E$ of degree one over $\Q$, $w\nmid\ell$, $A$ a place of
$\Omega$ above $w$ and $\phi\in G_\nu$ a Frobenius element at $A$ (in $\pgG$). Then $D_A\subseteq
\pgG$ (decomposition group in $G_F$), $\phi$ is a Frobenius element at $A$ for $F$, $D_A$ is
topologically generated by $I_A$ and $\phi$, and $D_A$ acts trivially on $M$.
\end{enumerate}
\end{lemma}
\begin{proof}
(i) For $g\in\pgG$, $\pgcyc(gc)=-\pgcyc(g)$ and $\eta(gc)=-1$, so $\pkucp=\pgcyc\eta$ on both cosets;
it is a continuous character. A complex conjugation $c'$ has $\pgcyc(c')=-1$ and $c'\notin\pgG$,
so $\eta(c')=-1$ and $\pkucp(c')=1$. An inertia group $I_B$ at a finite prime lies in $\pgG$ ($E/F$
is unramified at finite places), and there $\pkucp=\pgcyc$, which is unramified off $\ell$. If
$\ell=2$: $\mu_4\subseteq E$, so $\pgcyc\equiv1\bmod4$ on $\pgG$, and $\pkucp(gc)=\pgcyc(g)$.

(ii) $G_\nu=\ker\eta\cap\ker(\pkucp\bmod n)$, an open normal subgroup; it lies in $\pgG$, so contains
no complex conjugation. On $G_\nu$, $\pkucp\equiv1\bmod n$ and $\pgcyc\equiv1\bmod n$, so $G_\nu$
acts trivially on $M$, $\mu_n$ and $M^*$. The map $g\mapsto(\eta(g),\pkucp(g)\bmod n)$ embeds
$G_F/G_\nu$ in $\{\pm1\}\times(\Z/n)^\times$. If $\ell$ is odd, the $\ell$-Sylow subgroup of the
latter is cyclic. If $\ell=2$, the image is generated by the image of $\pgG$, contained in the
cyclic group $\{1\}\times(1+4\Z)/n$, and by the image of $c$; so it is generated by two elements,
and so is its $2$-Sylow subgroup (the whole group).

(iii) $F$ is the maximal real subfield of a CM field, so totally real, and $[E:\Q]\ge[\Q(\mu_{5040})
:\Q]=\varphi(5040)=1152$.

(iv) The prime $v$ of $F$ below $w$ has the same residue field $\mathbb F_p$ as $w$ and $E/F$ is
unramified at $v$; so $w$ is split over $v$, $D_A\subseteq\pgG$, and the decomposition groups of $A$
in $G_F$ and $\pgG$ coincide. In particular $\phi$ is a Frobenius element for $F$, and $D_A$ is
topologically generated by $I_A$ and $\phi$. On $I_A$, $\pkucp=\pgcyc=1$ ($w\nmid\ell$), and
$\pkucp(\phi)=\pgcyc(\phi)\equiv1\bmod n$.
\end{proof}

\begin{definition}[Kummer maps]\label{pku:def-kummer}
Let $x=(w_i)_{i\in I}$ be a finite family of degree-one places $w_i\nmid\ell$ of $E$ above distinct
primes $v_i$ of $F$, with places $A_i$ of $\Omega$ above $w_i$ and Frobenius elements $\phi_{w_i}
\in G_\nu$ at $A_i$. Let $S_0=S_0(x)$ be the set of primes of $F$ above $\ell$ or equal to some
$v_i$, and fix a place $B_v$ of $\Omega$ above each $v\in S_0$, with $B_{v_i}=A_i$.

(a) A \emph{test cocycle for $x$ and $P$}, where $P$ is a finite set of primes of $F$ not in
$S_0$, is a continuous cocycle $z\colon G_F\to M$ that vanishes on $I_B$ for every place $B$ over a
prime not in $S_0\cup P$; vanishes on $D_{B_{v_i}}$ for every $i$; vanishes on $D_{B_v}\cap\ker
\pkucp$ for every $v\mid\ell$; and vanishes at every complex conjugation.

(b) A \emph{split prime} is a prime $p\notin S_0$ of $F$ with a chosen place $B_p$ over it, a
Frobenius element $\phi_p\in G_\nu$ at $B_p$, and a normalised tame generator $\gamma_p$ at $B_p$.

(c) Let $T\supseteq S_0$ be a finite set of primes of $F$. A \emph{relation off $T$} is a function
$y\colon G_F\to M^\vee$ such that: $y$ is locally constant; $y=0$ outside $G_\nu$; $y(uv)=y(u)+y(v)$
for $u,v\in G_\nu$; $y(gug^{-1})(a)=\pgcyc(g)\,y(u)(g^{-1}a)$ for $g\in G_F$, $u\in G_\nu$, $a\in
M$; $y=0$ on $I_B$ for every place $B$ over a prime not in $T$; and
\begin{equation}\label{pku:eq-recip}
\sum_{p\in P}y(\phi_p)\bigl(z(\gamma_p)\bigr)=0
\end{equation}
for every finite set $P$ of split primes not in $T$ and every test cocycle $z$ for $x$ and $P$.

(d) A \emph{Kummer map of exponent $e$ for $x$ off $T$} is a function $\kappa\colon\pgG\to\Z/n$
such that $y(g)(a)=a\kappa(g)/n$ ($g\in\pgG$, $a\in M$) for some relation $y$ off $T$. Its
\emph{reduction} is $\pkured\kappa\colon G_\nu\to\mathbb F_\ell$, $g\mapsto\kappa(g)\bmod\ell$.
\end{definition}

\begin{remark}\label{pku:rem-def}
(1) A test cocycle vanishes on $gHg^{-1}$ for every $g\in G_F$ and every subgroup $H$ on which it
vanishes and which acts trivially on $M$: $z(ghg^{-1})=z(g)+gz(h)-ghg^{-1}z(g)=0$, as $ghg^{-1}$
acts trivially. The groups $D_{B_{v_i}}$ and $D_{B_v}\cap\ker\pkucp$ act trivially on $M$
(Lemma~\ref{pku:basic}(iv) and the definition of $\pkucp$), so the conditions do not depend on the
chosen places. (2) The equivariance law, with $\pkucp$ acting on $M$, reads
$y(gug^{-1})=\pgcyc(g)\pkucp(g)^{e}\,y(u)$; for $g\in\pgG$ it says $\kappa(gug^{-1})=
\pgcyc(g)^{e+1}\kappa(u)$, and for $g=c$ that $\kappa(cuc)=-\kappa(u)$ ($u\in G_\nu$).
\end{remark}

\begin{proposition}[Kummer maps]\label{pku:main}
Let $x=(w_i)_{i\in I}$ be as in Definition~\ref{pku:def-kummer}, $e\in\Z$, and $(f_i)_{i\in I}$ any
family of maps $G_\nu\to\mathbb F_\ell$ indexed by the same set $I$. Then there is a Kummer map
$\kappa$ of exponent $e$ for $x$ off $T=S_0(x)$ whose reduction is not in the $\mathbb F_\ell$-span
of the $f_i$. More precisely, the reductions of the Kummer maps of exponent $e$ for $x$ off $S_0(x)$
contain an $\mathbb F_\ell$-subspace of dimension at least $\#I+r_F-4\ge\#I+1$.
\end{proposition}

\subsubsection{Duality: the inputs}\label{pku:ss-duality}

Let $A$ be a finite $G_F$-module of $\ell$-power order killed by $n$, and $A^*=\pkuHom(A,\mu_n)$.
For a place $v$ of $F$ the cup product and the evaluation $A\times A^*\to\mu_n$ give
$H^1(F_v,A)\times H^1(F_v,A^*)\to H^2(F_v,\mu_n)$, followed by the local invariant into $\Q/\Z$
(for real $v$, $H^2(\R,\mu_n)=\mathrm{Br}(\R)[n]\subseteq\frac12\Z/\Z$). For a subgroup $L_v\subseteq
H^1(F_v,A)$, $L_v^\perp\subseteq H^1(F_v,A^*)$ is its annihilator. We use:

\smallskip\noindent(D1) \emph{Local duality.} For finite $v$ the pairing $H^1(F_v,A)\times H^1(F_v,
A^*)\to\Q/\Z$ is perfect; if $v\nmid\ell$ and $A$ is unramified at $v$, the unramified subgroups
$H^1_\pkuur(F_v,A)=\ker(H^1(F_v,A)\to H^1(I_B,A))$ and $H^1_\pkuur(F_v,A^*)$ are exact annihilators of
each other (\cite[Ch.~I, Corollary~2.3 and Theorem~2.6]{MilneADT}; \cite[Theorems~7.2.6 and 7.2.15]{NSW}).

\smallskip\noindent(D2) \emph{Global reciprocity.} For $a\in H^1(F,A)$ and $b\in H^1(F,A^*)$, the
local invariants of $a_v\cup b_v$ vanish for almost all $v$ and $\sum_v\pkuinv_v(a_v\cup b_v)=0$, the
sum over all places of $F$; each term may be computed at any place of $\Omega$ above $v$. (This is
the reciprocity law of the Brauer group, $\sum_v\pkuinv_v=0$ on $\mathrm{Br}(F)$,
\cite[Theorem~8.1.17]{NSW}.)

\smallskip\noindent(D3) \emph{The Greenberg--Wiles formula.} Let $\Sigma$ be a finite set of places of
$F$ containing the archimedean ones and those above $\ell$, with $A$ unramified outside $\Sigma$; let
$\mathcal L=(L_v)_{v\in\Sigma}$, $L_v\subseteq H^1(F_v,A)$, and $\mathcal L^*=(L_v^\perp)_{v\in\Sigma}$.
Put $H^1_{\mathcal L}(F,A)=\{a\in H^1(G_\Sigma,A):a_v\in L_v\ (v\in\Sigma)\}$, $G_\Sigma$ the Galois
group of the maximal extension of $F$ unramified outside $\Sigma$, and likewise $H^1_{\mathcal L^*}(F,
A^*)$. These groups are finite and
\begin{equation}\label{pku:eq-GW}
\frac{\#H^1_{\mathcal L}(F,A)}{\#H^1_{\mathcal L^*}(F,A^*)}=\frac{\#H^0(F,A)}{\#H^0(F,A^*)}
\prod_{v\in\Sigma}\frac{\#L_v}{\#H^0(F_v,A)} ,
\end{equation}
with ordinary $H^0$ at the real places (\cite[Theorem~2.19]{DDT}; \cite[Theorem~8.7.9]{NSW}; over $\Q$ and for $\ell$ odd, \cite[Proposition~1.6]{WilesFLT}). The proof is the Poitou--Tate sequence and Tate's global Euler
characteristic formula.

\smallskip\noindent(D4) \emph{Chebotarev's density theorem}: for a finite Galois extension $K/F$ and
$\sigma\in\pkuGal(K/F)$ there are infinitely many primes $p$ of $F$, unramified in $K$, with a place
of $K$ over $p$ whose Frobenius element is $\sigma$.

\smallskip
Statement (D2) is used only for $a$ unramified outside a finite set; the terms at the places where
$a$ and $b$ are both unramified and $v\nmid\ell$ vanish by (D1).

\begin{lemma}[The pairing at a split prime]\label{pku:pairing}
Let $p\nmid\ell$ be a prime of $F$, $B$ a place over it, $A$ a finite $G_F$-module killed by $n$
on which $D_B$ acts trivially, and suppose $\auNm p\equiv1\pmod n$ (so $D_B$ acts trivially on
$\mu_n$ and $A^*$). Let $\phi\in D_B$ be a Frobenius element and $\gamma$ a normalised tame generator
at $B$. For a homomorphism $z\colon D_B\to A$ and an unramified homomorphism $\xi\colon D_B\to A^*$,
\[
\pkuinv_B(z\cup\xi)=\epsilon\,\theta\bigl(\xi(\phi)(z(\gamma))\bigr),
\]
where $\epsilon\in\{\pm1\}$ depends only on the conventions for cup products, local invariants and
the local reciprocity map, not on $p$, $B$, $A$, $z$, $\xi$.
\end{lemma}
\begin{proof}
This is Lemma~\ref{pse:pairing}, whose proof is local and applies verbatim to the local field
$F_p$; we recall it. Both sides are biadditive and compatible with decompositions of $A$; so we may
take $A=\Z/\ell^s$ ($s\le\nu$) with trivial action, $A^*=\mu_{\ell^s}$, and identify
$\theta_s\colon\mu_{\ell^s}\to\frac1{\ell^s}\Z/\Z$ by $\zeta^{n/\ell^s}\mapsto1/\ell^s$ (compatible
with $\theta$). By Kummer theory $\xi$ is the class of some $b\in F_p^\times$, a unit up to
$\ell^s$-th powers because $\xi$ is unramified, and $\xi(\phi)=\phi(b^{1/\ell^s})/b^{1/\ell^s}=
\bar b^{(\auNm p-1)/\ell^s}$ (as a root of unity). By local class field theory,
$\pkuinv(z\cup\partial b)=\pm z(\mathrm{rec}(b))$ (\cite[XIV, \S1, Proposition~3, p.~205]{SerreLF}), and for a unit $b$ the element $\mathrm{rec}(b)$ acts on $F_p(\varpi^{1/\ell^s})$
by $\varpi^{1/\ell^s}\mapsto\bar b^{\mp(\auNm p-1)/\ell^s}\varpi^{1/\ell^s}$ (the tame symbol,
\cite[Ch.~XIV, \S3, Proposition~8 and its Corollary, pp.~210--211]{SerreLF}; see also \S2, Proposition~6, p.~208). A homomorphism $D_B\to\Z/\ell^s$ restricted to $I_B$ factors through
the cyclic quotient of tame inertia of order $\ell^s$, generated by $\gamma$; so $z(\mathrm{rec}(b))
=kz(\gamma)$ with $\zeta^{kn/\ell^s}=\xi(\phi)^{\mp1}$, i.e.\ $k/\ell^s=\mp\theta(\xi(\phi))$. Hence
$\pkuinv(z\cup\xi)=\pm\theta(\xi(\phi))\,z(\gamma)$, with a sign fixed by the conventions.
\end{proof}

\subsubsection{Selmer structures}\label{pku:ss-selmer}

Fix $x$, $S_0=S_0(x)$, the places $B_v$ ($v\in S_0$) and $e$ as in Definition~\ref{pku:def-kummer};
put $k=\#I$, the number of primes of $S_0$ not above $\ell$. For every other place $v$ of $F$ fix a
place $B_v$ of $\Omega$ above it; for split primes take the chosen $B_p$. Write $D_v=D_{B_v}$,
$I_v=I_{B_v}$.

\begin{definition}\label{pku:def-sel}
Let $A\in\{M,M[\ell]\}$ and let $P$ be a finite set of split primes. For each place $v$ let $H_v
\subseteq D_v$ be: $D_v\cap\ker\pkucp$ if $v\mid\ell$; $D_v$ if $v\in S_0$, $v\nmid\ell$, or $v$ is
real; $\{1\}$ if $v\in P$; $I_v$ otherwise. Each $H_v$ acts trivially on $A$
(Lemma~\ref{pku:basic}). Put
\[
L^P_v(A)=\ker\bigl(H^1(F_v,A)\to\pkuHom(H_v,A)\bigr),
\]
the restriction to $H_v$ being a homomorphism. Let $\Sigma_P=S_0\cup P\cup\{\text{real places}\}$.
Then $\mathcal L^P(A)=(L^P_v(A))_{v\in\Sigma_P}$, and $\mathcal L(A)$ is the family with the same
$\Sigma_P$ and the same $L_v$ for $v\notin P$, but $L_p=H^1_\pkuur(F_p,A)$ for $p\in P$.
\end{definition}

So $H^1_{\mathcal L^P}(F,A)$ consists of the classes unramified outside $S_0\cup P$ with the
conditions at $S_0$ and at the real places, and $H^1_{\mathcal L}(F,A)$ of those unramified outside
$S_0$; $H^1_{\mathcal L}$ does not depend on $P$, and neither does $H^1_{\mathcal L^*}(F,A^*)$
(at $p\in P$ the dual condition is $H^1_\pkuur(F_p,A^*)$ by (D1)). We write
\[
\pkuSel=H^1_{\mathcal L^*(M)}(F,M^*)
\]
for the classes unramified outside $S_0$ that lie in $L_v(M)^\perp$ at $v\in S_0$ (no condition at
$v\in S_0$, $v\nmid\ell$, and at the real places, where $L_v=0$). The class of a test cocycle for
$x$ and $P$ lies in $H^1_{\mathcal L^P}(F,M)$, by Remark~\ref{pku:rem-def}(1).

\begin{lemma}[Local terms]\label{pku:local}
Let $A\in\{M,M[\ell]\}$.
\begin{enumerate}
\item[(a)] If $v\in S_0$, $v\nmid\ell$, or $v$ is real, then $L_v(A)=0$ and $H^0(F_v,A)=A$.
\item[(b)] If $v\mid\ell$, then $\#L_v(A)\le\#H^0(F_v,A)$.
\item[(c)] If $p$ is a split prime, then $D_p$ acts trivially on $A$ and $A^*$, $\auNm p\equiv1
\pmod n$, the map $H^1(F_p,A)\to A$, $z\mapsto z(\gamma_p)$, is surjective with kernel
$H^1_\pkuur(F_p,A)$, and $H^1_\pkuur(F_p,A^*)\to A^*$, $\xi\mapsto\xi(\phi_p)$, is bijective. In
particular $\#H^1(F_p,A)=(\#A)^2=\#A\cdot\#H^0(F_p,A)$.
\end{enumerate}
\end{lemma}
\begin{proof}
(a) Here $H_v=D_v$ acts trivially on $A$ (Lemma~\ref{pku:basic}(iv) for $v=v_i$; $D_v=\{1,c'\}$ with
$\pkucp(c')=1$ for real $v$). So $H^0(F_v,A)=A$, $H^1(F_v,A)=\pkuHom(D_v,A)$ and the restriction to
$H_v=D_v$ is the identity; $L_v(A)=0$.

(b) As $H_v$ is normal in $D_v$ and acts trivially, inflation--restriction identifies $L_v(A)$ with
$H^1(D_v/H_v,A)$. Now $D_v/H_v\cong\Gamma=\pkucp(D_v)$, a closed subgroup of $\Z_\ell^\times$ ($\ell$
odd) or of $1+4\Z_2$ ($\ell=2$, Lemma~\ref{pku:basic}(i)), hence procyclic; let $\sigma$ be a
topological generator. A continuous cocycle on $\Gamma$ is determined by its value at $\sigma$ (the
cocycle rule determines it on the dense subgroup generated by $\sigma$), and the coboundaries take
the values $(\sigma-1)A$ at $\sigma$. So $\#H^1(\Gamma,A)\le\#A/(\sigma-1)A=\#\ker(\sigma-1)=\#A^\Gamma
=\#H^0(F_v,A)$.

(c) $D_p$ is topologically generated by $I_p$, on which $\pkucp$ and the action on $\mu_n$ are
trivial ($p\nmid\ell$), and $\phi_p\in G_\nu$, which acts trivially on $A$, $A^*$, $\mu_n$
(Lemma~\ref{pku:basic}(ii)); $\phi_p$ acts on $\mu_n$ by $\auNm p$-th powers, so $\auNm p\equiv1
\pmod n$. Hence $H^1(F_p,A)=\pkuHom(D_p,A)$. A homomorphism $D_p\to A$ kills wild inertia (pro-$p$),
and its restriction to $I_p$ factors through the maximal quotient of exponent $n$ of tame inertia,
which is cyclic of order $n$ generated by the image of $\gamma_p$ (it maps isomorphically to $\mu_n$
by $\tau\mapsto\tau(\varpi^{1/n})/\varpi^{1/n}$). So the kernel of $z\mapsto z(\gamma_p)$ is the set
of homomorphisms vanishing on $I_p$, i.e.\ $H^1_\pkuur(F_p,A)$. Surjectivity: as $\mu_n\subseteq F_p$,
$\chi(g)=g(\varpi^{1/n})/\varpi^{1/n}$ is a homomorphism $D_p\to\mu_n$ with $\chi(\gamma_p)=\zeta$;
for $a\in A$, $g\mapsto n\,\theta(\chi(g))\,a$ is a homomorphism taking the value $a$ at $\gamma_p$.
Finally $H^1_\pkuur(F_p,A^*)=\pkuHom(D_p/I_p,A^*)$ and $D_p/I_p\cong\widehat\Z$ is generated by the
image of $\phi_p$.
\end{proof}

\subsubsection{Dual Selmer classes are relations}\label{pku:ss-relations}

\begin{lemma}\label{pku:relation}
Let $\xi\colon G_F\to M^*$ be a continuous cocycle whose class lies in $\pkuSel$. Define $y_\xi
\colon G_F\to M^\vee$ by $y_\xi(g)=\theta\circ\xi(g)$ for $g\in G_\nu$ and $y_\xi(g)=0$ otherwise.
Then $y_\xi$ is a relation off $S_0$, and $\kappa_\xi(g)=n\,y_\xi(g)(1)$ is a Kummer map of exponent
$e$ for $x$ off $S_0$. It depends only on the class of $\xi$, and $\xi\mapsto y_\xi$ is additive.
\end{lemma}
\begin{proof}
A coboundary $g\mapsto gf-f$ vanishes on $G_\nu$, which acts trivially on $M^*$; so $y_\xi$
depends only on the class, and it is clearly additive in $\xi$. It is locally constant and vanishes
off $G_\nu$. For $u,v\in G_\nu$, $\xi(uv)=\xi(u)+u\xi(v)=\xi(u)+\xi(v)$. For $g\in G_F$ and $u\in
G_\nu$, the cocycle rule gives $\xi(gug^{-1})=\xi(g)+g\xi(u)-gug^{-1}\xi(g)=g\xi(u)$; and
$\theta\bigl((gf)(a)\bigr)=\theta\bigl(g(f(g^{-1}a))\bigr)=\pgcyc(g)\,\theta\bigl(f(g^{-1}a)\bigr)$.
So $y_\xi(gug^{-1})(a)=\pgcyc(g)\,y_\xi(u)(g^{-1}a)$. Let $B$ be a place over a prime $v\notin S_0$
and $\tau\in I_B$; if $\tau\notin G_\nu$ then $y_\xi(\tau)=0$; otherwise $I_B$ acts trivially on $M^*$
($v\nmid\ell$ and Lemma~\ref{pku:basic}(i)), the class of $\xi$ is unramified at $v$, so $\xi|_{I_B}$
is a coboundary for the trivial action, i.e.\ zero, and $y_\xi(\tau)=0$.

Reciprocity: let $P$ be a finite set of split primes (not in $S_0$) and $z$ a test cocycle for $x$
and $P$. Apply (D2) to $[z]$ and $[\xi]$. At $v\in S_0$ and at real $v$, $[z]_v\in L_v(M)$ and
$[\xi]_v\in L_v(M)^\perp$, so the term vanishes. At a finite $v\notin S_0\cup P$, $v\nmid\ell$, and
both classes are unramified, so the term vanishes by (D1). At $p\in P$, $D_p$ acts trivially on $M$,
$\auNm p\equiv1\bmod n$ (Lemma~\ref{pku:local}(c)), $z|_{D_p}$ is a homomorphism and $\xi|_{D_p}$ an
unramified homomorphism; by Lemma~\ref{pku:pairing} the term is $\epsilon\,\theta\bigl(\xi(\phi_p)(
z(\gamma_p))\bigr)=\epsilon\,y_\xi(\phi_p)(z(\gamma_p))$, as $\phi_p\in G_\nu$. Hence $\epsilon\sum_{p\in
P}y_\xi(\phi_p)(z(\gamma_p))=0$, which is~\eqref{pku:eq-recip}. Finally $y_\xi(g)(a)=a\,y_\xi(g)(1)
=a\kappa_\xi(g)/n$.
\end{proof}

\subsubsection{The values at Frobenius elements}\label{pku:ss-image}

Let $P$ be a finite set of split primes. Define
\[
\lambda_P\colon H^1_{\mathcal L^P}(F,M)\to M^P,\quad z\mapsto(z(\gamma_p))_{p\in P},\qquad
\mu_P\colon\pkuSel\to(M^\vee)^P,\quad\xi\mapsto(\theta\circ\xi(\phi_p))_{p\in P}.
\]
They are well defined ($\gamma_p$ acts trivially on $M$ and $\phi_p$ on $M^*$, so coboundaries vanish
there). For $b\in M^P$ and $f\in(M^\vee)^P$ write $\langle b,f\rangle=\sum_pf_p(b_p)$; this pairing is
perfect. For $B\subseteq M^P$ let $\pkuAnn(B)\subseteq(M^\vee)^P$ be its annihilator.

\begin{lemma}\label{pku:image}
$\mu_P(\pkuSel)=\pkuAnn\bigl(\lambda_P(H^1_{\mathcal L^P}(F,M))\bigr)$.
\end{lemma}
\begin{proof}
Put $V=\bigoplus_{p\in P}H^1(F_p,M)/H^1_\pkuur(F_p,M)$ and $V'=\bigoplus_{p\in P}H^1_\pkuur(F_p,M^*)$. By
(D1) the local pairings induce a perfect pairing $V\times V'\to\Q/\Z$, and by
Lemma~\ref{pku:local}(c) $z\mapsto(z(\gamma_p))$ and $\xi\mapsto(\theta\circ\xi(\phi_p))$ identify $V$
with $M^P$ and $V'$ with $(M^\vee)^P$; by Lemma~\ref{pku:pairing} the pairing becomes $\epsilon
\langle\ ,\ \rangle$. Let $\lambda'\colon H^1_{\mathcal L^P}(F,M)\to V$ and $\mu'\colon\pkuSel\to V'$ be
the localisation maps, which become $\lambda_P$, $\mu_P$. Their kernels are $H^1_{\mathcal L}(F,M)$
(classes unramified at $P$) and $H^1_{(\mathcal L^P)^*}(F,M^*)$ (classes of $\pkuSel$ vanishing at
$P$; the dual of the condition $L^P_p=H^1(F_p,M)$ is $0$).

The inclusion $\mu'(\pkuSel)\subseteq\pkuAnn(\lambda'(H^1_{\mathcal L^P}))$ is (D2) applied to $a\in
H^1_{\mathcal L^P}(F,M)$, $b\in\pkuSel$: the terms at $S_0$ and the real places vanish ($a_v\in L_v$,
$b_v\in L_v^\perp$), those at the other places off $P$ vanish by (D1), and the terms at $P$ add up to
the pairing of $\lambda'(a)$ and $\mu'(b)$.

Counting: $\#\mu'(\pkuSel)=\#\pkuSel/\#H^1_{(\mathcal L^P)^*}(F,M^*)$, and
$\#\pkuAnn(\lambda'(H^1_{\mathcal L^P}))=\#V/\#\lambda'(H^1_{\mathcal L^P})=\#V\cdot\#H^1_{\mathcal L}(F,M)/
\#H^1_{\mathcal L^P}(F,M)$. Apply~\eqref{pku:eq-GW} with $\Sigma=\Sigma_P$ to $\mathcal L^P(M)$ and to
$\mathcal L(M)$ and divide: all factors agree except those at $p\in P$, whose quotient is
$\#H^1(F_p,M)/\#H^1_\pkuur(F_p,M)$; so
\[
\frac{\#H^1_{\mathcal L^P}(F,M)}{\#H^1_{(\mathcal L^P)^*}(F,M^*)}=\frac{\#H^1_{\mathcal L}(F,M)}
{\#\pkuSel}\cdot\#V .
\]
Hence the two finite groups have the same order, and the inclusion is an equality.
\end{proof}

\subsubsection{The \texorpdfstring{$\ell$}{l}-torsion of the relaxed Selmer group}\label{pku:ss-torsion}

\begin{lemma}\label{pku:torsion}
There is a finite set $P$ of split primes with
\[
\#H^1_{\mathcal L^P}(F,M)[\ell]\le\ell^{\#P-k-r_F+4}.
\]
\end{lemma}
\begin{proof}
\emph{Step 1: choice of $P$.} Let $A=M[\ell]$, a group of order $\ell$, and $A^*=\pkuHom(A,\mu_n)
=\pkuHom(A,\mu_\ell)$. The group $Y=H^1_{\mathcal L^*(A)}(F,A^*)$ (classes unramified outside $S_0$
with the dual conditions at $S_0$ and at the real places) is finite by (D3). For each $\xi\in Y$ whose
restriction to $G_\nu$ is nonzero, choose a split prime $p_\xi$ with $\xi(\phi_{p_\xi})\ne0$,
the $p_\xi$ distinct: the restriction $\xi|_{G_\nu}$ (well defined on the class, as $G_\nu$ acts
trivially on $A^*$) is a nonzero continuous homomorphism; its kernel $K$ is open in $G_\nu$ and normal
in $G_F$, since $\xi(gug^{-1})=g\xi(u)$ for $u\in G_\nu$ as in the proof of Lemma~\ref{pku:relation};
pick $u\in G_\nu$ with $\xi(u)\ne0$; by (D4) applied to the fixed field of $K$ and the class of $u$,
there are infinitely many primes $p\notin S_0$ of $F$ with a place $B_p$ over them and a Frobenius
element $\phi_p\in uK\subseteq G_\nu$; then $\xi(\phi_p)=\xi(u)\ne0$, and $\xi(\phi_p)$ does not depend
on the choice of the Frobenius element at $B_p$ ($\xi$ is unramified at $p$). Choose also a
normalised tame generator $\gamma_p$ at $B_p$. Let $P$ be the set of the $p_\xi$.

\emph{Step 2: the dual group for $P$.} If $\xi\in H^1_{(\mathcal L^P)^*}(F,A^*)$, then $\xi\in Y$
(at $p\in P$ the condition is $\xi_p=0$, in particular $\xi$ is unramified there) and
$\xi(\phi_p)=0$ for $p\in P$ ($\xi|_{D_p}$ is a coboundary for the trivial action of $D_p$ on $A^*$,
hence zero). By the choice of $P$, $\xi|_{G_\nu}=0$. So $H^1_{(\mathcal L^P)^*}(F,A^*)$ lies in the
kernel of the restriction $H^1(F,A^*)\to H^1(G_\nu,A^*)$, which is $H^1(\Delta,A^*)$, $\Delta=G_F/
G_\nu$ (inflation--restriction; $G_\nu$ acts trivially on $A^*$). Write $\Delta=\Delta_\ell\times
\Delta'$ with $\Delta_\ell$ its $\ell$-Sylow subgroup and $\#\Delta'$ prime to $\ell$
(Lemma~\ref{pku:basic}(ii)). By inflation--restriction for the normal subgroup $\Delta_\ell$, the sequence $0\to H^1(\Delta',(A^*)^{\Delta_\ell})\to H^1(\Delta,A^*)\to H^1(
\Delta_\ell,A^*)$ is exact and the first group is $0$; $\Delta_\ell$ acts trivially on $A^*$ (an
$\ell$-group acting on a group of order $\ell$ through $\mathbb F_\ell^\times$); so $H^1(\Delta,A^*)$
embeds in $\pkuHom(\Delta_\ell,\mathbb F_\ell)$, of order at most $\ell^2$. Thus
$\#H^1_{(\mathcal L^P)^*}(F,A^*)\le\ell^2$.

\emph{Step 3: the Selmer group of $M[\ell]$.} Apply~\eqref{pku:eq-GW} to $A$, $\Sigma_P$ and
$\mathcal L^P(A)$. The global factor is $\#H^0(F,A)/\allowbreak \#H^0(F,A^*)\le\ell$. By
Lemma~\ref{pku:local}, the local factor $\#L_v/\#H^0(F_v,A)$ is $\ell^{-1}$ at each of the $k$ primes
of $S_0$ not above $\ell$ and at each of the $r_F$ real places, at most $1$ at $v\mid\ell$, and
$\#H^1(F_p,A)/\#H^0(F_p,A)=\ell$ at $p\in P$. With Step 2,
\[
\#H^1_{\mathcal L^P}(F,M[\ell])\le\ell^2\cdot\ell\cdot\ell^{\#P-k-r_F}=\ell^{\#P-k-r_F+3}.
\]

\emph{Step 4: from $M[\ell]$ to $M$.} Let $j\colon M[\ell]\to M$ be the inclusion, $\pi\colon M\to
M/M[\ell]$ the projection and $\iota\colon M/M[\ell]\to M$ the injective map $\pi(a)\mapsto\ell a$,
with cokernel $M/\ell M$; all are $G_F$-maps, and $\iota\circ\pi$ is multiplication by $\ell$. Work in
$H^1(G_{\Sigma_P},-)$. Let $X=H^1_{\mathcal L^P}(F,M)[\ell]$. For $a\in X$, $\iota_*\pi_*(a)=\ell a=0$,
so $\pi_*(a)$ lies in the image of the connecting map $H^0(G_{\Sigma_P},M/\ell M)\to H^1(G_{\Sigma_P},
M/M[\ell])$; as $\#(M/\ell M)=\ell$, $\#\pi_*(X)\le\ell$. Let $X_0=X\cap\ker\pi_*$, so $\#X\le\ell\,\#X_0$.
An $a\in X_0$ is $j_*(w)$ for some $w\in H^1(G_{\Sigma_P},M[\ell])$. For every $v\in\Sigma_P$,
$j_*(w_v)=a_v\in L^P_v(M)$, i.e.\ the restriction of $j\circ w$ to $H_v$ is $0$; since $H_v$ acts
trivially and $j$ is injective, the restriction of $w$ to $H_v$ is $0$, i.e.\ $w_v\in L^P_v(M[\ell])$.
So $X_0\subseteq j_*\bigl(H^1_{\mathcal L^P}(F,M[\ell])\bigr)$ and, by Step 3,
$\#X\le\ell\cdot\ell^{\#P-k-r_F+3}$.
\end{proof}

\subsubsection{Proof of Proposition~\ref{pku:main}}\label{pku:ss-proof}

Let $P$ be as in Lemma~\ref{pku:torsion}, $B=\lambda_P(H^1_{\mathcal L^P}(F,M))\subseteq M^P$ and $R=
\pkuAnn(B)=\mu_P(\pkuSel)$ (Lemma~\ref{pku:image}). Let $X=M[\ell]^P\subseteq M^P$, and let $r\colon
(M^\vee)^P\to\pkuHom(X,\Q/\Z)$ be $r(f)(b)=\langle b,f\rangle$.

(1) $r(R)=\pkuHom(X/(X\cap B),\Q/\Z)$. Indeed $r(R)$ kills $X\cap B$; conversely a character of $X$
trivial on $X\cap B$ is a character of $(X+B)/B$, which extends to $M^P/B$ ($\Q/\Z$ is divisible),
and such an extension is $b\mapsto\langle b,f\rangle$ for some $f\in R$. Since $X\cap B=B[\ell]$,
$\#r(R)=\ell^{\#P}/\#B[\ell]$.

(2) $\#B[\ell]\le\#H^1_{\mathcal L^P}(F,M)[\ell]$: for a surjection $Y\to B$ of finite abelian groups,
$\#B[\ell]=\#(B/\ell B)\le\#(Y/\ell Y)=\#Y[\ell]$. With Lemma~\ref{pku:torsion},
$\#r(R)\ge\ell^{k+r_F-4}$.

(3) Let $\rho\colon\pkuSel\to\mathrm{Map}(G_\nu,\mathbb F_\ell)$, $\rho(\xi)=\pkured\kappa_\xi$, with
$\kappa_\xi$ the Kummer map of Lemma~\ref{pku:relation}; $\rho$ is additive, so its image is an
$\mathbb F_\ell$-subspace consisting of reductions of Kummer maps of exponent $e$ for $x$ off $S_0$.
The composite $r\circ\mu_P$ factors through $\rho$: for $b\in X$ write $b_p=(n/\ell)m_p$; then
$r(\mu_P(\xi))(b)=\sum_py_\xi(\phi_p)(b_p)=\sum_p\kappa_\xi(\phi_p)m_p/\ell$, which depends only on the
values of $\pkured\kappa_\xi$ at the $\phi_p\in G_\nu$. Hence $\#\rho(\pkuSel)\ge\#r(\mu_P(\pkuSel))=
\#r(R)\ge\ell^{k+r_F-4}$, i.e.\ $\dim\rho(\pkuSel)\ge k+r_F-4\ge k+1$ by Lemma~\ref{pku:basic}(iii).

(4) The span of the $f_i$ ($i\in I$) has dimension at most $k=\#I<\dim\rho(\pkuSel)$, so some
$\rho(\xi)$ is not in it. \qed

\begin{remark}\label{pku:rem-source}
The proof uses Lemma~\ref{pku:image} (Greenberg--Wiles and global duality, with $T=S_0$) to produce
relations, and Lemma~\ref{pku:torsion} (Greenberg--Wiles for $M[\ell]$, Chebotarev, and the
$\ell$-torsion argument of Step 4) for the count; the many real places of $F$ give room for the
constant losses.
\end{remark}

\subsubsection{Detecting points, and the triangle}\label{pku:ss-triangle}

Let $h=r^8$, $t=5r$, and write $N_0=q_1d$ (the integer $N$ of \S\ref{pod:sec}); recall
$N_0+t+2\le h$ (proof of Corollary~\ref{pod:kummer-b}). The integer called $N$ in \S\ref{pri:sec} is
$N_{\mathrm{ch}}$, and is not used under that name here.

\begin{proposition}[Detecting points]\label{pku:b}
Let $g_0$ be a topological generator of $\pgG$ modulo $\pgG\cap\ker\pgcyc$.
\begin{enumerate}
\item[(i)] For every prime $\ell$ there is a finite set $s\subseteq\pgG\cap\ker\pgcyc$ such that, for
all $\auO$, $\nu$ and $e$, a locally constant cocycle $z\colon\pgG\to(\auO/\ell^\nu)(-e)$ unramified
outside $\ell$ with $z(g_0)=0$ and $z|_s=0$ is zero.
\item[(ii)] If $\ell\le h^2$, $s$ can be chosen with $\#s\le h^{20}$, and then $(t+1)(\#s+1)\le h^{21}$.
\end{enumerate}
\end{proposition}
\begin{proof}
(i) is the second statement of Lemma~\ref{pod:detect} with $y=\max(N_0,\lceil\ell^{1/2}\rceil)$. (ii)
is the same with $y=h$, allowed as $h\ge N_0$ and $\ell\le h^2$; then $\#s\le h^{20}$ and, as $t+1\le
h-1$, $(t+1)(\#s+1)\le(h-1)(h^{20}+1)=h^{21}-h^{20}+h-1\le h^{21}$.
\end{proof}

For $\ell>h^2$ no bound on $\#s$ is claimed or needed.

\begin{proposition}[Proof of Proposition~\ref{pri:triangle}]\label{pku:triangle}
Proposition~\ref{pri:triangle} holds; in its proof (\S\ref{pri:ss-triangle}): (1) in ``Test points'', $s$ is the set
of Proposition~\ref{pku:b} for $g_0=a^{-1}$, and the inequality $(\#J+\#s+1)hh_0\le h^{22}h_0$ is
trivial if $h_0=0$ and follows from Proposition~\ref{pku:b}(ii) if $h_0\ne0$, because then $\ell\le
h^2$ by hypothesis; (2) in ``Induction'', Proposition~\ref{pku:main} is applied with $x$ the earlier
places of exponent $e$, the $f_i$ their reductions and $T_j=S_0(x)$; (3) the loss at the new place
is $2\kappa_j(\phi_j)=4\kappa_0(u)$; (4) in ``Property (iii)'' the local conditions are those of
Definition~\ref{pku:def-kummer}(a), at $S_0(x)=T_j$.
\end{proposition}
\begin{proof}
We give the steps of the proof in \S\ref{pri:ss-triangle} in full.

\emph{Choice of $a$, $h_0$.} If $\ell>h^2$ then $\ell\nmid N_0$ ($N_0\le h^2$) and $t<\ell-1$; an inertia
element at $\ell$ whose cyclotomic character reduces to a generator of $\mathbb F_\ell^\times$ fixes
$E\subseteq\Q(\mu_{N_0})$, and $\pgcyc^{\,e}-1$ is a unit at it for $0<e<t$; as every $\pgcyc(g)^e-1$ is
a multiple of $\pgcyc(a)^e-1$ ($a$ generates $\pgG$ modulo $\ker\pgcyc$, a procyclic quotient as
$\pgcyc\equiv1\bmod4$ on $\pgG$), $h_0=0$ works. Otherwise $h_0=\prival_\ell(N_{\mathrm{ch}})$ works,
by the Frobenius element at a chosen place above $p_i\ne\ell$. In both cases $h_0\le\prival_\ell(
N_{\mathrm{ch}})$ and ($\ell\le h^2$ or $h_0=0$).

\emph{Separation (ii).} With $u=\pgcyc(a)$: $u^e-u^k=u^k(u^{e-k}-1)$ for $k<e$ (and the negative for
$k>e$), so each factor divides $\ell^{h_0}$ and the product over the $t'-1$ indices $k\ne e$ divides
$\ell^{(t'-1)h_0}$. Also $4\mid u^e-1\mid\ell^{h_0}$ for $e>0$ ($u\equiv1\bmod4$; if $\ell$ is odd,
$4$ is a unit), so $4(u^e-1)^2\mid\ell^{3h_0}$. Both divide $\ell^{hh_0}$ as $t'+2\le h$. For the
triangle for $a^{-1}$ these are the conditions on $\pgcyc(a^{-1})^{-e}=u^e$.

\emph{Test points (i).} $J=\{(e,g):0<e<t',\ g\in s\}$, ordered lexicographically with $e$ first,
$\prity(e,g)=e$; $J$ and $\prity$ do not depend on $\nu$, since $s$ does not. $\#J+\#s+1\le
t'\#s+1\le(t+1)(\#s+1)$. If $h_0=0$ both sides of $(\#J+\#s+1)hh_0\le h^{22}h_0$ vanish; if $h_0\ne0$,
then $\ell\le h^2$ and Proposition~\ref{pku:b}(ii) gives $(t+1)(\#s+1)h\le h^{22}$. Detection with
$a^{-1}$ is Proposition~\ref{pku:b}(i)--(ii) for $g_0=a^{-1}$.

\emph{Induction.} Let the places $w_i$, $i<j$, with Kummer maps $\kappa_i$ of exponent $\prity(i)$
be chosen, above distinct primes of $F$, with the reductions of the $\kappa_i$ of exponent
$e=\prity(j)$ linearly independent. Proposition~\ref{pku:main}, applied to the family $x$ of the
$w_i$ with $\prity(i)=e$ and to their reductions, gives a Kummer map $\kappa_0$ of exponent $e$ for
$x$ off $T_j=S_0(x)$ whose reduction is outside their span. The Kummer map attached to the new
place $w_j$ chosen below is $\kappa_j:=\kappa_0$, of exponent $e$ for $x$ off $T_j=S_0(x)$; this is
what Property (iii) uses. By Remark~\ref{pku:rem-def}(2) all these
maps are additive on $G_\nu$, negated by $g\mapsto g^c$, and conjugation by $g\in\pgG$ multiplies a
map of exponent $d$ by $\pgcyc(g)^{d+1}$; so conjugation by $a$ followed by the power
$\pgcyc(a)^{-1}\bmod\ell^\nu$ multiplies it by $\pgcyc(a)^d$. The earlier maps of exponent $e$ and
$\kappa_0$ have independent reductions, so together they map $G_\nu$ onto $(\Z/\ell^\nu)^{m}$
(Nakayama); take $g'\in G_\nu$ where $\kappa_0$ is $1$ and the earlier maps of exponent $e$ vanish,
and apply $\prod_{k\ne e,\,k<t'}(\text{conj}_a\text{ then power}-\pgcyc(a)^k)$: this gives $u\in
G_\nu$ with $\kappa_i(u)=0$ for all $i<j$ (for $\prity(i)\ne e$ the factor $k=\prity(i)$ vanishes) and
$\kappa_0(u)=\prod_{k\ne e}(\pgcyc(a)^e-\pgcyc(a)^k)$. With $g=g_j$ the test point of $j$ and
$\sigma=\delta^{\ell^\nu}gg^cu(u^c)^{-1}$ ($\delta$ given by the hypothesis on $\mathcal W$ for $\nu$
and $g''=gg^cu(u^c)^{-1}$), every
Kummer map of exponent $d$ takes the value $2\kappa(u)$ near $\sigma$, and every even class of
exponent $e$ unramified outside $\ell$ the value $2z(g)$, as in \S\ref{pri:ss-triangle}. Let $U\subseteq G_\nu$ be an open neighbourhood of $\sigma$ (this is possible: $\sigma\in G_\nu$,
since $\delta$, $g$, $g^c$ have trivial cyclotomic character modulo $\ell^\nu$ and $u\in G_\nu$, and
$G_\nu$ is open) on which the
$\kappa_i$ ($i<j$), $\kappa_0$ and all even classes of exponent $e$ unramified outside $\ell$ are
constant (after composing with additive maps to $\Z/\ell^\nu$, the latter are finitely many cocycles,
all factoring through one finite quotient of $\pgG$); by the hypothesis, $\mathcal W$ contains almost
all Frobenius data with Frobenius in a nonempty open subset of $U$, and Chebotarev gives infinitely
many of them. Choose $w_j
\in\mathcal W$ of degree one with Frobenius $\phi_j\in U$, not above $\ell$ and above a prime
of $F$ not below an earlier place; this prime is then outside every $T_i$, $i\le j$. Then
$\kappa_i(\phi_j)=2\kappa_i(u)=0$ for $i<j$, $z(\phi_j)=2z(g)$, $\phi_j\in G_\nu$, and the loss is
\[
2\kappa_j(\phi_j)=4\kappa_0(u)=4\prod_{k\ne e}(\pgcyc(a)^e-\pgcyc(a)^k)\ \Big|\ \ell^{(3+(t'-1))h_0}
\ \Big|\ \ell^{hh_0}
\]
in $\Z/\ell^\nu$ (using $4\mid\ell^{3h_0}$ if $j$ exists). The independence of the reductions of the
maps of exponent $e$ up to $j$ holds by construction.

\emph{Property (iii).} Let $z$ be an even class of exponent $e=\prity(j)$ for the $w_i$ with
$z(\tau_i)=0$ for $i<j$ and $z(\phi_i)=0$ for $i<j$ of exponent $e$. Then $z$ is unramified at $w_i$
and $\overline w_i$ for $i<j$. The values of $z$ lie in $\auO/\ell^\nu$; the additive maps
$f\colon\auO/\ell^\nu\to\Z/\ell^\nu$ separate points and $f\circ z$ is again an even class with values in
$(\Z/\ell^\nu)(-e)$ with the same vanishing properties, so we may and do assume $\auO=\Z_\ell$. Its
extension $z^+(gc)=z(g)$ to $G_F$ is a cocycle for
$(\pkucp)^{-e}$, and $2z^+$ is a test cocycle for $x$ (the earlier places of exponent $e$) and $P=
\{$primes below $w_i$, $i\ge j\}$ in the sense of Definition~\ref{pku:def-kummer}(a): it vanishes at
complex conjugations ($2z^+(c')=z^+(c'^2)=0$, as $\pkucp(c')=1$); on $D_{B_v}\cap\ker\pkucp$ for
$v\mid\ell$ (for such $g$, $2z^+(g)=z^+(g^2)$ and $g^2\in\pgG\cap\ker\pgcyc$ lies in a decomposition
group above $\ell$, where $z$ vanishes); on $D_{A_i}$ for the earlier places of exponent $e$ ($z$
vanishes on $I_{A_i}$ and at $\phi_i$, which generate it topologically); and it is unramified outside
$S_0(x)\cup P$. The primes of $P$ are split primes not in $T_j$ (distinct, by the choice of the
$w_i$), with Frobenius elements $\phi_i\in G_\nu$ (the $\phi_j$ of Definition~\ref{pri:def-tri} are
Frobenius elements in the arithmetic sense of \S\ref{pku:ss-statement}, as chosen above) and normalised tame
generators $\gamma_i$ at the places $A_i$ of the $w_i$. The reciprocity law~\eqref{pku:eq-recip} of $\kappa_j$ gives $\sum_{i\ge j}
\kappa_j(\phi_i)\,2z(\gamma_i)=0$ (up to a common sign); $\kappa_j(\phi_i)=0$ for $i>j$, so
$2\kappa_j(\phi_j)z(\gamma_j)=0$. The value of $z$ at any element of $I_{A_j}$, in particular at
$\tau_j$, is a multiple of $z(\gamma_j)$; hence every multiple of $2\kappa_j(\phi_j)$ in $\Z/\ell^\nu$,
in particular $\ell^{hh_0}$, kills $z(\tau_j)$.
\end{proof}

\begin{remark}\label{pku:rem-tri-source}
The triangle constants are $b=h\,h_0$, $c=h^{22}h_0$ with $h=r^8$, as in
Proposition~\ref{pri:triangle}; the size count $(\#J+\#s+1)\le(t+2)h^{20}\le h^{21}$ is used only
when $\ell\le h^2$, the case $h_0=0$ being trivial.
\end{remark}



\section{Auxiliary statements for Ribet's method}\label{C3:g7}
\providecommand{\auO}{\mathcal O}
\providecommand{\auQlb}{\overline{\Q}_\ell}
\providecommand{\auGL}{\mathrm{GL}}
\providecommand{\auF}{\mathbb F}
\providecommand{\auU}{\mathrm U}
\providecommand{\auNm}{\operatorname{Nm}}
\providecommand{\pgZb}{\overline{\Z}_\ell}
\providecommand{\pgG}{G_E}
\providecommand{\pgtr}{\operatorname{tr}}
\providecommand{\pgFrob}{\operatorname{Frob}}
\providecommand{\priX}{\mathcal X}
\providecommand{\priT}{\mathbb T}
\providecommand{\pgNch}{N_{\mathrm{ch}}}
\providecommand{\pgaWD}{\operatorname{WD}}
\providecommand{\auf}[1]{\texttt{#1}}
\providecommand{\paxch}{\operatorname{ch}}
\providecommand{\paxL}{\mathcal L}

\subsection{Auxiliary statements for Ribet's method}\label{pax:sec}

This section proves statements that \S\ref{pri:ss-inputs} uses without proof:
Amitsur's formula, on which Lemma~\ref{pri:amitsur} rests (\S\ref{pax:ss-amitsur}); the chain
lemma, Lemma~\ref{pri:chain} (\S\ref{pax:ss-chain}); the inputs on forms, Lemma~\ref{pri:forms}
(\S\ref{pax:ss-forms}); and the block-diagonal lattices, Lemma~\ref{pri:block}
(\S\ref{pax:ss-block}); the integrality of the pairing at a free level (\S\ref{pax:ss-pairing}); the
dictionary at an auxiliary place, Proposition~\ref{pri:dict} (\S\ref{pax:ss-dict}); the base level of a label,
Proposition~\ref{pri:BL} (\S\ref{pax:ss-BL}); the sign hypothesis of Proposition~\ref{pga:G123}
(\S\ref{pax:ss-sign}); and the convention check of \S\ref{pga:ss-M} (\S\ref{pax:ss-conv}). Conventions are those of \S\ref{pri:sec}.

\subsubsection{Amitsur's formula}\label{pax:ss-amitsur}

Let $x_1,\dots,x_k$ be commuting indeterminates. A \emph{word} is a finite sequence $u=u_1\cdots
u_m$ in $\{1,\dots,k\}$, of length $|u|=m\ge1$; for matrices $A_1,\dots,A_k$ of the same size put
$A_u=A_{u_1}\cdots A_{u_m}$ and $x^u=x_{u_1}\cdots x_{u_m}$. A word is \emph{primitive} if it is not
a power $v^d$, $d\ge2$, of a shorter word; the cyclic rotations of a primitive word $w$ are $|w|$
distinct words; a \emph{Lyndon word} is a primitive word which is lexicographically smallest among
its rotations. Every word is uniquely $w'^{\,d}$ with $w'$ primitive, and every primitive word is a
rotation of exactly one Lyndon word.

\begin{proposition}[Amitsur's formula]\label{pax:amitsur}
Let $\mathsf A$ be a commutative ring and $A_1,\dots,A_k\in M_\sigma(\mathsf A)$. In the ring of
formal power series $\mathsf A[[x_1,\dots,x_k]]$,
\begin{equation}\label{pax:eq-amitsur}
\det\Bigl(1-\sum_{i=1}^kx_iA_i\Bigr)=\prod_{w\ \mathrm{Lyndon}}\det\bigl(1-x^wA_w\bigr),
\end{equation}
the product being convergent (only finitely many factors differ from $1$ modulo each power of the
ideal $(x_1,\dots,x_k)$). Consequently, for every $j\ge0$ the coefficient of $T^j$ in $\det(1-T\sum_i
A_i)\in\mathsf A[T]$ is a polynomial with integer coefficients, independent of $\mathsf A$ and the
$A_i$, in the coefficients of the characteristic polynomials of the products $A_w$ with $|w|\le j$.
\end{proposition}

\begin{proof}
It suffices to treat the universal case: $\mathsf A_0=\Z[a_{i,pq}]$, the polynomial ring in the
entries of $k$ generic matrices; the general case follows by the ring map $\mathsf A_0\to\mathsf A$
sending the generic entries to those of the $A_i$, applied coefficientwise. Since $\mathsf A_0$ is a
domain of characteristic $0$, it suffices to prove \eqref{pax:eq-amitsur} in $\mathsf A_0\otimes
\Q[[x]]$, where both sides have constant term $1$ and the logarithm $\log(1+f)=\sum_{m\ge1}
(-1)^{m-1}f^m/m$ ($f$ without constant term) is injective, with $\log(fg)=\log f+\log g$. For a matrix
$M$ with entries in the ideal $(x)$, $\log\det(1-M)=\operatorname{tr}\log(1-M)=-\sum_{m\ge1}\pgtr
(M^m)/m$ (replace $M$ by $YM$ with an auxiliary variable $Y$: both sides vanish at $Y=0$, and their
derivatives in $Y$ agree, by Jacobi's formula $\frac{d}{dY}\det(1-YM)=-\det(1-YM)\pgtr\bigl(M(1-YM)^{-1}
\bigr)$ and $\pgtr(M(1-YM)^{-1})=\sum_{m\ge1}Y^{m-1}\pgtr M^m$). With $M=\sum_ix_iA_i$ we get
\[
\log\det\Bigl(1-\sum_ix_iA_i\Bigr)=-\sum_{u}\frac1{|u|}\,x^u\,\pgtr A_u,
\]
the sum over all words. Group the words by their primitive root and cyclic rotation: every word is
$u=w_1^{\,d}$ with $w_1$ a rotation of a unique Lyndon word $w$, and $x^u=x^{w}{}^{d}$, $\pgtr A_u=
\pgtr\bigl(A_{w_1}^{\,d}\bigr)=\pgtr\bigl(A_w^{\,d}\bigr)$ by cyclicity of the trace (a rotation of
$w_1^{\,d}$ by one letter is $w_2^{\,d}$ with $w_2$ the rotation of $w_1$). For fixed $w$ and $d$ there
are exactly $|w|$ such words $u$ (the rotations of $w$ are distinct), each of length $d|w|$. Hence
\[
-\sum_u\frac{x^u\pgtr A_u}{|u|}=-\sum_{w\ \mathrm{Lyndon}}\ \sum_{d\ge1}|w|\cdot\frac{(x^w)^d\,
\pgtr(A_w^{\,d})}{d\,|w|}=\sum_{w\ \mathrm{Lyndon}}\log\det(1-x^wA_w),
\]
by the same identity applied to the single matrix $x^wA_w$. Rearrangements are legitimate because for
each monomial only finitely many terms contribute. Exponentiating gives \eqref{pax:eq-amitsur}.

For the consequence, substitute $x_1=\dots=x_k=T$ (a continuous ring map $\mathsf A[[x]]\to\mathsf
A[[T]]$): $\det(1-T\sum A_i)=\prod_w\det(1-T^{|w|}A_w)$, and the coefficient of $T^j$ involves only the
factors with $|w|\le j$, whose coefficients are, up to sign, the coefficients of the characteristic
polynomials of the $A_w$ (as $\det(1-YA)=Y^\sigma\det(Y^{-1}-A)$).
\end{proof}

\begin{remark}\label{pax:rem-amitsur}
(1) This is the formula of Amitsur~\cite{Amitsur}; the proof above is the standard one by
logarithms (it is the proof of Reutenauer and Sch\"utzenberger~\cite{ReutenauerSchutzenberger}). It is self-contained.
(2) \emph{Use in Lemma~\ref{pri:amitsur}.} There $\rho(P)=\sum_qE_q$, and the characteristic polynomial
of $\rho(P)\in M_{m\sigma}(\mathsf L)$ is $X^{m\sigma}\det(1-X^{-1}\rho(P))$. By
Proposition~\ref{pax:amitsur} with $\mathsf A=\mathsf L$, its coefficients are integral polynomials in
the coefficients of the characteristic polynomials of the products $E_{q_1}\cdots E_{q_k}$. Writing
$E_q=\epsilon_q\pi_q$ with $\epsilon_q\in M_{m\sigma\times\sigma}$ the $q$-th block column of $1$ and
$\pi_q\in M_{\sigma\times m\sigma}$ the $q$-th block row of $\rho(P)$, so that $\pi_q\epsilon_{q'}=
\rho(P_{qq'})$, the product is $\epsilon_{q_1}\rho(P_{q_1q_2})\cdots\rho(P_{q_{k-1}q_k})\pi_{q_k}$, and
since $\det(X-AB)=X^{a-b}\det(X-BA)$ for $A$ of size $a\times b$, $B$ of size $b\times a$, its
characteristic polynomial is a power of $X$ times that of $\rho(P_{q_1q_2}\cdots P_{q_{k-1}q_k}
P_{q_kq_1})$, whose coefficients lie in the image of $\mathsf R$ by hypothesis. As the image of
$\mathsf R$ is a subring, the proof of Lemma~\ref{pri:amitsur} is complete.
\end{remark}

\subsubsection{The chain lemma}\label{pax:ss-chain}

\emph{Why Proposition~\ref{pga:L} does not suffice.} Proposition~\ref{pga:L} computes the Hecke
polynomial of a theta lift; its roots contain the chain $\alpha_wq^l$, $l<t$. Lemma~\ref{pri:chain}
concerns a \emph{weak lift}, an eigenform whose normalised Galois representation has the eigenvalue
$1$ everywhere; such a form is not known to be a theta lift (this is the point of the false lifts in
Lemma~\ref{pri:weak}). Only the first root of the chain comes from Galois representations; the others
come from the unitarity of the local component, as follows.

Let $F$ be a non-archimedean local field with residue field of order $q$, uniformiser $\varpi$,
$K=\auGL_N(\auO_F)$, $\nu=|\det|$, and let $T_i$ ($0\le i\le N$) be the Hecke operator of $K\,
\mathrm{diag}(\varpi1_i,1_{N-i})K$ as in (N3) of {\cite[\S6.1]{P49}}. For an irreducible unramified
representation $\pi$ of $\auGL_N(F)$ with Satake parameters $\alpha_1,\dots,\alpha_N$ (that is, $\pi$
is the spherical subquotient of the normalised induction $\chi_1\times\dots\times\chi_N$ from
unramified characters with $\chi_k(\varpi)=\alpha_k$), the eigenvalues $t_i$ of the $T_i$ on
$\pi^K$ satisfy
\begin{equation}\label{pax:eq-satake}
H_\pi(X):=\sum_{i=0}^N(-1)^iq^{i(i-1)/2}t_iX^{N-i}=\prod_{k=1}^N\bigl(X-q^{(N-1)/2}\alpha_k\bigr),
\end{equation}
which is Satake's formula for $\auGL_N$, used in the same normalisation in the proof of
Proposition~\ref{pga:G123}.

\begin{lemma}[Chains in unitary spherical representations]\label{pax:tadic}
Let $\pi$ be an irreducible unitary unramified representation of $\auGL_N(F)$, $c$ a root of
$H_\pi$ and $k\ge1$ an integer with $|c|^2q^{k}=q^N$. Then $c,cq,\dots,cq^{k-1}$ are roots of
$H_\pi$, and $\prod_{i<k}(X-q^i)$ divides $H_\pi(cX)$.
\end{lemma}

\begin{proof}
This is Corollary~\ref{pay:roots}; see \S\ref{pay:ss-cons}, where it is proved without the
classification of the unitary dual. We also give the deduction from Tadi\'c's classification.

\emph{Input} (Tadi\'c's classification~\cite{TadicENS}, Theorem~D; for the
unramified case see the discussion below). Every irreducible unitary representation of $\auGL_N(F)$
is an irreducible product $\tau_1\times\dots\times\tau_s$ of representations of the following two
kinds: $u(\delta,m)$, the Langlands quotient of $\delta\nu^{(m-1)/2}\times\dots\times\delta\nu^{-(m-1)/2}$
for an irreducible unitary square-integrable $\delta$; and $u(\delta,m)\nu^a\times u(\delta,m)\nu^{-a}$
with $0<a<\frac12$.

\emph{Unramified constituents.} If $\pi=\tau_1\times\dots\times\tau_s$ (irreducible, normalised
parabolic induction from a standard Levi $M$) is unramified, so is each $\tau_i$: by the Iwasawa
decomposition $\auGL_N(F)=P K$, restriction of functions to $K$ identifies the $K$-fixed vectors of
$\operatorname{Ind}_P(\tau_1\otimes\dots\otimes\tau_s)$ with the $K\cap M$-fixed vectors of
$\tau_1\otimes\dots\otimes\tau_s$. Next, $u(\delta,m)$ with $\delta$ a representation of $\auGL_d(F)$ is
unramified only if $d=1$ and $\delta$ is an unramified unitary character $\chi$: an unramified
irreducible representation has supercuspidal support consisting of unramified characters of $F^\times$,
which forces $\delta$ to be a twist $\mathrm{St}_d\otimes\chi$ with $\chi$ unramified; and the
unramified irreducible representation with a given cuspidal support $\{\chi_1,\dots,\chi_N\}$ is the
Langlands quotient of $\chi_{1}\times\dots\times\chi_{N}$ ordered by decreasing exponents, whose
Langlands data consist of characters of $\auGL_1$ only, whereas the Langlands data of $u(\mathrm{St}_d
\otimes\chi,m)$ are $\mathrm{St}_d\chi\nu^{(m-1)/2},\dots$; by the uniqueness in the Langlands
classification, $d=1$. Finally $u(\chi,m)=\chi\circ\det$ on $\auGL_m(F)$, a subquotient of $\chi\nu^{(m-1)/2}\times
\dots\times\chi\nu^{-(m-1)/2}$. Since $\nu^s(\varpi)=q^{-s}$, the Satake parameters of $\pi$ are the
union, with multiplicity, of blocks of the two forms
\[
\text{(i)}\ \ \zeta\,q^{\,b-(m-1)/2}\ (0\le b<m),\qquad\text{(ii)}\ \ \zeta\,q^{\,b-(m-1)/2\pm a}\
(0\le b<m),\quad|\zeta|=1,\ 0<a<\tfrac12 .
\]

\emph{The chain.} Write $c=q^{(N-1)/2}\alpha$ with $\alpha$ a Satake parameter. Then $|\alpha|^2=
q^{1-k}$, i.e.\ $|\alpha|=q^{-(k-1)/2}$, with $-(k-1)/2\in\frac12\Z$. If $\alpha$ lay in a block of
type (ii), $\log_q|\alpha|=b-(m-1)/2\pm a$ would lie in $\frac12\Z\pm a$, which is disjoint from
$\frac12\Z$ as $0<a<\frac12$. So $\alpha=\zeta q^{b_0-(m-1)/2}$ lies in a block of type (i), with
$b_0-(m-1)/2=-(k-1)/2$. Then $0\le b_0=(m-k)/2$, so $k\le m$ and $k\equiv m\pmod2$, and for $0\le i<k$,
$\alpha q^i=\zeta q^{(b_0+i)-(m-1)/2}$ with $0\le b_0+i\le b_0+k-1=(m+k)/2-1\le m-1$. Hence
$\alpha q^i$ is a Satake parameter of the same block, and $cq^i$ is a root of $H_\pi$ by
\eqref{pax:eq-satake}. These $k$ roots are distinct ($q>1$, $c\ne0$), so $\prod_{i<k}(X-cq^i)$ divides
$H_\pi(X)$, and $\prod_{i<k}(X-q^i)$ divides $H_\pi(cX)$.
\end{proof}

\begin{remark}\label{pax:rem-tadic}
Lemma~\ref{pax:tadic} can be proved directly, without the classification: for a nonzero spherical
vector $x$ with Hecke polynomial $P$, a root $c$ with $qc$ not a root satisfies $|c|^2>q^{N-2}$ (an
identity obtained by pairing two explicit translates of $x$, computed by counting hyperplanes, at
pro-$p$ Iwahori level), and the roots are symmetric under $y\mapsto q^{N-1}/\bar y$ (the adjoint of
$T_i$ is $T_{N-i}$ up to the centre); an induction along the chain then gives the lemma. This
direct proof is written in \S\ref{pay:sec}; the paper uses this direct proof.
\end{remark}

\begin{proof}[Proof of Lemma~\ref{pri:chain}]
Let $y$ be a nonzero weak lift in base level (the forms of weight $\lambda$ of $O$ on $\auU(V)$, $n=
\operatorname{rk}V$, $r'=n_W$, $t'=n-r'\ge1$), $w\nmid\ell$ split and off the level, $q=\auNm w$,
$v$ the place of $E^+$ below $w$. \emph{Unitarity.} The space of all algebraic modular forms of weight
$\lambda$ (all levels) is a smooth admissible representation of $G_f$ by right translation, and the
inner product of {\cite[\S14.3]{P46}}, renormalised by the volume of the level, is $G_f$-invariant and
positive definite on it. Hence the $G_f$-representation generated by $y$ is a finite direct sum of
irreducible unitary representations $\pi=\otimes'_u\pi_u$; as $y$ is an eigenform of the $T_{w,i}$,
its projection to a constituent $\pi$ is nonzero and has the same eigenvalues, so $\pi_v\simeq\pi_w$
(under $\auU(V)(E^+_v)\simeq\auGL_n(E_w)$) is irreducible, unitary and unramified, and $H_w(y;X)=
H_{\pi_w}(X)$ in the notation~\eqref{pax:eq-satake}.

\emph{The first root.} By (G1) (Proposition~\ref{pga:G123}), $\det(X-S_y(\pgFrob_w))=\iota H_w(y;X)$.
The normalised representation $S_y\otimes\eta_O$ has the eigenvalue $1$ at $\pgFrob_w$, so $S_y(
\pgFrob_w)$ has the eigenvalue $\eta_O(\pgFrob_w)^{-1}=\iota X_0$, and $X_0$ is a root of $H_w(y;X)$.

\emph{Its absolute value.} By the normalisation of $\eta_O$ (\S\ref{pga:ss-L}: $\eta_O(\pgFrob_w)=
\alpha_w^{-1}$), $X_0=\alpha_w=\chi_{W,w}^{(r')}=\chi_W(\varpi_w/\overline\varpi_w)\upsilon(\varpi_w
)^{-r'}q^{r'/2}$. The character $\chi_W$ of the compact group $\auU(1)(E^+)\backslash\auU(1)(\mathbb A
_{E^+})$ is unitary; the splitting character $\upsilon$ restricts to a character of finite order of
$\mathbb A_{E^+}^\times$, and writing $\upsilon=\upsilon_0|\cdot|_{\mathbb A_E}^s$ with $\upsilon_0$
unitary, the restriction of $|\cdot|^s_{\mathbb A_E}$ to $\mathbb A_{E^+}^\times$ is $|\cdot|^{2s}$,
of finite order only if $s=0$; so $\upsilon$ is unitary. Hence $|X_0|^2=q^{r'}$, i.e.\
$|X_0|^2q^{t'}=q^n$.

Lemma~\ref{pax:tadic} with $N=n$, $c=X_0$, $k=t'$ gives that $\prod_{i<t'}(X-q^i)$ divides $H_w(y;
X_0X)$, which is Lemma~\ref{pri:chain} (for $X_0$ the value at a geometric Frobenius of $\iota^{-1}
\eta_O^{-1}$). The statement holds at every split $w\nmid\ell$ off the level at which (G1) holds,
i.e.\ (Proposition~\ref{pga:G123}) at every such $w$ at which $y$ is spherical.
\end{proof}

\begin{remark}\label{pax:rem-chain-sign}
Lemma~\ref{pri:chain} uses (G1) (Proposition~\ref{pga:G123}, which holds without any sign
hypothesis by Proposition~\ref{pax:sign}); nothing else from \S\ref{pga:sec} is used. If in the normalisation of
{\cite[\S14.3]{P46}} the operator $T_{w,i}$ is defined with $K\,d_i^{-1}K$ or with the action through
$\tau\circ c$ (see (N2)), the polynomial $H_w$ is replaced by the one with Satake parameters
$\bar\alpha_k^{-1}$ (the contragredient, equal to the complex conjugate for unitary $\pi$); the set of
Satake parameters of a unitary unramified representation is stable under $\alpha\mapsto\bar\alpha^{-1}$
(the blocks (i), (ii) are), and the argument is unchanged.
\end{remark}

\subsubsection{Inputs on forms}\label{pax:ss-forms}

Let $J_O$ be the base level of Definition~\ref{prm:def-level}, with its level character $\chi_O$. At a
chosen place $J_{O,w}$ is not congruent to $1$ modulo $w$ (its $\auGL_{t'}$-block is arbitrary), so
{\cite[Lemma~14.5]{P46}} does not apply to $J_O$ itself. Put
\[
J_0=J_O\cap\{k:\ k_w\equiv1\bmod w\ \text{at every chosen }w\},\qquad H=J_O/J_0\simeq\prod_{w\ \mathrm{chosen}}
J_{O,w}/(1+\varpi_wM_n(\auO_w)),
\]
a normal subgroup ($1+\varpi_wM_n(\auO_w)$ is normal in $\auGL_n(\auO_w)$) and a finite quotient, $H$
being a subgroup of $\prod_{w\ \mathrm{chosen}}\auGL_n(\auF_w)$. Since $J_{O,v}\subseteq K_v$ for every $v$ (Definition~\ref{prm:def-level}), $J_O\subseteq K^\circ$, so
 $J_0$ consists of elements congruent to $1$ modulo every chosen place and is free by
{\cite[Lemma~14.5]{P46}}; we call $J_0$ the \emph{free sublevel} and write $\chi=\chi_O|_{J_0}$, a character of
finite order, trivial at the chosen places. Let $w_1,\dots,w_s$ be
places of $E$, split over $E^+$, at which $J_0$ is hyperspecial, $J_{0,w_j}=\auGL_n(\auO_{w_j})$, and
$\chi$ trivial. Let $J_1\subseteq J_0$ be the subgroup of the elements congruent to $1$ modulo each
$w_j$ (principal level at the $w_j$); it is normal, and $J_0/J_1\simeq\Gamma=\prod_j\auGL_n(
\auF_{w_j})$. Then $\Gamma$ acts on $\mathcal A_\lambda(J_1,\chi)$ by $(k\varphi)(g)=\varphi(g\tilde k)$,
$\tilde k\in\prod_j\auGL_n(\auO_{w_j})\subseteq J_0$ a lift of $k$ with components $1$ outside the $w_j$ (independent of such a lift since $\chi$ is trivial on the components at
the $w_j$ and $J_1$ is normal). An \emph{integral operator} is a linear endomorphism preserving, for
every $\iota$, the forms integral at $\iota$.

The hypothesis $n\le2n_m$ of {\cite[Lemma~14.5]{P46}} holds: $n=r'+t'\le r+t=6r$, while
$n_m=r+mt+1$ with $m=2\varrho+2\ge2$, so $2n_m>2mt\ge20r>n$.

\begin{lemma}[Freeness of the principal level]\label{pax:free}
Let $\mathcal C_0=\Gamma_{\mathrm{gl}}\backslash G_f/J_0$ and $\mathcal C_1=
\Gamma_{\mathrm{gl}}\backslash G_f/J_1$, $\Gamma_{\mathrm{gl}}=\auU(V)(E^+)$. Choose for each class of
$\mathcal C_0$ a representative $g$. Then $(g,k)\mapsto\Gamma_{\mathrm{gl}}g\tilde kJ_1$ is a bijection
from $\{\text{representatives}\}\times\Gamma$ onto $\mathcal C_1$, the action of $h\in\Gamma$ on
$\mathcal C_1$ by right translation is $(g,k)\mapsto(g,kh)$, and $\varphi\mapsto\bigl((g,k)\mapsto
\varphi(g\tilde k)\bigr)$ identifies $\mathcal A_\lambda(J_1,\chi)$, $\Gamma$-equivariantly, with
$\operatorname{Map}(\{\text{representatives}\}\times\Gamma,W_\lambda)$, $\Gamma$ acting by right
translation on the second variable. A form is integral at $\iota$ if and only if its values
$\varphi(g\tilde k)$ have gauge at most $1$ at $(\iota,g)$.
\end{lemma}

\begin{proof}
Surjectivity is clear. If $\Gamma_{\mathrm{gl}}g\tilde kJ_1=\Gamma_{\mathrm{gl}}g\tilde k'J_1$, then
$\gamma g\tilde k=g\tilde k'j$ with $\gamma\in\Gamma_{\mathrm{gl}}$, $j\in J_1$, so $\gamma\in\Gamma_{
\mathrm{gl}}\cap gJ_0g^{-1}=1$ by {\cite[Lemma~14.5]{P46}}, hence $\tilde k'^{-1}\tilde k\in J_1$ and $k=k'$.
The same lemma shows that the defining relation $\varphi(\gamma xj)=\chi(j)\gamma\varphi(x)$ imposes
no condition on the values at the points $g\tilde k$ (each point of $\mathcal C_1$ is represented once
and has trivial stabiliser), which gives the identification. For the gauge: $\tilde k$ may be chosen
with components $1$ outside the $w_j$ and in $\auGL_n(\auO_{w_j})$ at the $w_j$, so it is integral
with integral inverse above $\ell$, and (G2) of {\cite[\S6.1]{P49}} gives that the gauge at
$(\iota,g\tilde k)$ equals that at $(\iota,g)$.
\end{proof}

\emph{Stabilisers of the base level.} For $g\in G_f$ the stabiliser $\mathrm{Stab}_g=\Gamma_{\mathrm{gl}}
\cap gJ_Og^{-1}$ maps to $H$ by $\gamma\mapsto g^{-1}\gamma g\bmod J_0$, with kernel $\Gamma_{\mathrm{gl}}
\cap gJ_0g^{-1}=1$ ({\cite[Lemma~14.5]{P46}}). So $\#\mathrm{Stab}_g$ divides $\#H$, and $\#H$ divides
$\prod_{w\ \mathrm{chosen}}\#\auGL_n(\auF_w)$, which divides $N=\pgNch=\prod_{w\ \mathrm{chosen}}
\#\auGL_{rt}(\auF_w)$ since $n=r'+t'\le r+t=6r\le5r^2=rt$ ($r'\le r$, $t'\le t=5r$, $r\ge2$) and $\auGL_n(\auF_w)$ embeds in $\auGL_{rt}(\auF_w)$
(Lagrange). We take $c_0=\#H$: a common multiple of the orders of the stabilisers of the base level,
dividing $N$.

\begin{proof}[Proof of Lemma~\ref{pri:forms}]
(a) Let $J_O^{(1)}=J_O\cap J_1$ be the base level with principal level at the $w_j$; the forms of
principal level at the $w_j$ are $\mathcal A_\lambda(J_O^{(1)},\chi_O)\subseteq\mathcal A_\lambda(J_1,\chi)$.
We first build an integral operator $\Psi_1$ on $\mathcal A_\lambda(J_1,\chi)$ with $\sum_{k\in\Gamma'}
k^{-1}\Psi_1k=1$, and then average over $H$. Let $\Gamma'\subseteq\Gamma$
be a subgroup (in particular a subgroup of a factor, or $\Gamma$ itself), and $S\subseteq\Gamma$ a set
of representatives of the left cosets $k\Gamma'$. In the model of Lemma~\ref{pax:free} let $\Psi$ be
multiplication by the indicator function of $\{\text{representatives}\}\times S$:
$(\Psi\varphi)(g,k)=\mathbf 1_S(k)\varphi(g,k)$. It is integral by the last assertion of
Lemma~\ref{pax:free}. For $h\in\Gamma'$, $(h^{-1}\Psi h\varphi)(g,k)=(\Psi h\varphi)(g,kh^{-1})=
\mathbf 1_S(kh^{-1})\varphi(g,k)$, so
\[
\Bigl(\sum_{h\in\Gamma'}h^{-1}\Psi h\,\varphi\Bigr)(g,k)=\#\{h\in\Gamma':kh^{-1}\in S\}\cdot
\varphi(g,k)=\#(S\cap k\Gamma')\,\varphi(g,k)=\varphi(g,k).
\]
Call this operator $\Psi_1$ (it is the $\Psi$ just defined, on $\mathcal A_\lambda(J_1,\chi)$). Now let
$e_H=\sum_{h\in H}R(\tilde h)$, $(R(\tilde h)\varphi)(g)=\varphi(g\tilde h)$, with $\tilde h\in J_O$ a lift
of $h$ with components $1$ outside the chosen places (well defined on $\mathcal A_\lambda(J_1,\chi)$ since
$J_0$ is normal in $J_O$ and $\chi_O$ is trivial at the chosen places). Each $R(\tilde h)$ is integral:
$\tilde h$ has integral components with integral inverses at places not above $\ell$, so (G2) of
{\cite[\S6.1]{P49}} applies as in the proof of Lemma~\ref{pax:free}. The $R(\tilde h)$ commute with the
action of $\Gamma$ (translations at disjoint sets of places), and on $\mathcal A_\lambda(J_O^{(1)},\chi_O)$
each acts by $\chi_O(\tilde h)=1$, so $e_H$ acts there by $\#H$; moreover $e_H$ maps
$\mathcal A_\lambda(J_1,\chi)$ into $\mathcal A_\lambda(J_O^{(1)},\chi_O)$. Put $\Psi=e_H\Psi_1$, restricted
to $\mathcal A_\lambda(J_O^{(1)},\chi_O)$: an integral operator on the forms of principal level at the
$w_j$, and
\[
\sum_{k\in\Gamma'}k^{-1}\Psi k=e_H\sum_{k\in\Gamma'}k^{-1}\Psi_1k=e_H=\#H=c_0
\quad\text{on }\mathcal A_\lambda(J_O^{(1)},\chi_O).
\]
For $\Gamma'=\Gamma$ and $y$ of base level (i.e.\ fixed by $\Gamma$), $c_0y=\sum_{k}k^{-1}\Psi ky=
\sum_kk^{-1}\Psi y$, a norm $\sum_{k\in\Gamma}k\,(\Psi y)$.

(b) Let $y'$ be integral at every isomorphism and $y''$ integral at $\iota$; both are forms of level
$J_O$. The pairing of {\cite[\S14.3]{P46}} at level $J_O$ is
\[
\langle y',y''\rangle=\sum_{g\in\Gamma_{\mathrm{gl}}\backslash G_f/J_O}\#\mathrm{Stab}_g^{-1}\,
\langle y'(g),y''(g)\rangle_{W_\lambda}.
\]
Each term $\iota\langle y'(g),y''(g)\rangle_{W_\lambda}$ is integral: this is the pointwise computation
in the proof of Proposition~\ref{pax:pairing}, which uses only that $y'$ is integral at $\bar\iota=
\iota\circ c$ and $y''$ at $\iota$, and not the freeness of the level. Since $\#\mathrm{Stab}_g$ divides
$N$ (see above), $N\langle y',y''\rangle=\sum_g(N/\#\mathrm{Stab}_g)\langle y'(g),y''(g)\rangle\in\pgZb$.
For completeness we recall the pointwise computation, with $\varphi=y'$, $\varphi'=y''$ and a fixed
$g$: after $\iota$, the vectors $\rho(g_\ell)^{-1}\iota\varphi'(g)$ and the corresponding
vectors for $\varphi$ through $\iota\circ c$ (whose action of $g_\ell$ is through the transposes of
the inverses of the components for $\iota$, as $g$ is unitary; (N2)) are integral, the invariant form
of $W_\lambda$ has an integral Gram matrix $\mathsf G$ in the basis $(e_T)$, and $\mathsf G$ is
invariant under the simultaneous action of $g_\ell$ on the two sides (through the transposed inverse
on the first), so the two factors $\rho(g_\ell)^{\pm1}$ cancel and each term is the value of
$\mathsf G$ on two integral vectors.
\end{proof}

\begin{remark}\label{pax:rem-forms}
At the free sublevel $J_0$ all stabilisers are trivial and the operator $\Psi_1$ gives the
constant $1$; at the base level $J_O$ the stabilisers $\Gamma_{\mathrm{gl}}\cap gJ_Og^{-1}$ may be
nontrivial (they embed in $H$), and averaging over $H$ gives $c_0=\#H$. All uses of $c_0$ in
\S\ref{pri:ss-norm} only need $c_0\mid N$. Likewise the pairing at level $J_O$ is not the pairing
at $J_0$ (for $J'\subseteq J$ one has $\langle\,,\rangle_{J'}=[J:J']\langle\,,\rangle_J$), which is why
(b) is proved directly at level $J_O$ with the factor $N$. The invariance of the Gram matrix
used at the end of (b) is proved in \S\ref{pax:ss-pairing} (Lemma~\ref{pax:gram}).
\end{remark}

\subsubsection{Block-diagonal lattices}\label{pax:ss-block}

For a field $L$ of characteristic $0$, a group $G$ and finite-dimensional representations $\rho_1,
\rho_2$ of $G$ over $L$: \emph{if $\rho_1,\rho_2$ are semisimple and $\pgtr\rho_1(g)=\pgtr\rho_2(g)$
for all $g\in G$, then $\rho_1\simeq\rho_2$} (the theorem of Brauer and Nesbitt in characteristic $0$; see also
Curtis--Reiner~\cite{CurtisReiner}). Proof: let $A$ be the image of $L[G]$ in
$\operatorname{End}_L(\rho_1\oplus\rho_2)$, a finite-dimensional $L$-algebra spanned by the operators of the
elements of $G$. It has a faithful semisimple module, $\rho_1\oplus\rho_2$, so its radical, which
annihilates every simple module, is $0$, and $A$ is semisimple. For a simple $A$-module $M$ let $e_M\in A$ be
the central idempotent that is the identity on $M$ and $0$ on the other simple modules. If $M$ occurs
$m_i$ times in $\rho_i$, the trace of $e_M$ on $\rho_i$ is $m_i\dim_LM$. Now $e_M$ is an $L$-linear
combination $\sum_gc_g(\rho_1\oplus\rho_2)(g)$, and its trace on $\rho_i$ is $\sum_gc_g\pgtr\rho_i(g)$,
the same for $i=1,2$. Hence $m_1\dim_LM=m_2\dim_LM$ in $L$, and $m_1=m_2$ because $\dim_LM\ne0$ in
characteristic $0$. As both are semisimple, $\rho_1\simeq\rho_2$.

\begin{lemma}[Exact form of Lemma~\ref{pri:block}]\label{pax:block}
Let $\rho_x$ ($x\in Y$ finite) be continuous \emph{semisimple} representations $\pgG\to\auGL_n(\auQlb)$
and, for $x\in\priX\subseteq Y$, $\pgtr\rho_x=\psi\,\pgtr R_x+\sum_{e<t'}\chi^e$, with $R_x$ of rank
$r'$, $n=r'+t'$, $\psi$ a continuous character and $\chi=\varepsilon_\ell^{-1}$. Let $u'_x$ be
lattices in $R_x$ over a coefficient ring $\auO_1$ such that, for every large $\auO\supseteq\auO_1$,
the $u'_x(g)$ span $M_{r'}(\auO)$ over $\auO$ (property (M4) of the main summands). Then for every large
$\auO$ there are lattices $u_x$ in the $\rho_x$ (continuous homomorphisms $\pgG\to\auGL_n(\auO)$
conjugate over $\auQlb$ to $\rho_x$) with $u_x=\psi u'_x\oplus\bigoplus_{e<t'}\auO(-e)$ for $x\in
\priX$, and the $\psi u'_x(g)$ span $M_{r'}(\auO)$.
\end{lemma}

\begin{proof}
Take $\auO$ large enough to contain the values of $\psi$ (a continuous character of the compact group
$\pgG$ takes values in the units of the ring of integers of a finite extension of $\Q_\ell$), to
satisfy the spanning hypothesis, and such that every $\rho_x$, $x\notin\priX$, has a lattice over
$\auO$ (a continuous representation of the compact group $\pgG$ over a finite extension $K'$ stabilises
a lattice, the $\pgG$-orbit of any lattice being finite; and each $\rho_x$ is defined over a finite
extension, by compactness and Baire's theorem; standard). For $x\notin\priX$ take such a lattice.

Let $x\in\priX$ and $\rho'_x=\psi u'_x\oplus\bigoplus_{e<t'}\auO(-e)$, a continuous homomorphism
$\pgG\to\auGL_n(\auO)$ ($\auO(-e)$ is the character $\varepsilon_\ell^{-e}=\chi^e$). \emph{Spanning.}
The $\psi(g)u'_x(g)$ span $M_{r'}(\auO)$: each is a unit multiple of $u'_x(g)$, and multiplying the
members of a spanning family by units does not change their span. \emph{Semisimplicity.} Over
$\auQlb$ the summand $\psi u'_x$ is absolutely irreducible, since the span of its image is
$M_{r'}(\auQlb)$ (Burnside), and the other summands are characters; so $\rho'_x\otimes\auQlb$ is
semisimple. \emph{Traces.} $\pgtr\rho'_x=\psi\pgtr u'_x+\sum_e\chi^e=\psi\pgtr R_x+\sum_e\chi^e=
\pgtr\rho_x$, as $u'_x$ is a lattice in $R_x$. By Brauer--Nesbitt in characteristic $0$, $\rho'_x
\otimes\auQlb\simeq\rho_x$, i.e.\ $u_x:=\rho'_x$ is a lattice in $\rho_x$ of the required shape.
\end{proof}

\begin{remark}\label{pax:rem-block}
Lemma~\ref{pri:block} has two hypotheses that the proof needs: the \emph{semisimplicity} of the
$\rho_x$ (in characteristic $0$ equivalent to the non-degeneracy of the trace form on the span of the
image) and the \emph{spanning} of $M_{r'}(\auO)$ by the $u'_x(g)$ (property (M4) of the main
summands, which is the input of the conclusion ``the $\psi u'_x(g)$ span''). Both hold where
Lemma~\ref{pri:block} is used (Proposition~\ref{pri:columns}): $\rho_x=S_x\otimes\eta_O$ is
semisimple, $S_x$ being a direct sum of the semisimple representations $r_\iota(\cdot)$ twisted by
characters (Proposition~\ref{pga:G123}), and $R_x$ is a main summand, so (M4) holds. Conjugate
self-duality is not used.
\end{remark}

\subsubsection{Integrality of the pairing at a free level}\label{pax:ss-pairing}

This completes the proof of Lemma~\ref{pri:forms}(b), which used the pointwise integrality of
$\langle\varphi(x),\varphi'(x)\rangle$ proved below. Let $G_\lambda$ be the Gram matrix of the invariant hermitian form of
$W_\lambda$ in the basis $(e_T)$; it is a positive definite \emph{integer} matrix
({\cite[\S5.2]{P53}}: $W_\lambda(\Z)$ is the lattice of integral tensors and the invariant form is the
standard form of the tensor space restricted to $W_\lambda$, {\cite[(38)]{P53}}). Write the form
as $\langle a,b\rangle={}^t\bar a\,G_\lambda\,b$ on coordinate vectors (the variable that is
conjugated is the one that must be integral at $\bar\iota$, as in {\cite[\S14.3]{P46}}).

\begin{lemma}[Invariance of the Gram matrix]\label{pax:gram}
For every field $L$ of characteristic $0$ and every $g\in\prod_\tau\auGL_N(L)$,
${}^t\rho_\lambda(g)\,G_\lambda\,\rho_\lambda({}^tg^{-1})=G_\lambda$, where $\rho_\lambda(g)=
\bigotimes_\tau(\det g_\tau)^{e_\tau}S_{\kappa_\tau}(g_\tau)$ is the matrix of (N1) of
{\cite[\S6.1]{P49}}.
\end{lemma}
\begin{proof}
This is the computation of (P1) in the proof of {\cite[Lemma~5.26]{P53}}: invariance of the hermitian
form under the compact group $\prod\auU(N)$ is ${}^t\overline{\rho_\lambda(g)}G_\lambda\rho_\lambda(g)=
G_\lambda$; the entries of $\rho_\lambda(g)$ are integral polynomials in the entries of $g$ and in
$\det g^{\pm1}$, so $\overline{\rho_\lambda(g)}=\rho_\lambda(\bar g)=\rho_\lambda({}^tg^{-1})$ for $g$
unitary, and with $h=\bar g$ the identity reads ${}^t\rho_\lambda(h)G_\lambda\rho_\lambda({}^th^{-1})=
G_\lambda$ on $\prod\auU(N)$. Both sides are polynomial in the entries of $h$ and $\det h^{-1}$ with
integer coefficients, and $\prod\auU(N)$ is Zariski dense in $\prod\auGL_N$; hence the identity holds
in the polynomial ring and therefore over every field of characteristic $0$.
\end{proof}

\begin{proposition}[Free-level pairing]\label{pax:pairing}
Let the level $J$ be free ($\Gamma_{\mathrm{gl}}\cap gJg^{-1}=1$ for all $g$). If $\varphi$ is integral at
$\bar\iota=\iota\circ c$ and $\varphi'$ at $\iota$, then $\iota\langle\varphi,\varphi'\rangle\in
\auO_{\C_\ell}$.
\end{proposition}
\begin{proof}
The computation of each term below is pointwise and does not use freeness. At a free level $\langle\varphi,\varphi'\rangle=\sum_x\langle\varphi(x),\varphi'(x)\rangle$ over
representatives $x$ of $\Gamma_{\mathrm{gl}}\backslash G_f/J$ (all weights $\#\mathrm{Stab}^{-1}$ equal $1$). Fix
$x$. By definition of the gauge (N2), $\iota(\varphi'(x))=\rho_\iota(x_\ell)\beta$ with $\beta$
integral, where $\rho_\iota(x_\ell)=\bigotimes_\tau\rho_\tau(\hat\iota_\tau(x_{w_\tau}))$. For
$\bar\iota$ the primes $w_\tau$ are replaced by their conjugates; since $x$ is unitary, its component
at $\overline{w}$ is ${}^t\bar x_w^{-1}$, so $\rho_{\bar\iota}(x_\ell)=\rho_\iota({}^tx_\ell^{-1})$ in
the sense of Lemma~\ref{pax:gram} applied factorwise (with $L=\C_\ell$), and $\iota(\overline{
\varphi(x)})=\bar\iota(\varphi(x))=\rho_\iota({}^tx_\ell^{-1})\alpha$ with $\alpha$ integral. Hence
\[
\iota\langle\varphi(x),\varphi'(x)\rangle={}^t\alpha\;{}^t\rho_\iota({}^tx_\ell^{-1})\,G_\lambda\,
\rho_\iota(x_\ell)\,\beta={}^t\alpha\,G_\lambda\,\beta
\]
by Lemma~\ref{pax:gram} with $g={}^tx_\ell^{-1}$, and the right side is integral since
$G_\lambda$, $\alpha$, $\beta$ are.
\end{proof}

With Proposition~\ref{pax:pairing}, the proof of Lemma~\ref{pri:forms}(b) in \S\ref{pax:ss-forms} is
complete.

\subsubsection{The dictionary at an auxiliary place}\label{pax:ss-dict}

We prove Proposition~\ref{pri:dict}(a)--(c), and (d), using the local Lemma~\ref{pax:toproot}
proved in \S\ref{pay:ss-cons}. Notation is that of Proposition~\ref{pri:dict}: $w$ a split degree-one
place, $w\nmid\ell$, off the base level, $q=\auNm w\equiv1\pmod{\ell^\nu}$, $F=E_w$, $\Gamma_w=
\auGL_n(\auF_w)$, $\gamma=\mathrm{Art}_w(z)$ the tame generator, $\phi$ a geometric Frobenius with
$\mathrm{Art}_w(\varpi)=\phi$ on the maximal abelian quotient. For a system $x$ of principal level
at $w$ write $\rho_x=S_x\otimes\eta_O$; the traces $\pgtr[\,\cdot\,]$ act on the forms of $x$ by the
scalars $\pgtr\rho_x(\cdot)$. Each identity below is an identity of operators on the forms of
principal level at $w$; since the $\Gamma_w$-action, the operators $U$ and the elements of $\priT$
commute with the projections to systems, it suffices to check it on forms $y$ in a single system
$x$, which we fix, together with an irreducible constituent $\pi$ of the automorphic representation
generated by the components of $y$; let $\Pi=\pi_w$, an irreducible unitary representation of
$\auGL_n(F)$ with nonzero vectors fixed by $K(1)=1+\varpi M_n(\auO_F)$.

\emph{Two facts used throughout.} (F1) $\eta_O$ is unramified at $w$ (it is the character of the
algebraic Hecke character of Proposition~\ref{pga:trace}, built from $\chi_W$ and $\upsilon$, which
are unramified at the places off the level of the pair; Lemma~\ref{pay:eta}), so
$\rho_x|_{I_w}=S_x|_{I_w}$. (F2) By
Proposition~\ref{pga:G123} (G1) and its proof, with Caraiani's part of (GR),
$\operatorname{WD}(S_x|_{G_{E_w}})^{\mathrm{F\text{-}ss}}\simeq\operatorname{rec}(\Pi\otimes|\det|^{(1-n)/2})$;
if the monodromy of the right side vanishes, $S_x|_{I_w}$ is the restriction of
$\operatorname{rec}(\Pi)$ to $I_w$ (a semisimple representation of finite image), and $S_x(\phi)$ has
the Frobenius eigenvalues of $\operatorname{rec}(\Pi\otimes|\det|^{(1-n)/2})$, the decomposition
under $D_w$ being respected. We use the sign normalisation $\epsilon=1$ of \S\ref{pga:ss-local}.

\begin{proof}[Proof of Proposition~\ref{pri:dict}(a)]
Let $\Gamma_w$ act on $y$ by $\chi\circ\mathrm{ind}\circ\det$, and let $\tilde\chi$ be a tamely
ramified character of $F^\times$ with $\tilde\chi(u)=\chi(\mathrm{ind}(\bar u))$ for $u\in
\auO_F^\times$. Then $\Pi\otimes(\tilde\chi^{-1}\circ\det)$ has a nonzero $\auGL_n(\auO_F)$-fixed
vector, so it is unramified, and $\operatorname{rec}(\Pi)=\operatorname{rec}(\Pi\otimes
\tilde\chi^{-1})\otimes\tilde\chi$ has monodromy $0$ and inertia acting by the scalar
$\tilde\chi\circ\mathrm{Art}_w^{-1}$. By (F1), (F2), $\rho_x(\gamma)=\tilde\chi(z)\cdot1=\chi(1)\cdot1$,
as $\mathrm{ind}(\bar z)=1$. Hence $\pgtr\rho_x((\gamma-1)g\,p(\phi))=(\chi(1)-1)\pgtr\rho_x(g\,p(\phi))$.
On the other side $\langle\gamma\rangle=\mathrm{diag}(\bar z,1,\dots,1)$ acts on $y$ by
$\chi(\mathrm{ind}\det\langle\gamma\rangle)=\chi(1)$. Both sides of (a) act on $y$ by
$(\chi(1)-1)\pgtr\rho_x(g\,p(\phi))$.
\end{proof}

For (b), let $\mathcal P\subseteq\auGL_n(\auO_F)$ be the stabiliser of the line $\auF_we_1$ modulo
$\varpi$ (block upper triangular of type $(1,n-1)$ modulo $\varpi$), $P=MN$ the upper parabolic of
type $(1,n-1)$, $d=\mathrm{diag}(\varpi,1,\dots,1)$, and $U$ the operator of $I_1dI_1$, $I_1$ the
pro-$p$ Iwahori subgroup (upper unipotent modulo $\varpi$).

\begin{lemma}[$U$ on parahoric vectors]\label{pax:Ueig}
Let $\tilde\chi$ be a tamely ramified character of $F^\times$, nontrivial on $\auO_F^\times$, and
suppose $\Pi$ has a nonzero vector $v$ on which $\mathcal P$ acts by $k\mapsto\tilde\chi(k_{11})$.
Then $\Pi\simeq\Pi_1\times\Pi_2$ with $\Pi_1$ a character of $F^\times$ restricting to $\tilde\chi$
on $\auO_F^\times$ and $\Pi_2$ unramified; the space of such vectors is one-dimensional; it is
contained in $\Pi^{I_1}$; and $Uv=q^{(n-1)/2}\Pi_1(\varpi)\,v$.
\end{lemma}
\begin{proof}
The first assertion is Lemma~\ref{pga:parahoric} with $(n_1,n_2)=(1,n-1)$ and $\psi=\tilde\chi|_{\auO_F^\times}$
(viewed on $\auF_w^\times$), applied to the vector $\pi(t)v$, where $t=\mathrm{diag}(1,\varpi^{-1}1_{n-1})\in
\auGL_n(F)$. Indeed $t\begin{pmatrix}a&b\\c&D\end{pmatrix}t^{-1}=\begin{pmatrix}a&\varpi b\\\varpi^{-1}c&D
\end{pmatrix}$, so $t\mathcal Pt^{-1}$ is exactly the parahoric of the matrices that are block lower
triangular modulo $\varpi$ for the blocks $(1,n-1)$, the entry $k_{11}$ is unchanged, and
$\pi(tkt^{-1})\pi(t)v=\pi(t)\pi(k)v=\tilde\chi(k_{11})\pi(t)v$ for $k\in\mathcal P$; the representation
$\Pi$ itself is not changed. Realise
$\Pi$ as the normalised induction $\operatorname{Ind}_P(\Pi_1\otimes\Pi_2)$, functions with
$f(mn g)=\delta_P^{1/2}(m)(\Pi_1\otimes\Pi_2)(m)f(g)$. By the Iwasawa decomposition $G=PK$ and
the Bruhat decomposition of $\auGL_n(\auF_w)$ for the parabolic of type $(1,n-1)$, $G=P\mathcal P
\sqcup Ps_1\mathcal P$, with $s_1$ the transposition of $e_1$ and $e_n$. A vector $f$ of type
$(\mathcal P,\tilde\chi(k_{11}))$ is determined by $f(1)$ and $f(s_1)$. On $P\mathcal P$: $f(1)$ is
fixed by $1\times\auGL_{n-1}(\auO_F)$ and $\auO_F^\times\times1$ acts on it by $\tilde\chi$, consistent
with $\Pi_1|_{\auO_F^\times}=\tilde\chi$; the possible values form a line, as $\Pi_2$ is unramified. On
$Ps_1\mathcal P$: for $k\in\mathcal P$ with $k'=s_1ks_1^{-1}\in M$, $f(s_1k)=f(k's_1)$ shows that
$f(s_1)\in\Pi_1\otimes\Pi_2$ is an eigenvector of $M\cap s_1\mathcal Ps_1^{-1}$ with the character
$k'\mapsto\tilde\chi(k'_{nn})$, up to the factor $\delta_P^{1/2}\Pi_1(k'_{11})$. The $\auGL_{n-1}$-part
of this group contains the parahoric $\mathcal Q\subseteq\auGL_{n-1}(\auO_F)$ of the matrices whose
reduction stabilises the line $\auF_w e_n$ (indeed $(s_1ks_1^{-1})_{jn}=k_{s_1(j),1}\equiv0$ for
$j\neq n$; these are the matrices that are block lower triangular of type $(n-2,1)$ modulo $\varpi$),
acting by $\tilde\chi$ of the
last diagonal entry. By Lemma~\ref{pga:parahoric} (for $\auGL_{n-1}$ with blocks $(1,n-2)$, applied to the image of
the vector under $\Pi_2(t'w'_0)$, where $w'_0$ is the antidiagonal permutation matrix of $\auGL_{n-1}$,
which reverses the indices and carries $\mathcal Q$ to the parahoric that is block upper triangular of type
$(1,n-2)$ modulo $\varpi$ with the character on the first diagonal entry, and $t'=\mathrm{diag}(1,
\varpi^{-1}1_{n-2})$, which carries the latter onto the block lower triangular parahoric of type
$(1,n-2)$ as in the first step; the vector $f(s_1)$, as a vector of $\Pi_2$ up to the scalar action of
the first factor, is replaced by its image under $\Pi_2(t'w'_0)$) a nonzero such vector would force a character with ramified restriction $\tilde\chi$
in the supercuspidal support of $\Pi_2$, which consists of unramified characters. So $f(s_1)=0$, and
the space is one-dimensional, spanned by $f_0$ supported on $P\mathcal P$. As $I_1
\subseteq\mathcal P$ and $\tilde\chi(k_{11})=1$ on $I_1$, $f_0\in\Pi^{I_1}$. Now $I_1=(I_1\cap\bar N)
(I_1\cap M)(I_1\cap N)$ with $I_1\cap N=N(\auO_F)$, $d^{-1}(I_1\cap\bar N)d\subseteq I_1\cap\bar N$
and $d$ normalising $I_1\cap M$, so $I_1dI_1=\bigsqcup_{u\in N(\auO_F)/dN(\auO_F)d^{-1}}u\,d\,I_1$, with
$q^{n-1}$ cosets. $U$ preserves the vectors of type $(\mathcal P,\tilde\chi(k_{11}))$ in $\Pi^{I_1}$:
$\mathcal P=(\mathcal P\cap\bar N)(\mathcal P\cap M)(\mathcal P\cap N)$, where $\mathcal P\cap\bar N=
\bar N(\varpi\auO_F)\subseteq I_1$ fixes the $I_1$-fixed vector $Uf$ and has trivial character;
$\mathcal P\cap N=N(\auO_F)$ permutes the cosets $ud\,I_1$ and has trivial character; and the elements
$m$ of $\mathcal P\cap M=\auO_F^\times\times\auGL_{n-1}(\auO_F)$ commute with $d$ and normalise $N(\auO_F)$
and $dN(\auO_F)d^{-1}$, so $\pi(m)Uf=\sum_u\pi(mum^{-1}d)\pi(m)f=\tilde\chi(m_{11})Uf$. Hence
$Uf_0=cf_0$ with $c=(Uf_0)(1)/f_0(1)$, and
\[
(Uf_0)(1)=\sum_uf_0(ud)=q^{n-1}f_0(d)=q^{n-1}\delta_P^{1/2}(d)\,\Pi_1(\varpi)f_0(1)=
q^{(n-1)/2}\Pi_1(\varpi)f_0(1),
\]
using $f_0(ud)=f_0(d)$ ($u\in N$) and $\delta_P(d)=|\varpi|^{n-1}=q^{1-n}$.
\end{proof}

\begin{proof}[Proof of Proposition~\ref{pri:dict}(b)]
Let $\mathcal P$ act on $y$ by $\chi(\mathrm{ind}(k_{11}))$ with $\chi(n!)\ne1$; in particular
$\chi\ne1$, and $\tilde\chi$ (as in (a)) is ramified. By Lemma~\ref{pax:Ueig}, $\Pi=\Pi_1\times\Pi_2$,
so $\operatorname{rec}(\Pi\otimes|\det|^{(1-n)/2})=\lambda\oplus r_2$ with $\lambda=
\operatorname{rec}(\Pi_1|\cdot|^{(1-n)/2})$ a character, $\lambda|_{I_w}=\tilde\chi\circ
\mathrm{Art}_w^{-1}\ne1$, $\lambda(\phi)=q^{(n-1)/2}\Pi_1(\varpi)$, and $r_2$ unramified with
monodromy $0$. By (F1), (F2), $\rho_x|_{G_{E_w}}=L\oplus L'$ with $L$ a line on which $D_w$ acts by
$\lambda\eta_O$ and $I_w$ acting trivially on $L'$. Let $e$ be the projection onto $L$ along $L'$;
it commutes with $\rho_x(\phi)$, and $\rho_x(\gamma)-1=(\lambda(\gamma)-1)e$. With $\mu=
\lambda(\phi)\eta_O(\phi)$,
\[
\pgtr\rho_x\bigl((\gamma-1)g\,p(\phi)\bigr)=(\lambda(\gamma)-1)\pgtr\bigl(p(\rho_x(\phi))e\rho_x(g)
\bigr)=p(\mu)\,(\lambda(\gamma)-1)\pgtr(e\rho_x(g))=p(\mu)\pgtr\rho_x((\gamma-1)g).
\]
The same computation holds for every constituent $\pi$ meeting $y$, with the same $\lambda$ (it is
read from $\rho_x$: the unique ramified line of $\rho_x|_{D_w}$). By Lemma~\ref{pax:Ueig}, applied
to the components of $y$ (which lie in the one-dimensional parahoric-type spaces of the constituents),
$Uy=q^{(n-1)/2}\Pi_1(\varpi)y=\lambda(\phi)y$, so $\tilde Uy=\eta_O(\phi)Uy=\mu y$ and $p(\tilde U)y=
p(\mu)y$. Hence both sides of (b) act on $y$ by $p(\mu)\pgtr\rho_x((\gamma-1)g)$.
\end{proof}

\begin{proof}[Proof of Proposition~\ref{pri:dict}(c)]
For any form $z$ of principal level at $w$, $\mathrm{Nm}_wz=\sum_{k\in\Gamma_w}kz$ is fixed by
$\Gamma_w$, i.e.\ of base level at $w$, hence spherical at $w$; its components in the systems $x$
are spherical at $w$, so $S_x$ is unramified at $w$ by (G2), and with (F1) $\rho_x(\gamma)=1$, so
$\pgtr\rho_x((\gamma-1)g)=0$ for each system $x$ meeting $\mathrm{Nm}_wz$. As $\pgtr[(\gamma-1)g]$
commutes with $\Gamma_w$, $\mathrm{Nm}_w\pgtr[(\gamma-1)g]z=\pgtr[(\gamma-1)g]\mathrm{Nm}_wz=0$.
\end{proof}

\emph{Part (d).} Let $\delta\in\priT$, $p=(X-q^{t'-1})H$, $B(1)=p(\tilde U)L(1)$ with $L(1)=
\sum_{k\in\mathcal P/K(1)}k$ the plain norm over the line stabiliser (weight $1$). The
deduction of (d) uses the following local statement.

\begin{lemma}[The top root of a chain]\label{pax:toproot}
Let $\pi$ be an irreducible unitary unramified representation of $\auGL_n(F)$ with invariant inner
product $(\ ,\ )$, $x\ne0$ spherical, $X_0\ne0$, $k\ge1$, and suppose $H_\pi(X_0X)=\prod_{i<k}(X-q^i)
\cdot S(X)$ with $S\mid R$. Then $(x,\ ((X-q^{k-1})R)(X_0^{-1}U)\,z)=0$ for every $z\in\pi^{I_1}$.
\end{lemma}
\noindent Proved in \S\ref{pay:ss-cons}, from the spherical chain lemma,
Theorem~\ref{pay:chain}(ii).\par\smallskip

\begin{proof}[Proof of Proposition~\ref{pri:dict}(d)]
Put $Z=\mathrm{Nm}\,B(1)$ ($\mathrm{Nm}=\sum_{k\in\Gamma}k$); its values are $\Gamma$-fixed, i.e.\ of
base level. Let $u$ be a form of base level in a system $x$ with $\delta(x)\ne0$; by hypothesis
$H_w(x;X_0X)=\prod_{i<t'}(X-q^i)\cdot S$ with $S\mid H$, $X_0=\eta_O(\phi)^{-1}$, so $\tilde U=X_0^{-1}U$.
For every form $y$, $L(1)y$ is fixed by $I_1$ (as $I_1\subseteq\mathcal P$), and since the
components of $u$ in the constituents are spherical, Lemma~\ref{pax:toproot} (with $k=t'$, $R=H$),
applied in each constituent with the restriction of the global inner product (which is
$G_f$-invariant), gives $\langle u,B(1)y\rangle=0$. As $u$ is $\Gamma$-fixed and the inner product
$\Gamma$-invariant, $\langle u,Zy\rangle=\sum_k\langle k^{-1}u,B(1)y\rangle=\#\Gamma\,\langle u,B(1)y
\rangle=0$. Now $\delta Zy$ is of base level and lies in the span of the systems with $\delta\ne0$;
by normality of $\delta$ (adjoint in $\priT$, preserving the systems) $\langle\delta Zy,\delta Zy\rangle
=\langle\delta^*\delta Zy,Zy\rangle$, and $\delta^*\delta Zy$ is a sum of base-level forms in systems
with $\delta\ne0$, orthogonal to $Zy$ by the above. So $\delta Zy=0$.
\end{proof}

\subsubsection{Base level of a label}\label{pax:ss-BL}

We prove Proposition~\ref{pri:BL}, with the following precise meaning of ``base level'': the forms
of the weight $\lambda$ of $O$ on $\auU(V)$ of the level $J_O$ and level character $\chi_O$ of
Definition~\ref{prm:def-level} (at a chosen $w$, $k\equiv\left(\begin{smallmatrix}1_{r'}&*\\0&*
\end{smallmatrix}\right)$ modulo $w$ with trivial character; at $w\in Q$, $k$ block lower triangular
modulo $w$ with character $\psi_w^{-1}(\det k^{(11)})$; hyperspecial elsewhere), that are fixed by the
idempotent $e_c$ of the labels. Let $G_c=\prod_{w\ \mathrm{chosen}}\auGL_{r'}(\auF_w)$, acting on
these forms by $\varrho$ of Definition~\ref{prm:def-level} (lifts reducing to $\mathrm{diag}(g_w,
1_{t'})$), twisted by $\prod_w\psi_w(\det g_w)$ ($\psi_w$
the tame part at $w$), and $e_c\in\C[G_c]$ the central idempotent of the irreducible character named
by the labels.

\begin{proof}[Proof of Proposition~\ref{pri:BL}]
(a) holds by definition.

(b) $e_c=\sum_{g\in G_c}a_gg$ with $a_g=\dim(\sigma)\,\overline{\operatorname{ch}_\sigma(g)}/\#G_c$
(the central idempotent of the irreducible character $\sigma$), so $\#G_c\cdot a_g$ is an algebraic
integer. The twisting factors are roots of unity. The lift of $g$ supported at the chosen places has
integral components with integral inverses, so its translation preserves forms integral at $\iota$
((G2) of {\cite[\S6.1]{P49}}). Hence $\#G_c\,e_cy$ is integral when $y$ is, and so is $Ne_cy$, as $\#G_c$
divides $N=\pgNch=\prod_{w\ \mathrm{chosen}}\#\auGL_{rt}(\auF_w)$ (by Lagrange's theorem, since
$G_c=\prod_w\auGL_{r'}(\auF_w)$ and $\auGL_{r'}(\auF_w)$ embeds in $\auGL_{rt}(\auF_w)$ as
$r'\le r\le rt$).

(c) Let $x$ be a system occurring in base level with principal level at $U$, $y\ne0$ a form of $x$,
and $S_x=S_y$ as in Proposition~\ref{pga:G123}; put $\rho_x=S_x\otimes\eta_O$.
\emph{Conjugate self-duality.} By (GR), $S_x^c\simeq S_x^\vee\varepsilon_\ell^{1-n}$ (each summand
$r_\iota(\Sigma)\varepsilon_\ell^{e}$ of Proposition~\ref{pga:G123} is polarised in this way, with the
twist of Proposition~\ref{pax:sign}). The algebraic Hecke character
$\eta$ with $\ell$-adic character $\eta_O$ (Proposition~\ref{pga:trace}) satisfies $\eta\eta^c=
|\cdot|_{\mathbb A_E}^{\,r'}$: the factors $\chi_W(y/\bar y)$ cancel, $\upsilon(y)\upsilon(\bar y)=
\upsilon|_{\mathbb A_{E^+}}(\auNm y)=1$ ($\upsilon|_{\mathbb A_{E^+}^\times}$ is the quadratic character
of $E/E^+$), and $|y|^{-r'/2}|\bar y|^{-r'/2}=|y|^{-r'}$, inverted. So $\eta_O\eta_O^c=
\varepsilon_\ell^{r'}$ and $\rho_x^c\simeq\rho_x^\vee\varepsilon_\ell^{1-n+r'}=\rho_x^\vee\chi^{t'-1}$,
which is conjugate self-duality with string length $t'$ (the form (M1)).
\emph{Hodge--Tate data.} By (G3) the weights of $S_x$ are determined by $\lambda$, and those of
$\eta_O$ by the pair; let $J'$ be the resulting data (independent of $x$).
\emph{Ramification.} $y$ is fixed by a hyperspecial subgroup at every place off $\ell p_0p_1$,
$Q\cup\overline Q$, $U\cup\overline U$ (the level $J_O$ is hyperspecial there and $y$ has principal
level only at $U$), so $S_x$ is unramified there by (G2); $\eta_O$ is unramified off $\ell$ and the
places of the pair ((F1) of \S\ref{pax:ss-dict}; Lemma~\ref{pay:eta}).
\emph{Inertia at a chosen $w$.} By Lemma~\ref{prm:label}(b), $\Pi_w^{K(1)}$ contains an irreducible
representation of $\auGL_n(\auF_w)$ with cuspidal support $(L,\bigotimes_k\sigma_{\tilde f_k}\otimes
\mathbf 1^{\otimes t'})$, $\tilde c_w=\prod_k\tilde f_k$ the rescaled label. By the cuspidal-support
statement in the proof of
Lemma~\ref{pga:depth0} (Moy--Prasad), the supercuspidal support of $\Pi_w$ consists of the depth-zero
supercuspidals attached to the irreducible factors of $\tilde c_w$ and of $t'$ unramified characters;
so $\operatorname{rec}(\Pi_w)|_{I_w}$ (semisimplification) has the eigenvalues $\omega(\xi)$, $\xi$ the
roots of $\tilde c_w$, and $1$ ($t'$ times), at a tame generator. The characteristic polynomial of
$S_x(s)$, $s\in I_w$, equals that of $\operatorname{rec}(\Pi_w)(s)$ (by (GR) the Weil--Deligne
representations agree up to Frobenius semisimplification, and the nilpotent monodromy commutes with
$I_w$, so it does not change characteristic polynomials of inertia). As $\eta_O$ is unramified at $w$ (Lemma~\ref{pay:eta}), inertia at $w$ acts on $\rho_x$ through the roots of the rescaled
label times $(X-1)^{t'}$.
\emph{No monodromy at a chosen $w$.} We show that $\operatorname{rec}(\Pi_w)$ has monodromy $0$;
here $\pi$ is an irreducible constituent of the automorphic representation generated by $y$ on
which the orthogonal projection $y_\pi$ of $y$ is non-zero, $S_x=S_y$ is formed with this $\pi$ in
Proposition~\ref{pga:G123} (up to isomorphism $S_y$ depends only on $x$: it is semisimple and its
Frobenius characteristic polynomials at almost all split places are given by the Hecke eigenvalues
through (G1), so Chebotarev and Brauer--Nesbitt apply), and $\Pi_w=\pi_v$ is the component at $w$
((BC) at split places). Notation: $G_w=\auGL_n(\auF_w)$, $P=MU_P\subseteq G_w$ the block upper
triangular parabolic subgroup of type $(r',t')$ ($r',t'\ge1$, both being odd), $\sigma'=\sigma_{c_w}
\otimes(\psi_w^{-1}\circ\det)=\sigma_{\tilde c_w}$, as in Lemma~\ref{prm:label}.

\emph{Step 1: a vector of type $\sigma'\boxtimes\mathbf 1$.} By the proof of Lemma~\ref{prm:label}(a),
$y$ is fixed by $U_P$ and by $1\times\auGL_{t'}(\auF_w)$ and lies in the $\sigma'$-isotypic part of
$V_w$ for $g\mapsto\mathrm{diag}(g,1)$. The orthogonal projection onto the $\pi$-part commutes with
right translations, so $y_\pi$ has the same three properties, and by the proof of
Lemma~\ref{prm:label}(b) the $\pi$-part of $V_w$ is $\Pi_w^{K_w(1)}\otimes M_\pi$ with $G_w$ acting on
the first factor only. Writing $y_\pi=\sum_iv_i\otimes m_i$ with the $m_i$ linearly independent, each
$v_i$ has the three properties, and some $v_i\ne0$. So $\Pi_w^{K_w(1)}$ contains $v\ne0$ fixed by
$U_P$ and by $1\times\auGL_{t'}(\auF_w)$, in the $\sigma'$-isotypic part for $\auGL_{r'}(\auF_w)\times1$.

\emph{Step 2: the Jacquet module.} Let $\mathcal Q\subseteq\auGL_n(\auO_w)$ be the inverse image of
$P$, $\mathcal Q_1$ the inverse image of $U_P$, $P_w=M_wN_w$ the block upper triangular parabolic
subgroup of $\auGL_n(E_w)$ of type $(r',t')$, and $M_1=\mathcal Q_1\cap M_w=(1+\varpi_wM_{r'}(\auO_w))
\times(1+\varpi_wM_{t'}(\auO_w))$. Statement (J) of the proof of Lemma~\ref{pga:parahoric}
(Moy--Prasad~\cite[Proposition~6.7]{MoyPrasad}) holds for $\mathcal Q$, with the same verification of
its hypotheses: $\mathcal Q$ is a non-maximal parahoric subgroup $\mathcal G_x$ for a point $x$ of
the apartment of the diagonal torus, with pro-unipotent radical $\mathcal Q_1$ and attached Levi
subgroup $M_w$, and either of the two opposite parabolic subgroups with Levi $M_w$ may be used. So
the projection $\Pi_w\to\Pi_{w,N_w}$ restricts to an isomorphism $\Pi_w^{\mathcal Q_1}\to(\Pi_{w,N_w})^{
M_1}$, equivariant for $\mathcal Q/\mathcal Q_1=M\simeq\auGL_{r'}(\auF_w)\times\auGL_{t'}(\auF_w)$. Now
$v\in(\Pi_w^{K_w(1)})^{U_P}=\Pi_w^{\mathcal Q_1}$ lies in the $(\sigma'\boxtimes\mathbf 1)$-isotypic part
for $M$. Hence $(\Pi_{w,N_w})^{M_1}$ has a non-zero $(\sigma'\boxtimes\mathbf 1)$-isotypic part. Taking
isotypic parts for the compact group $M_w\cap\auGL_n(\auO_w)$ is exact and $\Pi_{w,N_w}$ has finite
length, so some irreducible subquotient $\tau=\tau_1\otimes\tau_2$ of $\Pi_{w,N_w}$ has this property:
$\tau_1^{K(1)}$ contains $\sigma'$, and $\tau_2^{K(1)}$ contains the trivial representation of
$\auGL_{t'}(\auF_w)$, that is, $\tau_2$ has a non-zero vector fixed by $\auGL_{t'}(\auO_w)$. (The
normalised Jacquet module differs from $\Pi_{w,N_w}$ by the unramified character
$\delta_{P_w}^{-1/2}$ of $M_w$, which changes neither property.)

\emph{Step 3: the supercuspidal support.} $\sigma'$ has cuspidal support $(L_1,\bigotimes_k\lambda_k)$,
$\lambda_k=\sigma_{\tilde f_k}$ (paragraph ``Twisted labels'' of \S\ref{prm:ss-level}); by the
cuspidal-support statement in the proof of Lemma~\ref{pga:depth0} (for $\auGL_{r'}$), $\tau_1$ has
supercuspidal support $\{\pi_1,\dots,\pi_s\}$, $\pi_k$ a depth-zero supercuspidal representation of
$\auGL_{d_k}(E_w)$ containing $\sigma_{\tilde f_k}$. The trivial representation of $\auGL_{t'}(\auF_w)$
has cuspidal support $(T,\mathbf 1)$ (it is a constituent of $R_T^{\auGL_{t'}}(\mathbf 1)$), so by the
same statement (for $\auGL_{t'}$, all $f_k=X-1$) $\tau_2$ has supercuspidal support $t'$ characters
$\chi_1,\dots,\chi_{t'}$ of $E_w^\times$ trivial on $\auO_w^\times$, i.e.\ unramified. The supercuspidal
support of $\Pi_w$ is that of $\tau$ (Bernstein--Zelevinsky, as in the proof of
Lemma~\ref{pga:parahoric}): $\{\pi_1,\dots,\pi_s,\chi_1,\dots,\chi_{t'}\}$.

\emph{Step 4: no linking, and the shape of $\Pi_w$.} A segment of $\pi_k$ can be linked to one of
$\chi_j$ only if $\pi_k\simeq\chi_j|\cdot|^u$ for some $u$; then $d_k=1$ and $\pi_k$ is trivial on
$\auO_w^\times$. But for $d_k=1$ the restriction of $\pi_k$ to $\auO_w^\times$ is the inflation of the
character $\sigma_{\tilde f_k}$ of $\auF_w^\times$, $\tilde f_k=X-\xi$, which is trivial only if
$\xi=1$; and no rescaled label has the root $1$ (lift setting, \S\ref{pri:sec}). Moreover two
$\pi_k$ are never inertially equivalent, since $\tilde c_w$ is squarefree (it is a label; proof of
Lemma~\ref{pga:depth0}, paragraph on depth-zero supercuspidals with a given $\sigma_f$). Hence, by
Zelevinsky's classification exactly as in the proof of Lemma~\ref{pga:parahoric},
$\Pi_w\simeq\Pi_1\times\Pi_2$ with $\Pi_1$, $\Pi_2$ irreducible of supercuspidal supports
$\{\pi_k\}$ and $\{\chi_j\}$; and $\Pi_1\simeq\pi_1\times\dots\times\pi_s$ by Zelevinsky's
irreducibility criterion~\cite[Theorem~4.2]{Zelevinsky}, as in the proof of Lemma~\ref{pga:depth0}.
By the geometric lemma of Bernstein--Zelevinsky, the composition factors of the Jacquet module of
$\Pi_1\times\Pi_2$ for $P_w$ are subquotients of $(\alpha_1\times\alpha_2)\otimes(\beta_1\times\beta_2)$,
where $\alpha_i\otimes\beta_i$ is a composition factor of the Jacquet module of $\Pi_i$ for a parabolic
subgroup of type $(a_i,b_i)$, $a_1+b_1=r'$, $a_2+b_2=t'$, $a_1+a_2=r'$, $b_1+b_2=t'$. The supercuspidal
support of $\tau_2$ consists of unramified characters, while that of $\beta_1$ is contained in
$\{\pi_k\}$, which contains none; so $b_1=0$, hence $a_1=r'$, $a_2=0$, $b_2=t'$, and the only such
factor is $\Pi_1\otimes\Pi_2$. Thus $\tau_2$ is $\Pi_2$ up to an unramified twist, and $\Pi_2$ has a
non-zero vector fixed by $\auGL_{t'}(\auO_w)$: it is spherical. So
\[
\Pi_w\simeq\pi_1\times\dots\times\pi_s\times\Pi_2,\qquad\Pi_2\ \text{an irreducible spherical
representation of }\auGL_{t'}(E_w);
\]
in particular no constituent of Steinberg type occurs in the unramified part.

\emph{Step 5: the parameter.} $\operatorname{rec}(\Pi_w)=\bigoplus_k\operatorname{rec}(\pi_k)\oplus
\operatorname{rec}(\Pi_2)$; each $\operatorname{rec}(\pi_k)$ is irreducible, with monodromy $0$, and
$\operatorname{rec}(\Pi_2)$ is unramified with monodromy $0$ ($\Pi_2$ spherical, as in the proof of
Lemma~\ref{pga:parahoric}). So $\operatorname{rec}(\Pi_w)$ has monodromy $0$, and by the last
assertion of Proposition~\ref{pga:G123} (which rests on the identity of (GR) including monodromy,
\cite[Thm.~2.1.1(2)]{BLGGT}, at $w\nmid\ell$) the monodromy of $\operatorname{WD}(S_x|_{G_{E_w}})$ vanishes. Then
$S_x|_{I_w}$ is the restriction to $I_w$ of the Weil group representation underlying this
Weil--Deligne representation, which is trivial on an open subgroup of $I_w$; so $S_x(s)$ has finite
order and is semisimple for $s\in I_w$, and so is $\rho_x(s)$ ($\eta_O$ is unramified at $w$). By the
previous paragraph its eigenvalues are roots of $o_w(X)(X-1)$, where $o_w(X)=\prod_\xi(X-\omega(\xi)^
\epsilon)$, $\xi$ running over the roots of $\tilde c_w$; hence
\[
o_w(\rho_x(s))\,(\rho_x(s)-1)=0\qquad(s\in I_w).
\]
\emph{Inertia above $Q\cup\overline Q$.} At $w\in Q$ the form $y$ is an eigenvector of the parahoric
of type $(n_W,t')$ with the character $\psi_{\mathcal P}$, $\psi=\psi_w\ne1$; by (G5)
(Proposition~\ref{pga:G45}), $(S_x(s)-\psi(s))(S_x(s)-1)=0$ for $s\in I_w$, and $\eta_O$ is unramified at
$w$. At $\bar w$ the same holds by conjugate self-duality, $\psi$ being anticyclotomic. This is (c).
\end{proof}

\begin{remark}\label{pax:rem-BL}
So Proposition~\ref{pri:BL} follows from Propositions~\ref{pga:G123}, \ref{pga:G45}, the
cuspidal-support part of Lemma~\ref{pga:depth0} and the definitions; {\cite[Lemma~14.5]{P46}} is not
needed for (a)--(c) (it enters through Lemma~\ref{pri:forms}). That $\eta_O$ is unramified at
$Q\cup\overline Q$ and at the places $U$ (indeed at every finite place not above $\ell$) is
Lemma~\ref{pay:eta}. The trivial block $\mathbf 1_{\auGL_{t'}}$ comes from the equivariance of the
single lift, Lemma~\ref{pay:product}; the property of the base level at the chosen places is
Lemma~\ref{prm:label} (Definition~\ref{prm:def-level}). The label at $w$ is not squarefree (the factor $(X-1)^{t'}$), so
Lemma~\ref{pga:depth0} itself does not apply to $\Pi_w$; it is used for the cuspidal support only.
The vanishing of the monodromy at a chosen $w$ is proved in the paragraph ``No monodromy at a chosen
$w$'' from that cuspidal support, statement (J) of the proof of Lemma~\ref{pga:parahoric}, Zelevinsky's
classification and the geometric lemma; it uses that no rescaled label has the root $1$. Without
that clause of the lift setting a factor $X-1$ of $\tilde c_w$ would contribute an unramified
character linked possibly to the $\chi_j$, and Steinberg-type monodromy could not be excluded.
\end{remark}

\subsubsection{The sign hypothesis of Proposition~\ref{pga:G123}}\label{pax:ss-sign}

\begin{proposition}\label{pax:sign}
In Proposition~\ref{pga:G123} the hypothesis of (GR) holds for every constituent, after replacing the
twist $|\det|^{(m_i-1)/2}$ by $|\det|^{s_i}$ with $s_i\in\frac12\Z$, $s_i\equiv(n_i-1-n+m_i)/2
\pmod1$; with this twist,
\[
S_y=\bigoplus_i\bigoplus_{a=0}^{m_i-1}r_\iota\bigl(\sigma_i|\det|^{s_i}\bigr)\otimes
\varepsilon_\ell^{\,e_{i,a}},\qquad e_{i,a}=\tfrac{m_i-1}2-a-s_i+\tfrac{n_i-n}2\in\Z ,
\]
and the conclusions of Proposition~\ref{pga:G123} hold without any sign hypothesis.
\end{proposition}
\begin{proof}
Fix $i$ and write $\sigma=\sigma_i$, $m=m_i$, $N=n_i$, $\Sigma=\sigma|\det|^{s}$, $s=s_i$.
\emph{Algebraicity.} $\Pi$ is cohomological for $\auGL_n$, so its infinitesimal character at each
$\tau$ has entries in $\Z+\frac{n-1}2$. The entries for $\operatorname{Speh}(\sigma,m)$ are those of
$\sigma$ shifted by $\frac{m-1}2-a$, $0\le a<m$; so those of $\sigma$ lie in $\Z+\frac{n-m}2$, and those
of $\Sigma$ in $\Z+\frac{n-m}2+s=\Z+\frac{N-1}2$. Regularity of $\Pi$ gives regularity of $\Sigma$. So
$\Sigma$ is regular algebraic cuspidal.
\emph{Polarisation.} By (BC) $\operatorname{Speh}(\sigma,m)$ is conjugate self-dual; as $\operatorname{Speh}
(\sigma,m)^c=\operatorname{Speh}(\sigma^c,m)$, $\operatorname{Speh}(\sigma,m)^\vee=\operatorname{Speh}
(\sigma^\vee,m)$ and $\operatorname{Speh}(\cdot,m)$ is injective on cuspidal representations (its
cuspidal support determines $\sigma$), $\sigma^c\simeq\sigma^\vee$. Hence $\Sigma^c\simeq\Sigma^\vee
\otimes|\det|^{2s}=\Sigma^\vee\otimes(\phi\circ\auNm_{E/E^+}\circ\det)$ with $\phi=|\cdot|_{\mathbb
A_{E^+}}^{2s}$, an algebraic character ($2s\in\Z$) with $\phi_v(-1)=1$ for every $v\mid\infty$. This is
the hypothesis of (GR), the sign being the same at all infinite places. If the further sign condition in
the definition of a polarised pair in \cite[\S2.1]{BLGGT} is not met, replace $\phi$ by $\phi\,
\delta_{E/E^+}$, $\delta_{E/E^+}$ the quadratic character of $E/E^+$: since $\delta_{E/E^+}\circ
\auNm_{E/E^+}=1$ the relation $\Sigma^c\simeq\Sigma^\vee\otimes(\phi\circ\auNm_{E/E^+}\circ\det)$ is
unchanged, as in \cite[\S2.1]{BLGGT}.
\emph{Galois representation.} At a finite $v\nmid\ell$ where everything is unramified, the Satake
parameters of $\operatorname{Speh}(\sigma,m)_v\otimes|\det|^{(1-n)/2}$ are those of $\Sigma_v\otimes
|\det|^{(1-N)/2}$ times $|\cdot|^{e_a}$ with $e_a=\frac{m-1}2-a-s+\frac{N-n}2$, an integer by the
choice of $s$; $|\cdot|^{e}$ corresponds under $\operatorname{rec}$ to $\varepsilon_\ell^{e}$ (geometric
Frobenius, $\varepsilon_\ell(\pgFrob_w)=\auNm w^{-1}$). So the summand $r_\iota(\Sigma)\varepsilon_\ell
^{e_a}$ has the local parameters of $\operatorname{Speh}(\sigma,m)$, and the rest of the proof of
Proposition~\ref{pga:G123} applies verbatim (that proof takes from the present one only what has been shown up to this point, namely the algebraicity, the polarisation and the local parameters, none of which uses it); in particular the ramified local--global
compatibility including the monodromy operator, used in (F2) of \S\ref{pax:ss-dict}, holds for the
polarised pair $(\Sigma,|\cdot|^{2s})$ by \cite{BLGGT}, as recalled in (GR), and does not require
$\Sigma$ to be conjugate self-dual. (The choice $s=\frac{m-1}2$
is admissible exactly when $n-n_i$ is even; this is the parity issue of Remark~\ref{pga:sign}: the
problem is algebraicity of the twist, not the sign, and it disappears with the choice of $s$ above.)
\end{proof}

\subsubsection{The convention check at the end of \S\ref{pga:ss-M}}\label{pax:ss-conv}

The check asks that the forms $\mathcal A_\lambda$ of {\cite[\S14.3]{P46}} have, in the sense of (G3),
the Harish-Chandra parameter $\varsigma=\lambda+\rho$ at each $\tau$, in the normalisation in which a
character of $\auU(1)$ of weight $c$ has parameter $c$.

\begin{lemma}\label{pax:conv}
Let $\varphi\in\mathcal A_\lambda(J,\chi)$ be nonzero, and $\pi$ an irreducible constituent of the
automorphic representation of $\auU(V)(\mathbb A_{E^+})$ formed by the functions $f_{\varphi,v'}(u,g)
=\langle v',u^{-1}\varphi(g)\rangle$ ($u\in\auU(V)(E^+\otimes\R)$, $v'\in W_\lambda^\vee$; this is the
reading of algebraic forms used in {\cite[Lemma~5.26]{P53}}'s sequel). Then at the archimedean place
$v$ below $\tau$, with $\auU(V)(E^+_v)\simeq\auU(n)$ through $\tau$, $\pi_v$ is the irreducible algebraic
representation of highest weight $\lambda_\tau$, so its parameter is $\lambda_\tau+\rho=\varsigma_\tau$.
For $n=1$ and $\lambda=e$ (the character $\det^e$) this is the character $z\mapsto\tau(z)^e$, of
parameter $e$.
\end{lemma}
\begin{proof}
By (N2), $\gamma\in\Gamma$ acts on $W_\lambda$ by $\rho_\lambda(\tau(\bar\gamma))=\rho_\lambda({}^t
\tau(\gamma)^{-1})$, and this extends to the real group: $u$ acts by $\rho_\lambda({}^t\tau(u)^{-1})=
\rho_\lambda(\overline{\tau(u)})=\overline{\rho_\lambda(\tau(u))}$ ($u$ unitary, $\rho_\lambda$ with real
coefficients). So $W_\lambda$, as a representation of $\auU(n)$ through $\tau$, is the complex
conjugate of $\rho_\lambda\circ\tau$, i.e.\ its contragredient (a unitary representation). Right
translation by $u_0\in\auU(V)(E^+_v)$ on $f_{\varphi,v'}$ gives $f_{\varphi,v'}(uu_0,g)=\langle v',u_0^{-1}
u^{-1}\varphi(g)\rangle=\langle u_0\cdot v',u^{-1}\varphi(g)\rangle$, so the functions $f_{\varphi,v'}$,
$v'$ varying, span a copy of $W_\lambda^\vee$, the dual of the contragredient of $\rho_\lambda\circ\tau$:
$\pi_v\simeq\rho_\lambda\circ\tau$, of highest weight $\lambda_\tau$ (the matrix $(\det g)^eS_\lambda(g)$
being the representation of highest weight $\lambda$ in the standard Borel by {\cite[(60)]{P49}}). The
case $n=1$ is the special case $S_\lambda=1$.
\end{proof}

So the dictionary is: Gross weight $\lambda$ $\leftrightarrow$ $\pi_\infty=\rho_\lambda\circ\tau$
$\leftrightarrow$ parameter $\varsigma=\lambda+\rho$, uniformly in $n$, including $n=1$; this is the
convention assumed in (G3), in Proposition~\ref{pga:weights}, and for the weight of $\mu^0_j$ (a
character of $\auU(1)$ read through the same (N2)). The relative sign used in
Proposition~\ref{pga:weights} is therefore justified. (Had the automorphic forms attached to
algebraic forms been taken to be the complex conjugates $\langle u^{-1}\varphi(g),v\rangle$, one would get the
contragredient at every place and $\varsigma$ would become $-w_0\varsigma$ uniformly in $n$, including
$n=1$; the relative sign would be unchanged.)



\providecommand{\auO}{\mathcal O}
\providecommand{\auGL}{\mathrm{GL}}
\providecommand{\auF}{\mathbb F}
\providecommand{\auU}{\mathrm U}
\providecommand{\auNm}{\operatorname{Nm}}
\providecommand{\pgtr}{\operatorname{tr}}
\providecommand{\pgFrob}{\operatorname{Frob}}
\providecommand{\payL}{\Lambda}
\providecommand{\payM}{\mathsf M}
\providecommand{\payU}{\mathbf U}
\providecommand{\payUd}{\mathbf U'}
\providecommand{\payN}{\mathcal N}

\subsection{The spherical chain lemma}\label{pay:sec}

This section proves, without using the classification of the unitary dual of $\auGL_N$, the local
statement behind Lemma~\ref{pax:tadic} (hence behind Lemma~\ref{pri:chain}) and the local
Lemma~\ref{pax:toproot}, used in the proof of Proposition~\ref{pri:dict}(d) in \S\ref{pax:ss-dict}. Both come from one theorem (Theorem~\ref{pay:chain}): for a nonzero
spherical Hecke eigenvector in a unitary representation of $\auGL_N(F)$, a root $c$ of its Hecke
polynomial with $qc$ not a root satisfies $|c|^2>q^{N-2}$, and of a chain of roots the operator $U$
of $\mathrm{diag}(\varpi,1,\dots,1)$ at pro-$p$ Iwahori level sees only the top one. The proof
consists of two counts of hyperplanes over the residue field (\S\ref{pay:ss-count}), which give
recursions for $U$ and for its adjoint on a chain of explicit vectors (\S\ref{pay:ss-rec}), and of
linear algebra in a positive definite hermitian space (\S\ref{pay:ss-lin}). The two facts used in
Remark~\ref{pax:rem-BL} are proved in \S\ref{pay:ss-BL}.

\subsubsection{Setting}\label{pay:ss-set}

Let $F$ be a non-archimedean local field, $\auO_F$ its ring of integers, $\varpi$ a uniformiser,
$\auF=\auO_F/\varpi$ the residue field, of order $q$, and $N\ge1$. Put $G=\auGL_N(F)$, $K=\auGL_N(
\auO_F)$, $e_1,\dots,e_N$ the standard basis of $F^N$ (column vectors), $d_i=\mathrm{diag}(\varpi
1_i,1_{N-i})$, $d=d_1=\mathrm{diag}(\varpi,1,\dots,1)$, and let $I_1\subseteq K$ be the pro-$p$
Iwahori subgroup: the elements of $K$ that are upper triangular unipotent modulo $\varpi$.

Let $(\rho,V)$ be a smooth representation of $G$ on a complex vector space with a positive definite
$G$-invariant hermitian form $(\ ,\ )$, linear in the first variable (convention (Cv1) of
{\cite[\S2.1]{P51}}); so $\rho(g)^*=\rho(g^{-1})$. For a vector $v$ fixed by $K$ put, as in (N3)
of {\cite[\S6.1]{P49}},
\[
T_iv=\sum_{h\in Kd_iK/K}\rho(h)v\qquad(0\le i\le N),
\]
and for $y$ fixed by $I_1$ let $Uy=\sum_{h\in I_1dI_1/I_1}\rho(h)y$, the operator $U$ of
Proposition~\ref{pri:dict} and of \S\ref{pax:ss-dict}. Let $\payM\colon V\to V^{I_1}$ be the mean
over $I_1$, $\payM v=\int_{I_1}\rho(k)v\,dk$ for the Haar measure of total mass $1$; as $v$ is fixed
by an open subgroup, this is a finite average. $\payM$ is the identity on $V^{I_1}$, and it is
self-adjoint: $(\payM a,b)=\int(a,\rho(k^{-1})b)\,dk=(a,\payM b)$, the measure being invariant under
$k\mapsto k^{-1}$.

If $x\in V^K$ satisfies $T_ix=t_ix$ for all $i$, its \emph{Hecke polynomial} is
\begin{equation}\label{pay:eq-P}
P_x(X)=\sum_{i=0}^N(-1)^iq^{\binom i2}t_iX^{N-i}\qquad(t_0=1).
\end{equation}
For $x$ a nonzero spherical vector in an irreducible unramified $\pi$ this is the polynomial
$H_\pi$ of~\eqref{pax:eq-satake}; and on forms, with the operators $T_i$ of (N3), it is the Hecke polynomial $H_w$ of Lemma~\ref{pri:chain} (the polynomial~\eqref{pax:eq-satake}).

\emph{Lattices.} For $0\le j\le N$ let $\payL_j$ be the set of $\auO_F$-lattices $L$ with $\varpi
\auO_F^N\subseteq L\subseteq\auO_F^N$ and $\dim_\auF(\auO_F^N/L)=j$, and $\payL^{(1)}_j=\{L\in
\payL_j:e_1\in L\}$. For a lattice $L$ fix $g_L\in G$ with $g_L\auO_F^N=L$; then $g_LK$ depends only on
$L$, and $gK\mapsto g\auO_F^N$ is a bijection from $G/K$ onto the set of lattices.

\begin{lemma}\label{pay:lat}
$Kd_jK/K$ corresponds under this bijection to $\payL_j$. Hence for $v\in V^K$, $T_jv=\sum_{L\in
\payL_j}\rho(g_L)v$. Moreover $\payL_0=\{\auO_F^N\}$, $\payL_N=\{\varpi\auO_F^N\}$,
$\payL^{(1)}_N=\emptyset$, $\payL^{(1)}_0=\{\auO_F^N\}$ and $\payL^{(1)}_{N-1}=\{\auO_Fe_1+\varpi\auO_F^N
\}$.
\end{lemma}
\begin{proof}
$d_j\auO_F^N\in\payL_j$, and $K$ preserves $\payL_j$; so $Kd_jK/K$ maps into $\payL_j$. Conversely,
$L\mapsto L/\varpi\auO_F^N$ is a bijection from $\payL_j$ onto the subspaces of dimension $N-j$ of
$\auF^N$, on which $K$ acts transitively through $K\to\auGL_N(\auF)$ (surjective); so every $L\in
\payL_j$ is $kd_j\auO_F^N$ with $k\in K$. The last assertions are immediate ($e_1\notin\varpi
\auO_F^N$; a subspace of dimension $1$ containing $\bar e_1$ is $\auF\bar e_1$).
\end{proof}

For $x\in V^K$ and $0\le j\le N$ put
\begin{equation}\label{pay:eq-xj}
x_j=\sum_{L\in\payL^{(1)}_j}\rho(g_L)x .
\end{equation}
So $x_0=x$, $x_N=0$, and $x_{N-1}=\rho(\mathrm{diag}(1,\varpi,\dots,\varpi))x$. Each $x_j$ is fixed by
$I_1$: an element $k\in K$ with $ke_1\in\auO_F^\times e_1+\varpi\auO_F^N$ (in particular $k\in I_1$)
permutes $\payL^{(1)}_j$ (it preserves $\payL_j$, and $ke_1\in L$ iff $e_1\in L$ as $\varpi\auO_F^N
\subseteq L$), and $\rho(k)\rho(g_L)x=\rho(g_{kL})x$ since $x$ is fixed by $K$.

\subsubsection{Two counts of hyperplanes}\label{pay:ss-count}

Let $\payN\subseteq K$ be the group of matrices that are equal to $1$ except in the first row, whose
off-diagonal entries are arbitrary in $\auO_F$ (and the diagonal entry is $1$); $\payN\simeq
\auO_F^{N-1}$ and $\payN\subseteq I_1$. For $u\in\auF^{N-1}$ fix $n_u\in\payN$ with first row reducing
to $(1,u)$.

\begin{lemma}[Count of hyperplanes]\label{pay:count}
Let $\sigma$ be a representation of $G$ on a vector space $W$, $v\in W$ fixed by $\sigma(K)$, and
$\psi\colon W\to W'$ a linear map with $\psi\circ\sigma(n)=\psi$ for all $n\in\payN$. Put $v_j=
\sum_{L\in\payL^{(1)}_j}\sigma(g_L)v$ and $v^\flat_j=\sum_{L\in\payL_j}\sigma(g_L)v$. Then for $0\le
j\le N-1$,
\[
q^{N-1}\,\psi\bigl(\sigma(d)v_j\bigr)=q^j\,\psi\bigl(v^\flat_{j+1}-v_{j+1}\bigr).
\]
\end{lemma}
\begin{proof}
Let $H=\{w\in\auO_F^N:w_1\in\varpi\auO_F\}=d\auO_F^N$. \emph{Step 1:} $L\mapsto dL$ is a bijection from
$\payL^{(1)}_j$ onto $\{M'\in\payL_{j+1}:M'\subseteq H\}$. Indeed, for $L\in\payL^{(1)}_j$: $dL
\subseteq d\auO_F^N=H$; $dL$ contains $de_1=\varpi e_1$ and $d(\varpi\auO_F^N)\supseteq\varpi
\sum_{l\ge2}\auO_Fe_l$, so $dL\supseteq\varpi\auO_F^N$; and $[\auO_F^N:dL]=[\auO_F^N:d\auO_F^N][d\auO_F^N:
dL]=q\cdot q^j$. Conversely, if $M'\in\payL_{j+1}$ and $M'\subseteq H$, then $d^{-1}M'\subseteq
d^{-1}H=\auO_F^N$, $d^{-1}M'\supseteq d^{-1}\varpi\auO_F^N\ni e_1$, $d^{-1}M'\supseteq\varpi\auO_F^N$,
and its colength is $j$.

\emph{Step 2:} for $M\in\payL_{j+1}$, the number of $u\in\auF^{N-1}$ with $n_uM\subseteq H$ is $0$
if $e_1\in M$ and $q^j$ if $e_1\notin M$. Indeed $n_uM\in\payL_{j+1}$, and $n_uM\subseteq H$ iff the
linear form $\lambda_u=(1,u)$ on $\auF^N$ (the reduction of the first row of $n_u$) vanishes on
$\bar M=M/\varpi\auO_F^N\subseteq\auF^N$, a subspace of dimension $N-j-1$. Every linear form $\lambda$
with $\lambda(\bar e_1)=1$ is $\lambda_u$ for exactly one $u$. The annihilator $\bar M^\perp$ has
dimension $j+1$, and evaluation at $\bar e_1$ is a linear form on it, which is zero iff $\bar e_1\in
\bar M$, i.e.\ iff $e_1\in M$. If it is nonzero, its fibre over $1$ is a coset of a subspace of
dimension $j$, with $q^j$ elements.

\emph{Step 3.} By Steps 1 and 2, the map $(u,L)\mapsto(u,n_u^{-1}dL)$ is a bijection from $\auF^{N-1}
\times\payL^{(1)}_j$ onto the pairs $(u,M)$ with $M\in\payL_{j+1}$ and $n_uM\subseteq H$, and for
fixed $M$ there are $q^j[e_1\notin M]$ such $u$. As $v$ is fixed by $K$, $\sigma(n_u^{-1}dg_L)v=
\sigma(g_{n_u^{-1}dL})v$. Hence
\[
\sum_{u\in\auF^{N-1}}\sigma(n_u^{-1}d)\,v_j=\sum_{M\in\payL_{j+1},\ e_1\notin M}q^j\sigma(g_M)v=q^j
(v^\flat_{j+1}-v_{j+1}).
\]
Apply $\psi$: as $n_u^{-1}\in\payN$, each term on the left gives $\psi(\sigma(d)v_j)$, and there are
$q^{N-1}$ of them.
\end{proof}

\begin{lemma}[Duality of lattices]\label{pay:dual}
Let $\theta(g)={}^tg^{-1}$, an automorphism of $G$ preserving $K$, and $\rho^\theta=\rho\circ\theta$.
Let $x\in V^K$ (so $x$ is fixed by $\rho^\theta(K)$), $c_0=\varpi^{-1}1_N$ and $c_1=\mathrm{diag}(1,
\varpi^{-1},\dots,\varpi^{-1})$. If $i+j=N$ then $\sum_{L\in\payL_j}\rho^\theta(g_L)x=\rho(c_0)
\sum_{L\in\payL_i}\rho(g_L)x$; if $1+i+j=N$ then $\sum_{L\in\payL^{(1)}_j}\rho^\theta(g_L)x=\rho(c_1)
x_i$.
\end{lemma}
\begin{proof}
For a lattice $L$ let $L^*=\{w\in F^N:{}^twL\subseteq\auO_F\}$; then ${}^tg_L^{-1}\auO_F^N=L^*$,
$(L^*)^*=L$, and $L\subseteq L'$ iff $L'^*\subseteq L^*$, with $[L^*:L'^*]=[L':L]$. If $\varpi
\auO_F^N\subseteq L\subseteq\auO_F^N$ has colength $j$, then $\auO_F^N\subseteq L^*\subseteq\varpi^{-1}
\auO_F^N$ with $[L^*:\auO_F^N]=q^j$. \emph{First case.} $L\mapsto\varpi L^*=c_0^{-1}L^*$ maps $\payL_j$
to $\payL_i$ ($[\auO_F^N:\varpi L^*]=[\varpi^{-1}\auO_F^N:L^*]=q^{N-j}$), with inverse $L'\mapsto
(\varpi^{-1}L')^*$. \emph{Second case.} $e_1\in L$ iff $w_1={}^twe_1\in\auO_F$ for all $w\in L^*$. For
$L\in\payL^{(1)}_j$ put $L'=c_1^{-1}L^*$: its first coordinates are those of $L^*$, in $\auO_F$, and
its other coordinates are $\varpi$ times those of $L^*$, in $\auO_F$; so $L'\subseteq\auO_F^N$, and
$L'\supseteq c_1^{-1}\auO_F^N=\auO_Fe_1+\varpi\auO_F^N$, so $e_1\in L'$ and $\varpi\auO_F^N\subseteq
L'$; finally $[\auO_F^N:L']=[\auO_F^N:c_1^{-1}\auO_F^N]/[L^*:\auO_F^N]=q^{N-1-j}=q^i$. Conversely, for
$L'\in\payL^{(1)}_i$, $c_1L'\supseteq c_1(\auO_Fe_1+\varpi\auO_F^N)\supseteq\auO_F^N$, $c_1L'\subseteq
c_1\auO_F^N$ has integral first coordinates and lies in $\varpi^{-1}\auO_F^N$, so $(c_1L')^*\in
\payL^{(1)}_j$ by the same computations; the two maps are inverse. In both cases ${}^tg_L^{-1}K=c_m
g_{L'}K$ (both send $\auO_F^N$ to $L^*=c_mL'$), and as $x$ is fixed by $K$, $\rho^\theta(g_L)x=\rho(
c_mg_{L'})x$. Summing gives the two identities.
\end{proof}

\subsubsection{The recursions}\label{pay:ss-rec}

Define endomorphisms of $V$
\[
\payU=q^{N-1}\,\payM\,\rho(d)\,\payM,\qquad\payUd=q^{N-1}\,\payM\,\rho(d^{-1})\,\payM .
\]
They are mutually adjoint: $(\payU a,b)=(a,\payUd b)$ and $(\payUd a,b)=(a,\payU b)$ for all $a,b\in V$,
since $\payM$ is self-adjoint and $\rho(d)^*=\rho(d^{-1})$.

\begin{lemma}\label{pay:UU}
On $V^{I_1}$, $U=\payU$; in particular $U$ preserves $V^{I_1}$.
\end{lemma}
\begin{proof}
Let $y\in V^{I_1}$ and $I'=I_1\cap dI_1d^{-1}$. The map $kI'\mapsto kdI_1$ is a bijection from
$I_1/I'$ onto $I_1dI_1/I_1$ ($kdI_1=k'dI_1$ iff $k^{-1}k'\in dI_1d^{-1}$), and the number of cosets is
$q^{N-1}$ (the decomposition $I_1dI_1=\bigsqcup_{u}u\,d\,I_1$, $u\in N(\auO_F)/dN(\auO_F)d^{-1}$, in
the proof of Lemma~\ref{pax:Ueig}). The vector $\rho(d)y$ is fixed by $dI_1d^{-1}\supseteq I'$, so
$k\mapsto\rho(kd)y$ is constant on the cosets $kI'$, each of measure $q^{1-N}$. Hence $\payM\rho(d)y=
q^{1-N}\sum_{k\in I_1/I'}\rho(kd)y=q^{1-N}Uy$, and $\payU y=q^{N-1}\payM\rho(d)\payM y=Uy$.
\end{proof}

\begin{proposition}[Recursions]\label{pay:rec}
Write $N=n+1$. Let $x\in V^K$ with $T_ix=t_ix$ ($0\le i\le N$), and $x_j$ as in~\eqref{pay:eq-xj}.
Then $t_0=1$ if $x\ne0$, and:
\begin{enumerate}
\item[(R1)] $\payU x_j=q^j\,(t_{j+1}x-x_{j+1})$ for $0\le j\le n$;
\item[(R2)] $\payUd x_{i+1}=q^{n-i-1}\bigl((t_{i+1}/q^n)\,\payUd x-x_i\bigr)$ for $0\le i<n$;
\item[(R3)] if $x\ne0$ then $t_N\ne0$, $\payUd x=q^nt_N^{-1}x_n$, and $\payU x_n=q^nt_Nx$.
\end{enumerate}
\end{proposition}
\begin{proof}
$T_0$ is the identity on $V^K$, so $t_0=1$ when $x\ne0$. As each $x_j$ and $x$ lie in $V^{I_1}$,
$\payM x_j=x_j$.

(R1) Apply Lemma~\ref{pay:count} with $\sigma=\rho$, $v=x$, $\psi=\payM$ (invariant under $\rho(I_1)
\supseteq\rho(\payN)$): $v_j=x_j$ and $v^\flat_{j+1}=T_{j+1}x=t_{j+1}x$ by Lemma~\ref{pay:lat}, so
$q^n\payM\rho(d)x_j=q^j(t_{j+1}x-x_{j+1})$, and the left side is $\payU x_j$.

(R2) Apply Lemma~\ref{pay:count} with $\sigma=\rho^\theta$, $v=x$, and $\psi=\payM\circ\rho(c_1^{-1})$.
The invariance: for $n'\in\payN$ with first row $(1,a)$, $\theta(n')$ is unipotent lower triangular with
entries $-a_l$ at $(l,1)$, $l\ge2$, and $0$ elsewhere off the diagonal; so $c_1^{-1}\theta(n')c_1$ has
entries $-\varpi a_l\in\varpi\auO_F$ at $(l,1)$ and lies in $I_1$, whence $\psi(\rho^\theta(n')w)=
\payM\rho(c_1^{-1}\theta(n')c_1)\rho(c_1^{-1})w=\psi(w)$. Note $\rho^\theta(d)=\rho(d^{-1})$. Fix
$0\le i<n$ and let $j=n-1-i\in[0,N-1]$. By Lemma~\ref{pay:dual}, $v_j=\rho(c_1)x_{i+1}$ (as $1+(i+1)+j=N$), $v^\flat_{j+1}=\rho(c_0)t_{i+1}x$ (as
$(i+1)+(j+1)=N$), and $v_{j+1}=\rho(c_1)x_i$ (as $1+i+(j+1)=N$). Lemma~\ref{pay:count} gives
\[
q^n\payM\rho(c_1^{-1}d^{-1}c_1)x_{i+1}=q^j\bigl(t_{i+1}\payM\rho(c_1^{-1}c_0)x-\payM x_i\bigr).
\]
Diagonal matrices commute and $c_1^{-1}c_0=d^{-1}$, so this reads $\payUd x_{i+1}=q^{n-1-i}(t_{i+1}
q^{-n}\payUd x-x_i)$, using $\payM\rho(d^{-1})x=q^{-n}\payUd x$.

(R3) By Lemma~\ref{pay:lat}, $T_Nx=\rho(\varpi1_N)x$ and $x_n=\rho(\mathrm{diag}(1,\varpi,\dots,
\varpi))x=\rho(d^{-1})\rho(\varpi1_N)x=t_N\rho(d^{-1})x$. As $\rho(\varpi1_N)$ is invertible,
$t_Nx=T_Nx\ne0$, so $t_N\ne0$. Then $\payUd x=q^n\payM\rho(d^{-1})x=q^nt_N^{-1}\payM x_n=q^nt_N^{-1}x_n$,
and $\payU x_n=q^n\payM\rho(d)x_n=q^nt_N\payM x=q^nt_Nx$.
\end{proof}

\begin{remark}\label{pay:rem-rec}
With the convention $x_N=0$, (R1) at $j=n$ is the last identity of (R3), and (R2) also holds for
$i=n$: by (R3) the right side is $q^{-1}(t_Nq^{-n}\cdot q^nt_N^{-1}x_n-x_n)=0=\payUd x_N$. We use (R2)
for $0\le i\le n$ in this extended form.
\end{remark}

\subsubsection{Linear algebra}\label{pay:ss-lin}

For $A=\sum_ig_iX^i\in\C[X]$ let $\bar A=\sum_i\bar g_iX^i$. If $U,U'$ are endomorphisms of a
hermitian space $(V,(\ ,\ ))$ (form linear in the first variable) with $(Ua,b)=(a,U'b)$ and $(U'a,b)=
(a,Ub)$ for all $a,b$, then by induction on the degree
\begin{equation}\label{pay:eq-adj}
(a,A(U)b)=(\bar A(U')a,b)\qquad(a,b\in V,\ A\in\C[X]).
\end{equation}

\begin{proposition}\label{pay:lin}
Let $(V,(\ ,\ ))$ be a complex vector space with a positive definite hermitian form, $U,U'$ mutually
adjoint endomorphisms as above, $q>1$ a real number and $n\ge0$ an integer. Let $x_0,\dots,x_{n+1}\in
V$ with $x:=x_0\ne0$ and $x_{n+1}=0$, and $t_0=1,t_1,\dots,t_{n+1}\in\C$ with $t_{n+1}\ne0$, such that
\begin{enumerate}
\item[(R1)] $Ux_j=q^j(t_{j+1}x-x_{j+1})$ for $0\le j\le n$;
\item[(R2)] $U'x_{i+1}=q^{n-i-1}\bigl((t_{i+1}/q^n)U'x-x_i\bigr)$ for $0\le i\le n$;
\item[(R3)] $U'x=q^nt_{n+1}^{-1}x_n$.
\end{enumerate}
Let $P=\sum_{j=0}^{n+1}(-1)^jq^{\binom j2}t_jX^{n+1-j}$. Then:
\begin{enumerate}
\item[(a)] if $P(c)=0$ and $P(qc)\ne0$, then $|c|^2>q^{n-1}$;
\item[(b)] if $\delta\in\C^\times$, $k\ge0$, $R\in\C[X]$, $P=\prod_{i<k}(\delta X-q^i)\cdot R$ and $P(q^k/\delta)=0$,
then $\bar R(U')x=0$.
\end{enumerate}
\end{proposition}

\begin{proof}
\emph{Step 1.} For $0\le m\le n+1$ let $p_m=\sum_{j\le m}(-1)^jq^{\binom j2}t_jX^{m-j}$, so $p_0=1$,
$p_{n+1}=P$ and $p_{m+1}=Xp_m+\kappa_{m+1}$ with $\kappa_{m+1}=(-1)^{m+1}q^{\binom{m+1}2}t_{m+1}$. Then
\begin{equation}\label{pay:eq-pm}
p_m(U)x=(-1)^mq^{\binom m2}x_m\qquad(0\le m\le n+1).
\end{equation}
For $m=0$ this is $x=x_0$; if it holds for $m\le n$, then by (R1)
$p_{m+1}(U)x=(-1)^mq^{\binom m2}q^m(t_{m+1}x-x_{m+1})+\kappa_{m+1}x=(-1)^{m+1}q^{\binom{m+1}2}x_{m+1}$,
as $\binom m2+m=\binom{m+1}2$. In particular $P(U)x=\pm q^{\binom{n+1}2}x_{n+1}=0$.

\emph{Step 2.} Fix $c\in\C$ and let $D_m=(p_m-p_m(c))/(X-c)\in\C[X]$ ($0\le m\le n+1$), and $D_{n+2}=
cD_{n+1}+P$. Then $D_0=0$ and $D_{m+1}=cD_m+p_m$ for $0\le m\le n+1$ (indeed $(X-c)(cD_m+p_m)=c(p_m-
p_m(c))+(X-c)p_m=Xp_m-cp_m(c)=p_{m+1}-p_{m+1}(c)$). We claim, for $0\le m\le n+1$,
\begin{equation}\label{pay:eq-Dm}
U'\bigl(D_{m+1}(U)x\bigr)=\frac{p_m(qc)}{q^m}\,U'x+q^{n-1}D_m(U)x .
\end{equation}
For $m=0$ both sides are $U'x$. Assume~\eqref{pay:eq-Dm} for some $m\le n$. Then $D_{m+2}(U)x=c\,
D_{m+1}(U)x+p_{m+1}(U)x$, and by~\eqref{pay:eq-pm}, (R2) with $i=m$, and~\eqref{pay:eq-pm} again,
\[
U'p_{m+1}(U)x=(-1)^{m+1}q^{\binom{m+1}2}q^{n-m-1}\Bigl(\frac{t_{m+1}}{q^n}U'x-x_m\Bigr)=
\frac{\kappa_{m+1}}{q^{m+1}}U'x+q^{n-1}p_m(U)x .
\]
Adding $c$ times~\eqref{pay:eq-Dm}, the coefficient of $U'x$ is $(qc\,p_m(qc)+\kappa_{m+1})/q^{m+1}=
p_{m+1}(qc)/q^{m+1}$, and the rest is $q^{n-1}(cD_m+p_m)(U)x=q^{n-1}D_{m+1}(U)x$. This
proves~\eqref{pay:eq-Dm} for $m+1$.

Put $z_c=D_{n+1}(U)x$. As $P(U)x=0$, \eqref{pay:eq-Dm} for $m=n+1$ reads
\begin{equation}\label{pay:eq-star}
c\,U'z_c=\frac{P(qc)}{q^{n+1}}\,U'x+q^{n-1}z_c .
\end{equation}

\emph{Step 3.} If $P(c)=0$ then $(U-c)z_c=(P-P(c))(U)x=P(U)x=0$, i.e.\ $Uz_c=cz_c$. By (R3) and (R1)
with $j=n$ ($x_{n+1}=0$), $UU'x=q^nt_{n+1}^{-1}\cdot q^nt_{n+1}x=q^{2n}x$, hence
$(U'a,U'x)=(a,UU'x)=q^{2n}(a,x)$ for all $a$. Pair~\eqref{pay:eq-star} with $U'x$ in the second
variable: the left side gives $c\,q^{2n}(z_c,x)$; on the right, $(U'x,U'x)=q^{2n}(x,x)$ and $(z_c,U'x)=
(Uz_c,x)=c(z_c,x)$. Hence, for every root $c$ of $P$,
\begin{equation}\label{pay:eq-E}
c\,(z_c,x)\,(q^{n+1}-1)=P(qc)\,(x,x).
\end{equation}

\emph{Step 4: proof of (a).} Let $P(c)=0$, $a=P(qc)\ne0$. By~\eqref{pay:eq-E}, $c\ne0$ and $z_c\ne0$.
Pair~\eqref{pay:eq-star} with $z_c$ in the second variable: the left side is $c(U'z_c,z_c)=c(z_c,Uz_c)
=|c|^2(z_c,z_c)$; on the right, $(U'x,z_c)=(x,Uz_c)=\bar c(x,z_c)$, and by the complex conjugate
of~\eqref{pay:eq-E}, $\bar c(x,z_c)=\bar a(x,x)/(q^{n+1}-1)$. Therefore
\[
\bigl(|c|^2-q^{n-1}\bigr)(z_c,z_c)=\frac{|a|^2\,(x,x)}{q^{n+1}(q^{n+1}-1)}>0 ,
\]
and $|c|^2>q^{n-1}$.

\emph{Step 5: two consecutive roots.} Let $c\ne0$ with $P(c)=P(qc)=0$. We show $(P/(X-c))(U)x=0$.
Here $P/(X-c)=D_{n+1}$, so the vector is $z_c$. By~\eqref{pay:eq-star}, $U'z_c=\mu z_c$ with $\mu=
q^{n-1}/c$, and by~\eqref{pay:eq-E}, $(z_c,x)=0$. For every $A\in\C[X]$, by~\eqref{pay:eq-adj},
$(z_c,A(U)x)=(\bar A(U')z_c,x)=\bar A(\mu)(z_c,x)=0$. With $A=D_{n+1}$ this gives $(z_c,z_c)=0$, so
$z_c=0$.

\emph{Step 6: cancellation of a chain.} Under the hypotheses of (b) we show $R(U)x=0$, by induction
on $k$. For $k=0$, $R=P$ and $P(U)x=0$. Let $k\ge0$ and $P=\prod_{i<k+1}(\delta X-q^i)\cdot R$ with
$P(q^{k+1}/\delta)=0$. Then $P=\prod_{i<k}(\delta X-q^i)\cdot R'$ with $R'=(\delta X-q^k)R$, and $P(q^k/\delta)=0$; by the
induction hypothesis $R'(U)x=0$, i.e.\ $UR(U)x=c\,R(U)x$ with $c=q^k/\delta\ne0$. As $P(c)=0=P(qc)$, Step 5
gives $(P/(X-c))(U)x=0$; but $P/(X-c)=\delta\,Q\,R$ with $Q=\prod_{i<k}(\delta X-q^i)$, so $0=Q(U)R(U)x=Q(c)R(U)x$,
and $Q(c)=\prod_{i<k}(q^k-q^i)\ne0$ as $q>1$. So $R(U)x=0$.

\emph{Step 7: proof of (b).} By Step 6, $R(U)x=0$, hence by~\eqref{pay:eq-pm} $R(U)x_j=\pm q^{-\binom
j2}p_j(U)R(U)x=0$ for $0\le j\le n+1$. Let $\mathcal S=\mathrm{span}(x_0,\dots,x_n)$. By (R3), $U'x_0\in \mathcal S$; by
(R2), $U'x_{i+1}\in\C U'x+\C x_i\subseteq \mathcal S$ for $i<n$. So $\mathcal S$ is $U'$-stable, and $B(U')x\in \mathcal S$ for
every $B\in\C[X]$; thus $R(U)B(U')x=0$. With $B=\bar R$ and~\eqref{pay:eq-adj},
$(\bar R(U')x,\bar R(U')x)=(x,R(U)\bar R(U')x)=0$, so $\bar R(U')x=0$.
\end{proof}

\subsubsection{The chain theorem}\label{pay:ss-thm}

\begin{theorem}[Spherical chain lemma]\label{pay:chain}
Let $(\rho,V)$ be as in \S\ref{pay:ss-set}, $x\in V^K$ nonzero with $T_ix=t_ix$ for $0\le i\le N$, and
$P=P_x$ its Hecke polynomial~\eqref{pay:eq-P}.
\begin{enumerate}
\item[(i)] If $P(c)=0$ and $P(qc)\ne0$, then $|c|^2>q^{N-2}$.
\item[(ii)] Let $\delta\in\C^\times$, $k\ge0$, $R\in\C[X]$, and suppose $P(\delta^{-1}X)=\prod_{i\le k}(X-q^i)
\cdot S(X)$ with $S$ dividing $R$. Then for every $y\in V^{I_1}$,
\[
\bigl(x,\ \Phi(U)\,y\bigr)=0,\qquad\Phi(X)=\bigl((X-q^k)R\bigr)(\delta X).
\]
\end{enumerate}
\end{theorem}

\begin{proof}
Write $N=n+1$ and let $x_j$ be as in~\eqref{pay:eq-xj}. By Proposition~\ref{pay:rec},
Remark~\ref{pay:rem-rec} and the adjointness of $\payU,\payUd$ (\S\ref{pay:ss-rec}), the hypotheses of
Proposition~\ref{pay:lin} hold for $(U,U')=(\payU,\payUd)$, the $x_j$ and the $t_j$; and the polynomial
$P$ there is $P_x$. (i) is Proposition~\ref{pay:lin}(a), as $q^{n-1}=q^{N-2}$.

(ii) Substituting $\delta X$ for $X$, $P(X)=\prod_{i\le k}(\delta X-q^i)\cdot S(\delta X)$. Write $R=SB$ and $A(X)=(\delta X-
q^k)S(\delta X)$; then $P=\prod_{i<k}(\delta X-q^i)\cdot A$ and $P(q^k/\delta)=0$. By Proposition~\ref{pay:lin}(b),
$\bar A(\payUd)x=0$. Now $\Phi(X)=B(\delta X)\,A(X)$. By Lemma~\ref{pay:UU}, $U$ preserves $V^{I_1}$ and equals
$\payU$ there, so $\Phi(U)y=\Phi(\payU)y$, and by~\eqref{pay:eq-adj}
\[
(x,\Phi(\payU)y)=\bigl(\bar\Phi(\payUd)x,y\bigr)=\bigl(\bar B(\bar \delta\,\payUd)\,\bar A(\payUd)x,\ y\bigr)=0 .
\qedhere
\]
\end{proof}

\begin{lemma}[Symmetry of the roots]\label{pay:sym}
In the situation of Theorem~\ref{pay:chain}, if $y,y'\in\C$ satisfy $y\,\bar y'=q^{N-1}$ and
$P(y')=0$, then $P(y)=0$.
\end{lemma}
\begin{proof}
Let $dg$ be the Haar measure on $G$ with $K$ of measure $1$; for $v\in V^K$, $T_iv=\int_{Kd_iK}\rho(g)v
\,dg$. For $v,v'\in V^K$, as $G$ is unimodular,
$(T_iv,v')=\int_{Kd_iK}(v,\rho(g^{-1})v')\,dg=\bigl(v,\int_{Kd_i^{-1}K}\rho(g)v'\,dg\bigr)$. Now
$d_i^{-1}=\varpi^{-1}\mathrm{diag}(1_i,\varpi1_{N-i})$ and $\mathrm{diag}(1_i,\varpi1_{N-i})\in
Kd_{N-i}K$ (conjugate by a permutation matrix), so $Kd_i^{-1}K=\varpi^{-1}Kd_{N-i}K$ and
$\int_{Kd_i^{-1}K}\rho(g)v'\,dg=\rho(\varpi^{-1}1_N)T_{N-i}v'$. With $v=v'=x$ and $z=(\rho(\varpi1_N)x,
x)$ this gives, for $0\le i\le N$,
\begin{equation}\label{pay:eq-ti}
t_i\,(x,x)=(x,\rho(\varpi^{-1}1_N)t_{N-i}x)=\bar t_{N-i}\,z .
\end{equation}
Now multiply $P(y)$ by $(x,x)$ and use~\eqref{pay:eq-ti} termwise:
$P(y)(x,x)=z\sum_i(-1)^iq^{\binom i2}\bar t_{N-i}y^{N-i}$. On the other hand, the term of index $b=N-i$
of $\overline{P(y')}$, multiplied by $y^N$, is
$(-1)^{N-i}q^{\binom{N-i}2}\bar t_{N-i}\,\bar y'^{\,i}y^N=(-1)^{N-i}q^{\binom{N-i}2+(N-1)i}\bar t_{N-i}
y^{N-i}$, using $(y\bar y')^i=q^{(N-1)i}$; and $\binom{N-i}2+(N-1)i=\binom N2+\binom i2$. Hence
\[
q^{\binom N2}\,(x,x)\,P(y)=(-1)^N\,z\,y^N\,\overline{P(y')} ,
\]
and $P(y')=0$ forces $P(y)=0$, as $(x,x)>0$.
\end{proof}

\begin{corollary}[Chains of roots]\label{pay:roots}
In the situation of Theorem~\ref{pay:chain}, let $c$ be a root of $P$ and $k\ge1$ an integer with
$|c|^2q^k=q^N$. Then $c,cq,\dots,cq^{k-1}$ are roots of $P$, and $\prod_{i<k}(X-q^i)$ divides $P(cX)$.
\end{corollary}
\begin{proof}
We show that $cq^i$ is a root for $0\le i<k$, by strong induction on $i$; write $k=i+1+j$ with $j\ge0$.
If $j<i$: $cq^j$ is a root by induction, and $(cq^i)\overline{(cq^j)}=|c|^2q^{i+j}=q^{N-k}q^{k-1}=q^{N-1}$,
so $cq^i$ is a root by Lemma~\ref{pay:sym}. If $j\ge i$ and $i=0$, $c$ is a root. If $j\ge i\ge1$ and
$cq^i$ were not a root, then $c'=cq^{i-1}$ is a root (induction) with $qc'$ not a root, so
$|c'|^2>q^{N-2}$ by Theorem~\ref{pay:chain}(i), i.e.\ $q^{N-k+2i-2}>q^{N-2}$, i.e.\ $2i>k=i+1+j$, so
$j<i-1$, a contradiction. As $c\ne0$ ($|c|^2=q^{N-k}$) and $q>1$, the $k$ roots $cq^i$ are distinct;
so $\prod_{i<k}(X-cq^i)$ divides $P(X)$, and $\prod_{i<k}(X-q^i)$ divides $P(cX)$.
\end{proof}

\subsubsection{Consequences for \S\ref{pax:sec}}\label{pay:ss-cons}

\begin{proof}[Proof of Lemma~\ref{pax:tadic} without the classification of the unitary dual]
Let $\pi$ be an irreducible unitary unramified representation of $\auGL_N(F)$, with invariant inner
product, and $x\ne0$ a spherical vector; $\pi^K$ is stable under the $T_i$, and $H_\pi$
in~\eqref{pax:eq-satake} is defined by their eigenvalues $t_i$ on $\pi^K$, so $T_ix=t_ix$ and $H_\pi=
P_x$ (both are the polynomial~\eqref{pay:eq-P}). Corollary~\ref{pay:roots} with $c$ a root of $H_\pi$
and $|c|^2q^k=q^N$ is the statement of Lemma~\ref{pax:tadic}.
\end{proof}

So the proof of Lemma~\ref{pri:chain} in \S\ref{pax:ss-chain} holds with Lemma~\ref{pax:tadic} proved
here: the input from Tadi\'c's classification, the discussion of unramified constituents and the
paragraph ``The chain'' of the proof in \S\ref{pax:ss-chain} are not needed. (The step ``Unitarity'' of the proof of Lemma~\ref{pri:chain} may even be shortened:
Corollary~\ref{pay:roots} applies to any smooth representation with a positive definite invariant
hermitian form, so it can be applied directly to the representation of $\auU(V)(E^+_v)\simeq
\auGL_n(E_w)$ on the forms of the given weight, with the global inner product, and to $x=y$, without
passing to an irreducible constituent.)

\begin{proof}[Proof of Lemma~\ref{pax:toproot}]
Let $\pi$, $x$, $X_0$, $k\ge1$, $S\mid R$ be as there, so $H_\pi(X_0X)=\prod_{i<k}(X-q^i)\cdot S(X)$, and
$H_\pi=P_x$ as in the previous proof. Apply Theorem~\ref{pay:chain}(ii) with $\delta=X_0^{-1}$ and $k-1$ in
place of $k$: then $P(\delta^{-1}X)=\prod_{i\le k-1}(X-q^i)\cdot S$, and $\Phi(X)=((X-q^{k-1})R)(X_0^{-1}X)$,
so $(x,((X-q^{k-1})R)(X_0^{-1}U)z)=0$ for every $z\in\pi^{I_1}$. (Orthogonality does not depend on the
order of the arguments of the inner product.)
\end{proof}

With this, the proof of Proposition~\ref{pri:dict}(d) given in \S\ref{pax:ss-dict} from
Lemma~\ref{pax:toproot} is complete, and Proposition~\ref{pri:dict} is proved in full (modulo the
published results cited in \S\ref{pax:ss-dict}).

\subsubsection{The two facts of Remark~\ref{pax:rem-BL}}\label{pay:ss-BL}

\begin{lemma}[$\eta_O$ is unramified away from $\ell$]\label{pay:eta}
Let $\eta$ be the algebraic Hecke character $y\mapsto(\chi_W(y/\bar y)\allowbreak \upsilon(y)^{-r'}\allowbreak |y|^{-r'/2})^{-1}$
of Proposition~\ref{pga:trace}, whose $\ell$-adic character is $\eta_O$. Then $\eta$ is trivial on
$\auO_{E,w}^\times$ for every finite place $w$ of $E$. Hence $\eta_O$ is unramified at every finite
place $w\nmid\ell$; in particular at the places of $Q\cup\overline Q$, at the places of $U\cup
\overline U$ in Proposition~\ref{pri:BL}(c), and at the chosen places not above $\ell$.
\end{lemma}
\begin{proof}
By \S\ref{pga:ss-L} and~{\cite[(1)]{P54}}, $\chi_W=a=\upsilon^{r_-}$ is a power of the splitting
character $\upsilon$. By {\cite[\S5.2.1]{P49}}, $\upsilon$ is unramified at every finite place
($\upsilon(x)=\prod_\tau x_\tau/|x_\tau|$ when the finite part of $x$ is a unit). Let $y\in
\auO_{E,w}^\times$ (embedded in $\mathbb A_E^\times$ with components $1$ elsewhere). Then $y/\bar y$ has
unit components at $w$ and $\bar w$ and $1$ elsewhere, so $\chi_W(y/\bar y)=1$; $\upsilon(y)=1$; and
$|y|=1$. So $\eta(y)=1$. For an algebraic Hecke character, the restriction of its $\ell$-adic character
to the inertia group at $w\nmid\ell$ is, through $\mathrm{Art}_w$, the restriction of $\eta_w$ to
$\auO_{E,w}^\times$ (the algebraic part only involves the components above $\ell$ and $\infty$); hence
$\eta_O|_{I_w}=1$. The places of $Q\cup\overline Q$ and of $U\cup\overline U$ do not lie above $\ell$
(they are disjoint from $\pgS(\ell)$, which contains the places above $\ell$, Definition~\ref{pri:def-aux}
and Proposition~\ref{pri:BL}(c)). The chosen places used in the proof of Proposition~\ref{pri:BL}(c) are those prime to $\ell$: if $\ell\in\{p_0,p_1\}$, the chosen places above $\ell$ are excluded, and at them the lemma asserts nothing.
\end{proof}

So (F1) of \S\ref{pax:ss-dict} holds at every finite place not above $\ell$, which is the first point
of Remark~\ref{pax:rem-BL}; it is used, for chosen $w\nmid\ell$, in the paragraph ``Inertia at a chosen $w$'' of the proof of
Proposition~\ref{pri:BL}(c). The rescaling of the label is by the tame part of the character of the
pair only.

\emph{Product labels.} The second point of Remark~\ref{pax:rem-BL} concerns how the group $\auG_O=
\prod_w\auGL_{r'}(\auF_w)$ of the chosen places, which acts on forms on $\auU(W)$ by right
translation (twisted by the tame part of $\mu$, as in {\cite[Proposition~14.27]{P46}}), acts on forms on
$\auU(V)$, $\operatorname{rk}V=n=r'+t'$. The following lemma shows that the single lift of
{\cite[\S2.1]{P54}} intertwines it with right translation by elements reducing to $\mathrm{diag}(u_w,
1_{t'})$; this is the ``product label'' $c_w\cdot(X-1)^{t'}$ used in the proof of
Proposition~\ref{pri:BL}(c).

\begin{lemma}[Equivariance of the single lift]\label{pay:product}
Let $u=(u_w)\in\auG_O$. Choose $k\in K_W$ with reduction $u_w$ at each chosen place $w$ and $1$ at
the places of $Q$, and $k'\in K_V$ with reduction $\mathrm{diag}(u_w,1_{t'})$ at each chosen place
and $1$ at the places of $Q$, and with components $k_v=1$, $k'_v=1$ at every place $v$ that is neither chosen nor in $Q$ (possible as the chosen places are split and reduction is onto). Let
$\Theta$ be the lift of {\cite[\S2.1]{P54}} with the block datum $\Phi_\vartheta$. Then for every form
$F$ on $\auU(W)$ and every $h$,
\[
(\Theta F)(hk')=\prod_{w\ \mathrm{chosen}}\psi_w(\det u_w)^{-1}\cdot\bigl(\Theta(F(\,\cdot\,k))\bigr)(h).
\]
\end{lemma}
\begin{proof}
At a chosen place $w$, by~{\cite[(1)]{P54}} (with $\eta_w$ unramified, hence trivial on $\auGL_{r'}(
\auO_w)\times\auGL_n(\auO_w)$, and $\chi_w(\det k_w)=\psi_w(\det u_w)$),
$\omega_\mu(k_w,k'_w)\Phi_{\vartheta,w}(T)=\psi_w(\det u_w)^{-1}\plone[k_w'^{-1}Tk_w\equiv\iota_0]=
\psi_w(\det u_w)^{-1}\plone[T\equiv k'_w\iota_0k_w^{-1}]$, and $k'_w\iota_0k_w^{-1}\equiv\iota_0$ since
$\mathrm{diag}(u_w,1_{t'})\iota_0=\iota_0u_w$. At the places of $Q$, $k$, $k'$ reduce to $1$, and
$\plone[T\text{ integral}]\psi_w(\det T'\bmod w)$ is invariant under $(k,k')$ with $k\equiv1$, $k'\equiv1$
(the upper block of $k'^{-1}Tk$ is congruent to $T'$), and the factor $\chi_w(\det k_w)^{-1}$ of~{\cite[(1)]{P54}} equals $1$, since $\chi_w$ is tamely ramified at $w\in Q$ (its restriction to $\auO_w^\times$ factors through $\psi_w$, i.e.\ through $\auO_w^\times\to(\auO_w/w)^\times$) and $\det k_w\equiv1\bmod w$. At all other places $k_v=k'_v=1$, so nothing changes there. Hence $\omega_\mu(k,
k')\Phi_\vartheta=\prod_w\psi_w(\det u_w)^{-1}\Phi_\vartheta$, i.e.\ $\omega_\mu(1,k')\Phi_\vartheta=
\prod_w\psi_w(\det u_w)^{-1}\omega_\mu(k^{-1},1)\Phi_\vartheta$. Insert this in $(\Theta F)(hk')=\int
\theta(\omega_\mu(g,hk')\Phi_\vartheta)F(g)\,dg$ and change variables $g\mapsto gk$ (Haar measure).
\end{proof}

Consequently, if the action of $\auG_O$ on forms on $\auU(V)$ is \emph{defined} by right translation
by such $k'$ (twisted by $\prod_w\psi_w(\det u_w)$), then $\Theta$ is $\auG_O$-equivariant, which is
the last sentence of Proposition~\ref{pga:L}, and the image of $e_c$-forms consists of $e_c$-forms
for this action; this is the action $\varrho$ of Definition~\ref{prm:def-level} used in
\S\ref{pax:ss-BL}, with the trivial block $\mathbf 1_{\auGL_{t'}}$ in the second position.

\begin{remark}[The base level at the chosen places]\label{pay:rem-level}
The proof of Proposition~\ref{pri:BL}(c), paragraph ``Inertia at a chosen $w$'', also uses a property
of the base level at the chosen places. The level is fixed in
Definition~\ref{prm:def-level} ($k\equiv\left(\begin{smallmatrix}1_{r'}&*\\0&*\end{smallmatrix}\right)$
modulo a chosen $w$), the lift lands in it (Lemma~\ref{prm:level}), and the property is
Lemma~\ref{prm:label}; see Corollary~\ref{prm:inertia}.
\end{remark}



\providecommand{\auO}{\mathcal O}
\providecommand{\auGL}{\mathrm{GL}}
\providecommand{\auF}{\mathbb F}
\providecommand{\auU}{\mathrm U}
\providecommand{\pgG}{G_E}
\providecommand{\pgtr}{\operatorname{tr}}
\providecommand{\priFitt}{\operatorname{Fitt}}
\providecommand{\prival}{\operatorname{val}}
\providecommand{\priAug}{\mathfrak I}
\providecommand{\priT}{\mathbb T}
\providecommand{\priX}{\mathcal X}
\providecommand{\priJ}{\mathcal J}
\providecommand{\prid}{\mathcal D}
\providecommand{\priN}{\mathcal N}
\providecommand{\prity}{\mathrm{ty}}
\providecommand{\plone}{\mathbf 1}
\providecommand{\prmL}{\Lambda}
\providecommand{\prmT}{\mathfrak T}
\providecommand{\prmK}{\mathfrak K}
\providecommand{\prmval}{\operatorname{vl}}
\providecommand{\prmgen}{\mathrm{gen}}
\providecommand{\prmn}{\mathfrak n}
\providecommand{\prmq}{\mathfrak q}

\subsection{Ribet's method: characters, the norm induction and the base level}\label{prm:sec}

This section proves the following steps of \S\ref{pri:sec}: the paragraph ``The characters'' of the
proof of Proposition~\ref{pri:ribet}, together with the gluing that follows it
(\S\ref{prm:ss-chars}); the induction in Proposition~\ref{pri:norm}(b) (\S\ref{prm:ss-norm}); and the
property of the base level at the chosen places that the proof of Proposition~\ref{pri:BL}(c) uses,
with the definition of that level (\S\ref{prm:ss-level}). \S\ref{prm:ss-blocks} lists the statements
of \S\S\ref{pri:sec}, \ref{pga:sec}, \ref{pse:sec} and \ref{pg:sec} whose proofs are given in later
sections, with the place where each is proved.

\subsubsection{The characters in Ribet's method}\label{prm:ss-chars}

We keep the hypotheses (1)--(7) of Proposition~\ref{pri:ribet} and the notation of its proof:
$\prmgen=(\delta_v)_v\sqcup s_0\sqcup s$ is the finite family of generators, $z_g\in\auO^{\prmgen}$
($g\in G$) are coordinates with $u_x(g)-1=\sum_iz_{g,i}(u_x(\prmgen_i)-1)$ for all $x\in Y$ and
$g\in G$, $c_\zeta=(\zeta_\pi-1)^{-1}$, $\nu_0=c_\zeta(\gamma_\pi-1)$, $\rho_0=\nu_0+(1-\nu_0)(a-1)$. Put
$\prmL=\auO[X_J]$. Each $u_x$ extends $\prmL$-linearly to a $\prmL$-algebra map $u_x\colon\prmL[G]\to
M_\sigma(\prmL)$, and each character $\xi$ of $G$ to a $\prmL$-algebra map $\eta_\xi\colon\prmL[G]\to
\prmL$; we write $\eta_e=\eta_{\xi_e}$. Let $\priAug\subseteq\prmL[G]$ be the augmentation ideal (the
kernel of $\eta_{\mathbf 1}$) and $\prmK=\bigcap_x\ker u_x$, a two-sided ideal.

\emph{Rows.} A row is a vector $r\in\prmL[G]^{\prmgen}$; its value is $\prmval(r)=\sum_ir_i(\prmgen_i-1)
\in\priAug$. The map $\prmval$ is left $\prmL[G]$-linear. For $m\in\auO^{\prmgen}$ let $\bar m\in
\prmL[G]^{\prmgen}$ be the row with entries $m_i\cdot1$ (the \emph{coordinate row}), and $\epsilon_i$
the $i$-th unit row. The \emph{tie ideal} $\prmT$ is the left ideal of $\prmL[G]$ generated by
$\prmK$, by $p\priAug$, by the tied relations $t_k=\tau_k-1-X_jH_j(\phi_k)(\phi_k-1)$, $k=(j,\pm)$,
and by $\priAug\rho_0$. The set $N$ of \emph{tied rows}, those with $\prmval(r)\in\prmT$, is a left
$\prmL[G]$-submodule of $\prmL[G]^{\prmgen}$.

\begin{lemma}[Tied rows]\label{prm:tied}
The following rows are tied:
{\rm(R1)} $\overline{z_{gh}}-\overline{z_g}-g\,\overline{z_h}$ for $g,h\in G$;
{\rm(R2)} $\overline{z_{\prmgen_i}}-\epsilon_i$ for every index $i$;
{\rm(R3)} $\overline{\sum_g\theta_gz_g}$ whenever $\theta\in\auO^{(G)}$ and $\sum_g\theta_g(u_x(g)-1)=0$
for all $x$;
{\rm(R4)} $p\,r$ for every row $r$;
{\rm(R5)} $\overline{z_{\tau_k}}-X_jH_j(\phi_k)\,\overline{z_{\phi_k}}$ for $k=(j,\pm)$;
{\rm(R6)} $(g-1)\bigl(c_\zeta\overline{z_{\gamma_\pi}}+(1-\nu_0)\overline{z_a}\bigr)$ for $g\in G$;
{\rm(R7)} $o_v(\delta_v)\,\epsilon_{v}$, $\epsilon_v$ the unit row of the coordinate of $\delta_v$.
\end{lemma}
\begin{proof}
For $m\in\auO^{\prmgen}$, $u_x(\prmval(\bar m))=\sum_im_i(u_x(\prmgen_i)-1)$; hence
$\prmval(\overline{z_g})-(g-1)\in\prmK$ for every $g$. (R1): modulo the two-sided ideal $\prmK$ the value
is $(gh-1)-(g-1)-g(h-1)=0$. (R2): the value is $\prmval(\overline{z_{\prmgen_i}})-(\prmgen_i-1)\in
\prmK$. (R3): $u_x$ of the value is $\sum_g\theta_g(u_x(g)-1)=0$ for all $x$. (R4): $\prmval(pr)=p\,
\prmval(r)\in p\priAug$. (R5): the value is $t_k+\bigl(\prmval(\overline{z_{\tau_k}})-(\tau_k-1)\bigr)-
X_jH_j(\phi_k)\bigl(\prmval(\overline{z_{\phi_k}})-(\phi_k-1)\bigr)\in\prmT$. (R6): the value of
$c_\zeta\overline{z_{\gamma_\pi}}+(1-\nu_0)\overline{z_a}$ is $\rho_0$ modulo $\prmK$, so the value of the row
is $(g-1)\rho_0$ plus an element of $\prmK$, and $(g-1)\rho_0\in\priAug\rho_0$. (R7): the value is
$o_v(\delta_v)(\delta_v-1)$, which every $u_x$ kills by hypothesis~(6).
\end{proof}

\emph{Reading rows on a character.} For a character $\xi$ let $E_\xi\colon\prmL[G]^{\prmgen}\to
\prmL^{s_0\sqcup s}$, $E_\xi(r)=(\eta_\xi(r_i))_{i\in s_0\sqcup s}$: read the row on $\xi$ and forget
the coordinates of the $\delta_v$. It is $\prmL$-linear, and $E_\xi(\lambda r)=\eta_\xi(\lambda)E_\xi(r)$
for $\lambda\in\prmL[G]$. The forms of \S\ref{pri:ss-fitting} are taken with $n=\prmgen$, the
function $z\colon g\mapsto z_g$, and the maps $\Phi$ below.

\begin{lemma}[Forms of tied rows are character cocycles]\label{prm:forms}
Let $\xi$ be a character of $G$ with $\xi(\gamma_\pi)=1$, let $A\colon\prmL^{s_0\sqcup s}\to\prmL^m$
be $\prmL$-linear, $K\subseteq\prmL^m$ a $\prmL$-submodule containing $A(E_\xi(N))$, and $\Phi=A\circ
E_\xi\circ\overline{(\cdot)}\colon\auO^{\prmgen}\to\prmL^m$. Then $\Phi(p\auO^{\prmgen})\subseteq K$, and
for every $\auO$-linear $\Psi\colon\auO^{\prmgen}\to\auO/p$ vanishing on $\Phi^{-1}(K)$ the function
$c=\Psi\circ z$ is a character cocycle for $\xi$ (relative to the $\delta_v$, $\gamma_\pi$ and $a$).
\end{lemma}
\begin{proof}
If $m\in\auO^{\prmgen}$ and some tied row $r$ has $\eta_\xi(r_i)=m_i$ for all $i\in s_0\sqcup s$, then
$E_\xi(r)=E_\xi(\bar m)$, so $\Phi(m)=A(E_\xi(r))\in K$ and $\Psi(m)=0$. By (R4), $\Phi(pm)\in K$.
Apply this to the rows of Lemma~\ref{prm:tied}. (R1), read on $\xi$, has entries $z_{gh,i}-z_{g,i}-
\xi(g)z_{h,i}$; so $c(gh)=c(g)+\xi(g)c(h)$, and $c$ is a cocycle for $\xi$. (R3) gives $\sum_g\theta_g
c(g)=0$ whenever $\sum_g\theta_g(u_x(g)-1)=0$ for all $x$: $c$ obeys the linear relations. (R2) for the
index of $\delta_v$: the unit row $\epsilon_v$ has no coordinate in $s_0\sqcup s$, so the row reads as
$\overline{z_{\delta_v}}$ and $c(\delta_v)=0$. (R6) read on $\xi$ has entries $(\xi(g)-1)\bigl(c_\zeta z_{
\gamma_\pi,i}+\eta_\xi(1-\nu_0)z_{a,i}\bigr)$, and $\eta_\xi(\nu_0)=c_\zeta(\xi(\gamma_\pi)-1)=0$; so
$(\xi(g)-1)(c_\zeta c(\gamma_\pi)+c(a))=0$ for all $g$, and if $c(\gamma_\pi)=0$ then $(\xi(g)-1)c(a)=0$.
\end{proof}

\begin{lemma}[Common unit rows]\label{prm:unit}
Assume, as in {\rm(2)} and {\rm(7)}, that $\xi_e(\gamma_\pi)=1$ and that $s$ detects every module of
character cocycles for $\xi_e$, for every $e<t'$. Then for every $i\in s_0$ there is a tied row
$\nu_i$ with $\eta_e(\nu_{i,j})=1$ if $j=i$ and $=0$ if $j\in s_0\smallsetminus\{i\}$, for all $e<t'$.
\end{lemma}
\begin{proof}
\emph{One character.} Fix $e$ and let $F\colon\prmL[G]^{\prmgen}\to\prmL^{s_0}$ be $E_{\xi_e}$ followed by
the projection to the coordinates in $s_0$, $K=F(N)$, $\Phi=F\circ\overline{(\cdot)}$. By
Lemma~\ref{prm:forms} (with $A$ the projection and $m=\#s_0$), $\Phi(p\auO^{\prmgen})\subseteq K$ and
the forms $\Psi\circ z$ ($\Psi$ vanishing on $\Phi^{-1}K$) make up a module of character cocycles for
$\xi_e$, which $s$ detects. For $g\in s$, the row (R2) of the index $g$ (in the copy $s$ of
$\prmgen$) is tied and $F(\epsilon_g)=0$, so $\Phi(z_g)\in K$ and every such $\Psi$ has $\Psi(z_g)=0$.
Hence $\Psi\circ z$ vanishes on $s$, so $\Psi\circ z=0$; in particular $\Psi(z_{i})=0$ for the element
$i\in s_0$. By Lemma~\ref{pri:fitt-basic}(c) (whose proof applies verbatim to $\prmL^{s_0}$),
$\Phi(z_i)\in K$. With the row (R2) of the index $i$, $F(\epsilon_i)=\Phi(z_i)-F(\overline{z_i}-
\epsilon_i)\in K$: some tied row $w$ has $\eta_e(w_j)=\delta_{ij}$ for $j\in s_0$.

\emph{All characters.} Let $E'$ be a set of indices $e$ for which a tied row $w$ with $\eta_{e}(w_j)=
\delta_{ij}$ ($j\in s_0$, $e\in E'$) has been found, and $e_1\notin E'$. By the first part choose tied
rows $m_{j}$ ($j\in s_0$) with $\eta_{e_1}(m_{j,j'})=\delta_{jj'}$ for $j'\in s_0$, and put
$w'=w+\sum_{j\in s_0}(\delta_{ij}-w_j)\,m_j$, a tied row ($N$ is a left module). For $j'\in s_0$,
$\eta_{e_1}(w'_{j'})=\eta_{e_1}(w_{j'})+\sum_j(\delta_{ij}-\eta_{e_1}(w_j))\delta_{jj'}=\delta_{ij'}$,
and for $e\in E'$, $\eta_e(\delta_{ij}-w_j)=0$, so $\eta_e(w'_{j'})=\eta_e(w_{j'})=\delta_{ij'}$.
Starting from $E'=\emptyset$ and $w=0$, after $t'$ steps we get $\nu_i$.
\end{proof}

Fix such rows $\nu_i$. For $e<t'$ let $A_e\colon\prmL^{s_0\sqcup s}\to\prmL^s$, $A_e(b)_g=b_g-\sum_{i\in
s_0}b_i\,\eta_e(\nu_{i,g})$ ($g\in s$), and $K_e=A_e(E_{\xi_e}(N))\subseteq\prmL^s$. Let $\prmn\subseteq
\prmL$ be the $\auO$-span of the monomials that do not divide $X_J$; it is an ideal (a multiple of such
a monomial is again such a monomial), and $f\in\prmn$ if and only if no monomial of $f$ divides $X_J$.
For $e<t'$ put $X_{<e}=\prod_{\prity(i)<e}X_i$, $X_{=e}=\prod_{\prity(i)=e}X_i$ and $\prmq_e=(\prmn:
X_{<e})=\{f:fX_{<e}\in\prmn\}$.

\begin{lemma}\label{prm:fittsum}
For a $\prmL$-submodule $K\subseteq\prmL^s$ and an ideal $\prmq$ of $\prmL$, $\priFitt(K+\prmq\prmL^s)
\subseteq\priFitt(K)+\prmq$.
\end{lemma}
\begin{proof}
A generator of the left side is $\det$ of a matrix with rows $k_l+q_l$, $k_l\in K$, $q_l\in\prmq
\prmL^s$. Expanding multilinearly in the rows, the term with all rows $k_l$ lies in $\priFitt(K)$, and
every other term has a row with entries in $\prmq$, so lies in $\prmq$.
\end{proof}

\begin{proposition}[The block of one character]\label{prm:block}
For every $e<t'$ there is $\Gamma_e\in\prmL$ with
$\prival(\Gamma_e(0))\le b_e+(\#\{\prity=e\}+\#s)\prival(C)$ and $\Gamma_eX_{=e}\in\priFitt(K_e+
\prmq_e\prmL^s)$.
\end{proposition}
\begin{proof}
Apply Proposition~\ref{pri:fitt-triangle} with $R=\prmL$, the set $s$, $N_e:=K_e+\prmq_e\prmL^s$ in the
role of $N$, $\Phi=A_e\circ E_{\xi_e}\circ\overline{(\cdot)}$, $z$ as above, $I=\{j:\prity(j)=e\}$,
$\tau_j=\tau_{(j,+)}$, $\tau'_j=\tau_{(j,-)}$ and likewise $\phi$. We check its hypotheses, writing
$y_g=\Phi(z_g)$. (i) $\Phi(p\auO^{\prmgen})\subseteq N_e$ by (R4). (ii) For $g\in s$, the row (R2) of
the index $g$ is tied and $A_e(E_{\xi_e}(\epsilon_g))=\epsilon_g$ (it has no coordinate in $s_0$), so
$y_g-\epsilon_g\in K_e$. (iii) By (3), $H_j(\xi_e(\phi_k))=1+pc$ with $c\in\auO$. The sum of the tied
rows (R5) and $p\,X_jc\,\overline{z_{\phi_k}}$ (R4) reads on $\xi_e$ as $E_{\xi_e}(\overline{z_{\tau_k}})
-X_j(1+pc)E_{\xi_e}(\overline{z_{\phi_k}})+pX_jcE_{\xi_e}(\overline{z_{\phi_k}})$, since $\eta_e(X_j
H_j(\phi_k))=X_jH_j(\xi_e(\phi_k))$; applying $A_e$, $y_{\tau_k}-X_jy_{\phi_k}\in K_e$ for $k=(j,\pm)$.
(iv) Let $i<j$ with $j\in I$. As $\prity$ is monotone, $\prity(i)\le e$; if $\prity(i)=e$ then $i\in I$;
otherwise $X_iX_{<e}$ is divisible by $X_i^2$, so it lies in $\prmn$, $X_i\in\prmq_e$ and
$X_i\prmL^s\subseteq N_e$. (v) By Lemma~\ref{prm:forms} (with $A=A_e$, $K=N_e\supseteq A_e(E_{\xi_e}(N))$
and $\xi_e(\gamma_\pi)=1$ by (2)), $\mathrm F(N_e)$ is a module of character cocycles for $\xi_e$, so by
(7) it has triangle cost $b_e$ along $I$ with constant $C$. Proposition~\ref{pri:fitt-triangle} gives
$\Gamma_e$ as stated.
\end{proof}

\begin{proposition}[Character rows]\label{prm:chars}
There are, for each $e<t'$, finitely many matrices $B_{e,i}\in M_s(\prmL)$ and weights $w_{e,i}\in
\prmL$ ($1\le i\le l_e$), elements $\zeta_e\in\prmq_e$ and $\Gamma_e$ as in
Proposition~\ref{prm:block}, such that
\[
d_e:=\sum_iw_{e,i}\det B_{e,i}=\Gamma_eX_{=e}-\zeta_e ,
\]
and, for each choice $\mathbf i=(i(e))_{e<t'}$, a square matrix $M_{\mathbf i}$ over $\prmL[G]$,
indexed by $\prmgen$, all of whose rows are tied, with
\[
\det\eta_e(M_{\mathbf i})=c_{\mathrm{ch}}\det B_{e,i(e)}\quad(e<t'),\qquad c_{\mathrm{ch}}=C^{\#s}\prod_vo_v(1).
\]
\end{proposition}
\begin{proof}
\emph{The determinants.} By Proposition~\ref{prm:block} and Lemma~\ref{prm:fittsum}, $\Gamma_eX_{=e}=
d_e+\zeta_e$ with $d_e\in\priFitt(K_e)$ and $\zeta_e\in\prmq_e$. By definition of the Fitting ideal,
$d_e=\sum_iw_{e,i}\det B_{e,i}$ with matrices $B_{e,i}$ all of whose rows lie in $K_e$.

\emph{Lifting the rows.} Let $b$ be a row of some $B_{e,i}$: $b=A_e(E_{\xi_e}(r))$ with $r\in N$. Put
$r'=r-\sum_{i'\in s_0}r_{i'}\nu_{i'}$, a tied row. For $j\in s_0$, $\eta_e(r'_j)=\eta_e(r_j)-
\sum_{i'}\eta_e(r_{i'})\delta_{i'j}=0$; for $g\in s$, $\eta_e(r'_g)=\eta_e(r_g)-\sum_{i'}\eta_e(r_{i'})
\eta_e(\nu_{i',g})=b_g$. So every row $b$ of $B_{e,i}$ is $E_{\xi_e}(r')|_s$ for a tied row $r'=:r'(e,i,b)$
whose coordinates in $s_0$ vanish on $\xi_e$.

\emph{Separating the characters.} By (4) and Lemma~\ref{prm:sep} below there are $q_e\in\auO[G]$ with
$\xi_{e'}(q_e)=C$ if $e'=e$ and $0$ otherwise.

\emph{The matrix.} Fix $\mathbf i$. Let $M_{\mathbf i}$ have: in the row of $\delta_v$, the row (R7)
$o_v(\delta_v)\epsilon_v$; in the row of $i\in s_0$, $\nu_i$; in the row of $g\in s$, $\sum_{e<t'}q_e\,
r'(e,i(e),b_{e,g})$, $b_{e,g}$ the row $g$ of $B_{e,i(e)}$. All rows are tied. Apply $\eta_e$ and order
the columns as $(\delta,s_0,s)$. The rows of the $\delta_v$ have the single entry $o_v(\xi_e(\delta_v))=
o_v(1)$ in the column of $\delta_v$ (by (6), $\xi_e(\delta_v)=1$). The rows of $s_0$ have the identity
in the columns of $s_0$. The row of $g\in s$ is $\sum_{e'}\xi_e(q_{e'})\eta_e(r'(e',\dots))=
C\,\eta_e(r'(e,i(e),b_{e,g}))$, which is $0$ in the columns of $s_0$ and $C\,b_{e,g}$ in those of $s$. So
$\eta_e(M_{\mathbf i})$ is block lower triangular with respect to $(\delta)\sqcup(s_0\sqcup s)$, with
diagonal blocks $\mathrm{diag}(o_v(1))$ and $\left(\begin{smallmatrix}1&*\\0&CB_{e,i(e)}\end{smallmatrix}
\right)$, and its determinant is $\prod_vo_v(1)\cdot C^{\#s}\det B_{e,i(e)}$.
\end{proof}

\begin{lemma}\label{prm:sep}
Let $\xi_0,\dots,\xi_{t'-1}$ be characters of $G$, $a\in G$, $e<t'$ and $C\in\auO$ a multiple of
$\prod_{e'\ne e}(\xi_e(a)-\xi_{e'}(a))$. Then $q_e=w\prod_{e'\ne e}(a-\xi_{e'}(a))\in\auO[G]$, where
$C=w\prod_{e'\ne e}(\xi_e(a)-\xi_{e'}(a))$, has $\xi_{e''}(q_e)=C$ for $e''=e$ and $0$ for $e''\ne e$.
\end{lemma}
\begin{proof}
$\xi_{e''}(q_e)=w\prod_{e'\ne e}(\xi_{e''}(a)-\xi_{e'}(a))$; for $e''\ne e$ the factor $e'=e''$ vanishes.
\end{proof}

\begin{lemma}[Coefficients below $X_J$]\label{prm:coeff}
Let $d_e,\zeta_e,\Gamma_e\in\prmL$ ($e<t'$) with $d_e+\zeta_e=\Gamma_eX_{=e}$ and $\zeta_e\in\prmq_e$,
and let $F\in\prmL$. Then $F\prod_ed_e\equiv F\prod_e\Gamma_e\cdot X_J\pmod\prmn$. Consequently no
monomial properly dividing $X_J$ occurs in $F\prod_ed_e$, and its coefficient of $X_J$ is the constant
term of $F\prod_e\Gamma_e$.
\end{lemma}
\begin{proof}
Order the $e$ increasingly. Telescoping, $\prod_e(d_e+\zeta_e)-\prod_ed_e=\sum_e\zeta_e\prod_{e'<e}(d_{e'}+
\zeta_{e'})\prod_{e'>e}d_{e'}$. In the $e$-th term, $\prod_{e'<e}(d_{e'}+\zeta_{e'})=\prod_{e'<e}
\Gamma_{e'}\cdot X_{<e}$ (the sets $\{\prity=e'\}$, $e'<e$, partition $\{\prity<e\}$), and $\zeta_eX_{<e}
\in\prmn$; so each term lies in the ideal $\prmn$. Also $\prod_e(d_e+\zeta_e)=\prod_e\Gamma_e\cdot X_J$.
Multiply by $F$. An element of $\prmn$ has no monomial dividing $X_J$; every monomial of $F\prod_e
\Gamma_e\cdot X_J$ is a multiple of $X_J$, and the only one dividing $X_J$ is $X_J$, with coefficient the
constant term of $F\prod_e\Gamma_e$.
\end{proof}

\begin{proof}[Proof of the paragraphs ``The characters'' and ``Gluing'' of Proposition~\ref{pri:ribet}]
Let $B_0$ be the matrix of tied rows of the paragraph ``The block'', with $\prival(\det\rho_\pi(B_0)
|_{X_J=0})\le b_\rho$ (hypothesis (5) of Proposition~\ref{pri:ribet}), and $B_{e,i}$, $w_{e,i}$, $\zeta_e$, $\Gamma_e$, $M_{\mathbf i}$ as in
Proposition~\ref{prm:chars}. Let $k\in\auO[G]$ be $1$ on every $\xi_e$ and $0$ on every $\rho_x$,
$x\in\priX$: for each $x$, $k_x=(\zeta_x-1)^{-1}(\zeta_x-\gamma_x)$ is $0$ on $\rho_x$ (where $\gamma_x$
is the scalar $\zeta_x$) and $1$ on the $\xi_e$ (where $\gamma_x$ is $1$), by (2); take $k=\prod_xk_x$. Let
$q_0$ be as in Lemma~\ref{prm:sep} for $e=0$. For each $\mathbf i$ apply
Proposition~\ref{pri:ribet-elt} (with $\auO$ replaced by $\prmL=\auO[X_J]$, each $u_x$ extended $\prmL$-linearly and $\mathsf L$ by $\mathsf L[X_J]$: its proof uses only Lemma~\ref{pri:amitsur} and the Cayley--Hamilton theorem, both valid over any commutative ring, and $S\supseteq\prmL$ since $\priJ$ is a trace module) to the matrix $(1-k)B_0+kM_{\mathbf i}$ (its rows are tied, $N$ being a
left module), with $\rho_1=kq_0$ and $\mu_0=C$: $kq_0$ acts as $C$ on $\xi_0$, as $0$ on the other
$\xi_e$ and as $0$ on every $\rho_x$, and $\rho_0$ acts as $0$ on $\xi_0=1$ ($\eta_0(\nu_0)=0$,
$\eta_0(a-1)=0$). It gives $f_{\mathbf i}\in S$ with $f_{\mathbf i}\det u_x(w)\in\priJ$ for $w\in
\priAug$ and, for $x\in\priX$, $f_{\mathbf i}(x)=C\det\rho_x(B_0)\prod_e\det\eta_e(M_{\mathbf i})$, because
on $\rho_x$ the matrix is $\rho_x(B_0)$ and on $\xi_e$ it is $\eta_e(M_{\mathbf i})$. Put $f=\sum_{\mathbf
i}\bigl(\prod_ew_{e,i(e)}\bigr)f_{\mathbf i}$ (the weights lie in $\prmL\subseteq S$). For $x\in\priX$,
\[
f(x)=C\det\rho_x(B_0)\sum_{\mathbf i}\prod_ew_{e,i(e)}\,c_{\mathrm{ch}}\det B_{e,i(e)}=C\,c_{\mathrm{ch}}^{t'}\det\rho_x(B_0)
\prod_{e<t'}d_e .
\]
By Lemma~\ref{prm:coeff} with $F=Cc_{\mathrm{ch}}^{t'}\det\rho_x(B_0)$, $f(x)$ has no monomial properly dividing
$X_J$ for every $x\in\priX$, and the coefficient of $X_J$ in $f(\pi)$ is
$d=C\,c_{\mathrm{ch}}^{t'}\det\rho_\pi(B_0)|_{X_J=0}\prod_e\Gamma_e(0)\in\auO$. Its valuation is at most
\[
\prival(C)+t'\Bigl(\#s\,\prival(C)+\sum_v\prival(o_v(1))\Bigr)+b_\rho+\sum_e\bigl(b_e+(\#\{\prity=e\}+\#s)
\prival(C)\bigr),
\]
which is the bound of Proposition~\ref{pri:ribet}, as $\sum_e\#\{\prity=e\}=\#J$. The last sentence of
the proof of Proposition~\ref{pri:ribet}, the identity $\det(1-u_x(g))=\det u_x(w)$ for $w=-(g-1)\in\priAug$, which does not involve the two paragraphs proved here, shows that $f$ is a common Ribet element.
\end{proof}

\begin{remark}\label{prm:rem-chars}
Two points of the paragraphs in the proof of Proposition~\ref{pri:ribet} deserve emphasis. (i) In
``Gluing'' the element $\rho_1$ must vanish on every $\rho_x$, as Proposition~\ref{pri:ribet-elt}
requires; $q_0$ alone (the element of Lemma~\ref{prm:sep} for $e=0$) does not, and $\rho_1=kq_0$ does. (ii) The vanishing of the character cocycles at the
$\delta_v$ comes from forgetting the coordinates of the $\delta_v$ (the row (R2)); the cost
$\prival(o_v(1))$ enters only through the rows (R7) that complete $M_{\mathbf i}$, not through
Lemma~\ref{pri:fitt-basic}(b).
\end{remark}

\subsubsection{The induction in Proposition~\ref{pri:norm}(b)}\label{prm:ss-norm}

We keep the notation of \S\ref{pri:ss-norm}. Let $R$ be the ring of integral operators on the forms
of principal level at the $w_j$, and $\sigma\colon R\to R'$ a ring map (in Proposition~\ref{pri:norm},
a reduction modulo $\ell^\nu/(2n!)$). For $j\in J$ let $R_j\subseteq R$ be the subring generated by
the operators made of translates at $w_j$; it contains $\Gamma_j$, $\mathrm{Nm}_j$ and the operator
$\partial_j$ of Lemma~\ref{pri:oneplace}. As recalled in the proof of Proposition~\ref{pri:norm},
$R_j$ and $R_{j'}$ commute for $j\ne j'$, and $\priT$ commutes with every $R_j$. For $S\subseteq J$ put
\[
\priN_S=\prod_{j\in S}\partial_j\prod_{j\notin S}\mathrm{Nm}_j ,
\]
the coefficient of $X_S$ in $\prod_j(\mathrm{Nm}_j+X_j\partial_j)$; the order of the factors does not
matter. We use three facts.
\begin{enumerate}
\item[(N1)] $\priN_\emptyset=\mathrm{Nm}$: expanding $\prod_j\sum_{k_j\in\Gamma_j}k_j$ gives $\sum_{k\in
\Gamma}k$, as $\Gamma=\prod_j\Gamma_j$ and translates at different places multiply as in $\Gamma$.
\item[(N2)] For $j\notin S$, $\priN_S=\priN'\mathrm{Nm}_j$ and $\priN_{S\cup\{j\}}=\priN'\partial_j$ with
$\priN'=\prod_{i\in S}\partial_i\prod_{i\notin S\cup\{j\}}\mathrm{Nm}_i$ (move the factor at $j$ to the
right).
\item[(N3)] $\delta'\in\priT$ commutes with every $\priN_S$, with $\delta=c_0^2\delta'$ and with every
$k\in\Gamma$, hence with $\mathrm{Nm}$.
\end{enumerate}

\begin{lemma}[From the factors to the group]\label{prm:invariant}
Let $\mathsf A\in R$ with $\sigma(\mathsf A(k-1))=0$ for every $k$ in every factor $\Gamma_j$, and let
$\Psi$ be as in Lemma~\ref{pri:forms}(a) for $\Gamma'=\Gamma$, so that $\sum_{k\in\Gamma}k^{-1}\Psi k=
c_0$. Then $\sigma(c_0\mathsf A)=\sigma(\mathsf A\Psi\,\mathrm{Nm})$.
\end{lemma}
\begin{proof}
The set $H=\{k\in\Gamma:\sigma(\mathsf A)\sigma(k)=\sigma(\mathsf A)\}$ contains every $\Gamma_j$ by
hypothesis and is closed under products: if $k,k'\in H$ then $\sigma(\mathsf A)\sigma(kk')=
\sigma(\mathsf A)\sigma(k)\sigma(k')=\sigma(\mathsf A)\sigma(k')=\sigma(\mathsf A)$. A subset of a finite
group closed under products and containing generators is the whole group, so $H=\Gamma$; in particular
$k^{-1}\in H$ for every $k$. Hence
$\sigma(c_0\mathsf A)=\sum_k\sigma(\mathsf A)\sigma(k^{-1})\sigma(\Psi k)=\sum_k\sigma(\mathsf A)\sigma(\Psi
k)=\sigma(\mathsf A\Psi\mathrm{Nm})$.
\end{proof}

\begin{proposition}[Proposition~\ref{pri:norm}(b)]\label{prm:normb}
Under the hypotheses of Proposition~\ref{pri:norm}, for every $S\subseteq J$ and every integer $i\ge
\#S$ there is $Y\in R$ with $\sigma(\delta^i\priN_S)=\sigma(Y\,\mathrm{Nm})$, and $Y=\delta^i$ if
$S=\emptyset$. In particular, with $i=\#J$, there are $Y_S$ with $\sigma(\delta^{\#J}\priN_S)=
\sigma(Y_S\mathrm{Nm})$ and $Y_\emptyset=\delta^{\#J}$.
\end{proposition}
\begin{proof}
Induction on $\#S$. For $S=\emptyset$, $\priN_\emptyset=\mathrm{Nm}$ by (N1); take $Y=\delta^i$.

Let $S\ne\emptyset$ and $i\ge\#S$; write $i=i'+1$, so $i'\ge\#S-1$. Put $\mathsf A=c_0\,\delta'\,
\delta^{i'}\priN_S$. We show $\sigma(\mathsf A(k-1))=0$ for $k$ in each factor $\Gamma_{j'}$.

\emph{Case $j'\notin S$.} By (N2), $\priN_S=\priN'\mathrm{Nm}_{j'}$, and $\mathrm{Nm}_{j'}k=\mathrm{Nm}_{j'}$
for $k\in\Gamma_{j'}$. So $\mathsf A(k-1)=0$.

\emph{Case $j'\in S$.} Let $S'=S\smallsetminus\{j'\}$. By (N2) for $S'$ and $j'$, $\priN_{S'}=\priN'
\mathrm{Nm}_{j'}$ and $\priN_S=\priN'\partial_{j'}$. Lemma~\ref{pri:oneplace} at $w_{j'}$ (whose
hypothesis is that of Proposition~\ref{pri:norm} at $w_{j'}$) gives $Z\in R$ with $\sigma(c_0\partial_{j'}
(k-1))=\sigma(\mathrm{Nm}_{j'}Z)$ and $\delta'\mathrm{Nm}Z=0$. Hence
\[
\sigma(\mathsf A(k-1))=\sigma(\delta'\delta^{i'}\priN')\,\sigma(c_0\partial_{j'}(k-1))=
\sigma(\delta'\delta^{i'}\priN'\mathrm{Nm}_{j'}Z)=\sigma(\delta'\delta^{i'}\priN_{S'}Z).
\]
By (N3), $\delta'\delta^{i'}\priN_{S'}=\delta^{i'}\priN_{S'}\delta'$. As $\#S'=\#S-1\le i'$, the induction
hypothesis gives $Y'$ with $\sigma(\delta^{i'}\priN_{S'})=\sigma(Y'\mathrm{Nm})$. Therefore, using (N3)
once more,
\[
\sigma(\mathsf A(k-1))=\sigma(Y')\,\sigma(\mathrm{Nm}\,\delta'Z)=\sigma(Y')\,\sigma(\delta'\mathrm{Nm}Z)=0 .
\]

So Lemma~\ref{prm:invariant} applies to $\mathsf A$: $\sigma(c_0\mathsf A)=\sigma(\mathsf A\Psi
\mathrm{Nm})$. Since $c_0\mathsf A=c_0^2\delta'\delta^{i'}\priN_S=\delta^{i}\priN_S$, the element $Y=
\mathsf A\Psi$ works.
\end{proof}

\begin{remark}\label{prm:rem-norm}
(i) The induction must be run for all $S$ simultaneously and, at each step, at every place
$j'\in S$, not only at one chosen $j$: this is why the claim is proved for every $i\ge\#S$ (the
exponent drops by one when one passes from $S$ to $S\smallsetminus\{j'\}$). (ii) The operators $Y_S$
need not commute with $\priT$ ($\Psi$ is not a Hecke operator); the consequence drawn in
Proposition~\ref{pri:norm} uses only that the $\priN_S$ commute with $\priT$, and that $\delta$
commutes with them; accordingly Proposition~\ref{pri:norm} asserts the commutation for the
$\priN_S$ only.
\end{remark}

\subsubsection{The base level at the chosen places}\label{prm:ss-level}

Let $O=(W,V,\mu,Q)$ be a pair with $r'=n_W$, $n=n_V$, $t'=n-r'$, lifted by the single lift $\Theta$ of
{\cite[\S2.1]{P54}} with the block datum $\Phi_\vartheta$. At a split place $w$ we identify $\auU(V)$
with $\auGL_n(E_w)$ in the coordinates of {\cite[\S2.1]{P54}}, in which $W$ corresponds to the first $r'$
basis vectors ($\iota_0=\left(\begin{smallmatrix}1_{r'}\\0\end{smallmatrix}\right)$ in the notation of
the proof of Lemma~\ref{pay:product}); the places of $Q$ and the chosen places are split. For a
matrix $k$ we write $k^{(11)}$ for its upper left $r'\times r'$ block, and $K_v$ for the stabiliser of
the standard lattice at $v$.

\begin{definition}[Base level]\label{prm:def-level}
The \emph{base level} of $O$ is the open compact subgroup $J_O=\prod_vJ_{O,v}$ of $\auU(V)(\mathbb A_f)$
with $J_{O,v}=K_v$ at the places not below a chosen place or a place of $Q$, and
\[
\begin{gathered}J_{O,w}=\Bigl\{k\in\auGL_n(\auO_w):k\equiv\begin{pmatrix}1_{r'}&*\\0&*\end{pmatrix}\bmod w\Bigr\}\
(w\text{ chosen}),\\
J_{O,w}=\Bigl\{k\in\auGL_n(\auO_w):k\equiv\begin{pmatrix}*&0\\ *&*\end{pmatrix}\bmod w\Bigr\}\ (w\in Q),\end{gathered}
\]
blocks of sizes $(r',t')$. Its \emph{level character} is $\chi_O(k)=\prod_{w\in Q}\psi_w\bigl(\det
k_w^{(11)}\bmod w\bigr)^{-1}$, $\psi_w$ the tame part of $\mu$ at $w$; it is a character because the
reductions at $w\in Q$ are block lower triangular, so the upper left block of a product is the product
of the upper left blocks. $\mathcal F(O)$ is the space of algebraic modular forms $y$ on $\auU(V)$ of
the weight of $O$ with $y(hk)=\chi_O(k)y(h)$ for $k\in J_O$ (the \emph{forms of base level}). For
$u=(u_w)\in\auG_O$ choose $k'_u\in\prod_vK_v$ reducing to $\mathrm{diag}(u_w,1_{t'})$ at each chosen
$w$ and to $1$ at the places of $Q$, and put $\varrho(u)y=\prod_{w\ \mathrm{chosen}}\psi_w(\det u_w)\,
y(\,\cdot\,k'_u)$.
\end{definition}

\begin{lemma}[The level and the action]\label{prm:level}
{\rm(a)} Each $k'_u$ normalises $J_O$ and $\chi_O(k_u'^{-1}kk'_u)=\chi_O(k)$; $\varrho(u)$ preserves
$\mathcal F(O)$, does not depend on the choice of $k'_u$, and $\varrho$ is a representation of $\auG_O$
on $\mathcal F(O)$.
{\rm(b)} For every form $F$ on $\auU(W)$ in the domain of $\Theta$, $\Theta F$ satisfies $(\Theta F)(hk)=
\chi_O(k)(\Theta F)(h)$ for $k\in J_O$; and for $u\in\auG_O$, $\varrho(u)(\Theta F)=\Theta\bigl(F(\,\cdot\,k_u)\bigr)$, $k_u\in K_W$
reducing to $u_w$ at the chosen places and to $1$ at $Q$; so $\Theta$ maps the $\sigma_c$-isotypic
part for (untwisted) right translation on $\auU(W)$ into the $\sigma_c$-isotypic part for $\varrho$.
\end{lemma}
\begin{proof}
(a) At a chosen $w$, $\mathrm{diag}(u_w,1)^{-1}\left(\begin{smallmatrix}1&B\\0&D\end{smallmatrix}\right)
\mathrm{diag}(u_w,1)=\left(\begin{smallmatrix}1&u_w^{-1}B\\0&D\end{smallmatrix}\right)$; at $w\in Q$,
$k'_u\equiv1$; elsewhere $k'_u\in K_v=J_{O,v}$. So $k'_u$ normalises $J_O$, and $\chi_O$, which only reads
the components at $Q$, where conjugation by $k'_u$ is trivial modulo $w$, is unchanged. If $y\in
\mathcal F(O)$ and $k\in J_O$, then $y(hkk'_u)=y(hk'_u\cdot k_u'^{-1}kk'_u)=\chi_O(k)y(hk'_u)$, so
$\varrho(u)y\in\mathcal F(O)$. Two choices of $k'_u$ differ by an element of $\prod_vK_v$ reducing to $1$
at the chosen places and at $Q$; it lies in $J_O$ and $\chi_O$ is $1$ on it, so the translates agree.
Finally $k'_uk'_{u'}$ is a valid choice for $uu'$, and the factors $\psi_w(\det u_w)$ are
multiplicative; so $\varrho(u)\varrho(u')=\varrho(uu')$.

(b) We use $(\Theta F)(hk)=\int\theta(\omega_\mu(g,hk)\Phi_\vartheta)F(g)\,dg$ and $\omega_\mu(g,hk)=
\omega_\mu(g,h)\omega_\mu(1,k)$, as in the proof of Lemma~\ref{pay:product}; it suffices to show
$\omega_\mu(1,k)\Phi_\vartheta=\chi_O(k)\Phi_\vartheta$ for $k\in J_O$, place by place. By
{\cite[(1)]{P54}}, for $k_w\in\auGL_n(\auO_w)$, $\omega_\mu(1,k_w)\Phi_{\vartheta,w}(T)=\Phi_{\vartheta,w}
(k_w^{-1}T)$ ($\eta_w$ is unramified, hence trivial on $\auGL_n(\auO_w)$). At a chosen $w$,
$\Phi_{\vartheta,w}=\plone[T\equiv\iota_0]$ (proof of Lemma~\ref{pay:product}); for $k_w\in J_{O,w}$ the
first $r'$ columns of $k_w$ reduce to those of $1_n$, so $k_w\iota_0\equiv\iota_0$ and $k_w^{-1}T\equiv
\iota_0\iff T\equiv k_w\iota_0\equiv\iota_0$: the datum is fixed. At $w\in Q$, $\Phi_{\vartheta,w}(T)=
\plone[T\text{ integral}]\,\psi_w(\det T^{(1)}\bmod w)$, $T^{(1)}$ the upper $r'\times r'$ block of
$T\in M_{n\times r'}$ (proof of Lemma~\ref{pay:product}); for $k_w\equiv\left(\begin{smallmatrix}a&0\\
c&d\end{smallmatrix}\right)$, $k_w^{-1}\equiv\left(\begin{smallmatrix}a^{-1}&0\\ *&d^{-1}\end{smallmatrix}
\right)$, so the upper block of $k_w^{-1}T$ is $a^{-1}T^{(1)}$ modulo $w$ and $\Phi_{\vartheta,w}(k_w^{-1}T)=
\psi_w(\det a)^{-1}\Phi_{\vartheta,w}(T)$. Elsewhere $k_w\in K_v$ fixes the vacuum ((W4) of
{\cite[Lemma~5.2]{P49}}). The product of the factors is $\chi_O(k)$. The equivariance: by Lemma~\ref{pay:product},
$(\Theta F)(hk'_u)=\prod_{w\ \mathrm{chosen}}\psi_w(\det u_w)^{-1}\,\Theta(F(\,\cdot\,k_u))(h)$, and the
factor $\prod\psi_w(\det u_w)$ in the definition of $\varrho$ cancels it, giving
$\varrho(u)(\Theta F)=\Theta(F(\,\cdot\,k_u))$; this does not depend on the choice of $k'_u$ by (a).
\end{proof}

So the space of forms produced by the lift is contained in $\mathcal F(O)$ of
Definition~\ref{prm:def-level}, and $e_c$ is the central idempotent of the family of labels $c$ for
$\varrho$; this is the precise form of the ``base level'' of Proposition~\ref{pri:BL} and of
\S\ref{pax:ss-BL}.

\emph{Twisted labels.} Let $w$ be chosen, $\auF_w=\auF_q$, and $c$ a label of degree $r'$ over
$\auF_w$ with factorisation $c=\prod_kf_k$ into irreducibles of degrees $d_k$. By
\S\ref{pga:ss-local}, $\sigma_c$ is a constituent of the Harish-Chandra induction from the cuspidal
representation $\bigotimes_k\sigma_{f_k}$ of $L_1=\prod_k\auGL_{d_k}(\auF_q)$. For a character $\phi$ of
$\auF_q^\times$ put $\lambda_k=\sigma_{f_k}\otimes(\phi\circ\det)$. Each $\lambda_k$ is irreducible
cuspidal (a vector fixed by the unipotent radical of a proper parabolic for $\lambda_k$ is fixed by it
for $\sigma_{f_k}$, as $\phi\circ\det$ is trivial on unipotent elements), and $\sigma_c\otimes(\phi
\circ\det)$ is a constituent of the Harish-Chandra induction of $\bigotimes_k\lambda_k$ (for a
character $\chi$ of the group, $\operatorname{Ind}_P(\tau)\otimes\chi\simeq\operatorname{Ind}_P(\tau
\otimes\chi|_P)$, and $\phi\circ\det$ restricted to a parabolic is inflated from its Levi). The
\emph{rescaled label} $\tilde c_w$ of $c_w$ is the label with $\sigma_{\tilde c_w}\simeq\sigma_{c_w}
\otimes(\psi_w^{-1}\circ\det)$; with $\phi=\psi_w^{-1}$, its irreducible factors $\tilde f_k$ have
$\sigma_{\tilde f_k}=\lambda_k$. (In terms of Deligne--Lusztig characters, $R_T(\theta)\otimes(\phi\circ
\det)=R_T(\theta\cdot\phi\circ\det|_T)$, so the roots of $\tilde c_w$ are those of $c_w$ multiplied by the
element of $\overline{\auF}_w^\times$ that names $\psi_w^{-1}$; this is the rescaling of the lift setting.
Only the cuspidal representations $\lambda_k$ are used below.)

\begin{lemma}[Product labels]\label{prm:label}
Let $c$ be a family of labels, $y\in\mathcal F(O)$ with $e_cy=y\ne0$, and $w$ a chosen place. Let
$P\subseteq G_w=\auGL_n(\auF_w)$ be the parabolic subgroup of block upper triangular matrices with
blocks $(r',t')$, $M\simeq\auGL_{r'}(\auF_w)\times\auGL_{t'}(\auF_w)$ its Levi subgroup, $\phi=
\psi_w^{-1}$, $\lambda_k$ as above for $c_w$, $L=\prod_k\auGL_{d_k}(\auF_w)\times\auGL_1(\auF_w)^{t'}$ and
$\lambda=\bigotimes_k\lambda_k\otimes\mathbf 1^{\otimes t'}$.
{\rm(a)} Let $V_w$ be the space of forms $y'$ of the weight of $O$ with $y'(hk)=\chi_O(k)y'(h)$ for $k$
in $J_O^{(w)}K_w(1)$, where $J_O^{(w)}$ is $J_O$ with the component at $w$ replaced by $1$ and $K_w(1)=
1+\varpi_wM_n(\auO_w)$. Then $y\in V_w$, $\auGL_n(\auO_w)$ acts on $V_w$ by right translation through
$G_w$, and every irreducible constituent of the $G_w$-module generated by $y$ has cuspidal support
$(L,\lambda)$ up to conjugacy.
{\rm(b)} For every irreducible automorphic representation $\Pi$ of $\auU(V)$ on which the orthogonal
projection of $y$ is non-zero, $\Pi_w^{K_w(1)}$ contains an irreducible representation of $G_w$ with
cuspidal support $(L,\lambda)$.
\end{lemma}
\begin{proof}
(a) $K_w(1)\subseteq J_{O,w}$ and $\chi_O$ is trivial on the component at $w$, so $y\in V_w$; $K_w(1)$
is normal in $\auGL_n(\auO_w)$ and acts trivially on $V_w$, so $\auGL_n(\auO_w)$ acts through
$G_w$. Let $U$ be the unipotent radical of $P$. The reduction of $J_{O,w}$ is $U\cdot(1\times
\auGL_{t'}(\auF_w))$, so $y$ is fixed by $U$ and by $1\times\auGL_{t'}(\auF_w)$.

For $g\in\auGL_{r'}(\auF_w)$, placed at $w$ (and $1$ at the other chosen places), $k'_g$ in
Definition~\ref{prm:def-level} may be taken with components in $K_v$ reducing to $1$ at the other chosen
places and at $Q$; those components lie in $J_O^{(w)}$ with $\chi_O=1$, so on $V_w$ the operator
$\varrho(g)$ is $\psi_w(\det g)$ times the action of $\mathrm{diag}(g,1)\in G_w$. Write $e_c=e_{c_w}e'$
with $e'$ the product of the idempotents at the other chosen places; these commute, so $e_{c_w}y=
e_{c_w}e_cy=e_cy=y$. Thus $y$ lies in the $\sigma_{c_w}$-isotypic part of $g\mapsto\psi_w(\det g)
\mathrm{diag}(g,1)$, that is, in the $\sigma'$-isotypic part of $V_w$ for $g\mapsto\mathrm{diag}(g,1)$,
where $\sigma'=\sigma_{c_w}\otimes(\phi\circ\det)$.

Let $X=\C[P]y$. As $U$ is normal in $P$ and fixes $y$, $U$ acts trivially on $X$, so $X=\C[M]y=\C[
\auGL_{r'}\times1]y$ ($1\times\auGL_{t'}$ fixes $y$ and commutes with $\auGL_{r'}\times1$). Hence $X$ lies
in the $\sigma'$-isotypic part for $\auGL_{r'}\times1$, and $1\times\auGL_{t'}$ acts trivially on it: as
a representation of $P$, $X$ is the inflation of $(\sigma'\boxtimes\mathbf 1)^{\oplus m}$ for some
$m\ge1$. The $G_w$-linear map $\C[G_w]\otimes_{\C[P]}X\to V_w$, $\alpha\otimes x\mapsto\alpha x$, has image
$\C[G_w]y$, and its domain is $\operatorname{Ind}_P^{G_w}X$, the Harish-Chandra induction of
$(\sigma'\boxtimes\mathbf 1)^{\oplus m}$. So every irreducible constituent of $\C[G_w]y$ is a constituent
of the Harish-Chandra induction $R_M^{G_w}(\sigma'\boxtimes\mathbf 1)$.

Now $\sigma'$ is a constituent of $R_{L_1}^{\auGL_{r'}}(\bigotimes_k\lambda_k)$ (twisted labels, above),
and the trivial representation of $\auGL_{t'}(\auF_w)$ is a constituent of $R_{T}^{\auGL_{t'}}(\mathbf 1)
=\C[\auGL_{t'}/B]$ (the constant functions), $T$ the diagonal torus and $B$ the Borel subgroup. So
$\sigma'\boxtimes\mathbf 1$ is a constituent of $R_L^M(\lambda)$, and by transitivity of Harish-Chandra
induction every constituent of $R_M^{G_w}(\sigma'\boxtimes\mathbf 1)$ is a constituent of
$R_L^{G_w}(\lambda)$, with $\lambda$ cuspidal. By Harish-Chandra's theorem (the irreducible constituents
of $R_L^{G}(\lambda)$, $\lambda$ cuspidal, have cuspidal support the $G$-conjugacy class of
$(L,\lambda)$; \cite[Proposition~6.3 and Theorem~6.4]{DigneMichel}; for $\auGL_n$ see \cite[\S2.3, Proposition; \S\S9.2--9.3]{ZelevinskyLNM}), every irreducible constituent of $\C[G_w]y$ has
cuspidal support $(L,\lambda)$.

(b) The space of forms of the weight of $O$ is the orthogonal direct sum of its $\Pi$-isotypic parts,
and the orthogonal projection $y\mapsto y_\Pi$ commutes with right translations; it maps $V_w$ into the
$\Pi$-part of $V_w$, on which $G_w$ acts through $\Pi_w^{K_w(1)}$ (tensored with a space on which $G_w$
acts trivially). So $\C[G_w]y_\Pi$, a non-zero quotient of $\C[G_w]y$, has an irreducible constituent,
which by (a) has cuspidal support $(L,\lambda)$ and is contained in $\Pi_w^{K_w(1)}$.
\end{proof}

\begin{corollary}[Inertia at a chosen place]\label{prm:inertia}
In the proof of Proposition~\ref{pri:BL}(c), paragraph ``Inertia at a chosen $w$'' of
\S\ref{pax:ss-BL}, the input ``$y$ is fixed by $K(1)$ at $w$ and lies in the isotypic part of the
representation induced from $\sigma_{\tilde c_w}\otimes\mathbf 1_{\auGL_{t'}}$'' holds in the form of
Lemma~\ref{prm:label}(b). With the cuspidal-support part of Lemma~\ref{pga:depth0} it gives: the
supercuspidal support of $\Pi_w$ consists of the depth-zero supercuspidal representations containing
the $\sigma_{\tilde f_k}$ and of $t'$ unramified characters (depth-zero characters of $E_w^\times$ whose
restriction to $\auO_w^\times$ inflates the trivial character of $\auF_w^\times$). Hence
Proposition~\ref{pri:BL} holds with the base level of Definition~\ref{prm:def-level}.
\end{corollary}
\begin{proof}
Lemma~\ref{prm:label}(b) provides an irreducible constituent of $\Pi_w^{K_w(1)}$ with cuspidal support
$(L,\lambda)$, $\lambda=\bigotimes_k\sigma_{\tilde f_k}\otimes\mathbf 1^{\otimes t'}$, which is the
hypothesis of the cuspidal-support paragraph of the proof of Lemma~\ref{pga:depth0}; the factor
$\auGL_1(\auF_w)$ with the trivial representation contributes a depth-zero supercuspidal of
$\auGL_1(E_w)$ containing the trivial character of $\auO_w^\times$, an unramified character. The rest of
the paragraph of \S\ref{pax:ss-BL} is unchanged ($\eta_O$ is unramified at $w$ by Lemma~\ref{pay:eta}).
Parts (a), (b) of Proposition~\ref{pri:BL} hold for $\mathcal F(O)$ by Definition~\ref{prm:def-level}
and Lemma~\ref{prm:level}(a) (the argument of \S\ref{pax:ss-BL}(b), with $\varrho$).
\end{proof}

\begin{remark}\label{prm:rem-level}
The level of Definition~\ref{prm:def-level} fixes the first $r'$ basis vectors modulo $w$ (the
vectors of $W$), and is the one the lift produces (Lemma~\ref{prm:level}(b)); a level of the form
$\{k\equiv\left(\begin{smallmatrix}*&*\\0&1_{t'}\end{smallmatrix}\right)\}$ is not. With it,
$\auG_O$ acts through $\mathrm{diag}(u_w,1_{t'})$, which normalises the level, and the trivial block
$\mathbf 1_{\auGL_{t'}}$ comes from the invariance under $1\times\auGL_{t'}$.
\end{remark}

\subsubsection{Where the inputs of \S\S\ref{pri:sec}, \ref{pga:sec}, \ref{pse:sec}, \ref{pg:sec} are proved}
\label{prm:ss-blocks}

The following list covers every statement of these four sections whose proof is given elsewhere,
with the place where it is proved.

\begin{enumerate}
\item Proposition~\ref{pri:BL} (base level of a label): proved in \S\ref{pax:ss-BL}, with
Lemma~\ref{pay:eta} ($\eta_O$ unramified away from $\ell$), Lemma~\ref{pay:product}, and
Definition~\ref{prm:def-level}, Lemmas~\ref{prm:level}, \ref{prm:label} and
Corollary~\ref{prm:inertia} (the base level and the product labels).
\item Lemma~\ref{pri:block} (block-diagonal lattices): Lemma~\ref{pax:block}, with the two hypotheses
(semisimplicity, spanning) made explicit there.
\item Lemma~\ref{pri:forms} (Higman operators, pairings): \S\ref{pax:ss-forms} and
\S\ref{pax:ss-pairing}.
\item Lemma~\ref{pri:chain} (chain lemma): \S\ref{pax:ss-chain}, with Lemma~\ref{pax:tadic} proved
directly in Corollary~\ref{pay:roots}.
\item Proposition~\ref{pri:odd} (odd classes): Proposition~\ref{pod:odd}, for every $\ell$ including
$\ell=2$ (Remark~\ref{pod:rem-two}).
\item Proposition~\ref{pri:kummer} (Kummer maps, detecting points): Proposition~\ref{pku:main} for (a)
and Proposition~\ref{pku:b} for (b).
\item Proposition~\ref{pri:dict} (dictionary at an auxiliary place): \S\ref{pax:ss-dict}, part (d)
through Lemma~\ref{pax:toproot}, proved in \S\ref{pay:ss-cons}.
\item The paragraphs ``The characters'' and ``Gluing'' of the proof of Proposition~\ref{pri:ribet}:
\S\ref{prm:ss-chars}.
\item The induction in Proposition~\ref{pri:norm}(b): Proposition~\ref{prm:normb}.
\item Speh constituents in Proposition~\ref{pga:G123} (Remark~\ref{pga:sign}):
Proposition~\ref{pax:sign} (the twist $|\det|^{s_i}$ is always of the right sign).
\item Theorem~\ref{pg:I}: the twist is Proposition~\ref{pax:sign}, the convention check is
Lemma~\ref{pax:conv}; the rest are published theorems, quoted with their hypotheses in
Theorem~\ref{pga:lit} and Lemmas~\ref{pga:depth0}, \ref{pga:parahoric}.
\item Proposition~\ref{pg:odd}: Proposition~\ref{pod:odd}.
\item Statement (IM) of \S\ref{pse:ss-stick} (integral measures): {\cite[Proposition~3.8]{P54}} for $z\ge2$
(at every $\ell\nmid\auNm Q$, $\ell=d$ included by {\cite[Lemma~3.6]{P54}}) and {\cite[Proposition~3.12]{P54}}
for $z=1$ (from the exchange identity); in (c) the place $v$ must satisfy $v,\bar v\notin
\mathsf M\cup Q\cup\bar Q$.
\item Proposition~\ref{pg:loss} (the local loss): proved in \S\ref{pll:sec}, the explicit
cup exponents being Theorem~\ref{pll:exponent}, proved in {\cite[\S2.3]{P56}}. Nothing in the
present section bears on it.
\end{enumerate}


\appendix

\section{Statements used from \texorpdfstring{\cite{P46}}{[P46]}, \texorpdfstring{\cite{P49}}{[P49]}, \texorpdfstring{\cite{P53}}{[P53]}, \texorpdfstring{\cite{P54}}{[P54]}, \texorpdfstring{\cite{P56}}{[P56]}}\label{C3:quoted}

For the convenience of the reader we reproduce here, without proofs, the statements of other papers that are used in this one, in the notation of the paper quoted. Each is headed by its place in that paper.

\begin{quotedthm}{\cite[Lemma~13.4]{P46}}{}
For admissible data the following hold.
\begin{enumerate}
\item[(S1)] $x^{10^{6}}\le x_1$, $x_1^{10^{6}}\le b$,
$4\log x_1\le(\log x_1)^{6}\le\log b$,
$(\log b)^{100}\le x$, $\log b\ge128$, and
$\log x\ge10^{14}(q_0+2d+|\aneu|)$.
\item[(S2)] $x^{11}\le2^{s}\le2^{22}x^{11}$; $l_0(s)\le2592\log x$;
$2m\le r$, so $r/2\le k\le r$.
\item[(S3)] $x^{11}\le T(s)=\#\anPa\le f\le q_1\le x^{12}$, and
$f=\varphi(q_0)\,\#\anPa$; moreover $1\le\#\anPa/\#\anPb\le\log b$.
\item[(S4)] $\anrd_E=\sqrt d\;q_1\prod_{p\mid q_0}p^{-1/(p-1)}$, and
$x^{10}\le\anrd_E\le x^{12}$.
\item[(S5)] $1\le xb/x_1\le\kappa(\tau)\le b\le a\le b^{2}$ for all $\tau$;
$b^{20}\le M\le b^{22}$;
$20\log b\le\log(M\anrd_E/2\pi)\le23\log b$; and
$m\log b\le\log N\le25\,m\log b$.
\item[(S6)] With $\Omega:=e^{\aneu}\cdot2\pi/\anrd_E$: $0<\Omega\le1$ and
$20\log b\le\log(M/\Omega)\le23\log b$.
\item[(S7)] With $n=m/2$ and $\varrho=n-1$ (so that $r=k+2\varrho+1$): $r$ is odd; every prime
dividing $q_1$ divides $q_0$; $q_0^{\,q_0}\le x$; $\varrho\ge10^{5}$;
$k\ge512000\,n^{9}$ and $k\ge1000\,(2\varrho+2)^{9}$; $k\ge2\varrho^{2}+3$; $18\varrho f\le b^{8}$;
$f\le b/1000$; $\anrd_E\le b$; and $N\le b^{60n}$. Moreover, with $\Delta=2r^{8}$:
\begin{itemize}
\item[(a)] $\varrho\ge10^{8}$ (so $n\ge100$); $n^{7}\le k$; $12n^{3}+500\le k$; $k-3n-6>10^{9}n^{3}$;
\item[(b)] $r\ge100$ (so $2r^{8}\ge500r^{3}+10r^{5}$); $3r^{3}f\le r^{7}$; $8nr\le b$; $kf\ge40$; $f\ge4000\,n$;
\item[(c)] $2f\le2^{2\varrho}$ and $2^{-\varrho}\le b^{-6}$;
\item[(d)] $(200\,f\,k\,\anrd_E)^{1000}\le b$ and $(25\,n^{3}f\,\anrd_E)^{1000}\le b$;
\item[(e)] $\sqrt d\le\anrd_E$, and $n,r,f,\sqrt d\le b$; $M\le b^{23}$;
\item[(f)] $\binom{a+k-1}{k-1}^{f}\le b^{2kf}$ (this bounds the dimension $d_\infty=\dim W_{\lambda^X}$ of
{\cite[\S14.3]{P46}}, a product of at most $f$ such binomial coefficients);
\item[(g)] $m\Delta+r\le\kappa(\tau)$ for all $\tau$; $k+2b+2m\Delta\le b^{2}$ (this bounds every
entry of $2\varsigma_\tau$, so that $L_\tau:=\max_i|2\varsigma_{\tau,i}|\le b^{2}$); and
$(r+1)^{4}(b^{2}+r+1+2n^{2})\le b^{5}$.
\end{itemize}
\end{enumerate}
\end{quotedthm}

\begin{quotedthm}{\cite[Lemma~14.5]{P46}}{ (Free levels)}
Let $n\le2n_m$ and let $J\subseteq K^\circ$ consist of elements congruent to $1$ modulo every
chosen place. Then $\Gamma\cap gJg^{-1}=1$ for all $g\in G_f$.
\end{quotedthm}

\begin{quotedthm}{\cite[Proposition~14.27]{P46}}{ (Equivariance of the lower chain)}
For every tame character $\vartheta$ and every $u\in\auG$ there is $\beta_\vartheta(u)$ in the
stabiliser $K^\circ$ of the lattice of $Y_m$ such
that $(\iota_\vartheta(\varrho(u)z))(g)=(\iota_\vartheta z)(g\beta_\vartheta(u))$ for all
$z\in\auAz$ and all $g$.
\end{quotedthm}

\begin{quotedthm}{\cite[Proposition~14.36]{P46}}{ (Equality of multiplicities)}
For $1\le j\le m$ and every element $h$ of the spherical Hecke algebra away from $M$ (the
compactly supported functions on the finite adelic group outside the chosen places that are
bi-invariant under the stabiliser of the lattice, at all places, split or not),
$\operatorname{tr}(h\,|\,e_0\mathcal A(W_j))=\operatorname{tr}(h\,|\,\auT)$. In particular
$\dim e_0\mathcal A(W_j)=\audT$, and each system of eigenvalues of the Hecke operators at the split
places prime to $M$ has the same multiplicity in $e_0\mathcal A(W_j)$ and in $\auT$. (The
hypothesis of {\cite[Proposition~4.6]{P51}} holds for the data: at every $\tau$, two entries of
$\varsigma_\tau$ on coordinates of different signs {\cite[(43)]{P51}} differ by at least
$\Delta=2r^8\ge500r^3+10r^5$.)
\end{quotedthm}

\begin{quotedthm}{\cite[Proposition~14.37]{P46}}{ (Factorisation of the chain determinant)}
For $1\le j\le m$ and a twist $\vartheta$ let $Z_j(\vartheta)=\det(u^*u)>0$, where
$u=\Theta_{j,\vartheta}$ is the theta lift from $e_0\mathcal A(W_j)$ to algebraic modular forms on
$\auU(V_j)$ with the character $\mu_j^0\vartheta$, normalised as in {\cite[\S2.4]{P53}} (joint
harmonic at infinity; at the chosen places the indicator of the inclusion; at $w\in Q_\vartheta$
the function $x\mapsto\psi_w(\det x^{\mathrm{top}})$; values in the algebraic modular forms of level
$J^V_\vartheta$ with their inner product). There are $e\in\Z$ and the ramified ratios $\rho_{ij}$,
$0<i<j\le m$, of {\cite[(9)]{P53}}, such that
for every $m$-fold twist $\vartheta$
\[
\frac{Z(1)}{Z(\vartheta)}=\Bigl((\auNm Q_\vartheta)^{e/2}\prod_{0<i<j\le m}\rho_{ij}(\vartheta)
\Bigr)^{-\audT}\cdot\prod_{j=1}^m\frac{Z_j(1)}{Z_j(\vartheta^j)} .
\]
\end{quotedthm}

\begin{quoteddef}{\cite[Definition~14.39]{P46}}{ (Regular families of twists)}
Put $\aubud=\aueu/(mk)$. A family $\theta\colon\auI\to\{m\text{-fold twists}\}$ is \emph{regular}
if:
\begin{enumerate}
\item[(R1)] $\ell(\iota)\nmid\auNm Q_{\theta_\iota}$ for all $\iota$;
\item[(R2)] for every ramified critical pair $(\nu,z)$, $1\le n\le m$ and sign $\pm$, the family
$\iota\mapsto|L(z,\nu^{-1})|/|L(z,(\nu\theta_\iota^{\pm n})^{-1})|$ has height at most $x_1$;
\item[(R3)] for $1\le j\le m$ there are $\beta_\ell\ge0$ with $\sum_{\ell\in T}\beta_\ell\le
\aubud$ for all finite $T$ such that, for every $\iota$ with prime $\ell$: (i) the twist
$\theta_\iota^{\,j}$ \emph{has denominators at most $\beta_\ell$ at $\iota$ for the $j$-th lift}
in the sense of {\cite[Definition~2.8]{P54}} (the false lifts being those of
{\cite[Definition~2.6]{P54}}: the forms in the image of $e_0$ that are orthogonal to the lifts and
are killed by the Hecke polynomial $P_w(X_w)$ of {\cite[Lemma~2.2(L3)]{P54}} at one split place $w$
off the level; a form is orthogonal to the false lifts if it is so for every such $w$); and (ii)
for every main pair $p$ for $j$ ({\cite[Definition~2.19]{P54}}) and all test forms $\Xi,\Xi'$ at
$\ell$ for $(p,\theta_\iota^{\,j})$ ({\cite[Definition~2.22]{P54}}), the inner product of the orthogonal
projections of $\Xi$ and $\Xi'$ to the line of the normalised lift
$\varphi_{p,\theta_\iota^{\,j}}$ has absolute value at most $e^{\beta_\ell}$ at $\iota$.
\end{enumerate}
\end{quoteddef}

\begin{quotedthm}{\cite[Proposition~14.40]{P46}}{ (Regular families exist)}
For all admissible data there is a regular family of $m$-fold twists.
\end{quotedthm}

\begin{quotedthm}{\cite[Proposition~14.41]{P46}}{ (One lift and Hecke values)}
Let $1\le j\le m$ and $d_j=\dim e_0\mathcal A(W_j)$. There are a constant $C_j\ne0$, for every twist
$\vartheta$ a number $\Delta_j(\vartheta)\ne0$ (the determinant of a \emph{comparison operator}
compressed to the space of lifts) and the positive number $H_j(\vartheta)$ of {\cite[(11)]{P54}}: for the three critical pairs $(\nu_i,z_i)$ of {\cite[\S2.7]{P54}} it is the product of $(\auNm\mathsf S_\vartheta)^{z_i^2}\prod_{1\le z\le z_i}|L^{\mathsf S_\vartheta}(z,(\nu_i\vartheta^{\pm1})^{-1})|^2$, such that
\[
Z_j(\vartheta)\cdot\Delta_j(\vartheta)=C_j\cdot\bigl(H_j(\vartheta)/H_j(1)\bigr)^{d_j},
\]
and
\begin{enumerate}
\item[(a)] $|\sigma\Delta_j(1)|\le\exp(3d_j\cdot f\,m\,b\,(\log b)^2)$ for all
$\sigma\in\auAut(\C)$;
\item[(b)] if $\ell(\iota)\nmid\auNm Q_\vartheta$ and $\vartheta$ has denominators at most $\beta$
at $\iota$ for the $j$-th lift, then $|\Delta_j(\vartheta)|_\iota\le\exp(d_j\,r\,\beta)$.
\end{enumerate}
\end{quotedthm}

\begin{quotedthm}{\cite[Proposition~14.42]{P46}}{ (Heights of twisted Hecke values)}
For a regular family $\theta$ and $1\le j\le m$ the family $\iota\mapsto H_j(1)/H_j(\theta_\iota^
{\,j})$ has height at most $10\,m\,x_1+k\,(6\,\aubud-\auhdir)$.
\end{quotedthm}

\begin{quotedthm}{\cite[Proposition~14.43]{P46}}{ (The cost of one lift)}
For a regular family $\theta$ and $1\le j\le m$ the family $\iota\mapsto Z_j(1)/Z_j(\theta_\iota^
{\,j})$ has height at most $\audT\bigl(f\,a\,(\log b)^4/m-k\,\auhdir\bigr)$.
\end{quotedthm}

\begin{quotedthm}{\cite[Lemma~5.2]{P49}}{ (Standard facts)}
Let $W$ be a hermitian space over $E$.
\begin{enumerate}
\item[(W1)] $\omega(g)\psrho(w)\omega(g)^{-1}=\psrho(gw)$. There is a functional $\theta=\theta_W$
on $\psHS(W)$, unique up to a scalar, with $\theta\circ\psrho(w)=\theta$ for $w\in W(E)$; it
satisfies $\theta\circ\omega(\gamma)=\theta$ for $\gamma\in\auU(W)(E^+)$.
\item[(W2)] Let $I\subset W$ be an isotropic $E$-subspace, $W_I=I^\perp/I$, and $P_I$ the stabiliser
of $I$ in $\auU(W)$. There is a continuous linear map $\psZ_I\colon\psHS(W)\to\psHS(W_I)$ (the
\emph{zero term}) with $\theta_{W_I}(\psZ_IF)=\int_{I(E)\backslash I(\auA)}\theta(\psrho(i)F)\,di$,
$\psZ_I\psrho(w)=\psrho(\bar w)\psZ_I$ for $w\in I^\perp(\auA)$ with image $\bar w$, and
$\psZ_I\omega(p)=\chi_I(p)\,\omega_{W_I}(\bar p)\psZ_I$ for $p\in P_I(\auA)$, where $\chi_I(p)=
(\upsilon|\cdot|^{-1/2})(\det p|_I)^{-1}=\upsilon(\det p|_I)^{-1}|\det p|_I|^{1/2}$, with
$\det p|_I\in\auA_E^\times$ the determinant on the $\auA_E$-module $I(\auA)$ and $|\cdot|$ the
idele norm of $E$; so $\chi_I=1$ on unipotent elements and on $P_I(E^+)$.
For isotropic $I\subset I'$ one has $\psZ_{I'}=\psZ_{I'/I}\psZ_I$, and $\psZ_I=\psZ_{I,\infty}
\otimes\psZ_{I,f}$.
\item[(W3)] For $F\in\psHS(W)$, $\theta(F)=\sum_{w\in W(E)/I^\perp(E)}\theta_{W_I}(\psZ_I\psrho(w)F)$,
the series converging absolutely. If $W=\auHom(V,V')$, $F=\Psi\otimes f$ with $\Psi$ smooth, and
$\mathcal C$ is a compact set of finite adelic points of $\auU(V)\times\auU(V')$, there are $C,A$ with
$\sum_w|\theta_{W_I}(\psZ_I\psrho(w)\omega(g,h)F)|\le C(\|g_\infty\|\|h_\infty\|)^A$ whenever
$(g_f,h_f)\in\mathcal C$.
\item[(W4)] Let $L\subset W(\auA_f)$ be a lattice that is its own dual for $[\,,]$. The vectors of
$\psHS_f(W)$ on which $\psrho(l)$, $l\in L$, acts by a fixed character form a line. If $\psvac_L$
spans the line of the character $\epsilon_L$ with $\epsilon_L(l)\epsilon_L(l')=\beta(l,l')
\epsilon_L(l+l')$ fixed below, and $\psrho^\circ(x)=\epsilon(x)\psrho(x)$ for a fixed function
$\epsilon$ with values in roots of unity extending $\epsilon_L^{-1}$, then
$\psrho^\circ(l)\psvac_L=\psvac_L$, $\psrho^\circ(x+l)\psvac_L=\zeta\psrho^\circ(x)\psvac_L$ with a
root of unity $\zeta$, the $\psrho^\circ(x)\psvac_L$, $x\in W(\auA_f)/L$, form a basis of
$\psHS_f(W)$, and $\omega(g)\psvac_L\in\C^\times\psvac_{gL}$ for $g\in\auU(W)(\auA_f)$. If
$g\in\auU(W)(\auA_f)$ stabilises $L$ then $\omega(g)\psvac_L=\psvac_L$.
\item[(W5)] At a place $v=w\aubar w$ of $E^+$ split in $E$, $\psHS_v(\auHom(V,V'))$ is the space
of Schwartz--Bruhat functions on $\auHom(V_w,V'_w)$, with $(\omega_\mu(g,h)\Psi)(x)=\mu_w(\det h_w)
\mu_w(\det g_w)^{-1}\upsilon'_w(g_w,h_w)\,\Psi(h_w^{-1}xg_w)$, where, for even ranks $n,n'$, $\mu_w\upsilon'_w$ is given by {\cite[(62)]{P49}}: the factor is
$\chi_w(\det h_w)|\det h_w|_w^{-n/2}\chi_w(\det g_w)^{-1}|\det g_w|_w^{n'/2}$ with $\chi$ the base
change of $\mu$; for arbitrary ranks the untwisted operator is
\[\Psi\mapsto\upsilon_w(\det h_w)^{-n}\upsilon_w(\det g_w)^{n'}|\det h_w|_w^{-n/2}|\det g_w|_w^{n'/2}\Psi(h_w^{-1}xg_w)\]
({\cite[Proposition~2.17]{P49}}), so that all factors multiplying $\Psi(h_w^{-1}xg_w)$ are built from $\mu_w$, $\upsilon_w$ and $|\cdot|_w$ evaluated at $\det g_w,\det h_w$; $\psrho(a_w,0)\Psi(x)=\Psi(x+a_w)$
and $\psrho(0,a_{\bar w})$ is multiplication by the character $x\mapsto[x,a_{\bar w}]$.
\item[(W0), (W6)] Two further facts are used; they are stated in full in {\cite[\S5.2.1]{P49}}.
(W0): $\theta$ is defined on the smooth vectors of the Fock space tensored with $\psHS_f(W)$ and is
continuous in the first factor. (W6): for hermitian spaces $V,V'$ of ranks $n,n'$ the map
$w\mapsto w^\dagger$ identifies $\psHS(\auHom(V,V'))$ with $\psHS(\auHom(V',V^-))$, carrying
$\omega_{\auHom(V,V')}(g,g')$ to $\upsilon(\det g)^{n'}\upsilon(\det g')^{-n}\,
\omega_{\auHom(V',V^-)}(g',g)$. It carries the self-dual lattice $\alpha\auHom(\psM,\psM')$ to
$\aubar\alpha\auHom(\psM',\psM)$, which differs from $\alpha\auHom(\psM',\psM)$ at the places above $q_1$
({\cite[Remark~7.18]{P49}}).
\end{enumerate}
\end{quotedthm}

\begin{quotedthm}{\cite[Theorem~3.2]{P53}}{ (Exchange identity)}
Let $W,V,V'$ be spaces of ranks $N=a+b$, $a$, $b$, with $a=b+2e$, $e\ge0$ (so $N$ is even). Let
$q,q_V,q_{V'}$ be their numbers of negative signs; assume $q=q_V+q_{V'}$ and $2q+2\le N$. Let
$\mu=(\mu_\tau)$ be integers with $N\le2\mu_\tau+2q$ for all $\tau$, and let $v$ be a vector of the
Fock space of $\auHom(V',W)\otimes\R$ with finite moments. Then there are $C>0$ and a vector $x$ of
the Fock space of $\auHom(W,V)\otimes\R$ with finite moments, such that: for every modulus
$\pcm\ne(1)$ and every tame character $\nu$ for $\pcm$ of weight $\mu$ there is $c\in\C$ with
\[
|c|^2=C\cdot L_*(\pcm,\nu,e)\ne0,
\]
such that for all automorphic characters $\alpha,\alpha'$ of $\auU(1)$ with
\[
\alpha\alpha'=\gamma:=\nu^{-1}\upsilon^{N/2},
\]
all functions $\Phi$ on $\pcMat_{a\times N}(\pcR_\pcm)$ and $\Psi$ on $\pcMat_{N\times b}(\pcR_\pcm)$
matched for $\chi$, and all $g\in\auU(W)(\auA)$,
\begin{equation*}\tag*{{\cite[(12)]{P53}}}
\int_{[\auU(V)]}\theta^{W,V}_{\alpha,\gamma}(g,h;x,\Phi)\,dh=c\int_{[\auU(V')]}
\theta^{V',W}_{\gamma,\alpha'}(h',g;v,\Psi)\,dh' .
\end{equation*}
Here $\theta(\cdot\,;x,\Phi)$ is the kernel of the datum $[x,\Phi]$ of {\cite[\S2.3]{P53}} (matrices
in the frames), and $C$, $x$ do not depend on $\pcm,\nu,\alpha,\alpha',\Phi,\Psi,g$.
\end{quotedthm}

\begin{quotedthm}{\cite[Lemma~5.26]{P53}}{ (The duality pairing)}
Let $\lambda=(\kappa,e)$ and $\lambda'=(\kappa',e')$ be weights of $\prod_\tau\auGL_N$ whose doubled
Harish-Chandra parameters are $l$ and $-l$. Then $\kappa'(h)=\kappa(N-h)$ for $1\le h<N$ and
$e+e'=-d$ at every $\tau$. Put $A=G_{\lambda'}J$, with $J$ the tensor product over $\tau$ of the
matrices of {\cite[Lemma~5.25]{P53}}, and $\langle v',v\rangle_{\pgsdu}={}^tv'Av$ for coordinate vectors
$v'$ of weight $\lambda'$ and $v$ of weight $\lambda$. Then:
\begin{enumerate}
\item[(P1)] $\langle\ ,\,\rangle_\pgsdu$ is invariant: $\langle\rho'(g)v',\rho(g)v\rangle_\pgsdu=
\langle v',v\rangle_\pgsdu$ for $g\in\prod_\tau\auGL_N(\C)$;
\item[(P2)] there is an integer matrix $C$ with $CA=h\cdot1$, where $h=H_{\lambda'}=\prod_\tau
H_{\mu'_\tau}$ is the hook number of $\lambda'$;
\item[(P3)] for every $v$ of weight $\lambda$ the vector $v^\flat=\overline{Jv}$ of weight $\lambda'$
satisfies $\|v^\flat\|=\|v\|$ and $\langle a,v\rangle_\pgsdu=\langle v^\flat,a\rangle$ for all $a$,
where $\langle p,q\rangle={}^t\bar pG_{\lambda'}q$ is the invariant hermitian form of weight
$\lambda'$.
\end{enumerate}
\end{quotedthm}

\begin{quotedthm}{\cite[Lemma~2.2]{P54}}{ (Standing facts)}
Let $\vartheta$ be a twist.
\begin{enumerate}
\item[(L2)] $\Theta_\vartheta\circ\varrho(s)=\varrho_V(s)\circ\Theta_\vartheta$ for $s\in\auG$.
Hence $\Theta_\vartheta(e_0\poA(W_j))\subseteq\varrho_V(e_0)\poA(V;\vartheta)$.
\item[(L3)] Let $w$ be off the level, with residue field of $q$ elements, $P_w(X)=\sum_{i\ge0}
(-X)^iq^{\binom i2}T_{w,i}$ the operator of (N3) on $\poA(V;\vartheta)$, and $\alpha_w$ the number
with $\omega_\mu(1,d_1)\plone=\alpha_w\,\plone(d_1^{-1}T)$ at $w$ (by {\cite[(1)]{P54}}, $\alpha_w=\eta_w(1,d_1)$ does not depend on $\vartheta$, and
$|\alpha_w|=q^{r/2}$). Put $X_w=(\alpha_wq^{t-1})^{-1}$. Then $P_w(X_w)\,\Theta_\vartheta y=0$ for
every $y\in\poA(W_j)$.
\end{enumerate}
\end{quotedthm}

\begin{quotedthm}{\cite[Lemma~3.6]{P54}}{ (Kernel integrality at $\ell=d$)}
Let $\ell=d$, $\iota\colon\C\to\C_d$, $\delta\in\auO_K\smallsetminus\{0\}$ and $f$ a standard datum of
level $\delta$ at $d$. Then $\Theta_f(d\,\aubar\delta s_0\,\cdot)$ is $\partial$-integral; hence for every
$k$ and every homogeneous polynomial $Q$ of degree $k$ with coefficients in $\pao$, $d^kQ(\partial)
\Theta_f(\aubar\delta s_0\,\cdot)(0)\in\pao$.
\end{quotedthm}

\begin{quotedthm}{\cite[Proposition~3.8]{P54}}{ (Statement (IM) for $z\ge2$)}
Let $p=(\nu,z)$ be a ramified critical pair with $z\ge2$. Put $I=\{z\}$, $m=0$ and
\[
F(\vartheta)=\mathsf N_z\,n_p^2\,\frac{\|\varphi_{p,\vartheta}\|^2}{(\auNm Q)^{2z-1}}\qquad
(\vartheta\text{ tame of modulus }Q).
\]
Then: (a) $F(1)\ne0$; (b) for every $Q$ there are finitely many $x_a\in\auU(1)(\auA_{E^+,f})$ and
$c_a\in\C$, with $c_a$ integral at every prime $\ell\nmid\auNm Q$ under every embedding $\C\to\C_\ell$,
such that $F(\vartheta)=\sum_ac_a\vartheta(x_a)$ for every tame $\vartheta$ of modulus $Q$; (c) if $v\notin Q$ and $\vartheta'$ is
$\vartheta$ regarded as of modulus $Q\cup\{v\}$, then $F(\vartheta')=F(\vartheta)\prod_\pm(1-\alpha_v^{\pm1}
q_v^{-z})$, $\alpha_v$ the value at a uniformiser of the base change of $\nu\vartheta$ at $v$; (d) for
every twist $\vartheta$ of modulus $Q$, $F(\vartheta)=F(1)\,\rho_p(\vartheta)$.
\end{quotedthm}

\begin{quotedthm}{\cite[Proposition~3.12]{P54}}{ (Statement (IM) for $z=1$)}
Let $p=(\nu,1)$ be a ramified critical pair, $I=\{0,1\}$, $m=1$, and $F(\vartheta)=|\mathsf N_1n_1
\Lambda(\vartheta)|^2$, with $\mathsf N_1=\#\auGL_{4f}(\auF_3)$ and $n_1\ge1$ an integer, depending only on $p$, making the kernel values below integral at all
$\ell\nmid\auNm Q$ (it exists by {\cite[Proposition~5.30]{P49}}, {\cite[Corollary~7.16]{P49}} and {\cite[Lemma~3.6]{P54}}). Then {\rm(a)}--{\rm(d)}
of {\rm(IM)} hold: $F(1)\ne0$; $F$ is an integral measure at every modulus; for $v,\aubar v\notin\mathsf M
\cup Q\cup\aubar Q$, $F(\vartheta')=F(\vartheta)\prod_{e\in\{0,1\}}\prod_\pm(1-\alpha_v^{\pm1}q_v^{-e})$; and
$F(\vartheta)=F(1)\,\auNm Q\,\rho_p(\vartheta)^2$ for every twist.
\end{quotedthm}

\begin{quoteddef}{\cite[Definition~2.4]{P56}}{ (Letters that read thickness)}
Let each embedding $\gamma$ be \emph{thick} or \emph{thin} and carry a sign $\varepsilon(\gamma)$, and
let $T$ (a set of letters) and $s$ (a sign on letters) be given. We say that \emph{the letters read
thickness} (for $Q$, $h$) if: $e\le h$; for every $\gamma$ that is thick, and for every $\gamma$ if
$\ell>Q$, one has [$\gamma$ thick $\Leftrightarrow\gamma(\zeta)\in T$] and
$\varepsilon(\gamma)=s(\gamma(\zeta))$; either $e=1$ and $\ell\ge7$, or $\ell\le Q$ and $s$ is
constant; and, if $\ell\le Q$, there is no letter $z$ with $z$, $z^{\ell^{a}}$, $z^{\ell^{a+a'}}$ all
in $T$ for some $a,a'\in[1,\pkilog_\ell h^3]$.
\end{quoteddef}

\begin{quotedthm}{\cite[Proposition~2.5]{P56}}{ (Generic primes)}
Let $e=1$, $\ell\ge5$, and $a,\eta,W\ge0$ with $\eta\ge1$, such that for every letter $z$:
$b_z+W+\eta\le(\ell-3)a$; $6(a+\eta+W)<b_z$; $\mathrm{NM}(\ell,6(a+\eta+W),b_z)$; $U_z\subseteq
\pkiWin(\eta,-\varepsilon_z,b_z,W)$ and $U'_z\subseteq\pkiWin(\eta,\varepsilon_z,b_z,W)$.
\begin{enumerate}
\item[(i)] If no element of any $U(\gamma)\cup U'(\gamma)$ lies in $[-a,0)$, or if every
$U'(\gamma)\subseteq[-\eta,0]$, then $\plldef_{\pkiC(U)}(L)=0$.
\item[(ii)] If $L$ has at most $\mu$ jumps in $[-\eta,\eta]$ at each $\gamma$, then
$\plldef_{\pkiC(U)}(L)\le11f\bigl(\mu(a+1)+\pllrk_\auO L\cdot(W+1)\bigr)$.
\end{enumerate}
\end{quotedthm}

\begin{quotedthm}{\cite[Proposition~2.10]{P56}}{ (The summands fit the windows)}
Let $h=r^8$, $Q=kt$, $G$ the weight at the places that are not $b$-places, $\eta=h^2+1$, $W=4t$, and
for $0\le j\le m$ let $\delta_j=j\Delta+h-2t$. Then (S1)--(S4) hold for $(Q,m,h,G,b,\delta_j,\eta,W)$.
Let $L_c$ and $\mathcal C_v$ be as in Proposition~\ref{pll:main} and let $U(\gamma)$ be the jump set
of $\mathcal C_v$, $U'(\gamma)$ that of $L_c^*$. There are $j\le m$ and, at each $\gamma$, a sign
$\varepsilon(\gamma)$, with $U(\gamma)\subseteq\pkiWin(\eta,-\varepsilon(\gamma),a_\gamma+\delta_j,W)$
and $U'(\gamma)\subseteq\pkiWin(\eta,\varepsilon(\gamma),a_\gamma+\delta_j,W)$, where $a_\gamma=b$ if
$\gamma$ is thick and $a_\gamma=G$ otherwise. Moreover $L_{\mathrm{main}}^*$ has few low jumps for
$(h,\lfloor r'm/k\rfloor,\eta)$, and for a character $L_e$ every $U'(\gamma)\subseteq[-\eta,0]$.
At $\ell=p_i$, $\mathcal C_v\subseteq\pkiC(U,K)$ and $L_c^*\in\pkiC(U',K^{-1})$ for a $K$ as in
{\cite[Proposition~2.8]{P56}} with $N=\max(r+t,2r)\,r$, $\varepsilon(\gamma)$ and $a_\gamma$ do not depend on $\gamma$ (write $a$ for the common value), and
$\ell>a+\delta_j+6(W+\eta)+r^2(N+1)^2+2$ for $a\in\{b,G\}$.
\end{quotedthm}

\begin{quotedthm}{\cite[Theorem~2.13]{P56}}{ (Uniform crystalline bound)}
Assume (S1)--(S4) for $(Q,m,h,G,b,\delta,\eta,W)$, that the letters read thickness, that
$U(\gamma)\subseteq\pkiWin(\eta,-\varepsilon(\gamma),a_\gamma+\delta,W)$ and
$U'(\gamma)\subseteq\pkiWin(\eta,\varepsilon(\gamma),a_\gamma+\delta,W)$ for all $\gamma$ (with
$a_\gamma$ as in {\cite[Lemma~2.12]{P56}}), that $L$ is crystalline with jumps in $U'$, and that
either $L$ has few low jumps for $(h,\mu,\eta)$, or every $U'(\gamma)\subseteq[-\eta,0]$. Let
$\lambda=\lambda_\ell(h,b)$ and $K=\lfloor2\lambda/Q\rfloor$. Then:
\begin{enumerate}
\item[(b)] if $\ell\le Q$, $\pkiC(U)$ has cup exponent $10^4K(\pkilog_\ell h^3+1)$ on $L$;
\item[(c)] if $\ell>Q$ and $\lambda\ge h^3$, it has cup exponent $10^4Km$;
\item[(d)] if $\ell>Q$ and $\lambda<h^3$, then $\plldef_{\pkiC(U)}(L)=0$ unless $\lambda\ge\lfloor
h/5\rfloor$ and $L$ has few low jumps; and always
$\plldef_{\pkiC(U)}(L)\le[E_v:\Q_\ell]\cdot10^5(\pllrk L\cdot Km+\mu\lambda)$.
\end{enumerate}
In all cases $\plldef_{\pkiC(U)}(L)\le[E_v:\Q_\ell]\cdot10^5(\pllrk L\cdot Km+\mu\lambda)$.
\end{quotedthm}

\begin{quoteddef}{\cite[Definition~7.1]{P56}}{ (Convention)}
For $g\in\pkjG$ let
\[
\sigma_g=\iota^{-1}\circ g^{-1}\circ e_v\quad\text{(reading (R$^-$))}\qquad\text{or}\qquad
\sigma_g=\iota^{-1}\circ g\circ e_v\quad\text{(reading (R$^+$))}.
\]
The paper uses reading (R$^-$); we prove everything for both readings, putting $\varpi=-1$ for (R$^-$)
and $\varpi=+1$ for (R$^+$), one of the two being fixed throughout a proof. Write
$P(g)\in\pkjPl$ for the place of $\sigma_g$ and $\varepsilon(g)=\varepsilon(\sigma_g)$; call $g$
\emph{thick} if $P(g)\in\pkjPb$. A coefficient ring $\auO\subset\auQlb$ is \emph{large} if it contains
the ring of integers $\auO_{E_v}$ of $E_v$. For such $\auO$, every $\Z_\ell$-algebra embedding $\gamma\colon\auO\to
\C_\ell$ has image in $\auQlb$ and is the restriction of some $g\in\pkjG$ (extend $\gamma$ $\Q_\ell$-linearly to the algebraic extension $\auO[1/\ell]$ of $\Q_\ell$, whose image is algebraic over $\Q_\ell$, and then to an automorphism of $\auQlb$); the embeddings of {\cite[Proposition~2.9]{P56}} are of this kind; the restriction of $g$ to
$e_v(E)\subset\auO[1/\ell]$ is determined by $\gamma$, hence so is $\sigma_g$ (in (R$^-$) because
$g^{-1}|_{e_v(E)}=(g|_{e_v(E)})^{-1}$). We put $\sigma_\gamma=\sigma_g$, and $\gamma$ is thick,
resp.\ has sign $\varepsilon(\gamma)$, if $g$ is thick, resp.\ $\varepsilon(\gamma)=\varepsilon(g)$.
This is ``the complex embedding of $E$ attached to $\gamma$ through $\iota$ and $v$'' of
{\cite[Proposition~2.9]{P56}}.
\end{quoteddef}

\begin{quotedthm}{\cite[Theorem~7.9]{P56}}{ ({\cite[Proposition~2.9]{P56}})}
Let $Q=kt$, $h=r^8$, and let thickness and $\varepsilon(\gamma)$ be as in {\cite[Convention~7.1]{P56}}
(either reading). For every prime $\ell$, every $v\mid\ell$, every large $\auO$ and every
$\zeta\in\auO$ of exact order $\ell^f-1$ (such $\zeta$ exist), there are a set $T$ of letters and a sign
$s$ on letters for which the letters read thickness for $(Q,h)$ ({\cite[Definition~2.4]{P56}}).
\end{quotedthm}

\begin{quoteddef}{\cite[Definition~9.2]{P56}}{ ($D_{\mathrm{pcris}}$, $\beta_s$, killing)}
Let $Y$ be a lattice and $V=Y[1/\ell]$. Put $\pkwD(Y)=\bigcup_H\pkwB(Y)^H$, the union over the open
subgroups $H$ of $\Gamma$. $Y$ is \emph{potentially crystalline} if $\pkwB(Y)^H$ spans $\pkwB(Y)$ over
$\pkwB$ for some open $H\subseteq\Gamma$ (by (F1), equivalently $V$ is crystalline on $H$). For such $Y$
and an inertial $s\in\Gamma$, $\beta_s=\beta_s^Y$ denotes the unique $\pkwB$-linear endomorphism of
$\pkwB(Y)$ that agrees with $s$ on $\pkwD(Y)$ ({\cite[Lemma~9.3(i)]{P56}}). For $P=\sum_kP_kX^k\in\auO[X]$,
``$P$ \emph{kills $s$ on $\pkwD(Y)$}'' means $P(\beta_s)=\sum_kP_k\beta_s^k=0$ on $\pkwB(Y)$,
equivalently ({\cite[Lemma~9.3(i)]{P56}}) $\sum_kP_ks^k=0$ on $\pkwD(Y)$.
\end{quoteddef}

\begin{quotedthm}{\cite[Theorem~9.9]{P56}}{ ((T2): descent)}
Fix a tame chart ({\cite[Definition~9.5]{P56}}) with ramification index $e$. Let $Y$ be a potentially
crystalline lattice such that $X^e-1$ kills every inertial $s\in\Gamma$ on $\pkwD(Y)$. Then $Y$ is
crystalline on $\Gamma'$. Its Hodge--Tate jumps are unchanged (they are those of $Y[1/\ell]$, which do not
depend on the open subgroup on which they are computed), so if $Y$ has jumps in $U$ it is a lattice with
jump sets over the chart ({\cite[Definition~3.4]{P56}}).
\end{quotedthm}

\begin{quotedthm}{\cite[Proposition~9.23]{P56}}{ (Uniform tame polynomials)}
Let $E_v=\Q_\ell$, $r,N_0\ge1$, $e'=\operatorname{lcm}_{1\le a\le r}(\ell^a-1)$. For each inertial $s\in\Gamma$ let
$M_s$ be a multiset of at most $N_0$ elements of $\auQlb$, each satisfying $y^{\ell^a-1}=1$ for some $1\le a\le r$,
and let $\mathcal T$ be a set of potentially crystalline lattices such that for every $T\in\mathcal T$ and every
inertial $s$, all eigenvalues of $\beta_s^T$ lie in $M_s$. (Here an eigenvalue of $\beta_s^T$ is a
$y\in\auQlb$ such that $\beta_s-y$ has nonzero kernel on $\pkwB(T)\otimes_{\Z_\ell}\Z_\ell[y]$; since $\beta_s$ has
finite order $n$ (see the proof), each such $y$ lies in $\mu_n$, and $\pkwB(T)\otimes_{\Z_\ell}\Q_\ell(\mu_n)$ is
the direct sum of the kernels of $\beta_s-y$, $y\in\mu_n$.) Put $q_s=\prod_{y\in M_s}m_y$, $m_y$ the minimal
polynomial of $y$ over $\Q_\ell$. Then $q_s\in\Z_\ell[X]$ is monic of degree $\le N_0r$, depends only on $s$, and
$q_s$ and $X^{e'}-1$ kill $s$ on $\pkwD(T)$ for every $T\in\mathcal T$, and on $\pkwD(T')$ for every lattice $T'$
with an injective equivariant map into a finite direct sum of members of $\mathcal T$.
\end{quotedthm}

\begin{thebibliography}{P46}
\bibitem[P46]{P46} D.~A.~Trejo Pizzo, \emph{The sign of the radial intersection form on the idele class group}, preprint, 2026, \url{https://github.com/dtrejopizzo/riemann/tree/main/04-papers/46-row-d-sign-radial-intersection-form}.
\bibitem[P49]{P49} D.~A.~Trejo Pizzo, \emph{The Short Forms Theorem: theta lifts, complex multiplication and divisibility}, preprint, 2026, \url{https://github.com/dtrejopizzo/riemann/tree/main/04-papers/49-short-forms-theorem-theta-lifts-complex-multiplication}.
\bibitem[P51]{P51} D.~A.~Trejo Pizzo, \emph{The Spanning Theorem: the second moment of a Poincar\'e map}, preprint, 2026, \url{https://github.com/dtrejopizzo/riemann/tree/main/04-papers/51-spanning-theorem-second-moment-poincare-map}.
\bibitem[P53]{P53} D.~A.~Trejo Pizzo, \emph{The Covolume Theorem: exchange identities and the determinant of a chain of theta lifts}, preprint, 2026, \url{https://github.com/dtrejopizzo/riemann/tree/main/04-papers/53-covolume-theorem-exchange-identities-chain-determinant}.
\bibitem[P54]{P54} D.~A.~Trejo Pizzo, \emph{The Covolume Theorem: one theta lift, Hecke $L$-values and test forms of least valuation}, preprint, 2026, \url{https://github.com/dtrejopizzo/riemann/tree/main/04-papers/54-covolume-theorem-one-theta-lift-hecke-l-values}.
\bibitem[P56]{P56} D.~A.~Trejo Pizzo, \emph{Explicit cup exponents for local Galois representations beyond the Fontaine--Laffaille range}, preprint, 2026, \url{https://github.com/dtrejopizzo/riemann/tree/main/04-papers/56-explicit-cup-exponents-local-galois-representations}.
\bibitem{Washington} L.~C.~Washington, \emph{Introduction to Cyclotomic Fields}, 2nd ed.,
Graduate Texts in Mathematics \textbf{83}, Springer, New York, 1997.
\bibitem{Finis} T.~Finis, \emph{Divisibility of anticyclotomic $L$-functions and theta functions
with complex multiplication}, Ann.\ of Math.\ (2) \textbf{163} (2006), 767--807.
\bibitem{HidaNonvanishing} H.~Hida, \emph{Non-vanishing modulo $p$ of Hecke $L$-values}, in:
Geometric Aspects of Dwork Theory, Vol.~II, de Gruyter, 2004, 735--784.
\bibitem{Labesse} J.-P.~Labesse, \emph{Changement de base CM et s\'eries discr\`etes}, in: On the
Stabilization of the Trace Formula, Int.\ Press, 2011, 429--470.
\bibitem{Prasanna} K.~Prasanna, \emph{Integrality of a ratio of Petersson norms and
level-lowering congruences}, Ann.\ of Math.\ (2) \textbf{163} (2006), 901--967.
\bibitem{Ribet1976} K.~A.~Ribet, \emph{A modular construction of unramified $p$-extensions of
$\mathbb Q(\mu_p)$}, Invent.\ Math.\ \textbf{34} (1976), 151--162.
\bibitem{SerreLF} J.-P.~Serre, \emph{Local Fields}, Graduate Texts in Mathematics 67, Springer,
New York, 1979.
\bibitem{Neukirch} J.~Neukirch, \emph{Algebraic Number Theory}, Grundlehren Math.\ Wiss.\
\textbf{322}, Springer, 1999.
\bibitem{BLGGT} T.~Barnet-Lamb, T.~Gee, D.~Geraghty and R.~Taylor, \emph{Potential automorphy and
change of weight}, Ann.\ of Math.\ (2) \textbf{179} (2014), 501--609.
\bibitem{BlochKato} S.~Bloch and K.~Kato, \emph{$L$-functions and Tamagawa numbers of motives},
in: The Grothendieck Festschrift, Vol.~I, Progr.\ Math.\ \textbf{86}, Birkh\"auser, 1990,
333--400.
\bibitem{FontaineLaffaille} J.-M.~Fontaine and G.~Laffaille, \emph{Construction de
repr\'esentations $p$-adiques}, Ann.\ Sci.\ \'Ecole Norm.\ Sup.\ (4) \textbf{15} (1982), 547--608.
\bibitem{NSW} J.~Neukirch, A.~Schmidt and K.~Wingberg, \emph{Cohomology of number fields}, 2nd
ed., Grundlehren Math.\ Wiss.\ \textbf{323}, Springer, 2008.
\bibitem{WilesFLT} A.~Wiles, \emph{Modular elliptic curves and Fermat's last theorem}, Ann.\ of
Math.\ (2) \textbf{141} (1995), 443--551.
\bibitem{MilneADT} J.~S.~Milne, \emph{Arithmetic Duality Theorems}, Academic Press, Boston, 1986 (1st ed.; all numbers cited refer to this edition), \url{https://www.jmilne.org/math/Books/adt1.pdf}.
\bibitem{Amitsur} S.~A.~Amitsur, \emph{On the characteristic polynomial of a sum of matrices},
Linear and Multilinear Algebra \textbf{8} (1979/80), 177--182.
\bibitem{LangCyclo} S.~Lang, \emph{Cyclotomic fields I and II}, combined 2nd ed., with an appendix
by K.~Rubin, Grad.\ Texts in Math.\ \textbf{121}, Springer, 1990.
\bibitem{RubinES} K.~Rubin, \emph{Euler systems}, Annals of Mathematics Studies \textbf{147},
Princeton University Press, 2000.
\bibitem{HarrisTaylor} M.~Harris and R.~Taylor, \emph{The geometry and cohomology of some simple
Shimura varieties}, Ann.\ of Math.\ Stud.\ \textbf{151}, Princeton Univ.\ Press, 2001.
\bibitem{Shin} S.~W.~Shin, \emph{Galois representations arising from some compact Shimura
varieties}, Ann.\ of Math.\ (2) \textbf{173} (2011), 1645--1741.
\bibitem{ChenevierHarris} G.~Chenevier and M.~Harris, \emph{Construction of automorphic Galois
representations, II}, Camb.\ J.\ Math.\ \textbf{1} (2013), 53--73.
\bibitem{CaraianiDuke} A.~Caraiani, \emph{Local-global compatibility and the action of monodromy
on nearby cycles}, Duke Math.\ J.\ \textbf{161} (2012), 2311--2413.
\bibitem{CaraianiANT} A.~Caraiani, \emph{Monodromy and local-global compatibility for $l=p$},
Algebra Number Theory \textbf{8} (2014), 1597--1646.
\bibitem{BLGGTlp1} T.~Barnet-Lamb, T.~Gee, D.~Geraghty and R.~Taylor, \emph{Local-global
compatibility for $l=p$, I}, Ann.\ Fac.\ Sci.\ Toulouse Math.\ (6) \textbf{21} (2012), 57--92.
\bibitem{BLGGTlp2} T.~Barnet-Lamb, T.~Gee, D.~Geraghty and R.~Taylor, \emph{Local-global
compatibility for $l=p$, II}, Ann.\ Sci.\ \'Ec.\ Norm.\ Sup\'er.\ (4) \textbf{47} (2014), 165--179.
\bibitem{MoyPrasad} A.~Moy and G.~Prasad, \emph{Jacquet functors and unrefined minimal
$K$-types}, Comment.\ Math.\ Helv.\ \textbf{71} (1996), 98--121.
\bibitem{MoyPrasad94} A.~Moy and G.~Prasad, \emph{Unrefined minimal $K$-types for $p$-adic
groups}, Invent.\ Math.\ \textbf{116} (1994), 393--408.
\bibitem{Zelevinsky} A.~V.~Zelevinsky, \emph{Induced representations of reductive $p$-adic
groups.~II. On irreducible representations of $\mathrm{GL}(n)$}, Ann.\ Sci.\ \'Ecole Norm.\ Sup.\
(4) \textbf{13} (1980), 165--210.
\bibitem{BHtame} C.~J.~Bushnell and G.~Henniart, \emph{The essentially tame local Langlands
correspondence, I}, J.\ Amer.\ Math.\ Soc.\ \textbf{18} (2005), 685--710.
\bibitem{MW89} C.~M\oe glin and J.-L.~Waldspurger, \emph{Le spectre r\'esiduel de $\mathrm{GL}(n)$},
Ann.\ Sci.\ \'Ecole Norm.\ Sup.\ (4) \textbf{22} (1989), 605--674.
\bibitem{BHeffective} C.~J.~Bushnell and G.~Henniart, \emph{To an effective local Langlands
correspondence}, arXiv:1103.5316v2 (published in Mem.\ Amer.\ Math.\ Soc.).
\bibitem{CasselmanNotes} W.~Casselman, \emph{Introduction to the theory of admissible
representations of $p$-adic reductive groups}, unpublished notes (1974/1995).
\bibitem{ZelevinskyLNM} A.~V.~Zelevinsky, \emph{Representations of finite classical groups. A Hopf
algebra approach}, Lecture Notes in Math.\ 869, Springer, Berlin, 1981.
\bibitem{RubinTexel} K.~Rubin, \emph{Kolyvagin's system of Gauss sums}, in: Arithmetic Algebraic
Geometry (Texel, 1989), Progr.\ Math.\ \textbf{89}, Birkh\"auser, 1991, 309--324.
\bibitem{DDT} H.~Darmon, F.~Diamond and R.~Taylor, \emph{Fermat's Last Theorem}, in: Elliptic
curves, modular forms \& Fermat's last theorem (Hong Kong, 1993), International Press, Cambridge,
MA, 1997, 2--140.
\bibitem{ReutenauerSchutzenberger} C.~Reutenauer and M.-P.~Sch\"utzenberger, \emph{A formula for the
determinant of a sum of matrices}, Lett.\ Math.\ Phys.\ \textbf{13} (1987), 299--302.
\bibitem{TadicENS} M.~Tadi\'c, \emph{Classification of unitary representations in irreducible
representations of general linear group (non-Archimedean case)}, Ann.\ Sci.\ \'Ecole Norm.\ Sup.\ (4)
\textbf{19} (1986), 335--382.
\bibitem{CurtisReiner} C.~W.~Curtis and I.~Reiner, \emph{Representation theory of finite groups and
associative algebras}, Interscience, 1962.

\bibitem{DigneMichel} F.~Digne and J.~Michel, \emph{Representations of finite groups of Lie type},
London Math.\ Soc.\ Student Texts \textbf{21}, Cambridge University Press, 1991.
\end{thebibliography}

\end{document}
